\documentclass[12pt,reqno]{amsart}

\usepackage[utf8]{inputenc}
\usepackage{listings}
\usepackage{amssymb}
\usepackage{amsmath}
\usepackage{amsthm}
\usepackage{amsfonts}
\usepackage{bbm}
\usepackage{bookmark}
\usepackage{hyperref}
\hypersetup{pdfstartview={FitH}}
\usepackage{color}
\usepackage{xcolor}
\usepackage{mathrsfs} 
\usepackage{enumitem}
\usepackage[normalem]{ulem}
\usepackage{tikz}
\usepackage{tikz-3dplot}
\usetikzlibrary{trees,shapes,backgrounds}
\usepackage{caption}
\usetikzlibrary{arrows, arrows.meta, calc, positioning}
\usetikzlibrary{shapes.multipart}
\usepackage{graphicx}
\usepackage{accents}
\usepackage{calc}
%\usetikzlibrary{external} 
%\tikzexternalize

\usepackage{fontawesome5} % This is for \faCar
\newcommand{\auto}{{({\tiny\faCar})}}

\addtolength{\hoffset}{-1.5cm}
\addtolength{\textwidth}{3cm}
\addtolength{\voffset}{-1cm}
\addtolength{\textheight}{1.2cm}

% \usepackage[maxbibnames=8,url=false]{biblatex}
% \addbibresource{arxivjul23.bib}

\usepackage{showkeys}
%\usepackage{refcheck}
 
\theoremstyle{plain}
\newtheorem{theorem}{Theorem}[section]
\newtheorem{lemma}[theorem]{Lemma}
\newtheorem{prop}[theorem]{Proposition}
\newtheorem{corollary}[theorem]{Corollary}
\newtheorem{proposition}[theorem]{Proposition}
%\theoremstyle{definition}
\newtheorem{definition}[theorem]{Definition}
\newtheorem{example}[theorem]{Example}
 
\numberwithin{equation}{section}
\newtheorem*{theorem*}{Theorem} 

\theoremstyle{remark}
\newtheorem{rem}{Remark}[section]


\newcommand{\Z}{{\mathbb Z}}
\newcommand{\R}{{\mathbb R}}
\newcommand{\Rm}{{\mathbb R}^m}
\newcommand{\N}{{\mathbb N}}
\newcommand{\C}{{\mathbb C}}
\newcommand{\D}{{\mathbb D}}
\newcommand{\T}{{\mathcal T}}
\newcommand{\Sp}{{\mathcal S}}
\newcommand{\I}{\mathcal{I}}
\newcommand{\K}{\mathcal{K}}
\newcommand{\M}{\mathbf{M}}
\newcommand{\F}{\mathbf{F}}


\def\wh{\widehat}
\def\det{\operatorname{det}}
\def\wt{\widetilde}
 \def\nct{{\boldsymbol \nu}}
\def\dker{{\boldsymbol \delta}}

\DeclareMathOperator{\II}{II}
\DeclareMathOperator{\dist}{dist}
\DeclareMathOperator{\supp}{supp}

\DeclareMathOperator{\size}{size}
\DeclareMathOperator{\ecc}{Ecc}
\DeclareMathOperator{\Mod}{Mod}

\providecommand{\la}[1]{ \langle #1 \rangle}

\providecommand{\s}[1]{\mathscr{#1}}
\providecommand{\trm}[1]{\textrm{#1}}
\providecommand{\rd}[1]{\textcolor{red}{#1}}
\newcommand{\comment}[1]{\marginpar{\scriptsize \textbf{comment:} #1}}


\newcommand{\pd}[1]{{\color{blue}{[P: #1]}}}
\newcommand{\ls}[1]{{\color{red}{[L: #1]}}}
\newcommand{\ct}[1]{{\color{magenta}{[C: #1]}}}
\newcommand{\jr}[1]{{\color{orange}{[J: #1]}}}
\newcommand{\cjr}{\tikz\fill[orange,scale=0.4](0,.35) -- (.25,0) -- (1,.7) -- (.25,.15) -- cycle;} 
\newcommand{\cpd}{\tikz\fill[blue,scale=0.4](0,.35) -- (.25,0) -- (1,.7) -- (.25,.15) -- cycle;} 
\newcommand{\cls}{\tikz\fill[red,scale=0.4](0,.35) -- (.25,0) -- (1,.7) -- (.25,.15) -- cycle;} 
\newcommand{\todo}[1]{{\color{green}{[TODO: #1]}}}

\DeclareFontFamily{U}{mathx}{\hyphenchar\font45}
\DeclareFontShape{U}{mathx}{m}{n}{
<5> <6> <7> <8> <9> <10>
<10.95> <12> <14.4> <17.28> <20.74> <24.88>
mathx10
}{}
\DeclareSymbolFont{mathx}{U}{mathx}{m}{n}
\DeclareFontSubstitution{U}{mathx}{m}{n}
\DeclareMathAccent{\widecheck}{0}{mathx}{"71}

\title[Norm-variation of multiple ergodic averages]{A blueprint for the formalization of norm-variation of multiple ergodic averages for commuting transformations}


 

\date{\today}


\begin{document}
\tdplotsetmaincoords{70}{110}

\begin{abstract}
\jr{Work in progress}
\end{abstract}
 
\maketitle

Conventions:
\begin{itemize}
\item All constants will be explicit and given names indexed by lemma number (as already done in Section \ref{sec:pfmainthm}). Constant definitions should generally be recursive in terms of previously defined constants. This ensures that typos in constants stay local, simplifies formalization and helps clarify dependencies. Dependencies of constants on parameters must explicitly appear in notation, write e.g. $C_{a,b,c}$ if $C$ depends on parameters $a,b,c$. This does not include dependence on previously defined constants.
When the constant is invoked, parameters should be explicitly instantiated as needed. 
%An exception to this rule are dependencies on the ambient dimension $n$ and the constant $C_0$ defined in Section \ref{sec:prelim red}. 
%These are project-wide constants and dependence should be left implicit in notation. E.g. simply write $C$ even if $C$ depends on $n$.
\item Generic smooth bump functions are denoted $\phi$, generic mean zero functions are denoted $\psi$.
\item After the introduction and overview sections we will use strict forward reasoning: no statement can be used before it is proved and all proofs must immediately follow statements. Exception: Machine-generated proofs are located in the supplemental section.
\item All theorems used must be referenced explicitly using {\texttt{ref}}. This will simplify automating creation of the metadata for the dependency graph.
\item We use strict separation of human-generated and machine-generated content. All machine-generated content must be in Section \ref{sec:machine generated}.
\end{itemize}


\section{Introduction}

\begin{definition}[Ergodic averages]\label{ergodic averages}
Let $(X,\Sigma,\mu)$ be a $\sigma$-finite measure space. We say that $T: X \to X$ is a measure preserving transformation if for every measurable $A\subset X$ the set $T^{-1}(A)$ is measurable and has the
same measure as $A$. 

Let $n\geq 1$ be an integer. Let $\mathbf{f}=(f_l)_{l\in [n)}$ be an $n$-tuple of complex-valued measurable functions  and let $(T_l)_{l\in [n)}$ be an $n$-tuple of measure preserving transformations on $X$. For an integer $N\geq 1$ and every $x\in X$, define the multiple ergodic averages
\begin{equation}\label{eq:averages-multiple}
M_{N}(\mathbf{f})(x) := \frac{1}{N}\sum_{i=0}^{N-1} \prod_{l=0}^{n-1} f_l(T_l^i x).
\end{equation}
\end{definition}

% \begin{definition}
% \[\| a \|_{V_{r}(I; B)}\]
% \end{definition}

The averages \eqref{eq:averages-multiple} are then also measurable functions $X\to \R$.

\begin{theorem}[Main ergodic theorem]\label{thm:ergodicthm}
Let $n\ge 2$ be an integer and let $r>2^{n-1}$, or $r\ge 2$ if $n=2$. There exists a $C \in (0,\infty)$ such that the following holds. 
Let $(X,\Sigma,\mu)$ be a $\sigma$-finite measure space,  $(T_l)_{l\in [n)}$ be an $n$-tuple of  mutually commuting measure preserving transformations
on $X$, and let $\mathbf{f}=(f_l)_{l\in [n)}$ be an $n$-tuple of complex-valued measurable functions on $X$ such that the Lebesgue seminorms $\|f_i\|_{L^{2n}(X)}$ are finite for all $i\in [n)$. 
Then
\begin{equation}
 \| M_{N} (\mathbf{f}) \|_{V_{r}(N\in\N_{\ge 1}; L^2(X))}
\le C 
 \prod_{l=0}^{n-1} \|f_l\|_{2n}.
\end{equation}
\end{theorem}
(See \S \ref{sec:variation-norm} for the definition of the $r$-variation norm.)

Since finite variation of a sequence implies convergence of the sequence and since $L^\infty(X)\subset L^{2n}(X)$ if $X$ is a probability space, Theorem \ref{thm:ergodicthm} immediately recovers Tao's theorem  \cite[Theorem 1.1]{Tao07}. %\jr{not $n=1$(?)}


The proof of this theorem can be reduced to the following {\em a priori} estimate over the real numbers. The proof of this latter estimate is the heart of the matter of this paper.

\begin{definition}[Twisted averages]
Let $n\geq 1$. For an $n$-tuple of real-valued Schwartz functions $\mathbf f = (f_0,f_1,\dots,f_{n-1})$ on $\R^n$, a bounded measurable function $\chi$ and $x\in\mathbb{R}^n$ denote
\begin{equation}\label{A_def}
A(\chi,\mathbf f)(x) = \int_{\mathbb{R}} \chi(s)\Big(\prod_{i\in [n)} f_i(x+se_i)\Big)\,ds.
\end{equation}
Also, $A_t(\chi,\mathbf{f})=A(\chi_{(t)},\mathbf{f})$, where $\chi_{(t)}(x)=t^{-1}\chi(t^{-1}x)$.
\end{definition}

\begin{theorem}[Main twisted theorem]\label{thm:nct main real}
    Let $n\ge 2$ be an integer. There exists a $C\in (0,\infty)$ such that the following holds. 
    For every positive integer $J$ and positive real numbers $t_0<t_1<\cdots < t_J$,  
    for every $n$-tuple of Schwartz functions ${\mathbf f}=(f_i)_{i \in [n)}$
    with \[ \|f_{i}\|_{2^{i+\min(n-i,2)}}=1 \] for $i\in [n)$, \begin{equation}\label{e:ncommuting_real} \sum_{j \in [J)} \|A_{t_{j + 1}}(\mathbf{1}_{[0,1]},\mathbf{f}) - A_{t_{j}}(\mathbf{1}_{[0,1]},\mathbf f)\|_{{L}^2(\R^n)}^2 \leq C J^{1-2^{-n+2}}. 
    \end{equation}
\end{theorem}

\jr{In the end we can insert an explicit constant in both the theorems}

\section{From real to ergodic}\label{sec:realtoergodic}\jr{working title}
The purpose of this section is to prove Theorem \ref{thm:ergodicthm} from Theorem \ref{thm:nct main real}.

Proposed plan: Theorem \ref{thm:nct main real} 

$\to$ permute exponents 

$\to$ density argument to extend to 
$L^p$ functions 

$\to$ apply multilinear interpolation 

$\to$ apply transference principle 

$\to$ ''generalize`` to complex-valued

$\to$ formulation with $r$-variation norm 

$\to$ specialize to symmetric exponents

Most of these steps will be implemented as propositions. A more precise variant of the final theorem will also be stated (general exponents, for fixed $J$).

There are two {\em external} inputs we will use in this reduction:

We do {\bf NOT} prove the transference principle, we only state it. Here we can cite, or depending on whether the statements align, at least refer to Felix Pernegger's thesis(?). The blueprint and formalization for that is a separate project.

We do {\bf NOT} prove the multilinear interpolation theorem, we only state it. We can cite an authoritative source, e.g. Bergh-L\"ofstrom \cite{BerghLofstrom}. The blueprint and formalization for that is a separate project.

\iffalse
\begin{theorem}\label{thm:nct main real-interpolation}
    Let $n\ge 2$ be an integer. There exists a $C\in \R_{>0}$ such that the following holds.  Let $(p_i)_{i\in [n)}$ be an $n$-tuple of exponents from the interval $[1,\infty]$ such that $\sum_{i\in [n)} p_i^{-1}=\frac{1}{2}$ and for every nonempty set $I \subset [n)$,
\[
\sum_{i\in I} p_i^{-1} \geq 2^{|I|-n-1},
\]
with the convention that $\infty^{-1}=0$.
    For every positive integer $J$ and positive real numbers $t_0<t_1<\cdots < t_J$,  
    for every $n$-tuple of measurable functions ${\mathbf f}=(f_i)_{i \in [n)}$
   such that $\|f_{l}\|_{p_l}$ is finite for $i\in [n)$, we have    \begin{equation}\label{e:ncommuting_real} \sum_{j \in [J)} \|A_{t_{j + 1}}(\mathbf{1}_{[0,1]},\mathbf{f}) - A_{t_{j}}(\mathbf{1}_{[0,1]},\mathbf f)\|_{{L}^2(\R^n)}^2 \leq C_{\ref{e:ncommuting_real}}J^{1-2^{-n+2}} \prod_{l=0}^{n-1} \|f_l\|^2_{L^{p_l}(\R^n)}. 
    \end{equation}
\end{theorem}

\begin{proof}
Since the Schwartz class is dense in $L^{k+\min(n-k,1)}(\R^n)$ for any $1\leq k \leq n$, inequality~\eqref{e:ncommuting_real} holds for any $n$-tuple of measurable functions $(f_i)_{i\in [n)}$ for which  $\|f_{k-1}\|_{2^{k+\min(n-k,1)}}=1$ for $1\le k\le n$.
\end{proof}
\fi


\begin{theorem}\label{thm:ergodicthm-interpolation}
Let $n\ge 2$ be an integer. There exists a $C \in \R_{>0}$ such that the following holds. Let $(p_i)_{i\in [n)}$ be an $n$-tuple of exponents from the interval $[1,\infty]$ such that $\sum_{i\in [n)} p_i^{-1}=\frac{1}{2}$ and for every nonempty set $I \subset [n)$,
\[
\sum_{i\in I} p_i^{-1} \geq 2^{|I|-n-1}.
\]
%with the convention that $\infty^{-1}=0$.
Let $(X,\Sigma,\mu)$ be a $\sigma$-finite measure space,  $(T_l)_{l\in [n)}$ be an $n$-tuple of  mutually commuting measure preserving transformations
on $X$, and let $(f_l)_{l\in [n)}$ be an $n$-tuple of measurable functions on $X$ such that the Lebesgue seminorms $\|f_{l}\|_{p_l}$ are finite for $l\in [n)$. Consider the averages \eqref{eq:averages-multiple}. 
For every positive integer $J$ and positive integers $N_0<N_1<\dots<N_J$, we have 
\begin{equation}
\sum_{j \in [J)}\|M_{N_{j+1}} - M_{N_{j}}\|_2^2 \leq CJ^{1-2^{2-n}} \prod_{l=0}^{n-1} \|f_l\|^2_{p_l}.
\end{equation}
\end{theorem}

\begin{theorem}\label{thm:ergodicthmJ}
Let $n\ge 2$ be an integer. There exists a $C \in \R_{>0}$ such that the following holds. 
Let $(X,\Sigma,\mu)$ be a $\sigma$-finite measure space,  $(T_l)_{l\in [n)}$ be an $n$-tuple of  mutually commuting measure preserving transformations
on $X$, and let $\mathbf{f}=(f_l)_{l\in [n)}$ be an $n$-tuple of measurable functions on $X$ such that the Lebesgue seminorms $\|f_i\|_{L^{2n}(X)}$ are finite for $0\le i< n$. Consider the averages \eqref{eq:averages-multiple}. 
For every positive integer $J$ and positive integers $N_0<N_1<\dots<N_J$, we have 
\begin{equation}
\sum_{j \in [J)}\|M_{N_{j+1}}(\mathbf{f}) - M_{N_{j}}(\mathbf{f})\|_2^2 \leq CJ^{1-2^{2-n}} \prod_{l=0}^{n-1} \|f_l\|^2_{2n}.
\end{equation}
\end{theorem}



\section{Overview}

\todo{To be written. This should be an informal overview of the ideas and overall proof structure, with explicit forward references to the precise statements}

\section{Preliminaries}\label{sec:prelim}

\jr{
Todo:

1. Print first part blue

1.1 Make sure whenever a proposition is used, it is referred to explicitly by reference. If a proposition is never referred to, either change that or remove/comment out the proposition.

1.2 Insert proofs where they are currently missing. Proofs should be blueprint-ready but err on the side of less details. If it is a self-contained "routine" fact, we can outsource it to AI. AI-written proofs must go in the supplements section at the end.

1.3 Insert the explicit recursive constants everywhere. We want all constants to be explicit.

\cjr 1.3.1. Constants in Section \ref{sec:prelim}

\cpd \cls 1.3.2. Constants in Section \ref{bump section}.

\cjr 1.3.3. Constants in Section \ref{sec:main argument}

1.3.4. Constants in Section \ref{sec:reduction}. Careful here: this was written with the convention that $n$-dependence doesn't have to be explicit. We don't want that anymore, all $n$ dependencies have to be explicit, because there is some hope we can get a very nice $n$ dependence (Theorem \ref{induct positive terms theorem} has $O(n^2)$) (PD) \pd{I started working on that, but it only makes sense to continue once the Gaussian domination has been completed,  the  section on windows has been revised, and $N$ is fixed once and for all}

1.3.5. Unfold the recursive constants to obtain an explicit constant depending on $n$ for Theorem \ref{thm:nct main real}

1.4 Enforce forward reasoning (should just be rearranging lemmas here and there)

2. Decide Gaussian domination refactor: unify? split? (JR)

% - Add $L^2$ convergence lemma

\cjr 3. Real vs complex valued

\cpd 3.1. remove remaining complex conjugates from \S \ref{sec:prelim red} onwards (check it's fine to remove and remove, also say explicitly that functions are real-valued where unclear)

% - Integrability/measurability/etc. conditions in separate Propositions? No

%- Insert appropriate $\Lambda$ to $\Lambda_1$ / $A$ to $\Lambda$ references 

%$\cpd$ 4. In Sections \ref{sec:pfmainthm}, \ref{sec:on-diag}, make appearance of functions $f$, $F$ explicit  %\pd{In Sec \ref{sec:pfmainthm} we fix $\mathbf{f}$ in the beginning and do not mention it in the notation. I didn't change that.}

%\cls 5. Add ergodic main theorem (statement only) 

%\cls 6. Interpolation of final result?

% $\cpd$ 7. In reduction argument fix $n=2$ issues 

\cjr 8. Change notation for Gaussian $g$ to $\mathfrak{g}$ or similar

\cpd 9. adapt Proof of Lemma \ref{special prism BL inequality} 

\cjr 10. explain the wlogs in proofs of Prop \ref{Cauchy-Schwarz at k}  and \ref{Cauchy-Schwarz at n-1}

\cpd 11. address comment in proof of Prop \ref{diagonal square root}

\cls 12. blueprint Prop. \ref{four scale Gaussian kernel}   

\cjr 13. explicit decay estimates for standard bump in Proposition \ref{standard bump properties}

\cjr 14. Consistently change notation in main argument: standard bump is not $\psi$ anymore, now it is $\Phi$. $\psi$ is for generic mean zero bump function.

15. Many definitions implicitly contain propositions relating to well-definedness (e.g. that something is integrable, measurable, or belonging to a space, or defined on the Fourier side). Here we split into Definition and proposition: the definition should not make any nontrivial mathematical claims (within reason). (LS) \ls{Done up to Section 5.}

16. Finish From real to ergodic section

\cjr 17. Two main theorems in introduction

18. Change any remaining uses of formulation `there is unique function $f$ such that $\widehat{f}=$` into `define $f=\mathcal{F}^{-1}(..)$` and then state well-definedness/properties
 

19. Check what from Sec \ref{sec:prelim red} should be moved to \ref{sec:prelim} (definitely everything that's used before to satisfy forward reasoning convention)

20. compare Sec \ref{sec:bump red} with Sec \ref{bump section} and decide what to do about the redundancy. (LS)

21. Finalize section structure and section naming. Maybe it is okay as is?

22. Add references, fill in discussion in introduction.

23. Write overview section

24. Clean up section on windows \S \ref{sec:windows} by addressing the comments 

25. fix/unify Schwartz vs test functions in the reduction section

26. Get rid of all occurrences of $\R_{>0}$ and similar subtypes of $\R$ (wirte out positivity assumption separately)

$\omega$. Move to git

$\omega+1$. Regenerate automatic section where needed (e.g. Gaussian domination)

$\omega+2$. Validation: proof correctness check and check for blueprint convention compliance

$\omega+3$. Decide on Lean metadata formatting and generate metadata for dependency graph.

}

\newcommand{\g}{\mathfrak{g}}

\subsection{Notation}
We use indexing starting at zero rather than starting at one. For example, if $x$ is a vector in $\R^n$ we write $x=(x_0,\dots,x_{n-1})$. By $\N$ we denote the natural numbers including zero.
For $n\in\N$ let
\[ [n)=\{0,\dots,n-1\}. \]
For a vector $x\in\R^n$ we denote by $|x|$ its standard Euclidean norm. For a multiindex $m\in\N^n$ we denote
\[ |m| = \sum_{i\in [n)}m_i. \]

For a subset $I\subset [n)$ (which inherits the natural order on $[m)$) and a $x\in \R^n$ we use notation
\[ x_I = (x_{i})_{i\in I} \in \mathbb{R}^{|I|} \]
and when $I\subset [n)$ and $i\in [n)$ we also write $I\setminus i$ for the set $I$ without the index $i$.

For $x\in \R^n$ let
\[\Sigma(x)=\sum_{i\in[n)} x_i.\]

For $y\in(\mathbb{R}^2)^n$ we write $y=(y^m_i)_{i\in [n), m\in [2)}$, $y^m=(y^m_i)_{i\in [n)}$ and similarly if $h\in [2)^n$, then write
\[ y^h = (y_i^{h_i})_{i\in [n)}. \]
Note there is no ambiguity in having $y^h$, $y^m$ mean different things, because the variables $m$, $h$ have different types.

For $i\in [n)$ we denote by $e_i$ the $i$th standard unit vector in $\R^n$, i.e. the vector with $(e_i)_{i'}=1$ if $i=i'$ and $=0$ otherwise.
Similarly, for $i\in [n), j\in [2)$ denote by $e_{i}^j\in (\R^2)^n$ the standard unit vector with $(e_{i}^j)_{i'}^{j'}=1$ if $i=i', j=j'$ and $=0$ otherwise.

The Fourier transform of an $L^1$ function on $\R^n$ will be defined by
\[\mathcal{F}f(\xi)= \widehat{f}(\xi) = \int_{\R^n} f(x) e^{-2\pi i x\cdot \xi} dx.\]
Here $x\cdot \xi=\sum_{i\in [n)} x_i \xi_i$ denotes the standard inner product on $\R^n$.
We have the usual properties of the Fourier transform. For example, $\widehat{f*g}=\widehat{f}\cdot \widehat{g}$.
The Fourier transform extends to an isometry on $L^2$: we have $\|\widehat{f}\|_2 = \|f\|_2$ for $f\in L^2$.
The inverse Fourier transform of an $L^1$ function on $\R^n$ is defined by
\[ \mathcal{F}^{-1}g(x)= \widecheck{g}(x) = \int_{\R^n} g(\xi) e^{2\pi i x\cdot \xi} d\xi.\]

If $\phi:\R^n\to \C, \psi:\R^m\to\C$ are functions, then $\phi\otimes \psi:\R^{n+m}\to\C$ denotes their tensor product
\[ (\phi\otimes \psi)(x,y) = \phi(x)\psi(y), \]
where we have identified $\R^{n+m}$ with $\R^n\times \R^m$. For $d\in\N$ also define
\[\chi^{\otimes d} = \underbrace{\chi \otimes \cdots \otimes \chi}_{d\;\text{times}}.\]

\begin{definition}\label{closed ball}
    For $n\in\N$ and $r>0$ and $x\in \R^n$ denote
\[ B_n(x,r) := \{ \xi\in\mathbb{R}^n\,:\, |\xi-x|\le r \} \]
\[ B_n(r) := B_n(0,r) \]

\end{definition}


\begin{definition}\label{gaussian}
The standard Gaussian on $\R$ is defined by
\[\g(x) = e^{-\pi x^2}\]
for $x\in\R^n$. This is chosen so that $\int_{\R^n} \g=1$.
\end{definition}
% The derivative of $\g$ is denoted by \jr{remove? needed?}
% \[h(x) = \g'(x) = -2\pi x \g(x)\]
% for $x\in\R$.

For $t>0$ and a function $\psi$ on $\R^n$ define the rescaled function
\begin{equation}
\psi_{(t)}(x)=t^{-n} \psi(t^{-1}x).
\end{equation}    
This is chosen so that $\int_{\R^n} \psi_{(t)} = \int_{\R^n} \psi$.

Denote by $\mathcal{S}(\R^n)$ the space of Schwartz functions on $\R^n$ and by $\mathcal{S}'(\R^n)$ the space of tempered distributions. The Fourier transform of a distribution $T\in\mathcal{S}'(\R^n)$ is defined as
\[ \mathcal{F}(T)(\phi) = T(\mathcal{F}\phi)\quad\text{for}\;\phi\in\mathcal{S}(\R^n). \]

Unless otherwise specified, all functions can be assumed real-valued.
We shall use freely the fact that if a function, or distribution $f$ on $\R^n$ is real-valued and even, then $\mathcal{F} f$ is also real-valued.



\begin{definition}[bracket bump]\label{bracket bump}
Define
\begin{equation}
    \left<x\right>:=(1+|x|)^{-1}\, .
\end{equation}    
We agree on the following parsing, which interprets the power first:
\begin{equation}
    \left<x\right>^N_{(s)}:=(\left< .\right>^N)_{(s)}(x)=s^{-1}(1+|s^{-1}x|)^{-N}\, .
\end{equation}  
\end{definition}


\subsection{The Wiener space \texorpdfstring{$W_0$}{}}
Let $\mathbb{K}$ be $\R$ or $\C$.
Functions in this section are $\mathbb{K}$-valued.


The  Wiener space $W$ is a space which locally behaves like $L^\infty$ and globally like $L^1$
\cite{Wiener32,Groechenig}. In this paper we will only use the space $W_0$, which is the closed subspace of continuous functions in $W$.
The exact definition follows below.

\begin{proposition}
Let $d\in \N$ and $f:\R^d\to \mathbb{K}$ be continuous.
Then for $r\in (0,\infty)$, the function $g:\R^d\to [0,\infty]$  defined by
\begin{equation}
    g(x)=\sup_{|x-y|\le r}|f(y)|
\end{equation} 
is finite and upper semi continuous and measurable.
\end{proposition}
\begin{proof}
 Pick $x\in \R^d$. Let $r>0$. 
 The function $f$ is continuous on the compact ball $B_d(x,2r)$ (Definition \ref{closed ball}) and thus bounded and uniformly continuous on that ball. It follows that $g$ is finite. By uniform continuity, for every $\lambda>g(x)$ there is an $\epsilon$ such that for $y\in B_d(x,r)$ and $z\in B_d(x,2r)$ with $|y-z|\le \epsilon$ we have 
 $|f(z)|<\lambda$. For $z\in B_d(x,r+\epsilon)$
 there is $y\in B(x,r)$ on the same ray emanating from $x$ as $z$ such that $|y-z|\le \epsilon$.
 Hence
 \begin{equation}
     \sup_{|y-x|\le r+\epsilon}|f(y)|\le \lambda \, .
 \end{equation}
 Hence, for  every $y\in B_d(x,\epsilon)$, we have
 $g(y)\le \lambda$. Upper semicontinuity of $g$ follows and then also measurability.
 \end{proof}

\begin{definition}
Let $d\in \N$.  Define 
\begin{equation}
\|f\|_{W_0,d,r}:=    \int_{\R^d} \sup_{|x-y|\le r} |f(y)|\,dx.
\end{equation}
As $d$ can be inferred from the type of $f$, we 
will omit it from the notation.
We also set 
\begin{equation}
    \|f\|_{W_0}:=\|f\|_{W_0,1}\, .
\end{equation}
Write $W_0(\R^d)$ or $W_0(d)$ for the set of continuous functions $f:\R^d\to\mathbb{K}$ so that
\[  \|f\|_{W_0}< \infty\, . \]
We also say that $f$ is in $W_0$ as again the type of $f$
determines $d$.
\end{definition}


\begin{proposition}\label{W_0 radius independence}
Let $d\in \N$. 
Let $0<r,s<\infty$ and set
\begin{equation}
    C_{\ref{W_0 radius independence},d,r,s}:=\max(1,(3s/r)^d)\, .
\end{equation}

Then for all continuous
$f:\R^d\to \mathbb{K}$ we have
\begin{equation}\label{W_0(r) property}
    \|f\|_{W_0,s}\le C_{\ref{W_0 radius independence},d,r,s}\|f\|_{W_0,r}\, .
\end{equation} 
In particular, for every $0<r<\infty$, a  continuous
$f:\R^d\to \mathbb{K}$ is in $W_0$ if and only if $\|f\|_{W_0,r}<\infty$.
\end{proposition}

\begin{proof}
If $0<s\le r<\infty$, the statement \eqref{W_0(r) property} follows with factor $1$  by nesting of the suprema. Now assume $0<r<s<\infty.$
Let $X$ be a set of points in $B_d(s)$ of
mutual distance at least $r$. As the open balls of radius $r/2$ about these points are disjoint and contained in $B_d(3s/2)$, volume comparison shows there
are at most $(3s/r)^d$ many points in $X$.
We may therefore pick a maximal set $X$ of such points. The closed balls of radius $r$ about these 
points cover $B_d(s)$, because any uncovered point would contradict maximality of $X$.

By translation, for every $x\in \R^d$, 
the closed ball $B_d(x,s)$ is covered by the closed balls
$B_d(x+z,r)$ with $z\in X$. It follows that
\begin{equation*}
    \|f\|_{W_0,s}\le \int_{\R^d} \sup_{z\in X} \left(\sup_{|x+z-y|\le r} |f(y)|\right)\,dx\, .
\end{equation*}
Estimating the supremum in $X$ by the sum in $X$ and pulling the sum out of the integral, we estimate the last display using translation invariance of the integral by
\begin{equation*}
         \sum_{z\in X} \int_{\R^d} \sup_{|x+z-y|\le r} |f(y)|\,dx \le \sum_{z\in X} \int_{\R^d} \sup_{|x-y|\le r} |f(y)|\,dx\le (3s/r)^d \|f\|_{W_0,r}\, .
\end{equation*}
This completes the proof of \eqref{W_0(r) property}. 


The characterization of $W_0(X)$ by arbitrary $r$ follows.



\end{proof}


\begin{proposition}\label{P:lp-embedding}
Let $d\in \N$ be given and set
\begin{equation*}
    C_{\ref{P:lp-embedding},d}:=\pi^{d/2}\Gamma(d/2+1)^{-1}\, .
\end{equation*}
If $f\in W_0(d)$ and $1\le p\le \infty$, then
\begin{equation}\label{Lp W_0}
\|f\|_p\le C_{\ref{P:lp-embedding},d}^{\frac 1p -1}\|f\|_{W_0}    \, .
\end{equation}
\end{proposition}

\begin{proof}
Let $f\in W_0(d)$. As $|f(x)|\le \sup_{|x-y|\le 1}|f(y)|$ for every $x\in \R^d$, integration yields
\begin{equation}
    \|f\|_1\le \|f\|_{W_0}\, .
\end{equation}
Moreover, for every $\lambda<\|f\|_\infty$  there is a point $x$ with $|f(x)|\ge \lambda$.
For every $y$ in  $B_d(x,1)$ (Definition \ref{closed ball}), we have
\begin{equation*}
    \sup_{|z-y|\le 1}
|f(z)|\ge \lambda \, .\end{equation*}
With $C_{\ref{P:lp-embedding},d}$  the volume of  $B_d(x,1)$, we obtain by integration
\begin{equation}
    \lambda C_{\ref{P:lp-embedding},d}\le \|f\|_{W_0}\, .
\end{equation}
Taking supremum over all $\lambda<
\|f\|_\infty$ gives
\begin{equation}
    \|f\|_\infty\le C_{\ref{P:lp-embedding},d}^{-1}\|f\|_{W_0}\, .
\end{equation}
Inequality \eqref{Lp W_0} follows by logarithmic convexity of the $L^p$ norms.
\end{proof}


\begin{proposition}\label{W_0 fiber integrals}
For every injective linear map $\pi:\R^{m}\times \R^{l}\to \R^n$ with $m,l\in\N$, $1\le m+l\le n$
there are $0<C_{\ref{W_0 fiber integrals},\pi,1},C_{\ref{W_0 fiber integrals},\pi,2}<\infty$ such that the following holds. For every $f\in W_0(n)$ and $u\in\R^m$, the function $h_u$ defined by
\begin{equation}
    h_u(v)=f(\pi(u,v))
\end{equation}
is in $W_0(l)$ with
\begin{equation}\label{W_0 on fiber}
    \|h_u\|_{W_0}\le C_{\ref{W_0 fiber integrals},\pi,1} \|f\|_{W_0}\, .
\end{equation}
The function $g:\Rm\to \C$ defined by
\[ g(u) := \int_{\R^l} f(\pi(u,v)) \,dv= \int_{\R^l} h_u(v) \,dv \]
is in $W_0(m)$ with
\begin{equation}\label{W_0 across fibers}
    \|g\|_{W_0}\le C_{\ref{W_0 fiber integrals},\pi,2}  \|f\|_{W_0}\, .
\end{equation}
\end{proposition}
 Note that $l,m,n$ are determined through the type of $\pi$, thus the two constants in the theorem
 may in particular depend on $l,m,n$.

\begin{proof}
The function $\pi$ is linear on finite dimensional spaces and thus has finite operator norm. The operator norm is non-zero as $\pi$ is injective. Then the function $(u,v)\to f(\pi(u,v))$ as composition of continuous maps is continuous
and thus continuous in each variable.

Boundedness of the operator norm  of $\pi$ in $u$ gives $\epsilon>0$, namely the inverse of the operator norm, such that $|u-u'|<\epsilon$ implies
\begin{equation}
|\pi(u',v)-\pi(u,v)|<1.    
\end{equation}

We have with $C_{\ref{P:lp-embedding},n}$ the volume of any ball of radius $1$ in $\R^n$,
introducing an artificial integral over a ball of radius $\epsilon$,
\begin{equation}\label{eq unif int}
    \int_{\R^l} \sup_{|u'-u|\le\epsilon}|f(\pi(u',v))| \,dv 
\end{equation}
\begin{equation*}
   \le \epsilon^{-n}C_{\ref{P:lp-embedding},n}^{-1}\int_{\R^l}\sup_{|u'-u|\le\epsilon}\int_{\R^n} 
    1_{|x-\pi(u',v)|\le 1}|f(\pi(u',v))| \,dxdv\, .
\end{equation*}
Consolidating the geometric conditions by the triangle inequality and using Tonelli, we estimate this by
\begin{equation}
    \le \epsilon^{-n}C_{\ref{P:lp-embedding},n}^{-1}\int_{\R^n}\int_{\R^l} 1_{|x-\pi(u,v)|\le  2} \sup_{|x-y|\le 1}|f(y)| \,dvdx\, .
\end{equation}
As $\pi$ is linear and injective, there exists $R>0$ such that for each $x$ there exists $v'$ such that
$|x-\pi(u,v)|\le 2$ implies $|v-v'|^\le R$. Namely, this statement is void if there is no $v'$ with
$|x-\pi(u,v')|\le 2$. If there is such a $v'$, pick any such $v'$. Then the statement follows from
\begin{equation}
    |\pi(u,v)- \pi(u,v')|\le 4
\end{equation}
with  $R$ being four times the square root of the smallest eigenvalue of $\pi^T\pi$.


Hence we can estimate the last display by
\begin{equation}
   \le \epsilon^{-n}C_{\ref{P:lp-embedding},n}^{-1}R^lC_{\ref{P:lp-embedding},l}\int_{\R^n}  \sup_{|x-y|<1}|f(y)| \,dx\, .
\end{equation}
This gives \eqref{W_0 on fiber} with 
\begin{equation}
    C_{\ref{W_0 fiber integrals},\pi,1}:=\epsilon^{-n}C_{\ref{P:lp-embedding},n}^{-1}R^lC_{\ref{P:lp-embedding},l}
\end{equation}
and thus $h_u\in W_0(\R^l)$.

Finiteness of \eqref{eq unif int} also
allows to apply Lebegue dominated convergence in the vicinity of $u$. As $f(\pi(u',v))\to f(\pi(u,v))$
for every $v$ and $u'\to u$, we thus see that  
\begin{equation}
    \lim_{u'\to u}g(u')=g(u)
\end{equation}
and thus $g$ is continuous.

Integrating \eqref{eq unif int} over $u$ and repeating similar steps as above gives
\begin{equation*}
   \|g\|_{W_0,\epsilon}\le \epsilon^{-n}C_{\ref{P:lp-embedding},n}^{-1}\int_{\R^m} \int_{\R^l}\sup_{|u'-u|<\epsilon}\int_{\R^n} 
    1_{|x-\pi(u',v)|\le 1}|f(\pi(u',v))| \,dxdvdu
\end{equation*}
\begin{equation*}
    \le \epsilon^{-n}C_{\ref{P:lp-embedding},n}^{-1}\int_{\R^n}\int_{\R^m}\int_{\R^l} 1_{|x-\pi(u,v)|\le 2} \sup_{|y-x|\le 1}|f(y)| \,dvdudx\, .
\end{equation*}
With the same $R$ as above, for each $x$ there exists $u',v'$ with 
$|x-\pi(u,v)|\le 2$ implies $|(u,v)-(u'v')|\le R$. Hence we estimate the last display by
\begin{equation*}
   \le \epsilon^{-n}C_{\ref{P:lp-embedding},n}^{-1}R^{m+l}C_{\ref{P:lp-embedding},m+l}\int_{x\in \R^n}  \sup_{|x,y|\le 1}|f(y)| \,dx\, <\infty.
\end{equation*}
This together with Proposition \ref{W_0 radius independence} proves \eqref{W_0 across fibers} with 
\begin{equation}
    C_{\ref{W_0 fiber integrals},\pi,2}:=C_{\ref{W_0 radius independence},d,\epsilon,1}\epsilon^{-n}C_{\ref{P:lp-embedding},n}^{-1}R^{m+l}C_{\ref{P:lp-embedding},m+l}
\end{equation}
 In particular, $g\in W_0$.
\end{proof}


\begin{proposition}\label{tensor Wiener}
    Let $J\in \N_{\ge 1}$. For $j\in [J)$, let $I_j\subset \N$ and assume the $I_j$ are pairwise disjoint.  
    Set $I:=\bigcup_{j\in [J)}I_j$. Let $f_j:\R^{|I_j|}\to \mathbb{K}$ in $W_0(|I_j|)$ and define
    \begin{equation}
        f(x_I)=\prod_{j\in [J)}f_j(x_{I_j})\, .
    \end{equation}
    Then $f\in W_0(|I|)$ and
    \begin{equation}\label{product W_0}
        \|f\|_{W_0}\le \prod_{j\in [J)}\|f_j\|_{W_0}\, .
    \end{equation}
\end{proposition}
\begin{proof}
    Define $g_j(x_I):=f_j(x_{I_j})$. Then $g_j$ is continuous as composition of the continuous
    projection $x_I\to x_{I_j}$ and the continuous map $f$. The function $f$ as product of
    continuous functions is also continuous.
We have
\begin{equation}
     \sup_{j\in [J)} |x_{I_j}-y_{I_j}|\le |x_I-y_I|\, ,
\end{equation}
and hence
\begin{equation}
    \sup_{|x_I-y_I|\le 1}|f(y_I)|\le \prod_{j\in [J)} \sup_{|x_{I}-y_{I}|\le 1}|f_j(y_{I_j})|\le 
    \prod_{j\in [J)} \sup_{|x_{I_j}-y_{I_j}|\le 1}|f_j(y_{I_j})|\, .
\end{equation}
It follows by Tonelli that
\begin{equation}
    \int_{\R^{|I|}}\sup_{|x_I-y_I|\le 1}|f(y_I)|dx_I\le 
    \prod_{j\in [J)} \int_{\R^{|I_j|}}\sup_{|x_{I_j}-y_{I_j}|\le 1}|f_j(y_{I_j})|dx_{I_j}\, ,
\end{equation}
    where each factor is finite by assumption and we have $f\in W_0$ and \eqref{product W_0}.
\end{proof}

\begin{proposition}\label{W_0 Brascamp Lieb}
 Let  $m,l\in\N$, $m+l\le n$ and assume $l\neq 0$.
Let $J\in \N_{\ge 1}$. For $j\in [J)$, let $l_j\in \N$ and let  $\Pi_j:\R^m\times \R^l \to \R^{l_j}$ 
be a linear map such that 
\begin{equation}\label{pi injective}
    \bigcap_{j\in [J)} {\rm ker}\Pi_j=\{0\}\, .
\end{equation}
Then there is $0<C_{\ref{W_0 Brascamp Lieb},(\Pi_j)_{j\in [J)}}<\infty$ (explicit below in the proof) such that the following holds.
For $j\in [J)$, let $f_j\in W_0(\R^{l_j})$. Then for each 
 $u\in \R^m$, the function 
 \begin{equation}
     v\to \prod_{j\in [J)}f_j(\Pi_{j}(u,v))
 \end{equation}
 is in $W_0(\R^l)$ and with $f$ defined by
    \begin{equation}
        f(u):=\int_{\R^l}\prod_{j\in [J)}f_j(\Pi_{j}(u,v))\, dv \, ,
    \end{equation}
    we have $f\in W_0(\R^{m})$ and
    \begin{equation}
        \|f\|_{W_0}\le C_{\ref{W_0 Brascamp Lieb},(\Pi_j)_{j\in [J)}} \prod_{j\in  [J)}\|f_j\|_{W_0}\, .
    \end{equation}
\end{proposition}
Note that $m=0$ is a case in the above proposition.
\begin{proof}
Let 
\begin{equation}
    I_j=[\sum_{h\le j}l_j)\setminus [\sum_{h<j}l_j)
\end{equation}
and $I=[\sum_{h<J} l_j)$.
Define 
$\Pi:\R^{m}\times \R^l\to \R^{|I|}$ by
\begin{equation}
    (\Pi(u,v))_{I_j}=(\Pi_{j}(u,v))\, .
\end{equation}
Then $\Pi$ is injective by \eqref{pi injective}.    

Define $g:\R^{|I|}\to \C$ by
\begin{equation}
    g(x_I)=\prod_{j\in [J)}f_j(x_{I_j})\, .
\end{equation}
By Proposition \ref{tensor Wiener}, $g\in W_0(\R^{|I|})$ with
\begin{equation}
    \|g\|_{W_0}\le \prod_{j\in[J)]}\|f_j\|_{W_0}\ .
\end{equation}
By Proposition \ref{W_0 fiber integrals}, $f\in W_0$ with
\begin{equation}
    \|f\|_{W_0}\le C_{\ref{W_0 fiber integrals},\Pi,2}\|g\|_{W_0}\, .
\end{equation}

This proves Proposition \ref{W_0 Brascamp Lieb} with 
\begin{equation}
    C_{\ref{W_0 Brascamp Lieb},(\Pi_j)_{j\in [J)}}:=C_{\ref{W_0 fiber integrals},\Pi,2}
\end{equation}.
\end{proof}

% \begin{proposition}
% The space $W_0(X)$ is a Banach space. If $f\in W_0(X)$ and $g$ is continuous on $X$ and $|g|\le |f|$ on $X$, then $g\in W_0(X)$.   
% \end{proposition}

\begin{proposition}\label{P:schwartz-into-wiener}
We have $\mathcal{S}(\R^n)\subset W_0(\R^n)$.
\end{proposition}



\subsubsection{Convolution along a vector}


\begin{definition}[Convolution along a vector]
For $\rho\in W_0(\R^n), \varphi\in W_0(\R)$ and a vector $\alpha\in\mathbb{R}^n$ define for $x\in\mathbb{R}^n$, 
\[ (\rho *_\alpha \varphi)(x) = \int_{\R} \rho(x-p\alpha) \varphi(p)\,dp \]
\end{definition}

\begin{proposition}[Properties of convolution along a vector]
\label{convolution vector}
For $\rho \in W_0(\R^n), \varphi\in W_0(\R), \alpha,\xi\in\mathbb{R}^n$ we have $\rho *_\alpha \varphi\in W_0(\R^n)$ and 
\[ \widehat{\rho *_\alpha \varphi}(\xi) = \widehat{\rho}(\xi)\widehat{\varphi}(\alpha\cdot\xi). \]
\end{proposition}

\begin{proof}
Apply Proposition \ref{W_0 Brascamp Lieb}. The left-hand side equals
\[\int_{\R^n} (\rho *_\alpha \varphi)(x)e^{-2\pi ix\cdot \xi} dx=\int_{\R^n}  \int_{\R} \rho(x-p\alpha) \varphi(p)\,dp\,e^{-2\pi ix\cdot \xi} dx\]
\[=\int_{\R^n}  \int_{\R} \rho(x) \varphi(p) e^{-2\pi i (x+p\alpha)\cdot \xi} dp\,dx,\]
which equals the right-hand side.
\end{proof}



\subsection{\texorpdfstring{$K$ kernels}{K kernels}}


\begin{definition}[normalized function tuples]\label{normalized function tuples}
 Let $\mathfrak{F}$ be the set of tuples 
$\F=(F_i)_{i\in [n)}$ of real-valued functions in $\mathcal{S}(\R^n)$
% satisfying for $1\le k\le n$ the normalization
%  \[ \|F_{k-1}\|_{2^{k+\min(n-k,1)}}=1 \]
% In other words,
%  \begin{equation}\label{f-normalization}
%   \|F_{k-1}\|_{2^{k+1}}=1\;\text{if}\;k\le n-1\;\text{and}\;\|F_{n-1}\|_{2^{n}}=1  \, . 
%  \end{equation} 
satisfying for $i\in [n)$ the normalization
 \[ \|F_{i}\|_{2^{i+\min(n-i,2)}}=1. \]
In other words,
 \begin{equation}\label{f-normalization}
  \|F_{i}\|_{2^{i+2}}=1\;\text{if}\;i\le n-2\;\text{and}\;\|F_{n-1}\|_{2^{n}}=1  \, . 
 \end{equation} 

\end{definition}



The word cube in the following definition stands for the parameter set $[2)^k$, which can be interpreted as a $k$ dimensional cube.

\begin{definition}[cube Brascamp--Lieb form]\label{cube Brascamp--Lieb}
Define for $k\in\N$ with $1 \le k\le  n$
and  $L\in W_0(\R^k)$ and $F\in W_0(\R^n)$
\[ \Theta_{{\rm cube}}(k,L)(F) = \int_{(\R^2)^k} \int_{\R^{n-k}} L(y^1-y^0) \Big(\prod_{h\in [2)^k} F(y^h, x_{[k,n)})\Big)^{2^{\min(n-k,1)}}\;dx_{[k,n)}\,dy.\]
\end{definition} 
Note that the integrand in $\Theta_{{\rm cube}}(k, L)(F)$ is in $W_0$ by Proposition \ref{W_0 Brascamp Lieb}.


\begin{proposition}[cube BL inequality]\label{cube BL inequality}
Let $k\in\N$ satisfy $1 \le k\le n$.  Let $L\in W_0(\R^k)$. Then, for all $F\in W_0(\R^n)$ with $\|F\|_{2^{k+\min(n-k,1)}}=1$,
\begin{equation}\label{cube-k-m-bound}
   |\Theta_{\rm cube}(k,L)(F)| \le\|L\|_1 .
\end{equation}
\end{proposition}

\begin{proof}
This follows from H\"older's inequality.
Changing variables $y^1\mapsto y^0 + u$ we have $y^h_i=y^0_i+h_i u_i$, so
\[ \Theta_{{\rm cube}}(k,L)(F) = \int_{\R^k} L(u) \Big(\int_{\R^k}\int_{\R^{n-k}} \Big(\prod_{h\in [2)^k} F((y^0_i+h_i u_i)_{i\in [k)}, x_{[k,n)})\Big)^{2^{\min(n-k,1)}}\;dx_{[k,n)}\,dy^0\Big)\,du \]
By the triangle inequality, $|\Theta_{{\rm cube}}(k,L)(F)|$ is no greater than $\|L\|_1$ times the supremum over $u\in\R^k$ of
\[ \int_{\R^n} \Big(\prod_{h\in [2)^k} |F(\pi_h(x))|\Big)^{2^{\min(n-k,1)}}\;dx, \]
where we have renamed $y^0$ to $x_{[0,k-1]}$ 
and set $\pi_h(x)=(x_i+h_i u_i)_{i\in [k)}, x_{[k,n)}$. By H\"older's inequality this is
\[ \le \Big(\prod_{h\in [2)^k} \|F\circ \pi_h\|_{2^{k+\min(n-k,1)}}\Big)^{2^{\min(n-k,1)}}, \]
which is equal to $1$ because each $\pi_h$ is an isometry and $\|F\|_{2^{k+\min(n-k,1)}}=1$.
\end{proof}


The word {\em prism} in what follows stands for the Cartesian product of a cube with a simplex.
The $2^k$ corners of the $k$ dimensional cubes are parameterized by $[2)^k$, 
while the $n-k+1$ corners of the simplex are
parameterized by integers in the interval $[k-1,n)$.
%\ls{$[k-1,n-1]$?}
Note that the simplex with two corners is an interval,
hence both cases $k=n-1$ and $k=n$ can be identified with a cube.



\begin{definition}[prism Brascamp--Lieb form]\label{prism Brascamp--Lieb}
Define for integers $1 \le k\le  n$ 
and  $K\in W_0(\R^{k+1})$ and $\F\in \mathfrak{F}$, 
% \[ \Theta_k(K)(\F) = \int_{(\R^{2})^k} \int_{\R^{n-k+1}} K(y^1-y^0, \Sigma(y^0) + \Sigma(x_{[k,n]})) 
% \prod_{\substack{h\in [2)^k,\\i\in [k,n]}} F_{i-1}( y^h, x_{[k,n]\setminus i})\;dx_{[k,n]}\;dy \]
% \jr{Use isntead (this is just renaming integration variables):}
\[ \Theta_k(K)(\F) = \int_{(\R^{2})^k} \int_{\R^{n-k+1}} K(y^1-y^0, \Sigma(y^0) + \Sigma(x_{[k-1,n)})) 
\prod_{\substack{h\in [2)^k,\\i\in [k-1,n)}} F_{i}( y^h, x_{[k-1,n)\setminus i})\;dx_{[k-1,n)}\;dy \]
\end{definition}

Note that the integrand in $\Theta_k(K)(\F)$ is in $W_0$ by Proposition \ref{W_0 Brascamp Lieb}.
The argument of $K$ in 
this definition is split as $\R^{k}\times \R$ 
while 
the argument of $F_{i-1}$  is split as $\R^{k}\times \R^{n-k}$. Recall that superscripts $0$ and $1$ as in $y^0$ and $y^1$ are not powers but just superscripts.

\begin{proposition}[prism BL inequality]\label{prism BL inequality}

Let $k\in\N$ satisfy $1\le k\le n$ and $K\in W_0(\R^{k+1})$ and $\F\in \mathfrak{F}$. Then
\begin{equation}\label{prism-k-bound}
   |\Theta_k(K)(\F)| \le \|K\|_1.
\end{equation}
\end{proposition}

\begin{proof}
This follows from H\"older's inequality.
Changing variables $(y^1,x_{k-1})\mapsto (u,s)$, where
\[
y^1=y^0+u,\qquad x_{k-1}=s-\Sigma(y^0)-\Sigma(x_{[k,n)}),
\]
we have $y_i^h=y_i^0+h_iu_i$ and
$\Sigma(y^0)+\Sigma(x_{[k-1,n)})=s$. For fixed $(u,s)$, rename
$(y^0,x_{[k,n)})$ as $x\in\R^n$ and set
\[
\pi_{h,i}(x)
=
\Big((x_l+h_lu_l)_{l\in [k)},
\big(s-\Sigma(x),x_{[k,n)}\big)_{[k-1,n)\setminus i}\Big).
\]
Then
\[
\Theta_k(K)(\F)
=
\int_{\R^{k+1}} K(u,s)
\int_{\R^n}
\prod_{\substack{h\in [2)^k\\ i\in [k-1,n)}}
F_i(\pi_{h,i}(x))
\,dx\,du\,ds.
\]
By the triangle inequality, $|\Theta_k(K)(\F)|$ is no greater than $\|K\|_1$ times the supremum over $(u,s)\in\R^{k+1}$ of
\[
\int_{\R^n}
\prod_{\substack{h\in [2)^k\\ i\in [k-1,n)}}
|F_i(\pi_{h,i}(x))|
\,dx.
\]
By H\"older's inequality this is no greater than
\[
\prod_{\substack{h\in [2)^k\\ i\in [k-1,n)}}
\|F_i\circ\pi_{h,i}\|_{2^{i+\min(n-i,2)}}.
\]
Each $\pi_{h,i}$ is invertible and measure preserving, so this is equal to one by the definition of $\mathfrak{F}$.
\end{proof}

The following proposition concerns the special case $k=n$.
\begin{proposition}[simplification 1 prism]\label{simplification 1 prism}
    Let $K\in W_0(\R^{n+1})$. If for all $z\in \R^{n+1}$ we have 
    \begin{equation}\label{1 prism mean zero assumption}
        \int_\R K(z+qe_n)\, dq=0\, ,
    \end{equation}
    then for all $\F\in \mathfrak{F}$,
    \begin{equation}
         \Theta_n(K)(\F)= 0\, .
    \end{equation}
\end{proposition}
\begin{proof}
The form is
\[ \Theta_{n}(K)(\F) = \int_{(\R^{2})^n} \int_{\R} K(y^1-y^0, \Sigma(y^0) + x_n) 
\prod_{\substack{h\in [2)^n}} F_{n-1}( y^h)\;dx_n\;dy. \]
Since the product $\prod_{\substack{h\in [2)^k}} F_{n-1}( y^h)$ does not depend on $x_n$, this is equal to
\[\Theta_n(K)(\F) = \int_{(\R^{2})^n} \Big( \prod_{\substack{h\in [2)^n}} F_{n-1}( y^h) \Big) \Big(\int_{\R} K(y^1-y^0, \Sigma(y^0) + x_n)\,dx_n\Big)\,dy. \]
Setting $z=(y^1-y^0, \Sigma(y^0))\in\R^{n+1}$ for every fixed $y\in (\R^2)^n$, we see that the inner integral over $x_n$ is equal to zero by the assumption \eqref{1 prism mean zero assumption}.
\end{proof}


In the next proposition we apply the Cauchy-Schwarz inequality as in the proof of Proposition \ref{prism BL inequality}, but this time carefully preserving in one of the terms on the right-hand side
a cancellation if the function $\phi_j$ has integral zero. 




\begin{proposition}[singly cancellative Cauchy-Schwarz]\label{single cancellative Cauchy-Schwarz}
Let $J\ge 1$ and $1\le k<n$ be integers. For
$(\rho_j)_{j\in [J)}\subset W_0(\R^{k+1})$, and
$(\phi_j)_{j\in [J)}\subset W_0(\R)$, define, for $u\in \R^{k+1}$,
\begin{equation}
    K(u):=\sum_{j\in[J)}\int_\R \rho_j(u_{[k)},u_k-q)\phi_j(q)\,dq\,.
\end{equation}
%$K(u)$ is well-defined, i.e. the integrand in the previous display is in $L^1(\R)$ for every $u\in\R^{k+1}$. 
Then $K\in W_0(\R^{k+1})$ by Proposition \ref{W_0 Brascamp Lieb}. 
For $u\in \R^{k+2}$, define
\begin{equation}
    \tilde{K}(u):=\sum_{j\in[J)}\int_\R
    |\rho_j|(u_{[k)},u_{k+1}-q)\phi_j(u_k+q)\phi_j(q)\,dq\, .
\end{equation}
Then $\tilde{K}\in W_0(\R^{k+2})$.
Moreover, for all $\F\in\mathfrak{F}$,
\begin{equation}
    |\Theta_k(K)(\F)|^2
    \le \Theta_{k+1}(\tilde{K})(\F)\sum_{j\in[J)}\|\rho_j\|_1\, .
\end{equation}
\end{proposition}

\begin{proof}
We compute, substituting $x_{k-1}$ by $x_{k-1}+q$ and then $q$ by $r-x_{k-1}$ in combination with suitable orders of integration in $dx_{k-1}$ and $dq$, $\Theta_k(K)(\F)$ equals
\begin{equation}\label{sccs0}
\begin{aligned}
\sum_{j\in[J)} & \int_{(\R^2)^k}\int_{\R^{n-k+1}}
\rho_j(y^1-y^0,\Sigma(y^0)+\Sigma(x_{[k-1,n)}))
\\
&\times
\Big(
\int_\R
\Big(
\prod_{\substack{h\in[2)^k\\ i\in[k,n)}}
F_{i}(y^h,r,x_{[k,n)\setminus i})
\Big)
\phi_j(r-x_{k-1})\,dr
\Big)
\\
&\times
\Big(
\prod_{h\in[2)^k}F_{k-1}(y^h,x_{[k,n)})
\Big)
dx_{[k-1,n)}\,dy\, .
\end{aligned}
\end{equation}
We estimate $\rho_j$ by $|\rho_j|$ and also the term in outmost brackets in \eqref{sccs0} by its absolute value $A$
and the last bracket in \eqref{sccs0} by its absolute value $B$.
We interpret the resulting expression as a pairing $\left< A , B \right>$ with respect to a positive measure.
Applying Cauchy-Schwarz, we estimate the modulus squared of \eqref{sccs0} by the product of two terms.
The first term $\left< A , A \right>$ is no greater than
\begin{equation}
\begin{aligned}
&\sum_{j\in[J)}\int_{(\R^2)^k}\int_{\R^{n-k+1}}
|\rho_j|(y^1-y^0,\Sigma(y^0)+\Sigma(x_{[k-1,n)}))
\\
&\times
\Big|
\int_\R
\Big(
\prod_{\substack{h\in[2)^k\\ i\in[k,n)}}
F_{i}(y^h,r,x_{[k,n)\setminus i})
\Big)
\phi_j(r-x_{k-1})\,dr
\Big|^2
dx_{[k-1,n)}\,dy,
\end{aligned}
\end{equation}
Using that the $F_i$ are real and renaming the integration variables $r$, which appear twice
through the square, once $y_k^0$ and once $y_k^1$,

\begin{equation}
\begin{aligned}
\sum_{j\in[J)}\int_{(\R^2)^{k+1}}\int_{\R^{n-k}}
\int_\R
|\rho_j|((y^1-y^0)_{[k)},\Sigma(y^0_{[k)})+x_{k-1}+\Sigma(x_{[k,n)}))
\\
\times
\Big(
\prod_{\substack{h\in[2)^{k+1}\\ i\in[k,n)}}
F_{i}(y^h,x_{[k,n)\setminus i})
\Big)
\phi_j(y_k^1-x_{k-1})\phi_j(y_k^0-x_{k-1})
\,dx_{k-1}\,dx_{[k,n)}\,dy
\end{aligned}
\end{equation}
This is equal to $\Theta_{k+1}(\tilde{K})(\F)$.
In the second identity, we substitute
$x_{k-1}$ by $x_{k-1}+y_k^0$ and then name the variable $x_{k-1}$ as $-q$.

The second term $\left< B , B \right>$, in Cauchy-Schwarz, after integrating the variable $x_{k-1}$ which only appears in the last argument of $\rho_j$,
is identified as
\begin{equation}
    \sum_{j\in[J)} \Theta_{\rm cube}(k,L_j)(F_{k-1})
\end{equation}
with
\begin{equation}
    L_j(u):=\int_\R |\rho_j(u,q)|\,dq
\end{equation}
and estimated by Proposition \ref{cube BL inequality}.
\end{proof}




If $k=n-1$, one may do an even more careful Cauchy-Schwarz, preserving cancellation in both terms on the right-hand-side.


\begin{proposition}[doubly cancellative Cauchy-Schwarz, $k=n-1$]\label{doubly cancellative Cauchy-Schwarz}
Let $J\ge 1$ be an integer. For $(\rho_j)_{j\in[J)}\subset W_0(\R^n)$ and 
$(\phi_j)_{j\in [J)}\subset W_0(\R)$ define for $u\in \R^n$,
\begin{equation}
    K(u):=\sum_{j\in [J)}\int_{\R^2} \rho_j(u_{[n-1)}, u_{n-1}-r-s) \phi_j(r) \phi_j(s) \, ds dr\, .
\end{equation}
Then $K\in W_0(\R^n)$.
For $u\in\R^{n+1}$ let
\begin{equation}
    \tilde{K}(u):=\sum_{j\in [J)}\int_\R |\rho_j|(u_{[n-1)}, u_{n}-q) \phi_j(u_{n-1}+q) \phi_j(q) \, dq\, .
\end{equation}
Then $\tilde{K}\in W_0(\R^{n+1})$ and for all $\F\in\mathfrak{F}$,
\begin{equation}
    |\Theta_{n-1}(K)(\F)| \le \sup_{\mathbf{\tilde{F}}\in\mathfrak{F}} |\Theta_{n}(\tilde{K})(\mathbf{\tilde{F}})|.
\end{equation}
\end{proposition}

\begin{proof}
We expand $\Theta_{n-1}(K)(\F)$, replacing $x_{n-2}$ by $x_{n-2}+r$ and $x_{n-1}$ by $x_{n-1}+s$ 
and then $r$ by $r-x_{n-2}$ and $s$ by $s-x_{n-1}$. We obtain that $\Theta_{n-1}(K)(\F)$ equals
\begin{equation}\label{ncs1}
\begin{aligned}
\sum_{j\in[J)} & \int_{(\R^2)^{n-1}} \int_{\R^2}
\rho_j(y^1-y^0, \Sigma(y^0)+\Sigma(x_{[n-2,n)}))  
\\
&\times \Big( \int_{\R} \Big(\prod_{h\in [2)^{n-1}} F_{n-1}(y^h,r)\Big)
\phi_j(r-x_{n-2})\,dr\Big)
\\
&\times \Big( \int_{\R} \Big(\prod_{h\in [2)^{n-1}} F_{n-2}(y^h,s)\Big)
\phi_j(s-x_{n-1})\,ds\Big)
    dx_{[n-2,n)}\,dy .
\end{aligned}
\end{equation}
We estimate $\rho_j$ by $|\rho_j|$ and also the two bracketed terms in \eqref{ncs1} by their absolute values.
We interpret the resulting expression as a pairing with respect to a positive measure.
Applying Cauchy-Schwarz, we estimate the modulus squared of \eqref{ncs1} by the product of two terms.

The first term is no greater than
\begin{equation}
\begin{aligned}
\sum_{j\in[J)} & \int_{(\R^2)^{n-1}} \int_{\R^2}
|\rho_j|(y^1-y^0, \Sigma(y^0)+\Sigma(x_{[n-2,n)}))
\\
&\times
\Big|
 \int_{\R} \Big(\prod_{h\in [2)^{n-1}} F_{n-1}(y^h,r)\Big)
\phi_j(r-x_{n-2})\,dr
\Big|^2
    dx_{[n-2,n)}\,dy .
\end{aligned}
\end{equation}
Using that the $F_i$ are real and renaming the integration variables $r$, which appear twice
through the square, once $y_{n-1}^0$ and once $y_{n-1}^1$, this is equal to
\begin{equation}
\begin{aligned}
\sum_{j\in[J)} & \int_{(\R^2)^n}\int_{\R^2}
|\rho_j|((y^1-y^0)_{[n-1)},\Sigma(y^0_{[n-1)})+x_{n-2}+x_{n-1})
\\
&\times
\Big(\prod_{h\in[2)^n}F_{n-1}(y^h)\Big)
\phi_j(y_{n-1}^1-x_{n-2})\phi_j(y_{n-1}^0-x_{n-2})
\,dx_{n-2}\,dx_{n-1}\,dy .
\end{aligned}
\end{equation}
In this identity, we substitute $x_{n-2}$ by $x_{n-2}+y_{n-1}^0$ and then name the variable $x_{n-2}$ as $-q$.
This identifies the last display as $\Theta_n(\tilde K)(\F)$.

The second term is no greater than
\begin{equation}
\begin{aligned}
\sum_{j\in[J)} & \int_{(\R^2)^{n-1}} \int_{\R^2}
|\rho_j|(y^1-y^0, \Sigma(y^0)+\Sigma(x_{[n-2,n)}))
\\
&\times
\Big|
 \int_{\R} \Big(\prod_{h\in [2)^{n-1}} F_{n-2}(y^h,s)\Big)
\phi_j(s-x_{n-1})\,ds
\Big|^2
    dx_{[n-2,n)}\,dy .
\end{aligned}
\end{equation}
Proceeding analogously, using that the $F_i$ are real and renaming the integration variables $s$, which appear twice
through the square, once $y_{n-1}^0$ and once $y_{n-1}^1$, this becomes
\begin{equation}
\begin{aligned}
\sum_{j\in[J)} & \int_{(\R^2)^n}\int_{\R^2}
|\rho_j|((y^1-y^0)_{[n-1)},\Sigma(y^0_{[n-1)})+x_{n-2}+x_{n-1})
\\
&\times
\Big(\prod_{h\in[2)^n}F_{n-2}(y^h)\Big)
\phi_j(y_{n-1}^1-x_{n-1})\phi_j(y_{n-1}^0-x_{n-1})
\,dx_{n-2}\,dx_{n-1}\,dy .
\end{aligned}
\end{equation}
In this identity, we substitute $x_{n-1}$ by $x_{n-1}+y_{n-1}^0$ and then name the variable $x_{n-1}$ as $-q$.
This identifies the last display as $\Theta_n(\tilde K)(\mathbf{\tilde F})$, where $\mathbf{\tilde F}$ arises from $\F$ by interchanging the functions $F_{n-2}$ and $F_{n-1}$.
Thus
\begin{equation}
    |\Theta_{n-1}(K)(\F)|^2
    \le \Theta_n(\tilde K)(\F)\Theta_n(\tilde K)(\mathbf{\tilde F}),
\end{equation}
which implies the claim since $\tilde{F}$ is again in $\mathfrak{F}$.
\end{proof}
    
The terms on the right-hand-side of Cauchy-Schwarz are by nature positive. This is 
stated independently in the following proposition.
\begin{proposition}[Positivity]\label{Positivity K}
 Let $1\le k\le n-1$. Let $\rho\in W_0(\R^{k+1})$ be non-negative and 
$\phi\in W_0(\R)$. Set for $u\in\R^{k+2}$ and $i\in [k+1)$
\begin{equation}
    K(u):=\int_\R \rho(u_{[k+1)\setminus i}, u_{k+1}-q) \phi(u_{i}+q) \phi(q) \, dq\, .
\end{equation}
Then $K\in W_0(\R^{k+2})$ and for all $\mathbf{F}\in\mathfrak{F}$
\begin{equation}
    \Theta_{k+1}(K)(\F)\ge 0\, .
\end{equation}
\end{proposition}

\begin{proof}
We have
\[
\begin{aligned}
 \Theta_{k+1}(K)(\F)
&= \int_{(\R^2)^{k+1}}\int_{\R^{n-k}}\int_\R
\rho((y^1-y^0)_{[k+1)\setminus i},\Sigma(y^0)+\Sigma(x_{[k,n)})-q)
\\
&\times
\phi(y_i^1-y_i^0+q)\phi(q)
\prod_{\substack{h\in [2)^{k+1}\\j\in[k,n)}}
F_j(y^h,x_{[k,n)\setminus j})
\,dq\,dx_{[k,n)}\,dy .
\end{aligned}
\]

Substituting \(q\) by \(y_i^0-q\) and integrating first in \(y_i^0,y_i^1\), this is equal to
\[
\begin{aligned}
&\int_{(\R^2)^{k}}\int_{\R^{n-k}}\int_\R
\rho((y^1-y^0)_{[k+1)\setminus i},
\Sigma(y^0_{[k+1)\setminus i})+q+\Sigma(x_{[k,n)}))
\\
&\times
\Big|
\int_\R
\Big(
\prod_{\substack{h\in [2)^{k+1}\\ h_i=0\\j\in[k,n)}}
F_j(y^h,x_{[k,n)\setminus j})
\Big)
\phi(y_i^0-q)\,dy_i^0
\Big|^2
\,dq\,dx_{[k,n)}\,dy_{[k+1)\setminus i},
\end{aligned}
\]
which is non-negative.
\end{proof}


\begin{proposition}[Monotonicity]\label{Monotonicity K}In the situation of Proposition \ref{Positivity K},
in particular, if $\rho\le \tilde{\rho}$ and 
\[ K_\rho(u):=\int_\R \rho(u_{[k)}, u_{k+1}-q) \phi(u_{k}+q) \phi(q) \, dq\]
then
\[ \Theta_{k+1}(K_\rho)(\F) \le \Theta_{k+1}(K_{\tilde{\rho}})(\F) \]
\end{proposition}
\begin{proof}
    Apply the previous proposition to $K_{\tilde{\rho}-\rho}$ and use linearity.
\end{proof}

\subsection{\texorpdfstring{$M$}{M} kernels}

\begin{proposition}[$M$ to $K$]\label{M to K}
Let $k\in\N$ satisfy $1\le k\le n$. Let $M\in W_0((\R^2)^k)$. 
Define for $z\in \R^{k+1}$,
\begin{equation}\label{kviam_2}
    K_k(M)(z)=\int_{\R^{k-1}}  M((z_{l}+p_{l},p_{l})_{l\in [k-1)}, z_{k-1}+z_{k}-\Sigma(p),z_{k}-\Sigma(p))\,dp.
\end{equation}
The integrand is in $W_0(\R^{k-1})$ for every $z\in\R^{k+1}$.
The argument of $M$ is identified with an element of $(\R^2)^{k}$.
Then $K_k(M)\in W_0(\R^{k+1})$ and
\[ \|K_k(M)\|_1 \le \|M\|_1. \]
\end{proposition}

\begin{proof}
The claims on membership in $W_0$ follow from 
Proposition \ref{W_0 Brascamp Lieb}.
The claimed inequality follows by the triangle inequality and Fubini-Tonelli's theorem.
\end{proof}

\begin{rem}
For example, for $k=1$, 
\[K_1(M)(z)=M(z_0+z_1,z_1).\]
\end{rem}

\begin{rem}
Note that $K_k(M)$ can also be written as
\begin{equation}\label{kviam_2'}
    K_k(M)(z)=\int_{\R^{k}}  M((z_{l}+p_{l},p_{l})_{l\in[k)}) \delta( z_{k}-\Sigma(p))\,dp.
\end{equation}
\end{rem}

\begin{definition}[prism form]
Let $1\le k\le n$. We define for $M\in W_0((\R^2)^k)$ and $\mathbf{F}\in\mathfrak{F}$,
\[\Lambda_k(M)(\mathbf{F}) := \Theta_k(K_k(M))(\mathbf{F}).\]
\end{definition}

\begin{proposition}[Positivity]\label{Positivity M}%\jr{may not need to keep both positivity statements}
Let $1\le k\le n$ and $i\in [k)$. Let $\tilde{M}\in W_0((\R^2)^{k-1})$ be nonnegative
and let $\phi\in W_0(\R)$ be real valued. Define for $y\in (\R^2)^k$,
\begin{equation}
    M(y)=\tilde{M}(y_{[k)\setminus i}) \phi^{\otimes 2}(y_{i}).
\end{equation}
Then $M\in W_0((\R^2)^k)$ and for all $\F\in {\mathfrak{F}}$,
\begin{equation}
    \Lambda_k(M)(\F)\ge 0.
\end{equation}
\end{proposition}

\begin{proof}
    Apply Proposition \ref{tensor Wiener} and Proposition \ref{Positivity K} after a change of variables.
\end{proof}


\begin{proposition}[Cauchy-Schwarz at $k$]\label{Cauchy-Schwarz at k} 
Let $k\in\N$ with $1\le k< n-1$ and  $J\in \N_{>0}$. For $j\in [J)$,  let $\rho_j\in W_0((\R^2)^k)$ and let $\varphi_j\in W_0(\R)$ be real valued with $\|\rho_j\|_1\le 1$.  
Let $i\in [k)$.
Let $M\in W_0(\R^k)$ and $\tilde{M}\in W_0(\R^{k+1})$ \pd{$M\in W_0((\R^2)^k)$ and $\tilde{M}\in W_0((\R^2)^{k+1})$?} be defined by
\begin{equation}\label{before CS}
M = \sum_{j\in [J)} \rho_j *_{e_{i}^0+e_{i}^1} \varphi_j,
\end{equation}
and for $y\in (\R^2)^{k+1}$
\begin{equation}
\tilde{M}(y)=\sum_{j\in [J)}|\rho_j(y_{[k)})| \varphi_j^{\otimes 2}(y_{k}).
\end{equation}
Then for all $\F\in\mathfrak{F}$,
\begin{equation}
|\Lambda_k(M)(\F)| \le |\Lambda_{k+1}(\tilde{M})(\F)|^{\frac12} J^{\frac12}.
\end{equation}
\end{proposition}
\begin{proof}
Without loss of generality we may set $i=k-1$ (replace $\rho_j$ with $\rho_j\circ \sigma$ where $\sigma:(\R^2)^k\to (\R^2)^k$ is the coordinate permutation switching the  $i$th pair with the $(k-1)$th pair, keeping the order inside the pair unchanged).

Thus from now on set $i=k-1$. Then
\[ K_k(M)(z) = \sum_{j\in [J)} \int_{\R^{k-1}} \int_\R \rho_j(((z_{l}+p_{l},p_{l})_{l\in [k-1)}, z_{k-1}+z_{k}-\Sigma(p) - q, z_{k}-\Sigma(p) - q) \varphi_j(q) dq\,dp \]
\[ = \sum_{j\in [J)} \int_{\R} K_k(\rho_j)(z_{[k)}, z_k-q) \varphi_j(q)\,dq. \]
Applying Proposition \ref{single cancellative Cauchy-Schwarz} with $K_k(\rho_j)$ in place of $\rho_j$, and applying Proposition \ref{Monotonicity K} with $\tilde{\rho_j}=K_k(|\rho_j|)\ge |K_k(\rho_j)|$ we obtain for $\F\in\mathfrak{F}$,  
\[ |\Lambda_k(M)(\F)| \le \Theta_{k+1}(\tilde{K})(\F)^{\frac12} J^{\frac12}, \]
where we used that $\|K_k(\rho_j)\|_1\le \|\rho_j\|_1\le 1$ and for $z\in\R^{k+2}$,
\[ \tilde{K}(z) = \sum_{j\in[J)}\int_\R
    K_k(|\rho_j|)(z_{[k)},z_{k+1}-q)\varphi_j(z_k+q)\varphi_j(q)\,dq \]
which equals $K_{k+1}(\tilde{M})(z)$.\end{proof}


\begin{proposition}[Cauchy-Schwarz at $n-1$]\label{Cauchy-Schwarz at n-1}
Let $J\in\N_{>0}$ and for $j\in [J)$, let 
$\rho_j\in W_0((\R^2)^{n-1})$  and $\varphi_j\in W_0(\R)$ be real valued.
Let $i \in [n-1)$.
Define
\begin{equation}\label{before CS 2}
    M=\sum_{j\in [J)} \rho_j *_{e_i^0+e_i^1} (\varphi_j * \varphi_j),
\end{equation}
and for $y\in (\R^2)^n$,
\begin{equation}
    \tilde{M}(y)=\sum_{j\in [J)} |\rho_j(y_{[n-1)})| \varphi_j^{\otimes 2}(y_{n-1}).
\end{equation}
Then $M\in W_0((\R^2)^{n-1})$ and $\tilde{M}\in W_0((\R^2)^n)$ and for all $\F\in\mathfrak{F}$
\begin{equation}
|\Lambda_{n-1}(M)(\F)| \le \sup_{\mathbf{\tilde{F}}\in\mathfrak{F}}|\Lambda_{n}(\tilde{M})(\mathbf{\tilde{F}})|.
\end{equation}
\end{proposition}


\begin{proof}Without loss of generality we may set $i=n-2$ (replace $\rho_j$ with $\rho_j\circ \sigma$ where $\sigma:(\R^2)^{n-1}\to (\R^2)^{n-1}$ is the coordinate permutation switching the  $i$th pair with the $(n-2)$th pair, keeping the order inside the pair unchanged).
Recall
\[ K_{n-1}(M)(z) = \int_{\R^{n-2}} M((z_l+p_l,p_l)_{l\in [n-2)}, z_{n-2} + z_{n-1} - \Sigma(p), z_{n-1} - \Sigma(p))\,dp. \]
For $z\in\R^n$,
% \[ M(y) = \sum_{j\in [J)} \int_\R \int_\R \rho_j(y_{[n-2)}, y_{n-2}^0 - q, y_{n-2}^1 - q) \varphi_j(q-r)\varphi_j(r)\,dr\,dq  \]
$K_{n-1}(M)(z)$ equals
\[ \sum_{j\in [J)} \int_{\R^{n-2}} \int_\R \int_\R \rho_j((z_l+p_l,p_l)_{l\in [n-2)}, z_{n-2} + z_{n-1} - \Sigma(p) - q, z_{n-1} - \Sigma(p) - q) \varphi_j(q-r)\varphi_j(r)\,dr\,dq\,dp.  \]
% Changing variables $q\mapsto z_{n-1}-\Sigma(p)-q$
% and renaming $q$ to $p_{n-2}$ and using Fubini's theorem, this equals
% \[ \sum_{j\in [J)} \int_{\R^{n-1}} \int_\R \rho_j((z_l+p_l,p_l)_{l\in [n-1)}) \varphi_j(z_{n-1}-\Sigma(p)-r)\varphi_j(r)\,dr\,dp. \]







Changing variables $q\mapsto q+r$ this is equal to
\[ \sum_{j\in [J)} \int_{\R^2} K_{n-1}(\rho_j)(z_{[n-1)}, z_{n-1} - q - r)\varphi_j(q)\varphi_j(r)\,d(q,r). \]

Applying Proposition \ref{doubly cancellative Cauchy-Schwarz} with $K_{n-1}(\rho_j)$ in place of $\rho_j$ gives  
\[ |\Lambda_{n-1}(M)(\F)| \le \sup_{\mathbf{\tilde{F}}\in\mathfrak{F}}|\Theta_n(\tilde{K})(\mathbf{\tilde{F}})|,  \]
where for $u\in \R^{n+1}$,
\[
\tilde{K}(u)=\sum_{j\in [J)}\int_\R |K_{n-1}(\rho_j)|(u_{[n-1)}, u_{n}-q)| \varphi_j(u_{n-1}+q) \varphi_j(q) \, dq.
\]
By Proposition \ref{Monotonicity K} we may replace $|K_{n-1}(\rho_j)|$ by $K_{n-1}(|\rho_j|)$ and then recover $K_n(\tilde{M})(u)$.
\end{proof}    

\subsection{Multiplicatively spaced monotone sequences}

\begin{definition}[Multiplicatively spaced monotone sequences]\label{multiplicatively spaced monotone sequences}
A sequence $a: \Z\to \R_{>0}$ is called multiplicatively spaced, if for each $j\in \Z$ 
\begin{equation}
    2a(j)\le a(j+1)\, .
\end{equation}
Let ${\rm A}$ be the set of such sequences.
\end{definition}



\begin{proposition}[Extension of sequences]\label{Extension of sequences}
Let $J\in \N_{>0}$ and $a:[J)\to \R_{>0}$ such that for all $ j\in [J-1)$ we have $a(j+1)\ge 2a(j)$.
Then there is a unique $b\in A$ such that for all $j\in [J)$ we have $a(j)=b(j)$ and for $j\ge J$ we have $b(j)=2^{j-J+1} a(J-1)$ and for $j<0$ we have $b(j)=2^j a(0)$\, .
\end{proposition}
For the proof see \S \ref{sec:auto multiplicative sequences}.


\begin{proposition}[Operations on spaced sequences]\label{Operations on spaced sequences}
    Let $a,b\in {\rm A}$. Then the following hold:

    (i) $\max(a,b)\in {\rm A}$
    
    (ii) $t\cdot a\in {\rm A}$ for every $t>0$.

    (iii) Let $n\in\mathbb{Z}$. If $c:\Z\to(0,\infty)$ is defined by $c(j)=a(j+n)$ for all $j\in\mathbb{Z}$, then $c\in {\rm A}$.

    (iv) If $c:\Z\to (0,\infty)$ is defined by $c(j)=\sqrt{a(j)^2+b(j)^2}$ for all $j\in\Z$, then $c\in {\rm A}$.
\end{proposition}

%\ls{For Definition~\ref{s multiplier}, we also need that $c(j):=\sqrt{a(j)^2+b(j)^2} \in A$ if $a\in A$ and $b\in A$.}
For the proof see \S \ref{sec:auto multiplicative sequences}.

\begin{definition}[Distance of spaced sequences]\label{Distance of spaced sequences}
Given $a,b\in {\rm A}$, define 
\begin{equation}
    \dist(a,b)=\min\{k\in\mathbb{Z}_{\ge 0}\,:\,\forall j, a(j-k)\le b(j)\le a(j+k)\}
\end{equation}
with the understanding that the value is $\infty$ if there is no such $k$.
\end{definition}

\begin{proposition}[Properties of distance of sequences]\label{Properties of distance of sequences}
    Let $a,b,c\in {\rm A}$.
    Then

    (i) If $\dist(a,b)=0$, then $a=b$.

    (ii) $\dist(a,b)=\dist(b,a)$
    
    (iii) $\dist(a,c)\le \dist(a,b)+\dist(b,c)$

    (iv) If $c(j)=b(j+1)$, then $\dist(a,c)\le \dist(a,b)+1$.

    (v) For all $t>0$, $\dist(t\cdot a,t\cdot b)=\dist(a,b)$
    
    (vi) For all $h\in\mathbb{Z}$, 
    $\dist(a,2^h a)\le |h|$.
\end{proposition}
For the proof see \S \ref{sec:auto multiplicative sequences}.

\begin{definition}[Closed balls in $\mathrm{A}$]\label{closed balls in A}
For $a\in\mathrm{A}$ and $r>0$ denote
\[B_{\mathrm{dist}}(a,r)=\{b\in\mathrm{A}\,:\,\dist(a,b)\le r\}.\]
\end{definition}

\subsection{Gaussians}

Recall that $\g(x)=e^{-\pi x^2}$ for $x\in\R$.

\begin{proposition}\label{square root one minus Gaussian}
    The function
    \begin{equation}
        f(x)=\sqrt{1-\g(x)}
    \end{equation}
    is well-defined
    using the nonnegative square root and is continuous on $\R$.
    For $|x|\le \frac 12$, we have
    \begin{equation}
        f(x)\ge \tfrac 12 |x|
    \end{equation}
    and for $|x|\ge \frac 12$ we have
 \begin{equation}
        1-\g(x)\le f(x)\le 1.
    \end{equation}
\end{proposition}
% For the  proof see \S \ref{sec:auto gaussian properties}. 

\begin{proof}%\jr{print it blue}
The function $f$ is well defined as the function $g$ takes values in $(0,1]$
and thus the square root is taken from a nonnegative number. 
Continuity follows from continuity of the square root and the gaussian. 
\end{proof}

\begin{proposition}\label{Gaussian bump decay}
For every $x\in\R$, $m,N\in\N$,
\[ |\g^{(m)}(x)|\le C_{\ref{Gaussian bump decay},m,N} \langle x\rangle^N, \]
where $C_{\ref{Gaussian bump decay},m,N}=(100(m+N+1))^{(m+N)/2}$. %10^{10^{10+m+n}}$. \jr{obviously way too much. AI can make it more rpecise}
\end{proposition}


For the  proof see \S \ref{sec:auto gaussian properties}.

\begin{proposition}[Elementary Gaussian properties]\label{Elementary Gaussian properties}

(i) $\g\in \mathcal{S}(\R)\subset W_0(\R)$

(ii) $\widehat{\g}=\g$. 

(iii) We have for $\lambda,\mu>0$
    \begin{equation}\label{eq:gaussconv}  \g_{(\lambda)}*\g_{(\mu)}=\g_{(\sqrt{\lambda^2+\mu^2})}\, .
    \end{equation}

(iv) We have for $\lambda>0$ 
 \begin{equation}\label{eq:scaleFT}
     \widehat{\g_{(\lambda)}}(\xi)=\g(\lambda \xi)
 \end{equation}
\end{proposition}
For the  proof see \S \ref{sec:auto gaussian properties}.

\begin{proposition}\label{poisson to abel}
Let
\begin{equation}
    p(x)=2(1+(2\pi x)^2)^{-1}.
\end{equation}
Then
\begin{equation}\label{eq: p hat}
    \widehat{p}(\xi)=e^{- |\xi|}.
\end{equation}
\end{proposition}
\begin{proof}[Proof of Proposition \ref{poisson to abel}]
For \(a\in\R\), set
\[
    I(a):=\int_0^\infty e^{-t}\cos(at)\,dt,
    \qquad
    J(a):=\int_0^\infty e^{-t}\sin(at)\,dt.
\]
Integration by parts, differentiating the trigonometric functions and integrating the exponential function gives
\[
    I(a)=1-aJ(a),
    \qquad
    J(a)=aI(a).
\]
Hence
\[
    I(a)=\frac{1}{1+a^2}.
\]
Using this for $a=2\pi x$ for \(x\in\R\)
\begin{equation}
    p(x)=2\int_0^\infty e^{-\xi}\cos(2\pi x\xi)\,d\xi  =\int_{\R}e^{-|\xi|}e^{2\pi i x\xi}\,d\xi    
\end{equation}
 The Fourier inversion formula gives \eqref{eq: p hat}. Namely, it is applicable because
 both $p$ and the right hand side of $\eqref{eq: p hat}$ are continuous and integrable.

\end{proof}





\begin{proposition}\label{auxiliary function B}
    With $p$ as in Proposition \ref{poisson to abel}, define $B:\R\to \R$ by
\begin{equation}
    B(\xi):=1-\sqrt{1-\g(\xi)}-\widehat{p}(\sqrt{\pi}\xi)\, .
\end{equation}
Then $B$ is smooth on $\R\setminus \{0\}$ and continuous on $\R$. The function $B'$ 
has a continuous extension to $\R$ which is integrable and the 
function $B''$, extended to be $0$ at $0$, is Borel measurable and integrable on $\R$. 
Moreover,
\begin{equation}\label{eq:B-explicit-norms}
    \|B\|_1\le 8,\qquad
    \|B'\|_1\le 20,\qquad
    \|B''\|_1\le 100.
\end{equation}
\end{proposition}

\begin{proof}[Proof sketch]
As $\g(\xi)\in (0,1]$ for all $\xi \in \R$, the square root
is taken from a nonnegative number and thus well defined. The function $B$ is continuous. It is smooth at $\xi\in \R$ unless possibly when 
$\g(\xi)=1$, which happens only at $\xi=0$. 

The function $p$ is integrable. Moreover, for $\g(\xi)<\frac 12$,
which happens for $|\xi|>1$, we have 
\begin{equation}
\tfrac 12 \le 1-\g(\xi)\le \sqrt{1-\g(\xi)}\le 1    
\end{equation}
and thus
\begin{equation}
|1-\sqrt{1-\g(\xi)}|\le \g(x)    
\end{equation}
As the right hand side is integrable, so is $B$.



For $\xi\neq 0$ we compute
\begin{equation}
    B'(\xi)=\frac{\g'(\xi)}{2\sqrt{1-\g(2^{-1}\pi^{-1/2}\xi)}} -\sqrt{\pi}p'(\sqrt{\pi}\xi)\, .
\end{equation}
Using Taylor expansion, we see that $B'$ extends continuously to $0$.
Using the lower bound on the square root seen above, we conclude integrability of this extension.
For $\xi\neq 0$ we compute
\begin{equation}
    B''(\xi)=\frac{\g''(\xi)}{2\sqrt{1-\g(\xi)}}+\frac{\g'(\xi)^2}{4(\sqrt{1-\g(\xi)})^3}-\pi p''(\sqrt{\pi}\xi)\, .
\end{equation}
Using Taylor expansion, we see that $B''$ remains bounded near $0$.
Using the lower bound on the square root seen above, we conclude integrability of $B''$.
\end{proof}
For more details and the proofs of the explicit bounds on $L^1$ norms see \S \ref{sec:auto gaussian properties}.

\begin{proposition}[square root of Gaussian decay]\label{square root of Gaussian decay}
For $x\in\R$ let
\[ \rho(x) = \mathcal{F}^{-1}(\xi \mapsto 1- \sqrt{1-\g(\xi)})(x). \]
This is a well-defined function in $W_0(\R)$ satisfying for all $x\in\R$,
\begin{equation} \label{sqr gauss C}
    0\le \rho(x) \le C_{\ref{square root of Gaussian decay}} \langle x\rangle^2,
\end{equation}  
where $C_{\ref{square root of Gaussian decay}}=100$.
\end{proposition}
\begin{proof}[Proof sketch]
$\rho$ is real-valued because it is the inverse Fourier transform of an even real-valued function.
By partial integration (for continuous and piecewise differentiable functions with integrable derivative) using the information about $B,B',B''$ in Proposition \ref{auxiliary function B}, we see that there is a  $C>0$ with ($\check B$ denoting the inverse Fourier transform)
\begin{equation}
    \sup_{x\in \R}|\check B(x)|, \; \sup_{x\in \R}|x\check B(x)|,\;
     \sup_{x\in \R}|x^2\check B(x)|\le C \, .
\end{equation}
The same holds for the function $p$ by Proposition \ref{poisson to abel}. Hence it holds for the function $\xi\mapsto 1-\sqrt{1-\g(\xi)}$.
By the Fourier inversion theorem, the desired estimate for $\rho$
holds.
\end{proof}
For more details on the proof  and the explicit constant see \S \ref{sec:auto gaussian properties}.

\subsection{Bumps and their estimates}\label{bump section}



\begin{lemma}
\label{lem:smoothdecay} 
Let $N\geq 2$ be an integer. Let $\zeta:\R \to \R$  be an even $N$ times continuously differentiable function with compact support. Then the function $\phi=\mathcal F^{-1} \zeta$ belongs to $W_0(\R)$ and we have for all $x\in \R$, $x\neq 0$ the estimate
%Let $N\ge 1$ be an integer and let $\phi$ be a $W_0(1)$  function with  $\widehat{\phi}$   compactly supported and  $N$ times continuously differentiable.
%$\phi:\R\to \R$ be a Schwartz function. 
%Then for $x\ne 0$, 
\[|\phi(x)|\le \min(\|\widehat{\phi}\|_1, \|\widehat{\phi}^{(N)}\|_1 (2\pi)^{-N} |x|^{-N}).\]
\end{lemma}

%\ct{this is unclear and needs to be clarified. Pick a proof that works.
%Where do you apply this and not \ref{smoothdecay 2}}

%\ls{I changed the assumption that $\phi$ is Schwartz to something weaker to make the lemma applicable in later proofs.}

\begin{proof}
By the assumptions on $\zeta$, we deduce that $\zeta \in L^1(\R)$. Thus, $\phi$ is a well-defined continuous function. Also, $\phi$ is real-valued since it is the inverse Fourier transform of an even real-valued function.

By definition, 
\begin{equation}
    \label{fourierinvphi}
    \phi(x) = \int_{\R} \zeta(\xi)e^{2\pi i x\xi} d\xi, 
\end{equation}
and  by the triangle inequality
\begin{equation}
    \label{xleone}
    |\phi(x)| = \Big | \int_{\R} \zeta(\xi)e^{2\pi i x\xi} d\xi \Big | \le \int_{\R} |\zeta(\xi)e^{2\pi i x\xi}| d\xi \le \|\zeta\|_1 .
\end{equation}


For the second estimate we integrate by parts  in \eqref{fourierinvphi} to obtain
\[\phi(x) = (2\pi i x)^{-N} \int_{\R} \zeta^{(N)}(\xi) e^{2\pi i x \xi} d\xi, \]
and thus by the triangle inequality, 
\[|\phi(x)| \le (2\pi)^{-N}|x|^{-N} \int_{\R}|\zeta^{(N)}(\xi) e^{2\pi i x \xi}| d\xi \]
\begin{equation}
    \label{xgeone}
    \le  (2\pi)^{-N} |x|^{-N} \|\zeta^{(N)}\|_1. 
\end{equation}
Combining \eqref{xleone} and \eqref{xgeone} gives that $\phi \in W_0(\R)$. By the Fourier inversion formula, we then conclude that $\zeta=\widehat{\phi}$.
\end{proof}

 
\begin{lemma}
    \label{lem: min and bracket} Let $N\ge 1$ be an integer, then
\[\min(1, |x|^{-N})\le 2^N   \langle x\rangle^{N}\, .\]
\end{lemma}
\begin{proof}
 We have, using that the mean is less than the maximum
    \[\min(1, |x|^{-N})=\max(1, |x|)^{-N}\le (2^{-1}(1+|x|))^{-N}=2^N  \langle x\rangle^{N}\, .\]
\end{proof}

We will mainly \ct{decide whether it is mainly or only, and then (with the anticipated decision) write only} \ls{In Lemma~\ref{L:second-gaussian-estimate}, we have a function with compactly supported Fourier transform, but we don't have a good control on the measure of the support, and thus need to use Lemma~\ref{lem:smoothdecay}.} use these lemmas when we have control on the $L^\infty$ norm and compact support of   the Fourier transform, so we state this case separately. 
 

\ct{I strengthened the following lemma by a factor $(2\pi)^N$ and improving the cosnatnt by a factor 2. The constant need no name, $2^N$ is shorter than the name.}\jr{I think the factor was probably omitted on purpose to simplify terms, see comment below}
\begin{lemma}
   \label{lem:smoothdecay2}
Let $N\geq 2$ be an integer. Let $\zeta:\R \to \R$  be an even $N$ times continuously differentiable function
 with compact support.
Then the function $\phi=\mathcal F^{-1} \zeta$ belongs to $W_0(\R)$ and we have for all $x\in \R$ the estimate 
\begin{equation}\label{E:smoothdecay2}
|\phi(x)|\leq C_{\ref{lem:smoothdecay2}, N} \max(\|\widehat{\phi}\|_\infty, (2\pi )^{-N}\|\widehat{\phi}^{(N)}\|_\infty)|\supp (\widehat{\phi})|   \langle x\rangle^{N}
\end{equation}
with $C_{\ref{lem:smoothdecay2}, N}=2^{N} $. 
\end{lemma}  


\begin{proof}
The fact that $\phi \in W_0(\R)$ and $\widehat{\phi}=\zeta$ follows from Lemma~\ref{lem:smoothdecay}. 
We have for any compactly supported bounded function $f:\R\to \R$ 
  \[\|f\|_1 \le \|f\|_\infty |\textup{supp}(f)| .\]
Using this twice and Lemma \ref{lem:smoothdecay} estimates $|\phi(x)|$ by 
\[\le \min(\|\widehat{\phi}\|_1, \|\widehat{\phi}^{(N)}\|_1 (2\pi)^{-N} |x|^{-N})
\le |\supp (\widehat{\phi})| \min(\|\widehat{\phi}\|_\infty, \|\widehat{\phi}^{(N)}\|_\infty (2\pi)^{-N} |x|^{-N})\, . \] 
  The claim now follows from Lemma \ref{lem: min and bracket}.
\end{proof}

 \jr{The constants in the inequality below don't conform to our conventions (maybe the conventions are not explicit enough). Every inequality should be $A\le C\cdot B$ where $A,B$ are target expressions reasonably free of explicit constants and $C$ is a numbered explicit constant (with explicit dependencies). It is worth reasonably minimizing the syntactic size of terms because the larger terms are the more cumbersome every step of the formalization. Attempting to faithfully track all the $2\pi$ factors only in the cases where they appear does not seem to be worth the trouble ($N$ will be a fixed number after all). 
 But since this lemma is used at most a couple of times and it is already written we can adopt a compromise and leave it as is here. But let us please not write more of these unnecessarily long terms in the future}
\begin{proposition}[mean value bump estimate]\label{mean value bump estimate 2}
Let $N\in\mathbb{N}_{\ge 1}$ and let $\rho$ be a $W_0(1)$ function with
$\widehat{\rho}$ supported in $[-1,1]$ and $\widehat{\rho}$ being $N$ times continuously differentiable.
%with \max_{0\le \nu\le N}\|\widehat{\rho}^{(N)}\|_\infty=1$ 
Then for all $x,y\in\mathbb{R}$
\[ |\rho(x+y)-\rho(x)| \le  2^{N+1}\max_{0\le \nu\le N}\|(2\pi)^{-\nu}\widehat{\rho}^{(\nu)}\|_\infty \min(1, 2\pi |y|) (\langle x+y\rangle^N + \langle x\rangle^N). \]



%$C_{\ref{mean value bump estimate},N}=2^{10N+10}$. 2
\end{proposition}
\begin{proof}
 Assume first the case that $2\pi |y| >1$. Then the claim holds by the triangle inequality and
Lemma \ref{lem:smoothdecay2}. Assume then the opposite case that $2 \pi |y|\leq 1$. For fixed $y$ consider $\tilde{\rho}(x)=\rho(x+y)-\rho(x)$. Then
\begin{equation}
    |\widehat{\tilde{\rho}}(\xi)|=|\widehat{\rho}(\xi)(e^{2\pi i\xi y}-1)|\le |\widehat{\rho}(\xi)||2 \pi y|
\end{equation}
and 
\begin{equation}
    |\widehat{\tilde{\rho}}^{(N)}(\xi)|\le |\widehat{\rho}^{(N)}(\xi)||2\pi y|+\sum_{\nu=0}^{N-1}|\widehat{\rho}^{(\nu )}(\xi)||2\pi y|^{N-\nu}
\end{equation}
and 
\begin{equation}
    (2\pi)^{-N}|\widehat{\tilde{\rho}}^{(N)}(\xi)|\le (2\pi)^{-N}|\widehat{\rho}^{(N)}(\xi)||2\pi y|+(\max_{0\le \nu <N}|(2\pi)^{-\nu}\widehat{\rho}^{(\nu )}(\xi)|)\sum_{\nu=0}^{N-1}|y|^{N-\nu}
\end{equation}
\begin{equation}
    \le (2\pi)^{-N}|\widehat{\rho}^{(N)}(\xi)||2\pi y|+(\max_{0\le \nu <N}|(2\pi)^{-\nu}\widehat{\rho}^{(\nu )}(\xi)|)|2\pi y|.
\end{equation}
The conclusion now follows from Lemma~\ref{lem:smoothdecay2}.
\end{proof}



\ct{ few comments: 

1) it appears a bit superfluous to compute Schwartz norms for compactly supported
functions (or FT thereof). You dont need two-parameter Schwartz norms, but
only sup norms of derivatives, i.e. a one parameter family of norms. See the above lemmas.


2) The construction using geometric decay is quite inefficient givin $2^{O(N^2)}$. 
One shoudl aim at $2^{O( N\log N)}$. Ofcourse for small $N$ the constants matter.
Presumably
less rapid decay is better. If one new upper
bound on the N needed, one could improve by convolving N indicators  of equal lentgh.
One woudl not hav Schwartz, but one never wanted that at the first palce
(one could make Scwart by concolving with very small further intervals. Screwing
up higher derivatives than N only.}





\begin{definition}[standard bump]\label{standard bump}
Define, for \(l\in \N_{>0}\),
\begin{equation}
 \Phi_l
 =
 \widehat{1_{[-3/4,3/4]}}
 \prod_{i\in [l)}
 2^{i+2}
 \widehat{1_{[-2^{-i-3},2^{-i-3}]}} .
\end{equation}
For $x\in\R$, define
\[\Phi(x)=\lim_{l\to\infty}\Phi_l(x).\]
(The limit exists for every $x\in\R$.)
\end{definition}


\begin{proposition}[standard bump]\label{standard bump properties}
The limit $\lim_{l\to \infty} \Phi_l$ exists in $L^\infty$  and $L^1$ sense and is a Schwartz function $\Phi$
whose Fourier transform takes values in $[0,1]$, is supported in $[-1,1]$, and constant $1$
on $[-1/2,1/2]$.

It satisfies for every $m,N\in\N$ and every $x\in\R$,
\begin{equation}
    \label{Phi_derbound}
    \|  (\widehat{\Phi^{(m)}})^{(N)}  \|_\infty \le \tilde{C}_{\ref{standard bump properties},m,N} 
\end{equation}
and 
\[ |\Phi^{(m)}(x)| \le C_{\ref{standard bump properties},m,N} \langle x\rangle^N, \]
where $\tilde{C}_{\ref{standard bump properties},m,N}=2^{4m + 2N^2 + 5N}$ and   $C_{\ref{standard bump properties},m,N}=2^{10m + 10N^2+10}.$
%\jr{$10$ is of course not optimal}
\end{proposition}

\begin{proof}

    The Fourier transform of the indicator function of an intervals is bounded by $C\langle x\rangle$.
    Hence each $\Phi_l$ is in $L^1$. Moreover, $|\Phi_l|$ is pointwise monotone decreasing in $l$,
    showing the existence of the limit in $L^\infty$ and $L^1$ sense. For any $m$ and sufficiently large
    $l$ one has that $\Phi_l$ is in $L^1$ with respect to the weight $|x|^m$ and the limit exists
    in that space a s well. This shows that $\widehat{\Phi}$ has any number of derivatives and thus is smooth.
    By induction, one proves
    that $\Phi_l$, a convolution product,  takes values in $[0,1]$, is supported in
    $[-2+2^{-l-1},2-2^{-l-1}]$ and is constant one on $[-1-2^{-l-1},1+2^{-l-1}]$. The desired properties
    of $\Phi$ follow.
Indeed, the Fourier support is in $[-1,1]$ because the Fourier transform of the product in the definition is supported on $[-a,a]$ with $a=\sum_{i=0}^\infty 2^{-i-3}=\frac14$.
Thus, $\widehat{\Psi}$ is supported on $[-b,b]$ with $b=\frac34+a=1$.
The Fourier support is one on the interval $[-c,c]$ with $c=\frac34-a=\frac12$.
For the decay estimates, note that
$m$ derivatives will give a factor like $(2\pi)^m$. To get $N$-decay for the $m$th derivative, we need to take $N$ derivatives on the Fourier side.
%which will make use of $N$ of the convolution factors each contributing a factor $2^{i+3}$, in total $\approx 2^{N(N+1)/2}$.
We compute
\[\frac{d}{d\xi} (f*2^{i+2}1_{[-2^{-i-3},2^{-i-3}]})(\xi) = 2^{i+2}(f(\xi+2^{-i-3})-f(\xi-2^{-i-3}))\]
with $f=1_{[-3/4,3/4]}$, so   the $L^\infty$ norm of this derivative  is bounded by $2^{i+3}$. Hence
\[\|(\widehat{\Phi})^{(N)}\|_\infty \le \prod_{i=0}^{N-1} 2^{i+3} = 2^{N(N-1)/2 +3N}\]
To compute the $N$-th derivative of $\widehat{\Phi^{(m)}}$ we use Leibniz rule: For $0\le k\le m$,
\[ \frac{d^N}{d\xi^N}\widehat{\Phi^{(m)}}(\xi)  = \frac{d^N}{d\xi^N}((2\pi i \xi)^m\widehat{\Phi}(\xi)) = \sum_{k=0}^N {N \choose k} \frac{m!}{(m-k)!}(2\pi i )^m \xi^{m-k} \widehat{\Phi}^{(N-k)}(\xi)\]
We obtain, using that $|\xi|\le 1$ in the support of $\widehat{\Phi}^{(N-k)}(\xi)$, and the bound  
\[\|  (\widehat{\Phi^{(m)}})^{(N)}  \|_\infty \le 2^{4m + 2N^2 + 5N}\]
Lemma \ref{lem:smoothdecay} then gives
\[ |\Phi^{(m)}(x)| \le C_{\ref{standard bump properties},m,N} \langle x\rangle^N. \]
\end{proof}
% \ct{we may only need the finite convolution product $\psi_l$. Length $l=4$ should be plenty.}\jr{it is cleaner to have it be a Schwartz function}

\begin{proposition}\label{compare brackets}
    Let $0<\lambda\le \mu$ with $1\le \mu$ and let $\lambda |y|\le \mu |x|$. Then for every $s,N>0$
\begin{equation}
    \mu^{-N}\langle x\rangle ^N_{(s)}\le \lambda^{-N}\langle y\rangle ^N_{(s)}
\end{equation}
\begin{proof}
We have, lowering both summands in the bracket on the left hand side,
\begin{equation*}
    s^{-1}(\mu+s^{-1}\mu |x|)^{-N}\le  s^{-1}(\lambda+s^{-1}|\lambda y|)^{-N} \, .
\end{equation*}
Pulling out the factors $\mu$ and  $\lambda$ on the two sides proves the proposition.
\end{proof}
\end{proposition}

\begin{proposition}[two bump estimate]\label{two bump estimate}
    Consider real numbers $x_i$ and  $s_i>0$ and $n_i>1$ for $i=0,1$.  Assume $s_0\ge s_1$ and define 
\begin{equation}\label{tb0}
     C_{\ref{two bump estimate},n_0,n_1}:=2^{1+\min(n_0,n_1)} [1+(\min(n_0,n_1)-1)^{-1}] 
\end{equation}    
  Then
\begin{equation}\label{two decays}
   | \int_{\R} (\prod_{i=0}^1 \left<x_i-p\right>^{n_i}_{(s_i)}) dp|\le C_{\ref{two bump estimate},n_0,n_1}
\left<x_0-x_1\right>^{\min(n_0,n_1)}_{(s_0)}\, .
\end{equation}
If $\min(n_0,n_1)\ge 2$, then 
\begin{equation}\label{constant two large decays}
C_{\ref{two bump estimate},n_0,n_1} \le 16.
\end{equation}
\end{proposition}
\begin{proof}
\eqref{constant two large decays} follows by arithmetic. Let us prove \eqref{two decays}.
By translating $p$ we may assume $x_0=-x_1$. 

If $|x_0-x_1|\le 2s_0$, we estimate the the first bump
by its $L^\infty$ norm, which is $s_i^{-1}$, and the second bump by its $L^1$ norm.
The latter is equal to the $L^1$ norm of $\langle x\rangle^{n_1}$, which in turn is estimated by splitting the integral
at $|x|<1$ and $|x|>1$ and thus by $2+2(n_1-1)^{-1}$. This gives the desired estimate.

If $|x_0-x_1|> 2s_0$, we split the integral into positive and negative axis
and do the favorable $L^1$ - $L^\infty$ estimate on either side. 
The $L^1$ norm is estimated as above, possibly with $n_0$ in place of $n_1$ if needed. The $L^\infty$ norm is
bounded by 
\begin{equation}
    s_i^{-1}(s_i^{-1}|x_0-x_1|/2)^{-\min(n_0,n_1)}
\end{equation}
for suitable $i$, which is however
maximized by $i=0$.
\end{proof}


\begin{proposition}\label{orthogonal domination}
Let     $\alpha_0 ,\alpha_1\in\mathbb{R}^2$ be linearly independent unit vectors. 
Then for every $x\in \R^2$ we have
 \begin{equation}
     |x\cdot \alpha_0^\perp|\le  2|\alpha_0\cdot \alpha_1^\perp|^{-1}| x\cdot \alpha_1|
 \end{equation}
or
 \begin{equation}
     | x\cdot \alpha_1^\perp|\le   2| \alpha_0\cdot \alpha_1^\perp|^{-1}|x\cdot \alpha_0|
 \end{equation}
\end{proposition}
\begin{proof}
As the desired inequalities do not change under replacing any of the orthogonal vectors by its negative,  
we may choose $\alpha_1^\perp$ so that  $\alpha_0\cdot \alpha_1^\perp$ is positive 
and $\alpha_0^\perp$ so that the second summands in \eqref{eq first alpha} and \eqref{eq second alpha} have the same sign.
We observe
\begin{equation}\label{eq first alpha}
     (x\cdot \alpha_1)( \alpha_1\cdot \alpha_0)+(x\cdot \alpha_1^\perp)( \alpha_1^\perp\cdot \alpha_0)=x\cdot  \alpha_0 \, ,
\end{equation}
 \begin{equation}\label{eq second alpha}
     (x\cdot \alpha_0)( \alpha_0\cdot \alpha_1)+(x\cdot \alpha_0^\perp)( \alpha_0^\perp\cdot \alpha_1)=x\cdot  \alpha_1\, .
\end{equation}
Adding and taking some terms to the other side and finally dividing by  $\alpha_0\cdot \alpha_1^\perp$ with the above sign conventions gives
\begin{equation*}
    |x\cdot \alpha_0^\perp|+|x\cdot \alpha_1^\perp|=\left|\frac{1-\alpha_0\cdot \alpha_1}{\alpha_0\cdot \alpha_1^\perp}x\cdot (\alpha_0+\alpha_1)
    \right|\le \frac{2}{|\alpha_0\cdot \alpha_1^\perp|}(|x\cdot \alpha_0|+|x\cdot \alpha_1|) \, .
\end{equation*}
The proposition follows.
\end{proof}
\begin{proposition}[orthogonal decay]\label{orthogonal decay}
Let $\alpha_0,\alpha_1\in\mathbb{R}^2$ be linearly independent unit vectors.
Let $\alpha_{m}^\perp$ be unit vectors orthogonal to $\alpha_m$ and note that $\alpha_0\cdot \alpha_1^\perp\neq 0$. Setting \begin{equation}
    C_{\ref{orthogonal decay}, \alpha_0,\alpha_1,n_0,n_1}:=\max( (2 |\alpha_0\cdot \alpha_1^\perp|^{-1})^{n_0},(2 |\alpha_0\cdot \alpha_1^\perp|^{-1})^{n_1}) \, ,
\end{equation}
then for all $n_0,n_1,s_0,s_1>0$ and $x\in\mathbb{R}^2$,
\begin{equation}\label{eq:two bump}
    \langle x\cdot \alpha_0\rangle_{(s_0)}^{n_0} \langle x\cdot \alpha_1\rangle_{(s_1)}^{n_1} 
\le C_{\ref{orthogonal decay},\alpha_0,\alpha_1,n_0,n_1} (\langle x\cdot \alpha_0\rangle_{(s_0)}^{n_0} \langle x\cdot \alpha_0^\perp\rangle_{(s_1)}^{n_1} +
\langle x\cdot \alpha_1^\perp\rangle_{(s_0)}^{n_0} \langle x\cdot \alpha_1\rangle_{(s_1)}^{n_1}).
\end{equation} 
\end{proposition}

\begin{proof}
This follows from Propositions \ref{compare brackets} and \ref{orthogonal domination}
\end{proof}

\begin{proposition}[bump triangle]\label{bump triangle}
Let $\alpha_0,\alpha_1\in\mathbb{R}^2$ and let $\alpha_2=c_0\alpha_0+c_1\alpha_1$ be a non-trivial linear combination of $\alpha_0,\alpha_1$, i.e. $c_0\not=0, c_1\not=0$. 
Let 
\[\tilde{C}_{\ref{bump triangle},c_0,c_1}=\max(2|c_0|, (2|c_0|)^{-1},2|c_1|, (2|c_1|)^{-1}),\]
\begin{equation}
    C_{\ref{bump triangle},c_0,c_1,n_0,n_1}=\max(\tilde{C}_{\ref{bump triangle},c_0,c_1}^{n_0}, \tilde{C}_{\ref{bump triangle},c_0,c_1}^{n_1}).
\end{equation}
Then
\[ \langle x\cdot \alpha_0\rangle_{(s_0)}^{n_0} \langle x\cdot \alpha_1\rangle_{(s_1)}^{n_1} 
\le C_{\ref{bump triangle},c_0,c_1,n_0,n_1}(\langle x\cdot \alpha_0\rangle_{(s_0)}^{n_0}\langle x\cdot \alpha_2\rangle_{(s_1)}^{n_1} + \langle x\cdot \alpha_2\rangle_{(s_0)}^{n_0}\langle x\cdot \alpha_1\rangle_{(s_1)}^{n_1}).
\]
\end{proposition}
\begin{rem}
It is not needed that $\alpha_0,\alpha_1$ are linearly independent, but in the linearly dependent case the proposition becomes trivial.
\end{rem}
\begin{proof}
We have by the triangle inequality and estimating the sum by twice the maximum
\begin{equation}
    |x\cdot \alpha_2|\le \max(2|c_0||x\cdot \alpha_0|,2|c_1||x\cdot \alpha_1|)\, .
\end{equation}
Assume the maximum is attained by one of the terms with index $i$.
If $2|c_i|\ge 1$, apply Proposition \ref{compare brackets} with $\lambda=1$ and $\mu=2|c_0|$.
If $2|c_i|< 1$, apply Proposition \ref{compare brackets} with $\lambda=2|c_0|$ and $\mu=1$.
\end{proof}


Recall the Gaussian $\g(x)=e^{-\pi x^2}$ for $x\in\R$.

\begin{proposition}[Gaussian domination]\label{Gaussian domination}
    Let 
    $N\in \R_{> 1}$ and $s\in \R_{>0}$. Then 
\begin{equation}
    \left< x\right>^{N}_{(s)}  \le C_{\ref{Gaussian domination}} 2^N   \sum_{m\in {\N}}  2^{(1-N)m} \g_{(2^ms)}(x),
\end{equation}
where $C_{\ref{Gaussian domination}}=e^\pi$.
\end{proposition}
 \begin{proof}
     The case for general $s$ follows from the case $s=1$ by scaling.  Consider $s=1$. For $|x|\le 1$, the summand $m=0$ on the right hand side suffices as it is at least $1$.
     
     For $|x|>1$ there is $m\in \N$ with $2^{m-1}<x\le 2^{m}$, and we estimate
     \begin{equation}
          \left< x\right>^{N}\le (2^{m-1})^{-N} \le (2^{m-1})^{-N} 2^me^\pi \g_{(2^m)}(x)\, . \,
     \end{equation}
     which is the $m$-th term on the right hand side.
 \end{proof}   


 
\begin{proposition}[diagonal square root]\label{diagonal square root}
For $t_0,t_1\in\R$ with $0<2t_0\le t_1$ define
\[ s = \mathcal{F}^{-1}(\xi\mapsto \sqrt{\g(t_0 \xi)- \g(t_1 \xi)}), \]
where the function under the square root is nonnegative. Then  
%s is positive and 
$s\in W_0(\R)$.
% there is a unique  function $s_{t_0,t_1}\in W_0(\R)$ such that
% \begin{equation}\label{E:square-root}
%     \widehat{s_{t_0,t_1}}(\xi)=\sqrt{\g(t_0 \xi)- \g(t_1 \xi)}\ge 0\, ,
% \end{equation}
% where function under the square root is nonnegative.
For all $N\in \mathbb N_{\geq 1}$ and $x\in \R$, the following estimate holds: %\jr{holds for every value of $8$ with constant depending on the value of $8$}
\begin{equation}\label{E:s-estimate}
    |s(x)|\le  C_{\ref{diagonal square root},N} ( \left<x\right>^N_{(t_0)}+ \left<x\right>^2_{(t_1)}),
\end{equation}
where %$C_{\ref{diagonal square root}}=2^{60}$
$C_{\ref{diagonal square root},N}=\sqrt{2}\max(C_{\ref{Gaussian bump decay},0,N}, C_{\ref{Gaussian bump decay},0,2}C_{\ref{square root of Gaussian decay}}C_{\ref{two bump estimate},2,2})$.
%$C_{\ref{diagonal square root}}=2^5C_{\ref{Gaussian bump decay},0,8}(1+C_{\ref{square root of Gaussian decay}}C_{\ref{two bump estimate},2,8})$.
\end{proposition}  

%\ls{Changed power $8$ to a general $N$.}

%\begin{rem}
%The conclusion also holds with $8$ replaced by any larger number $N$, with constant depending on $N$. 
%\end{rem}
% \ls{Should we also state this estimate for the derivative of $s$? This seems to be needed in the current (AI) proof of the Gaussian domination in the case $h<0$.}

\begin{proof}
The function $s$ is real-valued because it is the inverse Fourier transform of an even real-valued function.
By algebra, we obtain, 
\begin{equation}\label{dsr5}
    \widehat{s}(\xi) = \g(t_02^{-\frac12} \xi) \sqrt{1 - \g(r \xi)},
\end{equation}
where $r=(t_1^2-t_0^2)^{\frac 12}$.
Let $\rho$ be the function from Proposition \ref{square root of Gaussian decay} satisfying
\begin{equation}
    \widehat{\rho}(\xi)=1- \sqrt{1-\g(\xi)}.
\end{equation}
Then
\begin{equation}
\widehat{s}(\xi) = \g(\sigma\xi)(1-\widehat{\rho}(r\xi)),
\end{equation}
where $\sigma=t_0 2^{-\frac 12}$.

By Proposition \ref{square root of Gaussian decay},
\begin{equation}\label{dsr7}
    \rho(x)\le C_{\ref{square root of Gaussian decay}} \left<x\right>^2\, .
\end{equation}
It follows
\begin{equation}
    s=\g_{(\sigma)} - \g_{(\sigma)}*\rho_{(r)}.
\end{equation}
We use Proposition \ref{Gaussian bump decay} to obtain  
\[|\g_{(\sigma)}(x)|\le C_{\ref{Gaussian bump decay},0,N} \langle x \rangle_{(\sigma)}^N\]
and Proposition \ref{two bump estimate}, using $r>\sigma$ 
by assumptions on $t_i$, $i=0,1$, to obtain 
\[|(\g_{(\sigma)}*\rho_{(r)})(x)|\le C_{\ref{Gaussian bump decay},0,2}C_{\ref{square root of Gaussian decay}}C_{\ref{two bump estimate},2,2} \langle x \rangle_{(r)}^2.\]
 This gives 
\begin{equation}
    |s(x)|\le \max(C_{\ref{Gaussian bump decay},0,N}, C_{\ref{Gaussian bump decay},0,2}C_{\ref{square root of Gaussian decay}}C_{\ref{two bump estimate},2,2}) (\left<x\right>_{(\sigma)}^{N} +  \left<x\right>_{(r)}^2)\, .
\end{equation}
%with $C_{\ref{diagonal square root}} =C_{\ref{Gaussian bump decay},0,8}(1+C_{\ref{square root of Gaussian decay}}C_{\ref{two bump estimate},2,8})$. 
Since $\sigma=t_0 2^{-\frac12}$ and $\frac{\sqrt{3}}{2}t_1\le r\le t_1$, we obtain~\ref{E:s-estimate}. In particular, $s\in W_0(\R)$.
 \end{proof}


\begin{proposition}[derivative of diagonal square root]\label{derivative of diagonal square root}
Let $t_0,t_1,s$ be as in Proposition \ref{diagonal square root}.
Then for all $N\in\N_{\ge 1}$, $x\in\R$,
\begin{equation}\label{E:s-estimate}
    |s'(x)|\le  C_{\ref{derivative of diagonal square root},N} ( t_0^{-1} \left<x\right>^N_{(t_0)}+ t_1^{-1}\left<x\right>^2_{(t_1)}),
\end{equation}
where
$C_{\ref{derivative of diagonal square root},N}
=
2\max\bigl(
C_{\ref{Gaussian bump decay},1,N},
C_{\ref{Gaussian bump decay},0,2}
C_{\ref{Gaussian bump decay},1,2}
C_{\ref{two bump estimate},2,2}
\bigr).$
\end{proposition}

For the proof see \S \ref{sec:derivative of diagonal square root}.

\begin{lemma}\label{L:gaussian-estimate}
Let $N\in \N$. Then for $\xi \in \R$, $|\xi|\leq 1$, we have
\[
|(\g^{-1})^{(N)}(\xi)| \leq C_{\ref{L:gaussian-estimate},N},
\]
where $C_{\ref{L:gaussian-estimate},N}=e^{\pi} \sum_{l=0}^{\lfloor N/2 \rfloor} \frac{N!}{l!(N-2l)!} 2^{N-2l} \pi^{N-l}$.
\end{lemma}

\begin{proof}
We have $\g^{-1}(\xi)=e^{\pi\xi^2}$. The desired estimate then follows from the equality
\[
(\g^{-1})^{(N)}(\xi) 
=(\sqrt{\pi})^{N} e^{\pi  \xi^2} N!\sum_{l=0}^{\lfloor N/2\rfloor} \frac{2^{N-2l}}{l!(N -2l)!} (\sqrt{\pi} \xi)^{N-2l}.
\]
\end{proof}

\jr{We should normalize $c=1$ here and in the following lemmas.} \ls{The issue is that later on, $\widehat{\phi}$ will be equal to $1$ on $[-1/2,1/2]$, this may not be compatible with the normalization.}\jr{This cannot be an issue because the assumption $c=1$ is without loss of generality by homogeneity. As compromise I'm at least taking it out of the named constant}
\begin{lemma}\label{L:gaussian-bump-estimate}
Let $c>0$, $N \in \mathbb N$ and assume that $\phi$ is a $W_0(\R)$ function such that $\widehat{\phi}$ is supported in $[-1,1]$, $\widehat{\phi}$ is $N$ times continuously differentiable and $\max_{0\leq m \leq N} \|\widehat{\phi}^{(m)}\|_{\infty} \leq c$. 
For $\mu \in (0,1]$ and $\xi \in \R$, we define 
\[
Q(\xi)=(\g(\mu\xi))^{-1}\widehat{\phi}(\xi).
\]
Then $Q$ is $N$ times continously differentiable and 
\begin{equation}\label{E:estimate-Q}
\|(2\pi)^{-N}Q^{(N)}\|_{\infty} \leq c C_{\ref{L:gaussian-bump-estimate},N},
\end{equation}
where $C_{\ref{L:gaussian-bump-estimate},N}= (2\pi)^{-N}\sum_{l=0}^N {N \choose l} C_{\ref{L:gaussian-estimate},l}$.
\end{lemma}

\ls{In connection with windows, we are bounding $\max_{0\leq m\leq N}\|\phi^{(m)}\|_\infty$, while in connection with applying Proposition~\ref{mean value bump estimate 2}, we need to bound $\max_{0\leq m \leq N}(2\pi)^{-m}\|\phi^{(m)}\|_\infty$. Unify this?} \jr{We should avoid unnecessarily lengthy terms like $(2\pi)^{-N} Q^{(N)}$. It should just be $Q^{(N)}$}




\begin{proof}
%A combination of Lemma~\ref{L:gaussian-estimate} and of the fact that $\widehat{\Phi}$ is Schwartz and supported in $[-1,1]$ yields that $Q \in \mathcal{S}(\R)$. 
This follows from the Leibniz rule, Lemma~\ref{L:gaussian-estimate} and from the facts that $\mu \leq 1$ and $\widehat{\phi}$ is supported in the set where $|\xi|\leq 1$.
\end{proof}

\begin{lemma}\label{L:derivative-estimate-for-G}
Let $N\in \N$, $\nu\in [-1,0)$ and let $H(x)=(1-x)^{\nu}-1$ for $x\in [0,1)$. Then for $x\in [0,e^{-3\pi/16}]$, we have
\[
|H^{(N)}(x)| \leq C_{\ref{L:derivative-estimate-for-G},N},
\]
where $C_{\ref{L:derivative-estimate-for-G},N}= N! \left(1 - e^{-3\pi/16}\right)^{-(N+1)}$.
\end{lemma}

\begin{proof}
Assume that $N >0$. Then
\[
H^{(N)}(x)=\prod_{l\in [N)} (|\nu|+l)(1-x)^{\nu-N},
\]
and so
\[
|H^{(N)}(x)| \leq N! \left(1 - e^{-3\pi/16}\right)^{-(N+1)}.
\]
Also,
\[
|H(x)| \leq \frac{x}{1-x} \leq \left(1 - e^{-3\pi/16}\right)^{-1}.
\]
This yields the claim.
\end{proof}




\begin{lemma}\label{L:faa-di-bruno}
Let $N\in \N$, $\nu \in [-1,0)$, $t\geq \frac{\sqrt{3}}{2}$ and let $S(\xi)=(1-\g(t\xi))^{\nu}-1$ for $\xi \in \R$, $\xi \neq 0$. Then for $\xi \in \R$, $|\xi|\geq 1/2$, we have 
\begin{equation}\label{E:S}
|S^{(N)}(\xi)| \leq C_{\ref{L:faa-di-bruno},N} |\xi|^{-2},
\end{equation}
where
\[
C_{\ref{L:faa-di-bruno},N}=\tfrac{2^{N+1}}{\sqrt{3}} \sum_{p \in \Pi} C_{\ref{L:derivative-estimate-for-G},|p|} \prod_{B \in p} C_{\ref{Gaussian bump decay},|B|,N+2}.
\]
Here, $p$ runs through the set $\Pi$ of all partitions of the set $[N)$ and $B$ runs through the set of all blocks of the partition $p$. By $|p|$ we denote the number of blocks in the partition $p$ and by $|B|$ the size of the block $B$.
%\[
%B_N=N! \sum_{N} \prod_{m\in [N)} \frac{1}{k_m! ((m+1)!)^{k_m}}
%\]
%where $\sum_{N}$ stands for the sum over all $N$-tuples of nonnegative integers $(k_m)_{m\in [N)}$ such that $\sum_{m\in [N)} (m+1)k_m=N$.
\end{lemma}

\begin{proof}
Let us first assume that $t=1$. Let $H$ be the function from Lemma~\ref{L:derivative-estimate-for-G}. Then $S =H \circ \g$. By a combinatorial form of the Fa\`{a} di Bruno formula, we get
\[
S^{(N)}(\xi)=\sum_{p \in \Pi} H^{(|p|)}(\g(\xi)) \prod_{B \in p} \g^{(|B|)}(\xi),
\]
where $\Pi$, $p$ and $B$ are as in the statement of the lemma. Assume that $|\xi|\geq \sqrt{3}/4$. By Lemma~\ref{L:derivative-estimate-for-G} and Proposition~\ref{Gaussian bump decay}, we get
\begin{equation}\label{E:t=1}
|S^{(N)}(\xi)| \leq \sum_{p \in \Pi} C_{\ref{L:derivative-estimate-for-G},|p|} \prod_{B \in p} C_{\ref{Gaussian bump decay},|B|,N+2} \langle \xi \rangle^{N+2}
\leq \frac{\sqrt{3}}{2^{N+1}}C_{\ref{L:faa-di-bruno},N} \langle \xi \rangle^{N+2},
\end{equation}
where we estimated 
\[
\prod_{B \in p} \langle \xi \rangle^{N+2} \leq \langle \xi \rangle^{N+2}.
\]

Now, assume that $t\geq \sqrt{3}/2$ and $|\xi|\geq 1/2$. Applying the estimate~\eqref{E:t=1} with $\xi$ replaced by $t\xi$, we get
\[
|S^{(N)}(\xi)| \leq \frac{\sqrt{3}}{2^{N+1}} C_{\ref{L:faa-di-bruno},N} t^{N} \langle t\xi \rangle^{N+2}.
\]
Also,
\[
t^{N} \langle t\xi \rangle^{N+2} \leq \frac{t^{N}}{|t\xi|^{N}} \cdot \frac{1}{|t\xi|^2}
\leq \frac{2^{N+1}}{\sqrt{3}} |\xi|^{-2}.
\]
This yields~\eqref{E:S}.
\end{proof}

\begin{lemma}\label{L:second-gaussian-estimate}
Let $c>0$, $N\in \mathbb N$ and assume that $\phi$ is a $W_0(\R)$ function such that %$0\leq \widehat{\phi}(\xi) \leq 1$ for $\xi \in \R$, 
$\widehat{\phi}$ is supported in $[-1,1]$ and equal to $1$ on $[-1/2,1/2]$ and $\widehat{\phi}$ is $N$ times continuously differentiable and $\max_{0\leq m \leq N} \|\widehat{\phi}^{(m)}\|_{\infty} \leq c$. 
Assume that $0<\mu \leq \lambda \leq 1/2$, $t\geq \sqrt{3}/2$ and $\nu \in [-1,0)$. For $\xi \in \R$, we define 
\[
R(\xi)=(\g(\mu\xi))^{-1}\left(\widehat{\phi}(\lambda\xi)-\widehat{\phi}(\xi)\right)\left((1-\g(t\xi))^{\nu} -1\right).
\]
Then $R$ is $N$ times continuously differentiable and
\begin{equation}\label{E:estimate-R}
(2\pi)^{-N} \max(\|R^{(N)}\|_1,\|R^{(N)}\|_\infty) \leq cC_{\ref{L:second-gaussian-estimate},N},
\end{equation}
where $C_{\ref{L:second-gaussian-estimate},N}=8\sum_{l=0}^N {N \choose l} (2\pi)^{l-N} C_{\ref{L:gaussian-bump-estimate},l} C_{\ref{L:faa-di-bruno},N-l}$.

\end{lemma}


\begin{proof}
Let
\[
\widetilde{Q}(\xi):=(\g(\mu\xi))^{-1}\left(\widehat{\phi}(\lambda\xi)-\widehat{\phi}(\xi)\right).
\]
%A combination of Lemma~\ref{L:faa-di-bruno} and of the 
%fact that the function
%\[
%\widetilde{Q}(\xi):=(\g(\mu\xi))^{\nu}\left(\widehat{\Phi}(\lambda\xi)-\widehat{\Phi}(\xi)\right)
%\]
%is compactly supported and Schwartz (see Lemma~\ref{L:gaussian-bump-estimate}) yields that $R \in \mathcal{S}(\R)$.
%We next prove~\eqref{E:estimate-R}. 
By Lemma~\ref{L:gaussian-bump-estimate} and by the fact that $\mu \leq \lambda \leq 1$, we obtain for any $l\in \mathbb N$, $l\leq N$,
\begin{equation}\label{E:q-tilde-estimate}
\|(2\pi)^{-l}\widetilde{Q}^{(l)}\|_{\infty} \leq 2c C_{\ref{L:gaussian-bump-estimate},l}.
\end{equation}
We observe that $R$ is supported in $[1/2,1/\lambda]$. 
Using~\eqref{E:q-tilde-estimate}, Lemma~\ref{L:faa-di-bruno} and the Leibniz rule, we obtain for $|\xi|\geq 1/2$,
\[
|R^{(N)}(\xi)| \leq 2c\sum_{l=0}^N {N \choose l} (2\pi)^l C_{\ref{L:gaussian-bump-estimate}l} C_{\ref{L:faa-di-bruno},N-l} |\xi|^{-2}.
\]
This yields~\eqref{E:estimate-R}.
\end{proof}

\jr{Need this also for other than the standard bump function $\Phi$ (in reduction section)}
\ls{I changed the assumptions on $\Phi$ (now called $\phi$ since it is no longer the standard bump) to just reflect what is needed. It seems that in applications, $\phi$ is either the standard bump or a window.}
\ct{or maybe neither, as the stadnard bump definition may change
and likewise the window.}

\begin{lemma}\label{L:gaussian-bump-decomposition}
Let $\phi$ be a $W_0(\R)$ function such that %$0\leq \widehat{\phi}(\xi) \leq 1$ for $\xi \in \R$, 
$\widehat{\phi}$ is supported in $[-1,1]$ and equal to $1$ on $[-1/2,1/2]$.
Let $\mu_\pm, \lambda_\pm\in\mathbb{R}_{>0}$ satisfy
\begin{equation}\label{E:assumption-on-mu-lambda}
2\mu_-\le 2\lambda_-\le \lambda_+\le \mu_+,
\end{equation}
and let $\nu\in [-1,0)$.
%and assume that $\phi$ is a $W_0(\R)$ function such that $0\leq \widehat{\phi}(\xi) \leq 1$ for $\xi \in \R$, $\widehat{\phi}$ is supported in $[-1,1]$ and equal to $1$ on $[-1/2,1/2]$ and $\widehat{\phi}$ is $N$ times continuously differentiable. 
Define a function $\rho$ by 
\begin{equation}\label{E:definition-rho}
\rho = \mathcal F^{-1}\left(\xi \mapsto (\widehat{\phi}(\lambda_-\xi) - \widehat{\phi}(\lambda_+\xi)) (\g(\mu_- \xi)-\g(\mu_+\xi))^{\nu}\right). 
\end{equation}
Also, for $\xi \in \R$ we define
\begin{equation}
    \varrho_0(\xi)=(\g(\mu_{-}|\nu|^{1/2}\xi))^{-1}\widehat{\phi}(\lambda_{-}\xi)
\end{equation}
\begin{equation}
    \varrho_1(\xi)=-(\g(\mu_{-}|\nu|^{1/2}\xi))^{-1}\widehat{\phi}(\lambda_{+}\xi)
\end{equation}
\begin{equation}
    \varrho_2(\xi)=(\g(\mu_{-}|\nu|^{1/2}\xi))^{-1}\left(\widehat{\phi}(\lambda_{-}\xi)-\widehat{\phi}(\lambda_{+}\xi)\right)\left((1-\g(t\xi))^{\nu} -1\right)
\end{equation}
with $t=\sqrt{\mu_{+}^2-\mu_{-}^2}$.
Then $\rho$ is a well-defined continuous function and
\begin{equation}\label{E:decomposition-rho}
\rho=\sum_{l\in [3)} \mathcal F^{-1} (\varrho_l).
\end{equation}
\end{lemma}

\begin{proof}
Applying Lemmas~\ref{L:gaussian-estimate} and~\ref{L:second-gaussian-estimate} with $N=0$, we obtain that $\varrho_l$, $l\in [3)$, are continuous compactly supported functions. Therefore, their inverse Fourier transform is well-defined and continuous. The decomposition~\eqref{E:decomposition-rho} follows by algebra and elementary properties of $\g$. Thus, $\rho$ is a well defined continuous function.
Additionally, $\rho$ is real-valued because it is the inverse Fourier transform of an even real-valued function.
\end{proof}

\begin{proposition}[four scale Gaussian kernel estimate]\label{four scale Gaussian kernel}
Let $c>0$, $N\in \mathbb N_{\geq 2}$ and assume that $\phi$ is a $W_0(\R)$ function such that %$0\leq \widehat{\phi}(\xi) \leq 1$ for $\xi \in \R$, 
$\widehat{\phi}$ is supported in $[-1,1]$ and equal to $1$ on $[-1/2,1/2]$ and $\widehat{\phi}$ is $N$ times continuously differentiable and $\max_{0\leq m \leq N} \|\widehat{\phi}^{(m)}\|_{\infty} \leq c$. 
Let $\mu_\pm, \lambda_\pm\in\mathbb{R}_{>0}$ satisfy~\eqref{E:assumption-on-mu-lambda}, let $\nu\in [-1,0)$ and let $\rho$ be defined by~\eqref{E:definition-rho}.
Then $\rho \in W_0(\R)$ and for every $x\in\mathbb{R}$, 
\[ |\rho(x)| \le cC_{\ref{four scale Gaussian kernel},N} (\langle x\rangle^N_{(\lambda_-)} + \langle x\rangle^N_{(\lambda_+)}), \]
where
\[ C_{\ref{four scale Gaussian kernel},N} = 2C_{\ref{lem:smoothdecay2},N}\max(C_{\ref{L:gaussian-bump-estimate},0},C_{\ref{L:gaussian-bump-estimate},N}) +2^N\max(C_{\ref{L:second-gaussian-estimate},0},C_{\ref{L:second-gaussian-estimate},N}). \]
%= 4^{4^{N+m+2}}\]
%100^{100^{100+m+n+|\nu|}} 
%\] 
%\pd{$2^{2^{N+3}+30(m+N)^4}$ or $2^{2^{2^{N+m+2}}}$ or $4^{4^{N+m+2}}$ }
\end{proposition}

\begin{proof}
Let $\varrho_l$, $l\in [3)$, be as in Lemma~\ref{L:gaussian-bump-decomposition}. 
Rescaling the functions $\varrho_0$ and $\varrho_1$ so that $\lambda_{-}=1$ or $\lambda_{+}=1$, respectively, applying Lemmas~\ref{lem:smoothdecay2} and~\ref{L:gaussian-bump-estimate} and then rescaling back, we obtain
\begin{equation*}
     |\mathcal{F}^{-1}(\varrho_0)(x)|\le 2cC_{\ref{lem:smoothdecay2},N}\max(C_{\ref{L:gaussian-bump-estimate},0},C_{\ref{L:gaussian-bump-estimate},N}) \left<x\right>_{(\lambda_{-})}^N\,
\end{equation*}
and
\begin{equation*}
     |\mathcal{F}^{-1}(\varrho_1)(x)|\le 2cC_{\ref{lem:smoothdecay2},N}\max(C_{\ref{L:gaussian-bump-estimate},0},C_{\ref{L:gaussian-bump-estimate},N}) \left<x\right>_{(\lambda_{+})}^N\,.
\end{equation*}
Also, rescaling the function $\varrho_2$ so that $\lambda_{+}=1$, applying Lemmas~\ref{lem:smoothdecay}, \ref{lem: min and bracket} and~\ref{L:second-gaussian-estimate} and then rescaling back, we obtain
\begin{equation*}
     |\mathcal{F}^{-1}(\varrho_2)(x)|\le 2^Nc\max(C_{\ref{L:second-gaussian-estimate},0},C_{\ref{L:second-gaussian-estimate},N}) \left<x\right>_{(\lambda_{+})}^N.
\end{equation*}
Applying the decomposition~\eqref{E:decomposition-rho}, we get the desired estimate. In particular, since $N\geq 2$, we get $\rho \in W_0(\R)$.
\end{proof}

\begin{proposition}[mean value four scale Gaussian kernel estimate]\label{mean four scale Gaussian kernel}
Let $c>0$, $N\in \mathbb N_{\geq 2}$ and assume that $\phi$ is a $W_0(\R)$ function such that %$0\leq \widehat{\phi}(\xi) \leq 1$ for $\xi \in \R$, 
$\widehat{\phi}$ is supported in $[-1,1]$ and equal to $1$ on $[-1/2,1/2]$ and $\widehat{\phi}$ is $N$ times continuously differentiable and $\max_{0\leq m \leq N} \|\widehat{\phi}^{(m)}\|_{\infty} \leq c$. 
Let $\mu_\pm, \lambda_\pm\in\mathbb{R}_{>0}$ satisfy~\eqref{E:assumption-on-mu-lambda}. Assume, in addition, that $\lambda_{+}=2\lambda_{-}$. 
Let $\nu\in [-1,0)$ and let $\rho$ be defined by~\eqref{E:definition-rho}.
Then for all $x,y \in\mathbb{R}$, 
\begin{equation}\label{E:mean-value}
|\rho(x+y)-\rho(x)| \le c C_{\ref{mean four scale Gaussian kernel},N} \min(1,2\pi \lambda_{+}^{-1}|y|) (\langle x+y\rangle^N_{(\lambda_+)}+
\langle x\rangle^N_{(\lambda_+)}), 
\end{equation}
where
\[ C_{\ref{mean four scale Gaussian kernel},N} =  5\cdot 2^{N+1}\max_{0\leq l \leq N} C_{\ref{L:gaussian-bump-estimate},l} +2^{N+1}\max_{0\leq l \leq N} 2^lC_{\ref{L:second-gaussian-estimate},l}.\]
%= 4^{4^{N+m+2}}\]
%100^{100^{100+m+n+|\nu|}} 
%\] 
%\pd{$2^{2^{N+3}+30(m+N)^4}$ or $2^{2^{2^{N+m+2}}}$ or $4^{4^{N+m+2}}$ }
\end{proposition}

\begin{proof}
By homogeneity, we may assume $\lambda_{+}=1$. Then $\lambda_{-}=1/2$.
Let $\varrho_l$, $l\in [3)$, be as in Lemma~\ref{L:gaussian-bump-decomposition}. Rescaling the function $\varrho_0$ so that $\lambda_{-}=1$, applying Lemmas~\ref{mean value bump estimate 2} and~\ref{L:gaussian-bump-estimate} and then rescaling back, we obtain
\begin{equation*}
     |\mathcal{F}^{-1}(\varrho_0)(x+y)-\mathcal{F}^{-1}(\varrho_0)(x)|\le c2^{N+1}\max_{0\leq l \leq N} C_{\ref{L:gaussian-bump-estimate},l} 
     \min(1,4\pi|y|) (\langle x+y\rangle^N_{(1/2)}+
\langle x\rangle^N_{(1/2)})
\end{equation*}
\[
\leq c 2^{N+3} \max_{0\leq l \leq N} C_{\ref{L:gaussian-bump-estimate},l} \min(1,2\pi|y|) (\langle x+y\rangle^N+
\langle x\rangle^N).
\]
Similarly, applying Proposition~\ref{mean value bump estimate 2} and Lemma~\ref{L:gaussian-bump-estimate} to the function $\varrho_1$, we obtain
\begin{equation*}
     |\mathcal{F}^{-1}(\varrho_1)(x+y)-\mathcal{F}^{-1}(\varrho_1)(x)|\le c 2^{N+1}\max_{0\leq l \leq N} C_{\ref{L:gaussian-bump-estimate},l} 
     \min(1,2\pi|y|) (\langle x+y\rangle^N+
\langle x\rangle^N).
\end{equation*}

To estimate $\varrho_2$, we first consider its rescaled version $\widetilde{\varrho_2}:=\varrho_2(2\cdot)$, which is supported in $[-1,1]$. By Lemma~\ref{L:second-gaussian-estimate}, we have for $l\in \mathbb N$, $l\leq N$,
\[
\|(2\pi)^{-l}\widetilde{\varrho_2}^{(l)}\|_{\infty}
=2^l\|(2\pi)^{-l}\varrho_2^{(l)}\|_{\infty}
\leq c 2^lC_{\ref{L:second-gaussian-estimate},l}.
\]
By Proposition~\ref{mean value bump estimate 2},
\begin{equation*}
     |\mathcal{F}^{-1}(\varrho_2)(x+y)-\mathcal{F}^{-1}(\varrho_2)(x)|\le c 2^{N+1}\max_{0\leq l \leq N} 2^lC_{\ref{L:second-gaussian-estimate},l} 
     \min(1,2\pi|y|) (\langle x+y\rangle^N+
\langle x\rangle^N).
\end{equation*}
Using the decomposition~\eqref{E:decomposition-rho}, we get~\eqref{E:mean-value}.
\end{proof}









\section{The main argument}\label{sec:main argument}

The main theorem of this section is Theorem \ref{induct positive terms theorem} below.

\subsection{The sandwich kernel}

We first introduce  and discuss some geometric parameters useful to describe a family of kernels in Definition \ref{sandwich kernel}.

\begin{definition}[geometric parameters]\label{geometric parameters}
  Let $\Gamma$ be the set of triples $(k,u,a)$ with $k\in\N$, $1\le k\le n$, $u\in[2)^k$, $a\in ({\rm A}^{2})^k$  so that $\dist(a_{i}^0,a_{i}^1)<\infty$ for all $i\in [k)$.
  For a sequence pair $\alpha=(\alpha_0,\alpha_1)\in {\rm A}^2$ we write $\Delta(\alpha) = \mathrm{dist}(\alpha_{0}, \alpha_1)$
  and for $\gamma\in\Gamma$ let
  \[ \Delta_\gamma= 1+\sum_{i\in [k)} \Delta(a_i). \]
%\ct{maybe define +1, to avoid that this quantity gets zero)}\cjr
\end{definition}
% \pd{$i=1,\ldots k$ or $i\in [k)$?}

\begin{definition}[two unitary matrices]
For $u\in [2)$ define unitary matrices $W_u\in\R^{2\times 2}$ as follows:
Let $W_0$ be the $2\times 2$ identity matrix and $W_1=\frac{1}{\sqrt{2}}\begin{pmatrix}1 & 1\\-1 & 1\end{pmatrix}$.
\end{definition}



Next, we introduce some functions and functional concepts involved in the family of kernels of Definition \ref{sandwich kernel}.






\begin{definition}[double sequence of 2D functions]\label{double sequence of 2D functions}
  For $1\le k\le n$, let $\mathcal{X}_k$ be the set of
  double sequences
  \begin{equation}
      X=(X_{i,j})_{i\in [k),j\in \Z}
  \end{equation}
  with $X_{i,j}\in W_0(\R^2)$.
\end{definition}


The kernels of interest in Definition \ref{sandwich kernel} will be partial sums over a sequence of kernels.
We introduce this concept of kernel sequences.

\begin{definition}[kernel sequences]\label{kernel sequences}
Let $k\in\N$ with $1\le k\le n$. Let ${\rm M}(k)$ be the space of sequences $\M=(M_j)_{j\in \Z}$ such that for $j\in \Z$ we have
$M_j\in W_0((\R^2)^k)$. Define %\jr{$\rm{mBL}$ vs $\rm{pBL}$}
\begin{equation}
\|\M\|_{{\rm M}(k)}:=\sup_{J\in \N_{>0}, \F \in\mathfrak{F}}\min(1,J^{-1+2^{k-n+1}}) |\Lambda_k(\sum_{j\in [J)} M_j)(\F)| \, .
\end{equation}
\end{definition}

\begin{rem}
Note $\min(1,J^{-1+2^{k-n+1}})$ equals $J^{-1+2^{k-n+1}}$ if $k\le n-1$ and equals $1$ if $k=n$.
\end{rem}


\begin{definition}[2D Gaussians]\label{2D Gaussians}
For $q\in (\R_{>0})^2$ and $u\in[2)$, $v\in\mathbb{R}^2$ define
\[ G_{q,u}(v) = \g_{(q_0)}((W_u v)_0) \g_{(q_1)}((W_u v)_1). \]
For $\gamma=(k,u,a)\in \Gamma$, define
\[ (G_{\gamma})_{i,j} = G_{a_i(j),u(i)}. \]
\end{definition}




A sandwich kernel is a product, most factors of which are two dimensional Gaussians differently scaled and oriented 
using the data $\gamma$, while one
factor, which is  sandwiched between these Gaussians, is different and determined by an element of $\mathcal{X}_k$.
\begin{definition}[sandwich kernel]\label{sandwich kernel}
Let $\gamma=(k,u,a)\in \Gamma$, $X\in \mathcal{X}_k$, and $i\in [k)$. Define for  $j\in \Z$ and $y\in (\R^{2})^k$,
\[ (\M(\gamma,X,i))_j(y) =  \Big(\prod_{m\in [i)} (G_{\gamma})_{m,j}(y_m) \Big) X_{i,j}(y_i) \Big(\prod_{m=i+1}^{k-1} (G_{\gamma})_{m,j-1}(y_m) \Big).\] 
\end{definition}

\begin{rem}
When $k=1$, $i=0$, $\M(\gamma,X,1)$ doesn't depend on $u,a$.
Generally, $\M(\gamma,X,i)$ does not depend on $u_i, a_i$.    
\end{rem}

\begin{definition}[difference of 2D Gaussians]
Let $\gamma=(k,u,a)\in \Gamma$.  For $i\in [k)$, $j\in \Z$, define
\[ (Y_\gamma)_{i,j} = (G_\gamma)_{i,j-1} - (G_\gamma)_{i,j}\ .\]
\end{definition}

\begin{proposition}[telescoping terms]\label{telescoping terms}
Let $\gamma=(k,u,a)\in \Gamma$.  
Then $Y_\gamma\in \mathcal{X}_k$ and 
\begin{equation}
    \|\sum_{i\in [k)} \M(\gamma,Y_\gamma,i)\|_{{\rm M}(k)}\le 2.
\end{equation}
\end{proposition}

\begin{proof}
The fact that $Y_\gamma\in \mathcal{X}_k$ follows from Proposition~\ref{P:schwartz-into-wiener} since $(Y_\gamma)_{i,j}\in \mathcal S(\R^n)$ for $i\in [k), j\in \Z$.
For a given $J\in \N_{>0}$, 
\[
\sum_{j\in [J)} \sum_{i\in [k)} (\M(\gamma,Y_\gamma,i))_j
\]
telescopes into a difference of two boundary terms
\[
\prod_{m\in [k)} (G_{\gamma})_{m,-1}(y_m) -\prod_{m\in [k)} (G_{\gamma})_{m,J-1}(y_m).
\]
Estimate each with Propositions \ref{prism BL inequality} and \ref{M to K} and then make use of the fact that the $L^1$ norm of both these terms equals $1$.
\end{proof}



\begin{proposition}[positive terms]\label{positive terms}
Let $\gamma=(k,u,a)\in \Gamma$. Let $X\in \mathcal{X}_k$ and assume for each $i\in [k)$ and  $j\in \Z$ we have
 $X_{i,j}(u,v)=f_{i,j}(u)f_{i,j}(v)$ for some real valued $f_{i,j}\in W_0(\R)$. 
Then for each $i\in [k)$, $j\in\Z$, and for all $\F\in {\mathfrak{F}}$,
\begin{equation}\label{E:form-nonnegative}
    \Lambda_k(\M(\gamma,X,i)_j)(\mathbf{F})\ge 0.
\end{equation}
\end{proposition}

\begin{proof}
This follows from Proposition \ref{Positivity M}.
\end{proof}


%This follows from Proposition \ref{Positivity M} after rearranging variables.
%\ls{I don't think we need to rearrange variables.}

\subsection{\texorpdfstring{Multipliers $H$, $L$, $N$}{Multipliers H, L, N}}

\begin{definition}[square root Gaussian difference]\label{square root Gaussian difference}
Let $a\in A$ and $j\in\Z$.
Define the function $s(a,j):\R \to \R$ by 
\[s(a,j) = \mathcal{F}^{-1}(\xi \mapsto \sqrt{\g(a(j-1)\xi)-\g(a(j) \xi)}).\] 
\end{definition}

\begin{definition}[$s$ multiplier]\label{s multiplier}
    Let $\gamma=(k,u,a)\in \Gamma$. Let $i\in [k)$ and
    $j\in \Z$. Define  $(s_{\gamma})_{i,j}: \R \to \R$ by
\begin{equation}\label{E:b-definition}
   (s_{\gamma})_{i,j}=s(b,j)
\end{equation}
with 
\begin{equation}
   b(j)=\sqrt{a_{i}^0(j)^2 + a_{i}^1(j)^2}
\end{equation}
if $u_i=0$, and by
\begin{equation}
   (s_{\gamma})_{i,j}=s(\sqrt{2}a_{i}^1,j)
\end{equation}
if $u_i=1$.
% We also define\jr{maybe not needed}
% \[ S_{\gamma} = s_\gamma \otimes s_\gamma \]
\end{definition}

\begin{proposition}\label{square root Gaussian difference W0}
  For every $a\in\mathrm{A}$ and $j\in\mathbb{Z}$ 
  we have that $s(a,j)$ is a well-defined function in $W_0(\R)$. Consequently, if $\gamma \in \Gamma$ and $i\in [k)$, $j\in \Z$, then $(s_{\gamma})_{i,j}$ is a well-defined function in $W_0(\R)$.
\end{proposition}

\begin{proof}
The first part follows from Proposition~\ref{diagonal square root}. The second part follows from the first part and Proposition~\ref{Operations on spaced sequences}.
\end{proof}

Note that $\widehat{s(a,j)}$ is not smooth at the origin, due to the square root. The function $s(a,j)$ only decays with second power.
 







\begin{definition}[H multiplier]\label{H multiplier}
Let $\gamma=(k,u,a)\in \Gamma$. Define
    $H_{\gamma}=(H_\gamma)_{i\in [k),j\in \Z}$ by
\begin{equation}
(H_{\gamma})_{i,j}:=(s_{\gamma})_{i,j}\otimes (s_{\gamma})_{i,j}
  -(Y_{\gamma})_{i,j}
\end{equation}
for every $i\in [k), j\in\mathbb{Z}$.
\end{definition}

\begin{proposition}\label{H-in-X}
Let $\gamma=(k,u,a)\in \Gamma$. Then $H_{\gamma}\in \mathcal{X}_k$.
\end{proposition}
\begin{proof}
This follows from Propositions~\ref{tensor Wiener}, \ref{telescoping terms} and \ref{square root Gaussian difference W0}.
\end{proof}



\begin{proposition}[H vanishing]\label{H vanishing}
    For $\gamma=(k,u,a)\in \Gamma$, $i\in [k)$, $j\in \Z$, and $\xi\in \R$,
\begin{equation}
    \widehat{(H_{\gamma})_{i,j}}(\xi,-\xi)=0\, .
\end{equation}
\end{proposition}

\begin{proof}
Recall
\[ (Y_\gamma)_{i,j} = (G_\gamma)_{i,j-1} - (G_\gamma)_{i,j}\ \]
and with \eqref{eq:gaussconv} and \eqref{eq:scaleFT}, if $u_i=0$, then
\[ (\widehat{G_\gamma})_{i,j}(\xi, -\xi) = \g(\sqrt{(a_{i}^0(j))^2 + (a_{i}^1(j))^2}\cdot \xi) \]
and if $u_i=1$, then
\[ (\widehat{G_\gamma})_{i,j}(\xi, -\xi) = \g(\sqrt{2} a_{i}^1(j) \xi).\]
The proposition now follows with Definition \ref{square root Gaussian difference} of $s_\gamma$ and Proposition~\ref{square root Gaussian difference W0}.
\end{proof}


\begin{proposition}[H vanishing integral]\label{H vanishing integral}
    For $\gamma=(k,u,a)\in \Gamma$, $i\in [k)$, $j\in \Z$, and $z\in \R$,
\begin{equation}
   \int_{\R} (H_{\gamma})_{i,j}(z+p,p)\,dp=0.
\end{equation}
\end{proposition}

\begin{proof}
This follows from Proposition~\ref{H vanishing}.
\end{proof}







\iffalse
\begin{proposition}[Properties of L multipliers]\label{sum L multiplier convergence}
Let $\gamma=(k,u,a)\in \Gamma$, $i\in [k)$ and $\iota \in\mathcal{I}_{\gamma,i}$. Then
 $L_{\gamma,\iota}$ is a well-defined element of $\mathcal{X}_k$. 
 In addition, if $j\in \Z$ then
\begin{equation}\label{E:L2-convergence}
(H_{\gamma})_{i,j} = \sum_{\iota \in\mathcal{I}_{\gamma,i}} (L_{\gamma,\iota})_{i,j}, 
\end{equation}
where the sum on the right converges unconditionally in $L^2$. 
\end{proposition}

\begin{proof}
Each $(L_{\gamma,\iota})_{i,j}$ is real-valued because it is the inverse Fourier transform of a even real-valued function.
Additionally, for $v\in \R^2$, $(L_{\gamma,\iota})_{i,j}(v)$ can be written in the form 
\[
(L_{\gamma,\iota})_{i,j}(v)=\int_{\R} \rho(p) (H_\gamma)_{i,j}(v_1-p,v_2-p)\,dp,
\]
where $\rho \in \mathcal S(\R)$, and thus $\rho \in W_0(\R)$ by Proposition~\ref{P:schwartz-into-wiener}.  Proposition~\ref{H-in-X} yields that $(H_{\gamma})_{i,j} \in W_0(\R^2)$. Thus, $(L_{\gamma,\iota})_{i,j} \in W_0(\R^2)$ by Proposition~\ref{W_0 Brascamp Lieb}, i.e.\ $L_{\gamma,\iota} \in \mathcal{X}_k$.


Let $i\in [k), j\in\Z$. By Proposition~\ref{P:lp-embedding} and Plancherel's identity, $\widehat{(H_\gamma)_{i,j}}$ belongs to $L^2$. Let $\sigma: \mathbb N \rightarrow \mathcal{I}_{\gamma,i}$ be a bijection. Then  $\widehat{(H_\gamma)_{i,j}}-\sum_{l=0}^N \widehat{(L_{\gamma,\sigma(l)})_{i,j}}$ converges to $0$ a.e.\ as $N\to \infty$ and 
\[
|\widehat{(H_\gamma)_{i,j}}-\sum_{l=0}^N \widehat{(L_{\gamma,\sigma(l)})_{i,j}}|
\leq |\widehat{(H_\gamma)_{i,j}}|.
\]  
By the dominated convergence theorem, $\widehat{(H_\gamma)_{i,j}}-\sum_{l=0}^N \widehat{(L_{\gamma,\sigma(l)})_{i,j}}$ converges to $0$ in $L^2$. Using Plancherel's identity, we deduce~\eqref{E:L2-convergence}.
\end{proof}


\begin{lemma}\label{sandwich sums L2}
Let $\gamma=(k,u,a)\in \Gamma$.
Let $I$ be a countable index set.
Assume $X, X_\iota\in \mathcal{X}_k$ for $\iota\in I$ are such that for every $i\in [k)$, $j\in\Z$, $X_{i,j}=\sum_{\iota\in I} (X_\iota)_{i,j}$ holds unconditionally in $L^2$.
Then for every $i\in [k), j\in\Z$,
\[ \mathbf{M}(\gamma,X,i)_j = \sum_{\iota\in I} \mathbf{M}(\gamma,X_\iota,i)_j\quad\text{unconditionally in}\;L^2.\]
\end{lemma}

\begin{proof}
If the set $I$ is finite then the conclusion is a consequence of Definition~\ref{sandwich kernel}.
Assume that $I$ is infinite.
Let $i\in [k), j\in\Z$ and  let $\sigma: \mathbb N \rightarrow I$ be a bijection. Then
\begin{align}
&\mathbf{M}(\gamma,X,i)_j(y)-\sum_{l=0}^N \mathbf{M}(\gamma,X_{\sigma(l)},i)_j(y)\\
&=\Big(\prod_{m\in [i)} (G_{\gamma})_{m,j}(y_m) \Big) \Big(X_{i,j}(y_i)- \sum_{l=0}^N (X_{\sigma(l)})_{i,j}(y_i) \Big)\Big(\prod_{m=i+1}^{k-1} (G_{\gamma})_{m,j-1}(y_m) \Big).
\end{align}
Since $X_{i,j}-\sum_{l=0}^N (X_{\sigma(l)})_{i,j}$ converges to $0$ in $L^2$ as $N \rightarrow \infty$, we deduce that $\mathbf{M}(\gamma,X,i)_j-\sum_{l=0}^N \mathbf{M}(\gamma,X_{\sigma(l)},i)_j$ converges to $0$ in $L^2$ as $N \rightarrow \infty$ as well. 
\end{proof}

\begin{lemma}\label{prism sum le sum prism}
Let $I$ be a countable set.
If $M=\sum_{\iota\in I} M_\iota$ holds unconditionally in $L^2$ for $M,M_\iota\in W_0((\R^2)^k)$, then for every $\mathbf{F}\in\mathfrak{F}$, $1\le k\le n$,
\[ |\Lambda_k(M)(\mathbf{F})| \le \sum_{\iota\in I} |\Lambda_k(M_\iota)(\mathbf{F})|. \]
\end{lemma}

\begin{proof}
If $I$ is finite then the conclusion follows by the triangle inequality. Assume that $I$ is infinite. Let $\sigma: \mathbb N \rightarrow I$ be a bijection.
Let $E_N=M - \sum_{l=0}^N M_{\sigma(l)}$. By assumption, $\|E_N\|_2 \to 0$ as $N\to\infty$. Let $\mathbf{F}\in\mathfrak{F}$. By the triangle inequality, 
\[
|\Lambda_k(M)(\mathbf{F})| \le \sum_{l=0}^N |\Lambda_k(M_{\sigma(l)})(\mathbf{F})|
+|\Lambda_k(E_N)(\mathbf{F})|.
\]
We recall
\[
\Lambda_k(E_N)(\mathbf{F})=\Theta_k(K_k(E_N))(\F)
\]
\[= \int_{(\R^{2})^k} \int_{\R^{n-k+1}} \int_{\R^{k-1}} E_N((y_l^1-y_l^0+p_l,p_l)_{l\in [k-1)}, y_{k-1}^1-y_{k-1}^0+\Sigma(y^0)+\Sigma(x_{[k-1,n)})-\Sigma(p),
\]
\[
\Sigma(y^0)-\Sigma(x_{[k-1,n)})-\Sigma(p)) 
\prod_{\substack{h\in [2)^k,\\i\in [k-1,n)}} F_{i}( y^h, x_{[k-1,n)\setminus i})\; dp dx_{[k-1,n)}\;dy. 
\]
By the change of variables $p_l \mapsto p_l+y_l^0$, $l\in [k-1)$, this becomes
\begin{equation}\label{E:form-change-of-variables}
\int_{(\R^{2})^k} \int_{\R^{n-k+1}} \int_{\R^{k-1}} E_N((y_l^1+p_l,y_l^0+p_l)_{l\in [k-1)}, y_{k-1}^1+\Sigma(x_{[k-1,n)})-\Sigma(p),
\end{equation}
\[
y^0_{k-1}-\Sigma(x_{[k-1,n)})-\Sigma(p)) 
\prod_{\substack{h\in [2)^k,\\i\in [k-1,n)}} F_{i}( y^h, x_{[k-1,n)\setminus i})\; dp dx_{[k-1,n)}\;dy. 
\]
A further change of variables $y_l^i \mapsto y_l^i-p_l$, $l\in [k-1)$, $y_{k-1}^i \mapsto y_{k-1}^i +\Sigma(p)-\Sigma(x_{[k-1,n)})$, $i\in [2)$, translates the previous expression into
\[
\int_{(\R^{2})^k} E_N(y) \int_{\R^{n-k+1}} \int_{\R^{k-1}}
\]
\[ 
\prod_{\substack{h\in [2)^k,\\i\in [k-1,n)}} F_{i}( (y^h_l-p_l)_{l\in [k-1)},y^h_{k-1}+\Sigma(p)-\Sigma(x_{[k-1,n)}), x_{[k-1,n)\setminus i})\; dp dx_{[k-1,n)}\;dy. 
\]
We would like to bound the absolute value of this by
\[\le c_1\|E_N\|_{L^2},\] where $c_1$ is a constant that depends on $\mathbf{F}$. The best such constant satisfies
\[
c_1^2=\int_{(\R^{2})^k} \Big|\int_{\R^{n-k+1}} \int_{\R^{k-1}}
\]
\[
\prod_{\substack{h\in [2)^k,\\i\in [k-1,n)}} F_{i}( (y^h_l-p_l)_{l\in [k-1)},y^h_{k-1}+\Sigma(p)-\Sigma(x_{[k-1,n)}), x_{[k-1,n)\setminus i})\;
dp dx_{[k-1,n)}\Big|^2 \,dy.
\]
\ls{It seems that $c_1=\infty$ for $k>1$.}
\end{proof}

\fi

\begin{definition}
Let $\gamma=(k,u,a)\in \Gamma$, $i\in [k)$, $j\in \Z$. For $t>0$, we set
\[
(L_{\gamma,t})_{i,j} =(H_{\gamma})_{i,j} \ast_{(1,1)} \Phi_{(t)}.
\]  
\end{definition}

\begin{lemma}\label{L:F_t}
Let $\gamma=(k,u,a)\in \Gamma$ and $t>0$. Then $L_{\gamma,t} \in \mathcal{X}_k$ and for $i\in [k)$, $j\in \Z$, we have
\[
\lim_{t\to 0_+} (L_{\gamma,t})_{i,j} = (H_{\gamma})_{i,j} \quad\text{in}\;L^1.
\]
\end{lemma}

\begin{proof}
$L_{\gamma,t} \in \mathcal{X}_k$ by Propositions \ref{P:schwartz-into-wiener}, \ref{convolution vector}  and~\ref{H-in-X}.

Using the fact that $\int_{\R} \Phi_{(t)}(p)\,dp=1$, we write
\[
(L_{\gamma,t})_{i,j}(v)-(H_{\gamma})_{i,j}(v)
=\int_{\R} \Phi_{(t)}(p) ((H_\gamma)_{i,j}(v_1-p,v_2-p)-(H_\gamma)_{i,j}(v_1,v_2))\,dp. 
\]

Applying the previous expression and Fubini's theorem, we see that the $L^1$ norm of $(L_{\gamma,t})_{i,j}-(H_{\gamma})_{i,j}$ is bounded by 
\[
\int_{\R} |\Phi_{(t)}(p)| \omega(p)\,dp
=\int_{\R} |\Phi(q)| \omega(tq)\,dq
\]
where 
\[
\omega(p)=\int_{\R^2} |(H_\gamma)_{i,j}(v_1-p,v_2-p)-(H_\gamma)_{i,j}(v_1,v_2)|\,dv.
\]
By the continuity of the $L^1$ norm, we have $\lim_{t\to 0_+} \omega(tq)=0$ for every $q\in \R$. Additionally, $\omega(tq) \leq 2\|(H_{\gamma})_{i,j}\|_1$. Applying the Lebesgue dominated convergence theorem finishes the proof.
\end{proof}

\begin{lemma}\label{L:ft-infty}
Let $\gamma=(k,u,a)\in \Gamma$, $i\in [k)$, $j\in \Z$, $t>0$. Then 
\[
\lim_{t\to \infty} (L_{\gamma,t})_{i,j} = 0 \quad\text{in}\;L^1.
\]
\end{lemma}

\begin{proof}
Let $t>0$. We have
\begin{equation}\label{E:ft}
\|(L_{\gamma,t})_{i,j}\|_{1}
=\int_{\R^2} |(L_{\gamma,t})_{i,j}(x+y,y)|\,dx dy
=\int_{\R^2} \Big|\int_{\R} \Phi_{(t)}(p) (H_\gamma)_{i,j}(x+y-p,y-p)\,dp\Big| dx dy.
\end{equation}
By Proposition~\ref{H vanishing integral}, $H$ has integral zero in the diagonal direction. Using this and making the change of variables $q=y-p$, \eqref{E:ft} becomes
\begin{equation}\label{E:ft-rewritten}
\int_{\R^2} \Big|\int_{\R} (\Phi_{(t)}(y-q)-\Phi_{(t)}(y)) (H_\gamma)_{i,j}(x+q,q)\,dq\Big| dx dy.
\end{equation}
Performing a further change of variables $z=y/t$ and using the triangle inequality, \eqref{E:ft-rewritten} is bounded by 
\[
\int_{\R^2} \omega\Big(\frac{q}{t}\Big) |(H_\gamma)_{i,j}(x+q,q)|\,dx dq,
\]
where
\[
\omega(p)=\int_{\R}|\Phi(z-p)-\Phi(z)|\,dz.
\]
By the continuity of the $L^1$ norm, we have $\lim_{t\to \infty} \omega(q/t)=0$ for every $q\in \R$. Additionally, $\omega(q/t) \leq 2\|\Phi\|_1$. Applying the Lebesgue dominated convergence theorem finishes the proof.
\end{proof}

\begin{definition}[L multiplier]\label{L multiplier}
Let $\gamma=(k,u,a)\in \Gamma$.
 %Let $(H_\gamma)_{i,j}$ be as in Definition \ref{H multiplier}.
Define the index set 
\[\mathcal{I}_{\gamma} = \{(m,0)\,:\,m\in\mathbb{Z}, m\not=0\} \cup \{(0,l)\,:\,|l|\le \Delta_{\gamma}\} \subset \mathbb{Z}^2.\] 
%We equip $\mathcal{I}_{\gamma}$ with the lexicographic order, i.e. $(h,l)<(h',l')$ iff $h<h'$ or $h=h'$ and $l<l'$.
For every $\iota \in\mathcal{I}_{\gamma}$, we define  $L_{\gamma,\iota}=(L_{\gamma,\iota})_{i\in [k),j\in \Z}$ 
such that for $|l|\le \Delta_{\gamma}$,
\begin{equation}
  (L_{\gamma,(0,l)})_{i,j}:= (H_{\gamma})_{i,j} \ast_{(1,1)} (\Phi_{(a_i^1(j+l-1))}-\Phi_{(a_i^1(j+l))})\, ,
  %\mathcal F^{-1}((\xi,\eta) \mapsto\widehat{(\Phi_{(a_i^1(j+l-1))}}-\widehat{\Phi_{(a_i^1(j+l))}})(\xi+\eta)\widehat{(H_{\gamma})_{i,j}}(\xi,\eta))\,,
 \end{equation}
% (we also define $(L_{y,0})_{i,j} = \sum_{|l|\le d_i} (L_{y,0,l})_{i,j}$)
if $h>0$, then
\begin{equation}
(L_{\gamma,(h,0)})_{i,j}:= (H_{\gamma})_{i,j} \ast_{(1,1)} (\Phi_{(2^{h-1}a_i^1(j+\Delta_{\gamma}))}-\Phi_{(2^{h}a_i^1(j+\Delta_{\gamma}))})\, ,
%\mathcal F^{-1}((\xi,\eta) \mapsto \widehat{(\Phi_{(2^{h-1}a_i^1(j+d))}}-\widehat{\Phi_{(2^{h}a_i^1(j+d))}})(\xi+\eta)\widehat{(H_{\gamma})_{i,j}}(\xi,\eta))
\end{equation}
and if $h<0$, then
\begin{equation}
 (L_{\gamma,(h,0)})_{i,j}(\xi,\eta):= (H_{\gamma})_{i,j} \ast_{(1,1)} (\Phi_{(2^{h}a_i^1(j-\Delta_{\gamma}-1))}-\Phi_{(2^{h+1}a_i^1(j-\Delta_{\gamma}-1))})\, .
 %\mathcal  F^{-1}((\xi,\eta) \mapsto (\widehat{\Phi_{(2^{h}a_i^1(j-d-1))}}-\widehat{\Phi_{(2^{h+1}a_i^1(j-d-1))}})(\xi+\eta)\widehat{(H_{\gamma})_{i,j}}(\xi,\eta))
\end{equation}
\end{definition}

\begin{rem}
Note that we could have equivalently used $a_{i}^0$ in place of $a_i^1$.
\end{rem}


\begin{definition}[Summation over $\mathcal I_{\gamma}$]\label{summation-definition}
Let $X$ be a normed $\R$-vector space.
Let $\gamma=(k,u,a)\in \Gamma$ and let $D_\iota \in X$ for $\iota \in \mathcal{I}_{\gamma}$. We define
\[
\sum_{\iota \in \mathcal{I}_{\gamma}} D_{\iota}:=\lim_{N\to \infty} \sum_{\iota \in\mathcal{I}_{\gamma}, ~|\iota| \leq N} D_{\iota},
\]
where the limit is taken in $X$.
\ls{Do we want a pointwise definition, or rather a definition of the $L^1$ sum?}\jr{it seems to be used in the $L^1$ sense, but we should define it in any normed space, changed}
\end{definition}

\begin{proposition}[Properties of L multipliers]\label{sum L multiplier convergence-L1}
Let $\gamma=(k,u,a)\in \Gamma$ and $\iota \in\mathcal{I}_{\gamma}$. Then
 $L_{\gamma,\iota}$ is a well-defined element of $\mathcal{X}_k$. In addition, if $i
\in [k)$ and $j\in \Z$ then
\begin{equation}\label{E:L2-convergence}
(H_{\gamma})_{i,j} = \sum_{\iota \in\mathcal{I}_{\gamma}} (L_{\gamma,\iota})_{i,j}, 
\end{equation}
where the sum on the right converges in $L^1$. 
\end{proposition}

\begin{proof}
Each  
$(L_{\gamma,\iota})_{i,j}$ belongs to $W_0(\R^2)$ by Propositions \ref{P:schwartz-into-wiener}, \ref{convolution vector} and~\ref{H-in-X}. Thus, $L_{\gamma,\iota}\in \mathcal{X}_k$. 

Let $N\in \N$ be such that $N >d$. Then
\[
\sum_{\iota \in\mathcal{I}_{\gamma}, ~|\iota| \leq N} (L_{\gamma,\iota})_{i,j}
=(L_{\gamma,2^{-N} a_i^1(j-d-1)})_{i,j}-(L_{\gamma,2^{N} a_i^1(j+d)})_{i,j}.
\]
An application of Lemmas~\ref{L:F_t} and~\ref{L:ft-infty} yields~\eqref{E:L2-convergence}.
\end{proof}

\begin{lemma}\label{sandwich sums L1}
Let $\gamma=(k,u,a)\in \Gamma$.
Assume $X, X_\iota\in \mathcal{X}_k$ for $\iota\in \mathcal{I}_{\gamma}$ are such that for every $i\in [k)$, $j\in\Z$, $X_{i,j}=\sum_{\iota\in \mathcal{I}_{\gamma}} (X_\iota)_{i,j}$ holds in $L^1$.
Then for every $i\in [k), j\in\Z$,
\[ \mathbf{M}(\gamma,X,i)_j = \sum_{\iota\in \mathcal{I}_{\gamma}} \mathbf{M}(\gamma,X_\iota,i)_j\quad\text{in}\;L^1.\]
\end{lemma}

\begin{proof}
Let $N\in \N$ be such that $N >d$ and $y\in (\R^2)^k$.
Then
\begin{align}
&\mathbf{M}(\gamma,X,i)_j(y)-\sum_{\iota \in\mathcal{I}_{\gamma}, ~|\iota| \leq N} \mathbf{M}(\gamma,X_{\iota},i)_j(y)\\
&=\Big(\prod_{m\in [i)} (G_{\gamma})_{m,j}(y_m) \Big) \Big(X_{i,j}(y_i)- \sum_{\iota \in\mathcal{I}_{\gamma}, ~|\iota| \leq N} (X_{\iota})_{i,j}(y_i) \Big)\Big(\prod_{m=i+1}^{k-1} (G_{\gamma})_{m,j-1}(y_m) \Big).
\end{align}
Since $X_{i,j}-\sum_{\iota \in\mathcal{I}_{\gamma}, ~|\iota| \leq N}(X_{\iota})_{i,j}$ converges to $0$ in $L^1$ as $N \rightarrow \infty$, we deduce that $\mathbf{M}(\gamma,X,i)_j-\sum_{\iota \in\mathcal{I}_{\gamma}, ~|\iota| \leq N} \mathbf{M}(\gamma,X_{\sigma(l)},i)_j$ converges to $0$ in $L^1$ as $N \rightarrow \infty$ as well. 
\end{proof}


\begin{lemma}\label{prism sum le sum prism-L1}
Let $\gamma=(k,u,a)\in \Gamma$. 
If $M=\sum_{\iota\in \mathcal{I}_\gamma} M_\iota$ holds in $L^1$ for $M,M_\iota\in W_0((\R^2)^k)$, then for every $\mathbf{F}\in\mathfrak{F}$,
\[ |\Lambda_k(M)(\mathbf{F})| \le \sum_{\iota\in \mathcal{I}_\gamma} |\Lambda_k(M_\iota)(\mathbf{F})|. \]
\end{lemma}


\begin{proof}
% If $I$ is finite then the conclusion follows from the triangle inequality. Assume that $I$ is infinite and that the summation over $I$ is defined via a bijection $\sigma:\N \rightarrow I$, i.e.
% \[
% \sum_{\iota \in I} D_{\iota} =\lim_{N\to \infty} \sum_{l=0}^\infty D_{\sigma(l)}.
% \]
Let $E_N=M - \sum_{\iota \in\mathcal{I}_{\gamma}, ~|\iota| \leq N} M_{\sigma(l)}$. By assumption, $\|E_N\|_2 \to 0$ as $N\to\infty$. Let $\mathbf{F}\in\mathfrak{F}$. By the triangle inequality, 
\[
|\Lambda_k(M)(\mathbf{F})| \le \sum_{\iota \in\mathcal{I}_{\gamma}, ~|\iota| \leq N} |\Lambda_k(M_{\sigma(l)})(\mathbf{F})|
+|\Lambda_k(E_N)(\mathbf{F})|.
\]
By Propositions~\ref{prism BL inequality} and~\ref{M to K},
\[
|\Lambda_k(E_N)(\mathbf{F})| \leq \|E_N\|_1.
\]
Letting $N \to \infty$ then finishes the proof.
\end{proof}



\begin{definition}[N multiplier]\label{N multiplier}
  Let $\gamma=(k,u,a)\in \Gamma$ and assume $k\le n-1$.  
Let $\nu=1$ if $k<n-1$ and $\nu=2$ if $k=n-1$ and let $i\in [k)$, $j\in \Z$.
For $\iota\in\mathcal{I}_{\gamma}$ define functions $\sigma_{\gamma,\iota,i,j}:\R\to\R$ so that
if $|l|\le \Delta_{\gamma}$,
\[ \sigma_{\gamma,(0,l),i,j} = s(a_i^1(\cdot+l), j) \]
and for $h>0$,
\[ \sigma_{\gamma,(h,0),i,j} = s(2^h a_i^1(\cdot+\Delta_{\gamma}), j) \]
and for $h<0$,
\[ \sigma_{\gamma,(h,0),i,j} = s(2^h a_i^1(\cdot-\Delta_{\gamma}), j). \]
For every $\iota\in\mathcal{I}_{\gamma}$ we define $N_{\gamma,\iota}= (N_{\gamma,\iota})_{i\in [k),j\in \Z}$ such that
\begin{equation}
 (N_{\gamma,\iota})_{i,j} = 
 \mathcal F^{-1}((\xi,\eta) \mapsto \widehat{\sigma_{\gamma,\iota,i,j}}(\xi+\eta)^{-\nu}\widehat{(L_{\gamma,\iota})_{i,j}}(\xi,\eta))\, .
\end{equation}
\end{definition}

\begin{proposition}
Let $\gamma=(k,u,a)\in \Gamma$ and assume $k\le n-1$. Then for every $\iota\in\mathcal{I}_{\gamma}$, the double sequence $N_{\gamma,\iota}$ is well-defined and belongs to $\mathcal{X}_k$.
\end{proposition}

\begin{proof}
Each $(N_{\gamma,\iota})_{i,j}$ is real-valued because it is the inverse Fourier transform of an even real-valued function.
Additionally, $(N_{\gamma,\iota})_{i,j}$ can be written in the form
\[
(N_{\gamma,\iota})_{i,j}=(H_\gamma)_{i,j} \ast_{(1,1)} \rho
\]
for some function $\rho: \R \to \R$. 
Thanks to Proposition~\ref{four scale Gaussian kernel}, we have $\rho \in W_0(\R)$.  Proposition~\ref{H-in-X} yields that $(H_{\gamma})_{i,j} \in W_0(\R^2)$. Thus, $(N_{\gamma,\iota})_{i,j} \in W_0(\R^2)$ by Proposition~\ref{convolution vector}, i.e.\ $N_{\gamma,\iota} \in \mathcal{X}_k$.
\end{proof}





\subsection{Gaussian domination}

\begin{proposition}[Kernel estimate for Gaussian domination]\label{H kernel estimate Gaussian domination}
Let $\gamma=(k,u,a)\in \Gamma$ and let $i\in [k)$.
There exists a finite multiset $\mathcal{P}\subset B_{\mathrm{dist}}(a_i^1,\Delta_\gamma)^2\times [2)$ with $\#\mathcal{P}= 6$ and for every $v\in \R^2$ and $j\in\Z$,
\[|(H_\gamma)_{i,j}|(v) \le C_{\ref{H kernel estimate Gaussian domination}} \sum_{(t_0,t_1,u)\in\mathcal{P}} \langle (W_u v)_0\rangle_{(t_0(j))}^2 \langle (W_u v)_1\rangle_{(t_1(j))}^2, \] 
where 
$C_{\ref{H kernel estimate Gaussian domination}}
=
4C_{\ref{diagonal square root},2}^{\,2}
+
C_{\ref{Gaussian bump decay},0,2}^{\,2}.$
\end{proposition}
% Multiset because the sequences could coincide and then we want to count with multiplicity
\begin{proof}
For $m,\ell\in[2)$, $j\in\Z$, let
$t_{m,\ell}(j)=a_i^\ell(j+m-1).$
Thus $t_{0,\ell}(j)=a_i^\ell(j-1)$ and
$t_{1,\ell}(j)=a_i^\ell(j)$. By Propositions
\ref{Operations on spaced sequences} and
\ref{Properties of distance of sequences}, we have
$t_{m,\ell}\in{\rm A}$ and
$\dist(a_i^1,t_{m,\ell})
%\le \Delta(a_i)+1
\le \Delta_\gamma$.

Assume first that $u(i)=0$. For $m\in[2)$, let
$t_m^+=\max(t_{m,0},t_{m,1})\in \mathrm{A}$, and set
\[
\mathcal P
=
\left\{
(t_m^+,t_{m'}^+,0):m,m'\in[2)
\right\}
\cup
\left\{
(t_{m,0},t_{m,1},0):m\in[2)
\right\}.
\]
By Proposition~\ref{Operations on spaced sequences} and Proposition~\ref{Properties of distance of sequences}, if $t$ is one of $t_m^+,t_{m,0},t_{m,1}\in\mathrm{A}$, then 
\[\dist(a_i^1,t)\le 1 + \Delta(a_i)\le \Delta_\gamma,\]
so $\mathcal P\subset{B_{\dist}(a_i^1,\Delta_\gamma)}^2\times[2)$. Also by definition
$\#\mathcal P=6$. 

Fix $j\in\mathbb Z$, and write
$s=(s_\gamma)_{i,j}$ and $H=(H_\gamma)_{i,j}$. If $b$ is the
sequence in \eqref{E:b-definition}, then
\[
t_m^+(j)\le b(j+m-1)\le\sqrt2\,t_m^+(j).
\]

By Definitions~\ref{H multiplier} and \ref{2D Gaussians}, and the
nonnegativity of the Gaussians,
\[
|H(v)|
\le
|s(v_0)s(v_1)|
+
\sum_{m=0}^1
\g_{(t_{m,0}(j))}(v_0)
\g_{(t_{m,1}(j))}(v_1).
\]
Proposition~\ref{diagonal square root}, applied with $N=2$,
gives
\[
|s(x)|
\le
2C_{\ref{diagonal square root},2}
(\langle x\rangle_{(t_0^+(j))}^2+
\langle x\rangle_{(t_1^+(j))}^2).
\]
By Proposition~\ref{Gaussian bump decay} and scaling,
$\g_{(r)}(x)\le
C_{\ref{Gaussian bump decay},0,2}\langle x\rangle_{(r)}^2$. The claim follows.

Assume now that $u(i)=1$, and set
\[
\mathcal P
=
\left\{
(t_{m,1},t_{m',1},0):m,m'\in[2)
\right\}
\cup
\left\{
(t_{m,0},t_{m,1},1):m\in[2)
\right\}.
\]
Again, $\#\mathcal P=6$, and the required distance condition
holds.

In this case $s=(s_\gamma)_{i,j}=s(\sqrt2a_i^1,j)$.
Moreover,
\[
|H(v)|
\le
|s(v_0)s(v_1)|
+
\sum_{m=0}^1
\g_{(t_{m,0}(j))}\bigl((W_1v)_0\bigr)
\g_{(t_{m,1}(j))}\bigl((W_1v)_1\bigr).
\]
Comparing the scales $\sqrt2t_{m,1}(j)$ with $t_{m,1}(j)$ and
applying Proposition~\ref{diagonal square root} with $N=2$, we get
\[
|s(x)|
\le
2C_{\ref{diagonal square root},2}
(
\langle x\rangle_{(t_{0,1}(j))}^2
+
\langle x\rangle_{(t_{1,1}(j))}^2
).
\]
Together Proposition~\ref{Gaussian bump decay} we obtain the claim.
\end{proof}

\begin{proposition}[Kernel derivative estimate for Gaussian domination]\label{H kernel derivative estimate Gaussian domination}
Let $\gamma=(k,u,a)\in \Gamma$ and let $i\in [k)$.
There exists a finite multiset $\mathcal{P}\subset B_{\mathrm{dist}}(a_i^1,\Delta_\gamma)^2\times [2)$ with $\#\mathcal{P}\le 6$ and for every $v\in \R^2$ and $j\in\Z$,
    \[|(\partial_0+\partial_1)(H_\gamma)_{i,j}(v)| \le C_{\ref{H kernel estimate Gaussian domination}} \sum_{(t_0,t_1,u)\in\mathcal{P}} (t_0(j)^{-1} + t_1(j)^{-1}) \langle (W_u v)_0\rangle_{(t_0(j))}^2 \langle (W_u v)_1\rangle_{(t_1(j))}^2, \] 
where $
C_{\ref{H kernel derivative estimate Gaussian domination}}
=
4C_{\ref{diagonal square root},2}
C_{\ref{derivative of diagonal square root},2}
+
2C_{\ref{Gaussian bump decay},0,2}
C_{\ref{Gaussian bump decay},1,2}.
$ 
\end{proposition}

\begin{proof}[Proof \auto]
Let $\mathcal P$ be the multiset constructed in the proof of
Proposition \ref{H kernel estimate Gaussian domination}; it has the
required cardinality and distance properties. Fix $j\in\Z$, and write
$s=(s_\gamma)_{i,j}$ and $H=(H_\gamma)_{i,j}$.

The same scale comparisons used there, together with Propositions
\ref{diagonal square root} and
\ref{derivative of diagonal square root}, give
\[
|(\partial_0+\partial_1)(s\otimes s)(v)|
\le
4C_{\ref{diagonal square root},2}
C_{\ref{derivative of diagonal square root},2}
\sum_{(t_0,t_1,0)\in\mathcal P}
(t_0(j)^{-1}+t_1(j)^{-1})
\langle v_0\rangle_{(t_0(j))}^2
\langle v_1\rangle_{(t_1(j))}^2.
\]
By Proposition \ref{Gaussian bump decay}, scaling, and the chain rule,
\[
|(\partial_0+\partial_1)(Y_\gamma)_{i,j}(v)|
\le
2C_{\ref{Gaussian bump decay},0,2}
C_{\ref{Gaussian bump decay},1,2}
\sum_{(t_0,t_1,u)\in\mathcal P}
(t_0(j)^{-1}+t_1(j)^{-1})
\langle (W_uv)_0\rangle_{(t_0(j))}^2
\langle (W_uv)_1\rangle_{(t_1(j))}^2.
\]
Here the factor $2$ also covers the factor $\sqrt2$ arising when
$u=1$, since
\[
(\partial_0+\partial_1)(W_1v)_0=\sqrt2,
\qquad
(\partial_0+\partial_1)(W_1v)_1=0.
\]
The claim follows from Definition \ref{H multiplier}.
\end{proof}

\begin{proposition}[Gaussian domination, combined]\label{Gaussian domination combined}
Let
\[ C_{\ref{Gaussian domination combined},0} = 30,  C_{\ref{Gaussian domination combined},1} = 2\]
Let $\gamma=(k,u,a)\in \Gamma$ and assume $k\le n-1$. Let $i\in [k)$.  
For every $\iota\in\mathcal{I}_{\gamma}$ there exists a finite set $\mathcal{B}=\mathcal{B}_{\gamma,i,\iota}$ with 
\[\#\mathcal{B}\le C_{\ref{Gaussian domination combined},0},\]
and for every $b\in\mathcal{B}$ a $u_b\in[2)$ and
for every $(b,m)\in\mathcal{B}\times\mathbb{N}^2$ a sequence pair $p_{b,m}\in {\rm A}^2$
such that for every $b\in\mathcal{B}, m\in\mathbb{N}^2$,
\begin{equation}
\Delta(p_{b,m})\le C_{\ref{Gaussian domination combined},1} (\Delta_{\gamma}+|\iota_0| + |m|), \end{equation}
and for each $j\in \Z$, and $v\in \R^2$
\begin{equation}\label{Gaussian domination main estimate}
    |(N_{\gamma,\iota})_{i,j}(v)|\le C 2^{-\frac{|\iota_0|}{2}} \sum_{m\in\mathbb{N}^2} 2^{-\frac{|m|}{2}}\sum_{b\in \mathcal{B}}     
    G_{p_{b,m}(j), u_b}(v)
\end{equation}
and the non-negative series on the right-hand side converges for every $v\in\R^2$ and has finite $L^1$ norm, and
\[ C = C_{\ref{Gaussian domination combined},2} = 100^{100} \]
\end{proposition}
The proof of this proposition is split into three cases: Proposition \ref{Gauss domination case 1}, Proposition \ref{Gauss domination case 2}, Proposition \ref{Gauss domination case 3}, below. 

\begin{proposition}\label{Gauss domination case 1}
The conclusion of Proposition \ref{Gaussian domination combined} holds in the case
$\iota=(h,0)$, $h>0$ with the constant $C$ in \eqref{Gaussian domination main estimate} being
\[ C = C_{\ref{Gauss domination case 1}} =
2^7\pi e^{2\pi}
\tilde{C}_{\ref{standard bump properties},0,2}
C_{\ref{mean four scale Gaussian kernel},2}
C_{\ref{H kernel estimate Gaussian domination}}
\max(
C_{\ref{two bump estimate},\frac32,\frac32}+4,
4C_{\ref{bump triangle},-\frac12,\frac12,\frac32,\frac32}
C_{\ref{two bump estimate},\frac32,\frac32}
). \]
\end{proposition}

\begin{proof}
Let $\iota=(h,0)\in\mathcal{I}_{\gamma}$ with $h>0$.
Let $\rho$ be the Schwartz function defined as
\[\rho=\mathcal{F}^{-1}\Big(\frac{\widehat{\Phi_{(2^{h-1}a_i^1(j+\Delta_{\gamma}))}}-\widehat{\Phi_{(2^{h}a_i^1(j+\Delta_{\gamma}))}}}{\widehat{s(2^h a_i^1(\cdot+\Delta_{\gamma}), j)}^{\nu}}\Big).\]
Let $\lambda=\lambda(j)=2^h a_i^1(j+\Delta_{\gamma})$.
Note $\lambda\in{\rm A}$ by Proposition~\ref{Operations on spaced sequences}. We will omit the argument $j$ where clear from context.
Moreover, letting $F = (N_{\gamma,\iota})_{i,j}$ and $H=(H_\gamma)_{i,j}$, we have for $v\in\mathbb{R}^2$,
\[ F(v) = \int_{\mathbb{R}} \rho(p) H(v_0-p, v_1-p)\, dp. \]
By Proposition \ref{H vanishing integral}, $H$ has mean zero in the diagonal direction.
Setting $v_0=w_0-w_1$ and $v_1=w_0+w_1$ (then $w_m = 2^{-\frac12} (W_1 v)_m$) and changing variables in the integral and subtracting zero gives
\begin{equation}\label{eqn:Gaussianbumppf1}
F(v) = \int_{\mathbb{R}} \Big(\rho(w_0+p)-\rho(w_0)\Big) H(-w_1-p, w_1-p)\,dp.
\end{equation}
By Proposition \ref{mean four scale Gaussian kernel} and Proposition \ref{standard bump properties} (each applied with $N=2$),
\begin{equation}\label{eqn:Gaussianbumppf1rhomvt}
|\rho(w_0+p)-\rho(w_0)| \leq 2\pi \tilde{C}_{\ref{standard bump properties},0,2}
C_{\ref{mean four scale Gaussian kernel},2}\min(1, \lambda^{-1}|p|) (\langle w_0+p\rangle^2_{(\lambda)}+
\langle w_0\rangle^2_{(\lambda)})
\end{equation}
By Proposition \ref{H kernel estimate Gaussian domination} we obtain $\mathcal{P}\subset B_{\mathrm{dist}}(a_i^1,\Delta_\gamma)^2\times [2)$ with $\#\mathcal{P}\le 6$ 
such that for all $v\in\R^2$,
\[ |H(v)| \le C_{\ref{H kernel estimate Gaussian domination}} \sum_{(t_0,t_1,u)\in\mathcal{P}} \langle (W_u v)_0\rangle_{(t_0)}^2 \langle (W_u v)_1\rangle_{(t_1)}^2, \]
where we omitted the argument $j$.
Note that by Proposition \ref{Properties of distance of sequences}, if $t\in B_{\dist}(a_i^1,\Delta_\gamma)$, then
\begin{equation}\label{eqn:Gaussianbumppf1rfac}
\frac{t}{\lambda} \le 2^{-h}.
\end{equation}
In the following we will use the inequality $\min(1,a)\le a^\frac12$, valid for all $a\ge 0$. For every $p,w_1\in\R$ we have
\[ \min(1, \lambda^{-1}|p|) \le \sum_{\pm} \min(1, \tfrac12\lambda^{-1}|p\pm w_1|) \le \sum_{\pm} (\lambda^{-1}|p\pm w_1|)^{\frac12}\]
and if $t$ satisfies \eqref{eqn:Gaussianbumppf1rfac}, then the right-hand side is
\begin{equation}\label{Gaussianbumppf1u0}
\le 2^{-\frac{h}{2}} \sum_{\pm} (1+ t^{-1} |p\pm w_1|)^{\frac12}.
\end{equation}
Altogether, the integrand in \eqref{eqn:Gaussianbumppf1} is no greater than
\[ C_1 2^{-\frac{h}{2}} \sum_{\alpha\in[2)} \langle w_0+\alpha p\rangle_{(\lambda)}^2 \Big( 2\sum_{(t_0,t_1,0)\in\mathcal{P}} \langle -w_1-p\rangle_{(t_0)}^{\frac32} \langle w_1-p\rangle_{(t_1)}^{\frac32} +\sum_{(t_0,t_1,1)\in\mathcal{P}} \langle p\rangle^{\frac32}_{(t_0)} \langle w_1\rangle_{(t_1)}^2  \Big),\]
where we have used \eqref{Gaussianbumppf1u0} for the terms with $u=0$ and $\min(1,\lambda^{-1}|p|)\le 2^{-\frac{h}{2}}(1+t_0^{-1}|p|)^{\frac12}$ for the terms with $u=1$, and
$C_1 = 2\pi \tilde{C}_{\ref{standard bump properties},0,2}
C_{\ref{mean four scale Gaussian kernel},2} C_{\ref{H kernel estimate Gaussian domination}}.$
Integrating over $p\in\R$ we obtain that $|F(v)|$ is no greater than $2C_1 2^{-\frac{h}{2}}$ times
\begin{equation}\label{eqn:Gaussianbumppf1bump}
\sum_{(t_0,t_1,1)\in\mathcal{P}}\langle w_0\rangle_{(\lambda)}^2 \langle w_1\rangle^2_{(t_1)} \int_{\R}\langle p\rangle_{(t_0)}^{\frac32}\,dp
\end{equation}
\begin{equation}\label{eqn:Gaussianbumppf2bump}
 +\langle w_0\rangle_{(\lambda)}^2 \sum_{(t_0,t_1,0)\in\mathcal{P}} \int_{\R} \langle w_1+p\rangle_{(t_0)}^{\frac32} \langle w_1-p\rangle_{(t_1)}^{\frac32}\,dp  + \sum_{(t_0,t_1,1)\in\mathcal{P}} \langle w_1\rangle^2_{(t_1)} \int_\R \langle w_0+p\rangle_{(\lambda)}^2 \langle p\rangle_{(t_0)}^{\frac32}  \, dp
\end{equation}
\begin{equation}\label{eqn:Gaussianbumppf3bump}
+ \sum_{(t_0,t_1,0)\in\mathcal{P}}\int_{\R} \langle w_0+p\rangle_{(\lambda)}^2 \langle w_1+p\rangle_{(t_0)}^{\frac32} \langle w_1-p\rangle_{(t_1)}^{\frac32}\,dp
\end{equation}
%Define the multisets $\tilde{\mathcal{M}}_{m,u} = \{ t_m \,:\,(t_0,t_1,u)\in \mathcal{P}\}$ for $m,u\in [2)$.
By the two bump estimate, Proposition \ref{two bump estimate}, we may bound \eqref{eqn:Gaussianbumppf1bump} and \eqref{eqn:Gaussianbumppf2bump} by 
\[ 
C_{2} \sum_{r\in \mathcal{M}_1} \langle (W_1 v)_0\rangle_{(\lambda)}^{\frac32}\langle (W_1 v)_1\rangle_{(r)}^{\frac32},
\]
where $\mathcal{M}_1=\{\max(t_0,t_1)\,:\,(t_0,t_1,0)\in \mathcal{P}\}\cup \{ t_1 \,:\,(t_0,t_1,1)\in \mathcal{P}\}$ as multiset, and 
$C_2 = 4(C_{\ref{two bump estimate},\frac32,\frac32}+ \int_{\R}\langle p\rangle^{\frac32}\,dp).$

To estimate \eqref{eqn:Gaussianbumppf3bump} we first use Proposition \ref{bump triangle} to obtain
\[ \langle w_1+p\rangle_{(t_0)}^{\frac32} \langle w_1-p\rangle_{(t_1)}^{\frac32} 
\leq C_{\ref{bump triangle},-\frac12,\frac12,\frac32,\frac32}(\langle w_1\rangle_{(t_0)}^{\frac32} \langle w_1-p\rangle_{(t_1)}^{\frac32} + \langle w_1+p\rangle_{(t_0)}^{\frac32} \langle w_1\rangle_{(t_1)}^{\frac32}). \]
Together with Proposition \ref{two bump estimate} this allows us to bound \eqref{eqn:Gaussianbumppf3bump} by  
\[ C_3 \sum_{(t_0,t_1,0)\in\mathcal{P}} \langle w_0+w_1\rangle_{(\lambda)}^{\frac32} \langle w_1\rangle_{(t_0)}^{\frac32} + \langle w_0-w_1\rangle_{(\lambda)}^{\frac32} \langle w_1\rangle_{(t_1)}^{\frac32},   \]
where $C_3=C_{\ref{bump triangle},-\frac12,\frac12,\frac32,\frac32} C_{\ref{two bump estimate},\frac32,\frac32}$.
By Proposition \ref{orthogonal decay} the previous is no greater than
\[
16C_3\sum_{(t_0,t_1,0)\in\mathcal P}
\langle v_0\rangle_{(t_0)}^{\frac32}
\langle v_1\rangle_{(\lambda)}^{\frac32} +
\langle v_0\rangle_{(\lambda)}^{\frac32}
\langle v_1\rangle_{(t_1)}^{\frac32} +
\langle (W_1v)_0\rangle_{(\lambda)}^{\frac32}
\langle (W_1v)_1\rangle_{(t_0)}^{\frac32}+
\langle (W_1v)_0\rangle_{(\lambda)}^{\frac32}
\langle (W_1v)_1\rangle_{(t_1)}^{\frac32}
\]
which we can write as
\[ 16C_3 \sum_{(r_0,r_1,u)\in\mathcal{M}_2}\langle (W_u v)_0\rangle_{(r_0)}^{\frac32}
\langle (W_u v)_1\rangle_{(r_1)}^{\frac32} \]
with $\mathcal{M}_2
=
\{(t_0,\lambda,0):(t_0,t_1,0)\in\mathcal{P}\}
\cup
\{(\lambda,t_1,0):(t_0,t_1,0)\in\mathcal{P}\}
\cup
\{(\lambda,t_0,1):(t_0,t_1,0)\in\mathcal{P}\}
\cup
\{(\lambda,t_1,1):(t_0,t_1,0)\in\mathcal{P}\}$ as a multiset.
To summarize, we showed that 
\[ |F(v)|\le C_4 2^{-\frac h2} \sum_{b\in\mathcal{B}}\langle (W_{u_b} v)_0\rangle_{(p_{b,0})}^{\frac32}
\langle (W_{u_b} v)_1\rangle_{(p_{b,1})}^{\frac32},  \]
where $\mathcal{B}$ is a finite set in bijection with the multiset $\mathcal{M}_1\sqcup \mathcal{M}_2$ and $u_b=1$, $p_{b}=(\lambda, r)\in\mathrm{A}^2$ when $b$ corresponds to $r\in \mathcal{M}_1$, and $u_b=u$, $p_b=(r_0,r_1)$ when $b$ corresponds to $(r_0,r_1,u)\in\mathcal{M}_2$, and
$C_4 = 2C_1 \max(C_2,16C_3).$
Finally we use Gaussian domination, Proposition \ref{Gaussian domination}, to estimate the right hand side of the previous display by
\[ 8 C_{\ref{Gaussian domination}}^2 C_4  2^{-\frac h2}
\sum_{b\in\mathcal B}
\sum_{m\in\mathbb{N}^2}
2^{-\frac{|m|}{2}}
\g_{(2^{m_0}p_{b,0})}\big((W_{u_b}v)_0\big)
\g_{(2^{m_1}p_{b,1})}\big((W_{u_b}v)_1\big). \]
\end{proof}

\begin{proposition}[Gaussian domination, far case]\label{Gauss domination case 2}
The conclusion of Proposition \ref{Gaussian domination combined} holds in the case
$\iota=(h,0)$, $h<0$ with the constant $C$ in \eqref{Gaussian domination main estimate} being
\[ C = C_{\ref{Gauss domination case 2}} =
3\cdot 2^7 e^{2\pi}
\tilde{C}_{\ref{standard bump properties},0,3}
C_{\ref{four scale Gaussian kernel},3}
C_{\ref{H kernel derivative estimate Gaussian domination}}
C_{\ref{bump triangle},1,1,2,2}
C_{\ref{two bump estimate},2,2} \]
\end{proposition}

\begin{proof}[Proof \auto]
Let $\iota=(h,0)\in\mathcal I_\gamma$ with $h<0$, and let
\[
\rho=\mathcal F^{-1}\Bigg(
\frac{
\widehat{\Phi_{(2^h a_i^1(j-\Delta_\gamma-1))}}
-
\widehat{\Phi_{(2^{h+1}a_i^1(j-\Delta_\gamma-1))}}
}{
\widehat{s(2^h a_i^1(\cdot-\Delta_\gamma),j)}^\nu
}
\Bigg).
\]
Set $\lambda(j)=2^{h+1}a_i^1(j-\Delta_\gamma-1)$. Then $\lambda\in{\rm A}$ by Proposition \ref{Operations on spaced sequences}. As before, write
$F=(N_{\gamma,\iota})_{i,j}$ and $H=(H_\gamma)_{i,j}$, so that
\[
F(v)=\int_\R \rho(p)H(v_0-p,v_1-p)\,dp.
\]

By Propositions \ref{four scale Gaussian kernel} and
\ref{standard bump properties}, applied with $N=3$,
\[
|\rho(p)|
\le
\tilde{C}_{\ref{standard bump properties},0,3}
C_{\ref{four scale Gaussian kernel},3}
\bigl(
\langle p\rangle_{(\lambda/2)}^3+
\langle p\rangle_{(\lambda)}^3
\bigr)
\le
3\tilde{C}_{\ref{standard bump properties},0,3}
C_{\ref{four scale Gaussian kernel},3}
\langle p\rangle_{(\lambda)}^3.
\]
The multiplier defining $\rho$ vanishes in a neighborhood of the origin, hence $\int_\R\rho=0$. Consequently,
\[
\left|\int_{-\infty}^p\rho(q)\,dq\right|
\le
\frac32
\tilde{C}_{\ref{standard bump properties},0,3}
C_{\ref{four scale Gaussian kernel},3}
\lambda\langle p\rangle_{(\lambda)}^2.
\]
Integration by parts gives
\[
F(v)=
\int_\R
\left(\int_{-\infty}^p\rho(q)\,dq\right)
(\partial_0+\partial_1)H(v_0-p,v_1-p)\,dp.
\]

Let $\mathcal P$ be the multiset from Proposition
\ref{H kernel derivative estimate Gaussian domination}, and omit the argument $j$ from the scales. Setting
$v_0=w_0-w_1$ and $v_1=w_0+w_1$, Propositions
\ref{bump triangle}, \ref{two bump estimate}, and
\ref{orthogonal decay} bound, for $(t_0,t_1,0)\in\mathcal P$,
\[
\int_\R
\langle p\rangle_{(\lambda)}^2
\langle v_0-p\rangle_{(t_0)}^2
\langle v_1-p\rangle_{(t_1)}^2\,dp
\]
by
\[
8C_{\ref{bump triangle},1,1,2,2}
C_{\ref{two bump estimate},2,2}
\Bigl(
\langle v_0\rangle_{(t_0)}^2
\langle v_1\rangle_{(\max(\lambda,t_1))}^2
+
\langle v_0\rangle_{(\lambda)}^2
\langle v_1\rangle_{(\max(t_0,t_1))}^2
+
\langle (W_1v)_0\rangle_{(\lambda)}^2
\langle (W_1v)_1\rangle_{(\max(t_0,t_1))}^2
\Bigr).
\]
For $(t_0,t_1,1)\in\mathcal P$, the corresponding integral is bounded by
\[
4C_{\ref{two bump estimate},2,2}
\langle (W_1v)_0\rangle_{(\max(\lambda,t_0))}^2
\langle (W_1v)_1\rangle_{(t_1)}^2.
\]

If $t\in B_{\dist}(a_i^1,\Delta_\gamma)$, then
\[
\frac{\lambda}{t}\le 2^h.
\]
Thus
\[
\lambda(t_0^{-1}+t_1^{-1})\le 2^{h+1}.
\]
Let $\mathcal B$ index the bracket products displayed above, after collecting repetitions, and let $u_b$ and $p_b$ denote their orientations and sequence pairs. Using the explicit multiset $\mathcal P$ constructed in the proof of Proposition
\ref{H kernel estimate Gaussian domination}, we have
$\#\mathcal B\le14$. We obtain
\[
|F(v)|
\le
24
\tilde{C}_{\ref{standard bump properties},0,3}
C_{\ref{four scale Gaussian kernel},3}
C_{\ref{H kernel derivative estimate Gaussian domination}}
C_{\ref{bump triangle},1,1,2,2}
C_{\ref{two bump estimate},2,2}
2^{h/2}
\sum_{b\in\mathcal B}
\langle (W_{u_b}v)_0\rangle_{(p_{b,0})}^2
\langle (W_{u_b}v)_1\rangle_{(p_{b,1})}^2.
\]

For $m\in\N^2$, set
\[
p_{b,m}=(2^{m_0}p_{b,0},2^{m_1}p_{b,1}).
\]
Propositions \ref{Operations on spaced sequences} and
\ref{Properties of distance of sequences} give
\[
\Delta(p_{b,m})
\le
2(\Delta_\gamma+|h|+|m|).
\]
Finally, Proposition \ref{Gaussian domination}, applied with $N=2$ to both factors, gives
\[
|F(v)|
\le
C_{\ref{Gauss domination case 2}}
2^{-\frac{|h|}{2}}
\sum_{m\in\N^2}2^{-\frac{|m|}{2}}
\sum_{b\in\mathcal B}
G_{p_{b,m}(j),u_b}(v).
\]
The series converges pointwise and has finite $L^1$ norm by comparison with the product of two geometric series.
\end{proof}

\begin{proposition}[Gaussian domination, intermediate case]\label{Gauss domination case 3}
The conclusion of Proposition \ref{Gaussian domination combined} holds in the case
$\iota=(0,l)$, $|l|\le\Delta_\gamma$ with the constant $C$ in \eqref{Gaussian domination main estimate} being
\[ C = C_{\ref{Gauss domination case 3}} = 2^7e^{2\pi}
\tilde{C}_{\ref{standard bump properties},0,2}
C_{\ref{four scale Gaussian kernel},2}
C_{\ref{H kernel estimate Gaussian domination}}
C_{\ref{bump triangle},1,1,2,2}
C_{\ref{two bump estimate},2,2} \]
\end{proposition}

\begin{proof}[Proof \auto]
Let $\iota=(0,l)\in\mathcal I_\gamma$ with $|l|\le\Delta_\gamma$, and let
\[
\rho=\mathcal F^{-1}\Bigg(
\frac{
\widehat{\Phi_{(a_i^1(j+l-1))}}
-
\widehat{\Phi_{(a_i^1(j+l))}}
}{
\widehat{s(a_i^1(\cdot+l),j)}^\nu
}
\Bigg).
\]
Set
\[
\lambda_-(j)=a_i^1(j+l-1),
\qquad
\lambda_+(j)=a_i^1(j+l).
\]
Then $\lambda_-,\lambda_+\in{\rm A}$ by Proposition
\ref{Operations on spaced sequences}. Writing
$F=(N_{\gamma,\iota})_{i,j}$ and $H=(H_\gamma)_{i,j}$, we have
\[
F(v)=\int_\R\rho(p)H(v_0-p,v_1-p)\,dp.
\]
By Propositions \ref{four scale Gaussian kernel} and
\ref{standard bump properties}, applied with $N=2$,
\[
|\rho(p)|
\le
\tilde{C}_{\ref{standard bump properties},0,2}
C_{\ref{four scale Gaussian kernel},2}
\bigl(
\langle p\rangle_{(\lambda_-)}^2+
\langle p\rangle_{(\lambda_+)}^2
\bigr).
\]

Let $\mathcal P$ be the multiset from Proposition
\ref{H kernel estimate Gaussian domination}, and omit the argument $j$ from the scales. Set $v_0=w_0-w_1$ and $v_1=w_0+w_1$. For each
$\lambda\in\{\lambda_-,\lambda_+\}$ and
$(t_0,t_1,0)\in\mathcal P$, Propositions
\ref{bump triangle}, \ref{two bump estimate}, and
\ref{orthogonal decay} bound
\[
\int_\R
\langle p\rangle_{(\lambda)}^2
\langle v_0-p\rangle_{(t_0)}^2
\langle v_1-p\rangle_{(t_1)}^2\,dp
\]
by
\[
8C_{\ref{bump triangle},1,1,2,2}
C_{\ref{two bump estimate},2,2}
\Bigl(
\langle v_0\rangle_{(t_0)}^2
\langle v_1\rangle_{(\max(\lambda,t_1))}^2
+
\langle v_0\rangle_{(\lambda)}^2
\langle v_1\rangle_{(\max(t_0,t_1))}^2
+
\langle (W_1v)_0\rangle_{(\lambda)}^2
\langle (W_1v)_1\rangle_{(\max(t_0,t_1))}^2
\Bigr).
\]
For $(t_0,t_1,1)\in\mathcal P$, the corresponding integral is bounded by
\[
4C_{\ref{two bump estimate},2,2}
\langle (W_1v)_0\rangle_{(\max(\lambda,t_0))}^2
\langle (W_1v)_1\rangle_{(t_1)}^2.
\]

Let $\mathcal B$ index the resulting bracket products after collecting repetitions, and let $u_b$ and $p_b$ denote their orientations and sequence pairs. From the explicit form of $\mathcal P$ in the proof of Proposition
\ref{H kernel estimate Gaussian domination}, for each fixed $\lambda$ the first four terms of $\mathcal P$ produce at most eight elements of $\mathcal B$. The remaining two terms produce at most six elements if $u(i)=0$ and at most two elements if $u(i)=1$. Hence
\[
\#\mathcal B\le28.
\]
We obtain
\[
|F(v)|
\le
8
\tilde{C}_{\ref{standard bump properties},0,2}
C_{\ref{four scale Gaussian kernel},2}
C_{\ref{H kernel estimate Gaussian domination}}
C_{\ref{bump triangle},1,1,2,2}
C_{\ref{two bump estimate},2,2}
\sum_{b\in\mathcal B}
\langle (W_{u_b}v)_0\rangle_{(p_{b,0})}^2
\langle (W_{u_b}v)_1\rangle_{(p_{b,1})}^2.
\]

For $m\in\N^2$, set
\[
p_{b,m}=(2^{m_0}p_{b,0},2^{m_1}p_{b,1}).
\]
Propositions \ref{Operations on spaced sequences} and
\ref{Properties of distance of sequences}, together with
$|l|\le\Delta_\gamma$, give
\[
\Delta(p_{b,m})
\le
2(\Delta_\gamma+|m|).
\]
Finally, Proposition \ref{Gaussian domination}, applied with $N=2$ to both factors, gives
\[
|F(v)|
\le
C_{\ref{Gauss domination case 3}}
\sum_{m\in\N^2}2^{-\frac{|m|}{2}}
\sum_{b\in\mathcal B}
G_{p_{b,m}(j),u_b}(v).
\]
The series converges pointwise and has finite $L^1$ norm by comparison with the product of two geometric series.
\end{proof}

\begin{proposition}\label{Gauss domination constant}
We have
\[ \max(C_{\ref{Gauss domination case 1}},C_{\ref{Gauss domination case 2}},C_{\ref{Gauss domination case 3}})\le C_{\ref{Gaussian domination combined}}. \]
That is, the conclusion of Proposition \ref{Gaussian domination combined} holds with the claimed constant $C=C_{\ref{Gaussian domination combined}}$ in \eqref{Gaussian domination main estimate}.
\end{proposition}

\begin{proof}
\jr{fill in}
\end{proof}

(This is a purely arithmetic verification.)

\subsection{Main induction}

\newcommand{\InductPositiveTerms}[1]{{\hyperref[induct positive terms]{{\rm InductPositiveTerms}(#1)}}}
\begin{definition}[induct positive terms]\label{induct positive terms}
Let $k\in \N$ with $1\le k\le n$ and $C\in [1,\infty)$.
We say that $\InductPositiveTerms{k,C}$ holds if 
%there is a constant $C\in (0,\infty)$ such that
for all $\gamma=(k,u,a)\in \Gamma$, $i\in [k)$,
\begin{equation}
    \|\M(\gamma,s_{\gamma} \otimes s_{\gamma},i)\|_{{\rm M}(k)}\le C \Delta_\gamma^{2-2^{k-n+1}}.
\end{equation}
\end{definition}

\newcommand{\VanishingDiagonal}[1]{{\hyperref[vanishing diagonal]{{\rm VanishingDiagonal}(#1)}}}
\begin{definition}[vanishing diagonal]\label{vanishing diagonal}
Let $k\in \N$ with $1\le k\le n$ and $C\in [1,\infty)$.
We say that $\VanishingDiagonal{k,C}$ holds if 
%there is a constant $C\in (0,\infty)$ such that
for all $\gamma=(k,u,a)\in \Gamma$, $i\in [k)$,
\begin{equation}
    \|\M(\gamma,H_\gamma,i)\|_{{\rm M}(k)}
    \le C \Delta_\gamma^{2-2^{k-n+1}} .
\end{equation}
\end{definition}

\newcommand{\DiagonalBand}[1]{{\hyperref[diagonal band]{{\rm DiagonalBand}(#1)}}}
\begin{definition}[diagonal band]\label{diagonal band}
Let $k\in \N$ with $1\le k\le n-1$ and $C\in [1,\infty)$.
We say that $\DiagonalBand{k,C}$ holds if 
%there is a constant $C\in (0,\infty)$ such that
for all $\gamma=(k,u,a)\in \Gamma$, $i\in [k)$,
\begin{equation}
   \sum_{\iota\in\mathcal{I}_{\gamma}}
   \|\M(\gamma,L_{\gamma,\iota},i)\|_{{\rm M}(k)}
   \le C\Delta_\gamma^{2-2^{k-n+1}} .
\end{equation}
\end{definition}


\newcommand{\IncreaseData}[1]{{\hyperref[increase data]{{\rm IncreaseData}(#1)}}}
\begin{definition}[increase data]\label{increase data}
Let $k\in \N$ with $1\le k\le n-1$ and $C\in [1,\infty)$.
We say that $\IncreaseData{k,C}$ holds if 
%there is a constant $C\in (0,\infty)$ such that 
the following holds.

Let $\gamma=(k,u,a)\in \Gamma$, let $i\in [k)$, and let
$\iota\in\mathcal{I}_{\gamma}$. 
For $j\in\mathbb Z$ and $y\in(\mathbb R^2)^{k+1}$ define
\begin{equation}
    (\M_{i,\iota})_j(y):=
    \Big(\prod_{m\in [i)} (G_\gamma)_{m,j}(y_m)\Big)
    |(N_{\gamma,\iota})_{i,j}(y_i)|
    \Big(\prod_{m=i+1}^{k-1} (G_\gamma)_{m,j-1}(y_m)\Big)
    \bigl(\sigma_{\gamma,\iota,i,j}^{\otimes 2}\bigr)(y_{k}) .
\end{equation}
Then $\M_{i,\iota}\in {\rm M}(k+1)$ and
\begin{equation}
  \|\M_{i,\iota}\|_{{\rm M}(k+1)}
  \le
  C 2^{-\frac{|\iota_1|}{2}}(1+|\iota_1|)^2
  \Delta_\gamma^{2-2^{k-n+2}} .
\end{equation}
\end{definition}
We will prove that for $1\le k\le n-1$ the following implications hold:
\[ \InductPositiveTerms{k+1} \Longrightarrow \IncreaseData{k} \Longrightarrow \DiagonalBand{k} \]
\[ \Longrightarrow \VanishingDiagonal{k} \Longrightarrow \InductPositiveTerms{k} \]
(Here we have omitted the explicit constants for clarity.)
The last implication also holds for $k=n$. After observing that $\VanishingDiagonal{n}$ holds,
Theorem \ref{induct positive terms theorem} below follows by induction on $k$.

\begin{proposition}
\label{vanishing diagonal implies induct positive terms}
Let $k\in\N$ with $1\le k\le n$ and $C\in [1,\infty)$.
If $\VanishingDiagonal{k,C}$ holds, then $\InductPositiveTerms{k,kC+2}$ holds.
\end{proposition}

\begin{proof}
Fix $\gamma$. Let $C'=kC+2$.
By  positivity, Proposition \ref{positive terms}, it suffices to prove
\begin{equation}
    \|\sum_{i\in [k)} \M(\gamma,s_\gamma \otimes s_\gamma,i)\|_{{\rm M}(k)}\le C' \Delta_\gamma^{2-2^{k-n+1}}.
\end{equation}
Adding and subtracting $\M(\gamma, Y_\gamma, i)$ and using the triangle inequality, the left-hand side is
\[
\le \sum_{i\in [k)} \|\M(\gamma, H_\gamma,i)\|_{{\rm M}(k)} + \Big\|\sum_{i\in [k)} \M(\gamma, Y_\gamma,i)\Big\|_{{\rm M}(k)},
\]
which by $\VanishingDiagonal{k,C}$ and Proposition \ref{telescoping terms} is
$\le (kC+2) \Delta_\gamma^{2-2^{k-n+1}}$.
\end{proof}

\begin{proposition}
\label{diagonal band implies vanishing diagonal}
Let $k\in\N$ with $1\le k\le n-1$ and $C\in [1,\infty)$.
If $\DiagonalBand{k,C}$ holds, then $\VanishingDiagonal{k,C}$ holds.
\end{proposition}

\begin{proof}
By Proposition \ref{sum L multiplier convergence-L1},
\[ \|\M(\gamma,H_\gamma,i)\|_{{\rm M}(k)} \le \sum_{\iota \in \mathcal{I}_{\gamma}} \|\M(\gamma,L_{\gamma,\iota},i)\|_{{\rm M}(k)}, \]
where we used Lemma \ref{prism sum le sum prism-L1}, Lemma \ref{sandwich sums L1} and that $H_\gamma=\sum_{\iota\in \mathcal{I}_{\gamma}} L_{\gamma,\iota}$ converges in $L^1$.
The claim follows from $\DiagonalBand{k,C}$.
\end{proof}

\begin{proposition}[vanishing kernel integral]\label{vanishing kernel integral}
Let $\gamma=(n,w,a)\in \Gamma$.
For every $i\in [n)$,
\begin{equation}
    \|\M(\gamma,H_\gamma,i)\|_{{\rm M}(n)}=0.
\end{equation}
In particular, $\VanishingDiagonal{n,1}$ holds.
\end{proposition}

\begin{proof}
The claim follows by Proposition \ref{simplification 1 prism} if we can verify the assumption \eqref{1 prism mean zero assumption}.
For $z\in \R^{n+1}$,
\begin{equation}
\int_\R K_n((\M(\gamma,H_\gamma, i))_j(z+qe_n)\, dq
\end{equation}
\begin{equation}
=\int_\R \int_{\R^{n-1}}  M((z_{l}+p_{l},p_{l})_{l\in [n-1)}, z_{n-1}+z_{n}+q-\Sigma(p),z_{n}+q-\Sigma(p))\,dp\,dq,
\end{equation}
where we abbreviated $M=(\M(\gamma,H_\gamma, i))_j$.
By Fubini's theorem, changing variables $q\mapsto q - z_n + \Sigma(p)$ and renaming $q$ to $p_{n-1}$, this is equal to
\begin{equation}
\int_{\R^{n}}  M((z_{l}+p_{l},p_{l})_{l\in [n)}) 
 \,dp.
\end{equation}
We claim that this integral vanishes. This follows by unpacking the definition of $M$, separating the $i$th integral and 
observing by Proposition \ref{H vanishing integral} that
\begin{equation}
        \int_{\R}  (H_\gamma)_{i,j}(z_{i}+p_{i},p_{i}) 
     \,dp_i=0.
\end{equation}
\end{proof}

\begin{proposition}
\label{increase data implies diagonal band}
Let $k\in\N$ with $1\le k\le n-1$ and $C\in [1,\infty)$. Suppose that $\IncreaseData{k,C}$ holds. Then:

(i) If $k<n-1$, then $\DiagonalBand{k, 2^{10} C^{\frac12} }$ holds.

(ii) If $k=n-1$, then
$\DiagonalBand{n-1, 2^{10} C}$ holds.

\end{proposition}


\begin{proof}
Fix $\gamma,i,\iota$.
We will apply Proposition \ref{Cauchy-Schwarz at k} if $k<n-1$ and Proposition \ref{Cauchy-Schwarz at n-1} if $k=n-1$.
By definition,
\[ (L_{\gamma,\iota})_{i,j} = (N_{\gamma,\iota})_{i,j} *_{(1,1)} \sigma_{\gamma,\iota,i,j}^{*\nu}. \]


Recall that $\nu=1$ if $k<n-1$ and $\nu=2$ if $k=n-1$.
With this, the sandwich kernel $\M(\gamma, L_{\gamma,\iota},i)$ takes the form \eqref{before CS} if $k<n-1$, resp. \eqref{before CS 2} if $k=n-1$.
Here $\varphi_j = \sigma_{\gamma,\iota,i,j}$ and for $y\in (\mathbb{R}^2)^k$,
\[ \rho_j(y) = \Big(\prod_{m\in [i)} (G_{\gamma})_{m,j}(y_m) \Big) (N_{\gamma,\iota})_{i,j}(y_i) \Big(\prod_{m=i+1}^{k-1} (G_{\gamma})_{m,j-1}(y_m) \Big). \]

If $k<n-1$ we then obtain for every $\F\in\mathfrak{F}$,
\[ |\Lambda_k (\sum_{j\in [J)} \M(\gamma, L_{\gamma,\iota},i)_j )(\F)| \le |\Lambda_{k+1}(\tilde{M}_\iota)(\F)|^{\frac12} J^{\frac12}, \]
whereas if $k=n-1$, then for every $\F\in\mathfrak{F}$,
\[ |\Lambda_{n-1} (\sum_{j\in [J)} \M(\gamma, L_{\gamma,\iota},i)_j )(\F)| \le \sup_{\mathbf{\tilde{F}}\in\mathfrak{F}}|\Lambda_n(\tilde{M}_\iota)(\mathbf{\tilde{F}})|.\]
Here in either case for $y\in(\mathbb{R}^2)^{k+1}$,
\[ \tilde{M}_\iota(y) = \sum_{j\in [J)} |\rho_j((y_l)_{l\in [k)})| \varphi_j^{\otimes 2}(y_{k}).\]
That is, $\tilde{M}=\sum_{j\in [J)} (\M_{i,\iota})_j$ with $\M_{i,\iota}$ as in Definition \ref{increase data}.
Applying $\IncreaseData{k,C}$ we obtain for every $F\in\mathfrak{F}$,
\[ |\Lambda_{k+1}(\tilde{M}_\iota)(\F)|\le C \max(1,J^{1-2^{k-n+2}}) 2^{-\frac{|\iota_1|}{2}} (1+|\iota_1|)^2 %{2-2^{k-n+2}}
  \Delta_\gamma^{2-2^{k-n+2}}. \]
Thus in the case $k<n-1$,
\[ |\Lambda_k (\sum_{j\in [J)} \M(\gamma, L_{\gamma,\iota},i)_j )(\F)|\le C^{\frac12} J^{1-2^{k-n+1}} 2^{-\frac{|\iota_1|}{4}} (1+|\iota_1|) \Delta_\gamma^{1-2^{k-n+1}} \]
and in the case $k=n-1$,
\[ |\Lambda_{n-1} (\sum_{j\in [J)} \M(\gamma, L_{\gamma,\iota},i)_j )(\F)|\le C 2^{-\frac{|\iota_1|}{2}} (1+|\iota_1|)^2. \]
Summing over $\iota\in\mathcal{I}_{\gamma}$ we lose a factor of $2\Delta_{\gamma}+1\le 3 \Delta_\gamma$, which gives the claim.
Indeed, if $k\le n-1$ the new constant is $C^{\frac12}$ times
$3 + 2 \sum_{h=1}^\infty 2^{-h/4}(1+h) \le 2^{10},$
where we used $1+h\le 8\cdot 2^{h/8}$ for $h\ge 1$ and summed the geometric series.

If $k=n-1$, the new constant is $C$ times
$3 + 2 \sum_{h=1}^\infty 2^{-h/2} (1+h)^2\le 2^{10},$
where we used $(1+h)^2\le 2^{h/4+5}$ for $h\ge 1$ and summed the geometric series.
\end{proof}

\begin{proposition}
\label{induct positive terms imply increase data}
Let $k\in\N$ with $1\le k\le n-1$ and $C\in [1,\infty)$.
Assume that $\InductPositiveTerms{k+1,C}$ holds. Then $\IncreaseData{k,C\cdot C_{\ref{induct positive terms imply increase data}}}$ holds, where
\[C_{\ref{induct positive terms imply increase data}} = 2^{10} C_{\ref{Gaussian domination combined},0} (1+C_{\ref{Gaussian domination combined},1})^2 C_{\ref{Gaussian domination combined},2}.\]
\end{proposition}


\begin{proof}
Estimating the functions $|N_{\gamma,\iota}|$ by Gaussian domination using Propositions \ref{Gaussian domination combined}, Proposition \ref{Monotonicity K} and applying the dominated convergence theorem,
\begin{equation}\label{eqn:increasedatapf1}
\|\M_{i,\iota}\|_{{\rm M}(k+1)} \le C_{\ref{Gaussian domination combined},2} 2^{-|\iota_1|/2} \sum_{m\in\mathbb{N}^2} 2^{-|m|/2} \sum_{b\in\mathcal{B}_{i,\gamma,\iota}} \|\M(\widetilde{\gamma}_{b,m}, s_{\widetilde{\gamma}_{b,m}}^{\otimes 2}  ,k)\|_{{\rm M}(k+1)},    
\end{equation} 
and for each $(b,m)\in\mathcal{B}\times \mathbb{N}^2$ the new data $\widetilde{\gamma}_{b,m}$ is $(k+1,\widetilde{u}_{b},\widetilde{a}_{b,m})\in \Gamma$ with
\[ \widetilde{u}_b(i') = \left\{\begin{array}{ll}
u(i') & \text{if}\;i'=0,\dots,i-1,\;\text{or}\;i'=i+1,\dots,k-1,\\
u_b & \text{if}\;i'=i\\
0 & \text{if}\;i'=k,
\end{array}\right.
\]
\[ (\widetilde{a}_{b,m})_{i'} = \left\{\begin{array}{ll}
a_{i'} & \text{if}\;i'=0,\dots,i-1,\\
p_{b,m} & \text{if}\;i'=i\\
a_{i'}(\cdot -1) & \text{if}\;i'=i+1,\dots,k-1\\
\alpha_{\iota,i} & \text{if}\;i'=k
\end{array}\right.
\]
where $u_b\in [2), p_{b,m}\in {\rm A}^2$ is produced by Proposition \ref{Gaussian domination combined} and 
$\alpha_{\iota,i}\in {\rm A}^2$ can be inferred. To be explicit,
$\alpha_{(h,0),i}(j) = (2^{h-1/2} a_i^1(j+\Delta_{\gamma}),2^{h-1/2} a_i^1(j+\Delta_{\gamma}))$ if $h>0$,
$\alpha_{(h,0),i}(j) = (2^{h-1/2} a_i^1(j-\Delta_{\gamma}),2^{h-1/2} a_i^1(j-\Delta_{\gamma}))$ if $h<0$ and
$\alpha_{(0,l),i}(j) = (2^{-1/2} a_i^1(j+l),2^{-1/2} a_i^1(j+l))$ if $|l|\le \Delta_{\gamma}$.
In particular, $\Delta(\alpha_{\iota,i})=0$ and 
\[ \Delta_{\widetilde{\gamma}_{b,m}}=\Delta_\gamma-\Delta(a_i)+\Delta(p_{b,m}) \le  (1+C_{\ref{Gaussian domination combined},1})(\Delta_\gamma+|\iota_1|+|m|). \]
Also note $(s_{\widetilde{\gamma}_{b,m}})_{k,j}=\sigma_{\gamma,\iota,i,j}$.
We may apply $\InductPositiveTerms{k+1,C}$ and sum over $b\in \mathcal{B}_{i,\gamma,\iota}$ to obtain for each $\iota\in\mathcal{I}_{\gamma}$:
\[ \|\M_{i,\iota}\|_{{\rm M}(k+1)} \le C\; C_{\ref{Gaussian domination combined},0} (1+C_{\ref{Gaussian domination combined},1})^2 C_{\ref{Gaussian domination combined},2} 2^{-\frac{|\iota_1|}{2}} \sum_{m\in\mathbb{N}^2} 2^{-\frac{|m|}{2}} (\Delta_\gamma + |\iota_1|+|m|)^{2-2^{k-n+2}}. \]
Summing over $m$ we obtain the claim.
Indeed, estimate
\[ \sum_{m\in\mathbb{N}^2} 2^{-\frac{|m|}{2}} (\Delta_\gamma + |\iota_1|+|m|)^{2-2^{k-n+2}} \le \Delta_\gamma^{2-2^{k-n+2}} (1+|\iota_1|)^{2-2^{k-n+2}} \sum_{m\in\N^2} 2^{-\frac{|m|}{2}} (1+|m|)^{2-2^{k-n+2}} \]
\[
\leq \Delta_\gamma^{2-2^{k-n+2}} (1+|\iota_1|)^{2} \sum_{m\in\N^2} 2^{-\frac{|m|}{2}} (1+|m|)^{2}
\]
and $\sum_{m\in\N^2} 2^{-\frac{|m|}{2}} (1+|m|)^2\le 2^{10}$.
\end{proof}

\begin{proposition}\label{P:C_k-induction}
For every $k\in\N$ with $1\le k\le n$, $\InductPositiveTerms{k,C_k}$ holds, where
\[ C_k = \left\{\begin{array}{ll}
2+n, & \text{if}\;k=n,\\
2+(n-1)2^{10}C_{\ref{induct positive terms imply increase data}} C_n, & \text{if}\;k=n-1,\\
2+k 2^{10} (C_{\ref{induct positive terms imply increase data}}C_{k+1})^{\frac12}, & \text{if}\;1\le k\le n-2.
\end{array}
\right. \]
\end{proposition}

\begin{proof}
This is proved by (reverse) induction on $k$ with the base case being $k=n$.
The claim holds for $k=n$ by Proposition \ref{vanishing diagonal implies induct positive terms} with $k=n$ and Proposition \ref{vanishing kernel integral}.
Next assume $1\le k\le n-1$ and assume that $\InductPositiveTerms{k+1,C_{k+1}}$ holds.
Then $\InductPositiveTerms{k,C_k}$
follows by applying Proposition \ref{induct positive terms imply increase data}, Proposition \ref{increase data implies diagonal band}, Proposition \ref{diagonal band implies vanishing diagonal} and finally Proposition \ref{vanishing diagonal implies induct positive terms}.
\end{proof}

\ls{I have improved the constants, so that the final result holds with a constant independent of $n$.}
\jr{Oh, of course, because we only lose $k$ in the $k$th step and this is killed by the square root! So it was there already, I just inefficiently estimated the inducted constant, my bad! Very nice}
\begin{proposition}\label{P:better-induction}
\label{induct positive terms theorem}
For every $k\in\N$ with $1\le k\le n-1$, \\$\InductPositiveTerms{k,(2^{10}(k+2)C_{\ref{induct positive terms imply increase data}}^{\frac{1}{2}}+2^{\frac{1}{2}})^2}$ holds.
\end{proposition}

\begin{proof}
%If $n=2$ then the claim follows directly from Proposition~\ref{P:C_k-induction}. Assume that $n\geq 3$. 
Let $C_k$ be the constant from Proposition~\ref{P:C_k-induction}, it suffices to show that
\begin{equation}\label{E:induction-statement}
C_k \leq (2^{10}(k+2)C_{\ref{induct positive terms imply increase data}}^{\frac{1}{2}}+2^{\frac{1}{2}})^2.
\end{equation}
We prove the claim by reverse induction on $k$ with the base being $k=n-1$.

The claim holds for $k=n-1$ as
\[
2+(n-1)(n+2)2^{10}C_{\ref{induct positive terms imply increase data}} 
\leq (2^{10}(n+1)C_{\ref{induct positive terms imply increase data}}^{\frac{1}{2}}+2^{\frac{1}{2}})^2.
\]
Next, assume that $1\leq k \leq n-2$ and~\eqref{E:induction-statement} holds with $k$ replaced by $k+1$. By Proposition~\ref{P:C_k-induction},
\begin{equation}\label{E:k+1-to-k}
C_k=2+k 2^{10} (C_{\ref{induct positive terms imply increase data}}C_{k+1})^{\frac12}
\leq
2+k 2^{10} C_{\ref{induct positive terms imply increase data}}^{\frac12}(2^{10}(k+3)C_{\ref{induct positive terms imply increase data}}^{\frac{1}{2}}+2^{\frac{1}{2}}).
\end{equation}
The right-hand side of~\eqref{E:k+1-to-k} is bounded by
\[
(2^{10}(k+2)C_{\ref{induct positive terms imply increase data}}^{\frac{1}{2}}+2^{\frac{1}{2}})^2,
\]
which yields the conclusion.
\end{proof}

\begin{rem}
We only need to use the previous proposition for $k=2$. The result then holds with a constant independent of $n$. 
\end{rem}

\begin{theorem}\label{induct positive terms theorem}
$\InductPositiveTerms{2,C_{\ref{induct positive terms theorem}}}$ holds,
where $C_{\ref{induct positive terms theorem}}=(2^{12} \sqrt{C_{\ref{induct positive terms imply increase data}}}+\sqrt{2})^2$.
\end{theorem}

\begin{proof}
This follows from Proposition~\ref{P:better-induction} if $n\ge 3$ and from Proposition~\ref{P:C_k-induction} if $n=2$.
\end{proof}



\section{Reduction to the main argument}\label{sec:reduction}
In this section we prove Theorem \ref{thm:nct main real} from Theorem \ref{induct positive terms theorem}.

\subsection{Further preliminaries for the reduction}\label{sec:prelim red}

\subsubsection{Distributional prism form}
% \jr{WIP}
% \jr{Need: 
% $A$: original averages in terms of $f_i$

% $\Lambda$: prism form allowing distributional kernels in terms of $F_i$

% $\Lambda_1$: $k=1$ prism form in general induction only defined on kernels that are functions


% (1) Lemma $A$ to $\Lambda_1$

% (2) Lemma $\Lambda_1$ to $\Lambda$

% Then: 

% - Final reduction \S \ref{sec:pfmainthm} will call Lemma $A$ to $\Lambda_1$
% and invoke \S \ref{sec:on-diag}

% - Reduction on diag to off diag in \S \ref{sec:on-diag} is where we need distributional kernels: will call Lemma $\Lambda_1$ to $\Lambda$ and invoke 
% \S \ref{sec:off-diag}

% - \S \ref{sec:off-diag} proves the off diagonal estimates for $\Lambda_1$ using the main argument, Theorem \ref{induct positive terms theorem} for $k=2$: use main argument's $\Lambda_1$, $M(1), M(2)$ notation
% }


In the reduction argument we will sometimes need to refer to $\Lambda_1(M)$ when $M$ is a distribution, specifically when $M$ is $\delta_0$, or $\varphi\otimes \delta_0$, or $\delta\pm \varphi$.
For this reason we introduce a variant definition $\Lambda(M)$ that accepts distributional kernels $M$.

\begin{definition}
Let $\mathbf{F}\in\mathfrak{F}$.
Define for $y\in\R^2$,
\begin{equation}
\psi_{\Lambda}(\mathbf{F})(y) = \int_{\R^n} \prod_{h\in [2), i\in [n)} F_{i}(y^h-\Sigma(x), x_{[n)\setminus i})\,dx.
\end{equation}
\end{definition}

\begin{proposition}
\label{prop:Fx}
Let $\mathbf{F}\in\mathfrak{F}$.
For every $y\in\R^2$, the function $\R^n\to \R$ defined by
\[ x
\mapsto \prod_{h\in [2), i\in [n)} F_{i}(y^h-\Sigma(x), x_{[n)\setminus i}) \]
is measurable and integrable, and
$\psi_{\Lambda}(\mathbf{F})\in \mathcal{S}(\R^2)$. \pd{It is Schwartz in $x$, no? want this for Lemma $\ref{lem:form_pos}$}
\end{proposition}

\begin{definition}
Define for $M\in\mathcal{S}'(\R^2)$ and $\mathbf{F}\in\mathfrak{F}$,
\begin{equation}\label{eqn:Lambda-def}
\Lambda(M)(\mathbf{F}) = M(\psi_\Lambda(\mathbf{F}))
\end{equation}
\end{definition}

We also need a version of Proposition \ref{prism BL inequality}.

\begin{lemma}
\label{special prism BL inequality}  
    Let $\varphi\in\mathcal{S}(\R)$ with $\|\varphi\|_1=1$. Let $\delta_0$ denote the Dirac delta distribution at the origin.
    Then for all $\mathbf{F}\in\mathfrak{F}$,

    (i) $|\Lambda(\varphi\otimes \delta_0)(\mathbf{F})|\le 1$

    (ii) $|\Lambda(\delta_0\otimes \varphi)(\mathbf{F})|\le 1$

    (iii) $|\Lambda(\delta_0)(\mathbf{F})|\le 1$

    (iv) If $\phi\in\mathcal{S}(\R^2)$, then $|\Lambda(\phi)(\mathbf{F})|\le \|\phi\|_1$.
\end{lemma}
\begin{proof}
\begin{enumerate}
    \item[(i)] We have \[\Lambda(\varphi \otimes \delta_0)(\F) = \int_{\R}\varphi(y^0)\int_{\R^n}\Big(  \prod_{i\in [n)} F_{i}(y^0-\Sigma(x), x_{[n)\setminus i}) F_{i}(-\Sigma(x), x_{[n)\setminus i}) \Big) \,dx\, dy^0\]
    Changing variables $x_i\to -(x_i-x_{[n)\setminus i})$ gives
    \[\int_{\R}\varphi(y^0)\int_{\R^n}\Big(  \prod_{i\in [n)} F_{i}(y^0+x_i, x_{[n)\setminus i}) F_{i}(x_i, x_{[n)\setminus i}) \Big) \,dx\, dy^0\]
    Applying H\"older's inequality and then shifting by $y^0$ in the Lebesgue norms, 
    this is in absolute value
    \[\le \int_{\R}|\varphi(y^0)| dy^0 \le 1 .\]
    %\[\Lambda(\varphi \otimes \delta_0) = \int_{\R^n} \prod_{h\in [2), i\in [n)} F_{i}(y^h-\Sigma(x), x_{[n)\setminus i})\,dx.\]
        \item[(ii)] This follows from $\Lambda(\delta_0\otimes \varphi)(\F) = \Lambda(\varphi\otimes \delta_0)(\F)$ and (i).
    \item[(iii)] By definition, 
    \[\Lambda(\delta_0)(\mathbf{F}) = \delta_0(\psi_\Lambda(\F)) = \psi_\Lambda(\F)(0) = \int_{\R^n} \prod_{h\in [2), i\in [n)} F_{i}(-\Sigma(x), x_{[n)\setminus i})\,dx\]
    We change  variables $x_i\to -(x_i-x_{[n)\setminus i})$ and apply  H\"older's inequality in $x$, which  gives 
    \[|\Lambda(\delta_0)(\mathbf{F})|\le 1\]
    \item[(iv)] 
    By Proposition \ref{prism BL 
    inequality}, Proposition \ref{Lambda1 to Lambda}, and Proposition \ref{M to K},
    \[ |\Lambda(\phi)(\mathbf{F})|=|\Lambda_1(\phi)(\mathbf{F})|=|\Theta_1(K_1(\phi))(\mathbf{F})|\le \|K_1(\phi)\|_1\le \|\phi\|_1. \] \end{enumerate}
\end{proof}

\begin{lemma}[$\Lambda_1$ to $\Lambda$]\label{Lambda1 to Lambda}
Let $M\in W_0(\R^2)$ and $\mathbf{F}\in \mathfrak{F}$.
Then
\[ \Lambda(M)(\mathbf{F}) =  \Lambda_1(M)(\mathbf{F}). \]
\end{lemma}

\begin{proof}
By definition,
\[ \Lambda_1(M)(\F) = \int_{\R^2} \int_{\R^n} M(y_1+\Sigma(x),y_0 + \Sigma(x)) \prod_{h\in [2), i\in [n)} F_{i}(y_h, x_{[n)\setminus i}) dx\,dy. \]
By Fubini's theorem and changing variables $y_0\mapsto y_0-\Sigma(x)$, $y_1\mapsto y_1-\Sigma(x)$, this equals
\[ \int_{\R^2} M(y) \int_{\R^n} \prod_{h\in [2), i\in [n)} F_{i}(y_h-\Sigma(x), x_{[n)\setminus i}) dx\,dy = \int_{\R^2} M(y) \psi_\Lambda(\mathbf{F})(y)\,dy, \]
which is $\Lambda(M)(\mathbf{F})$.
\end{proof}


\begin{lemma}[$A$ to $\Lambda_1$]\label{A to Lambda}
Let $\phi\in \mathcal{S}(\R)$.
Let $\mathbf{f}=(f_0,\dots,f_{n-1})$ be functions in $\mathcal{S}(\R^n)$.
Then
\[
\|A(\phi, {\mathbf f})\|_2^2 = \Lambda_1(\phi^{\otimes 2})(\mathbf{F}),
\]
where for every $i\in [n)$, and $x\in\R^n$,
\[ F_i(x) = f_i(-x_{[1,i]},\Sigma(x),-x_{[i+1,n)}). \]
\end{lemma}

\begin{proof}
Expanding the $L^2$ norm and the square inside the $L^2$ integral, the left hand side equals
\[
 \int_{\R^n} \int_{\mathbb{R}^2}
 \Big(\prod_{(i,h)\in [n)\times [2)} f_i(u+y_h e_i) \Big)
 (\phi^{\otimes 2})(y)\, dy\, du .
\]
Changing variables $u=-x$, this becomes
\[
 \int_{\R^n} \int_{\mathbb{R}^2}
 \Big(\prod_{(i,h)\in [n)\times [2)} f_i(-x+y_h e_i) \Big)
 (\phi^{\otimes 2})(y)\, dy\, dx .
\]
Now change variables $y_h\mapsto y_h+\Sigma(x)$ for $h\in[2)$. We obtain
\[
 \int_{\R^n} \int_{\mathbb{R}^2}
 \Big(\prod_{(i,h)\in [n)\times [2)}
 f_i(-x+(y_h+\Sigma(x))e_i) \Big)
 \phi(y_0+\Sigma(x))\phi(y_1+\Sigma(x))\, dy\, dx .
\]
By the definition of $F_i$, for every $i\in[n)$ and $h\in[2)$,
\[
 F_i(y_h,x_{[n)\setminus i})
 =
 f_i\bigl(-x_{[0,i)},y_h+\Sigma(x)-x_i,-x_{[i+1,n)}\bigr)
 =
 f_i(-x+(y_h+\Sigma(x))e_i).
\]
Hence the last display is equal to
\[
 \int_{\R^2}\int_{\R^n}
 (\phi^{\otimes 2})(y_1+\Sigma(x),y_0+\Sigma(x))
 \prod_{h\in[2),\,i\in[n)}
 F_i(y_h,x_{[n)\setminus i})\,dx\,dy ,
\]
which is $\Lambda_1(\phi^{\otimes 2})(\mathbf{F})$.
\end{proof}



% \begin{lemma}[Change of variables]
% \label{lem:cov}
%     We have \ct{alert: this uses letters without quantifying them...}
%     \[\Lambda_{n}(K) = \int_{\R^n} \int_{\mathbb{R}^2} \Big(\prod_{(i,j)\in [n]\times [2)} f_i(x+u_j e_i) \Big) K(u)\, du\, dx. \]
% \end{lemma}

% \begin{proof}
% Recalling \eqref{eqn:Lambdalast_eq}, the right-hand side equals $\Lambda_n(K)$ after changing variables
% $u_j \mapsto x_n^j - \sum_{l\in[n]} x_l$ for $j=0,1$. 
% Indeed,
% %\[ \Pi(x, x_n^0, x_n^1) = \Big(x_n^0 - \sum_{l\in [n]} x_l, x_n^1 - \sum_{l\in [n]} x_l\Big),\;\text{and} \]
% \[ f_i(x+(x_n^j-\sum_{l\in[n]} x_l)e_i) = 
% F_i(x + (x_n^j - x_i)e_i).  \]
% % This is the only change of variables in $u$ that works (up to adding a constant) if we want the projections $\Pi_s$; this forces the definition of $\Pi$ (and of $\tau$).
% \end{proof}







% \begin{lemma}[Expanding the norm]\label{lem:ecv}Let $\phi\in W_0(\R)$.
% Let $\mathbf{f}=(f_0,\dots,f_{n-1})$ be functions in $W_0(\R^n)$.
% Then
% \[
% \|A(\phi, {\mathbf f})\|_2^2 = \Lambda_1(\phi^{\otimes 2})(F),
% \]
% where $f_i(x)=F_i(\sigma_i(x+(-x_i-\sum_{\ell\in [n)} x_\ell)e_i))$ and $\sigma_i$ is the coordinate permutation with $\sigma_i(x)=x_i e_0 + \sum_{j\in [i)} x_j e_{j+1} + \sum_{j\in [i+1,n-1]} x_j e_j=(x_i,(x_\ell)_{\ell\not=i})$.
% \end{lemma}
% \begin{proof}
% Expanding the $L^2$ norm, the left hand side equals
% \[  \int_{\R^n} \int_{\mathbb{R}^2} \Big(\prod_{(i,j)\in [n)\times [2)} f_i(x+u_j e_i) \Big) (\phi^{\otimes 2})(u)\, du\, dx.  \]
% The claim follows from the definitions of $\Lambda_1$, $K_1$. %Indeed,
% %\[ \Lambda_1(\phi^{\otimes 2}))(F) = \int_{\mathbb{R}^2} \int_{\mathbb{R}^n} \phi^{\otimes 2}(u_0,u_1) \prod_{i\in [n)} \prod_{h\in[2)} f_i(x+u_h e_i)\,dx\,d(u_0,u_1) \]
% \end{proof}



\subsubsection{Variation seminorms}\label{sec:variation-norm}

Let $B$ be a seminormed space, $I\subset\mathbb{R}$ a subset and $a:I\to B$ and $J,k\in\mathbb{N}$, $r\ge 1$. The \emph{$r$-variation seminorm} (with $J$ jumps) is 
\[ \|a\|_{V_{r,J}(I; B)} = \sup_{\substack{t_0<\cdots<t_J,\\t_j\in I\;\text{for all}\;j\in [J+1)}} \Big(\sum_{j\in[J)} \|a(t_{j+1})- a(t_{j})\|^r\Big)^{1/r}. \]

One usually considers $V_r(a)=\sup_{J\ge 1} V_{r,J}(a)$.
The $k$-th short variation can be expressed as
$\|a\|_{V_{r,J}([2^k, 2^{k+1}]; B)}$
and long variations as
$\|a\|_{V_{r,J}(2^\mathbb{Z}; B)}.$
% Corresponding \emph{short $r$-variation seminorms} are defined by
% \[ S_{J,r,k}(a) = \sup_{2^k\le t_0<\cdots<t_J\le 2^{k+1}} \Big(\sum_{j=1}^J \|a(t_j)- a(t_{j-1})\|^r\Big)^{1/r} \]
% and the \emph{long $r$-variation seminorms} by
% \[ L_{J,r}(a) = \sup_{k_0<\cdots<k_J} \Big(\sum_{j=1}^J \|a(2^{k_j})- a(2^{k_{j-1}})\|^r\Big)^{1/r}, \]
% where supremum is over integers $k_0<\cdots<k_J$.

%Is t_j meant to be nonnegative here?
\begin{lemma}\label{lem:shortlongjumps}
For every increasing sequence $(t_j)_{j\in [J+1)}$,
\[ \Big(\sum_{j\in [J)} \|a(t_{j+1})- a(t_{j})\|^r\Big)^{1/r} \le 2 \Big(\sum_{k\in \kappa} \|a\|_{V_{r,J}([2^k, 2^{k+1}]; B)}^r\Big)^{1/r} +\|a\|_{V_{r,J}(2^\mathbb{Z}; B)}, \]
where $\kappa$ is the set of $k\in \mathbb{Z}$ such that $2^{k}\le t_j<2^{k+1}$ for some $j\in [J+1)$.
Consequentially,
\[ \|a\|_{V_{r,J}(\mathbb{R}; B)} \le 2 \Big(\sum_{k\in\mathbb{Z}} \|a\|_{V_{r,J}([2^k, 2^{k+1}]; B)}^r\Big)^{1/r} + \|a\|_{V_{r,J}(2^\mathbb{Z}; B)}, \]
\end{lemma}
%\jr{One can replace the $r$-norm and the norm on $B$ by any pseudometrics; also nothing special about the sequence $(2^k)$}
\begin{proof}
For every real $t\in\mathbb{R}$ let $k(t)$ be the unique $k\in\mathbb{Z}$ such that $t\in [2^k,2^{k+1})$.
%Fix a sequence $t_0<\cdots<t_J$ and 
Set $k_j=k(t_j)$.
%Define $\mathcal{J}_\mathrm{short}$ as the set of $j=1,\dots,J$ such that $k_{j-1}=k_j$ and $\mathcal{J}_\mathrm{long}$ the complement.
% Then the left hand side is bounded by
% \[ \Big(\sum_{j\in\mathcal{J}_\mathrm{short}} \|a(t_j)- a(t_{j-1})\|^r\Big)^{1/r} + \Big(\sum_{j\in\mathcal{J}_\mathrm{long}} \|a(t_j)- a(t_{j-1})\|^r\Big)^{1/r} \]
%The first summand is no greater than $ \Big(\sum_{k} S_{r,J,k}(a)^r\Big)^{1/r}$.
By the triangle inequality the left hand side is
\[\le \Big(\sum_{j\in[J)} \|a(t_{j+1})- a(2^{k_{j+1}})\|^r\Big)^{1/r} + \Big(\sum_{j\in[J)} \|a(2^{k_{j}})- a(t_{j})\|^r\Big)^{1/r}\]
\[+ \Big(\sum_{j\in[J)} \|a(2^{k_{j+1}})- a(2^{k_{j}})\|^r\Big)^{1/r}. \]
Each of the first two sums is no greater than $(\sum_{k} \|a\|_{V_{r,J}([2^k, 2^{k+1}]; B)}(a)^r)^{1/r}$ and the third is $\le \|a\|_{V_{r,J}(2^\mathbb{Z}; B)}$.
%Sort into long and short jumps and use triangle inequality.
\end{proof}


\begin{lemma}\label{lem:ftccs-R}   Let $0<\alpha < \beta$ and $a:[\alpha,\beta]\to \R$ be continuously differentiable.
    Then 
\begin{equation}\label{ftccs1-R}
\|a\|_{V_{2,J}([\alpha, \beta]; \R)} \le (\tfrac{\beta-\alpha}{\alpha})^{1/2} \|ta'(t)\|_{L^2(t\in [\alpha, \beta],\frac{dt}{t})}, \end{equation}
\begin{equation}\label{ftccs2-R}
\|a\|_{V_{2,J}([\alpha, \beta]; \R)}^2 \le 8 \|a\|_{L^2([\alpha, \beta],\frac{dt}{t})}\|ta'(t)\|_{L^2(t\in [\alpha, \beta],\frac{dt}{t})}.
\end{equation}
\end{lemma}
\begin{proof}
First we prove \eqref{ftccs1-R}. Let $t_0\,\ldots,t_J$ satisfy    $\alpha \le t_0<\cdots<t_J \le \beta$. By the fundamental theorem of Calculus and the triangle inequality on $B$, 
\[|a(t_{j+1})- a(t_{j})|^2 = \Big | \int_{t_j}^{t_{j+1}} ta'(t) \frac{dt}{t} \Big |^2\le \Big ( \int_{t_j}^{t_{j+1}} | ta'(t)| \frac{dt}{t} \Big )^2 \]
By  the Cauchy-Schwarz inequality in $t$,  this is
\[ \le   (t_{j+1}-t_j)  \int_{t_j}^{t_{j+1}} |ta'(t)|^2 \frac{dt}{t^2}   \]
Using the crude bound $t_{j+1}-t_j\le \beta-\alpha$, and the fact that the integrand is non-negative we obtain 
\[\le   (\beta-\alpha)  \int_{t_j}^{t_{j+1}} |ta'(t)|^2 \frac{dt}{t^2} \le \tfrac{\beta-\alpha}{\alpha}  \int_{t_j}^{t_{j+1}} |ta'(t)|^2 \frac{dt}{t} = \tfrac{\beta-\alpha}{\alpha} \|ta'(t)\|^2_{L^2(t\in [t_j,t_{j+1}],\frac{dt}{t})}   \] 
Thus,
\[\|a\|_{V_{2,J}([\alpha,\beta]; \R)} = \sup_{\alpha \le t_0<\cdots<t_J \le \beta} \Big(\sum_{j\in[J)} |a(t_{j+1})- a(t_{j})|^2\Big)^{1/2}
\]
\[\le ( \tfrac{\beta-\alpha}{\alpha})^{1/2} \|ta'(t)\|_{L^2(t\in [\alpha,\beta],\frac{dt}{t})}  \]
This yields \eqref{ftccs1-R}. To show \eqref{ftccs2-R}, we first prove that for any   $\alpha\le t_0<\cdots<t_J \le \beta$ and an absolutely continuous non-negative  $a:\R\to \R$,  
\begin{equation}
    \label{ftccs2-pos}
    |a(t_{j+1}) - a(t_j)|^2 \leq 2 \|a\|_{L^2([t_j,t_{j+1}],\frac{dt}{t})} \|ta'(t)\|_{L^2(t\in [t_j,t_{j+1}],\frac{dt}{t})} 
\end{equation}
Once \eqref{ftccs2-pos} is shown,  we split
$a= a_+ - a_-$, where $a_+ = \max(a,0)$ and $a_- = -\min(a,0)$. 
By \eqref{ftccs2-pos} we obtain for $\ell\in\{+,-\}$, 
\[|a_\ell(t_{j+1}) - a_\ell(t_j)|^2 \le 2 \|a_\ell\|_{L^2([t_j,t_{j+1}],\frac{dt}{t})} \|ta_\ell'(t)\|_{L^2(t\in [t_j,t_{j+1}],\frac{dt}{t})}\]\[ \le 2  \|a\|_{L^2([t_j,t_{j+1}],\frac{dt}{t})} \|ta'(t)\|_{L^2(t\in [t_j,t_{j+1}],\frac{dt}{t})} \]
Therefore,    
\[\sum_{j\in [J)} |a_\ell(t_{j+1}) - a_\ell(t_j)|^2 \le 2   \sum_{j\in [J)} \|a\|_{L^2([t_j,t_{j+1}],\frac{dt}{t})}\|ta'(t)\|_{L^2(t\in [t_j,t_{j+1}],\frac{dt}{t})}.\] 
\[\le   2   \Big( \sum_{j\in [J)} \|a\|^{2}_{L^2([t_j,t_{j+1}],\frac{dt}{t})}\Big)^{1/2} \Big( \sum_{j\in[J)}\|ta'(t)\|^{1/2}_{L^2(t\in [t_j,t_{j+1}],\frac{dt}{t})}\Big)^{1/2} \]
\[ \le 2  \|a\|_{L^2([\alpha,\beta],\frac{dt}{t})} \|ta'(t)\|_{L^2(t\in [\alpha,\beta],\frac{dt}{t})} \]
where we used the Cauchy-Schwarz inequality for the first estimate in this display.
Thus,
\[\|a_\ell\|_{V_{2,J}([\alpha,\beta]; \R)} \le 2^{1/2} \|a\|^{1/2}_{L^2([\alpha,\beta],\frac{dt}{t})}\|ta'(t)\|^{1/2}_{L^2(t\in [\alpha,\beta],\frac{dt}{t})}.\]
By the triangle inequality for the variation norm,  
\[\|a\|_{V_{2,J}([\alpha,\beta]; \R)}  \le  \|a_+\|_{V_{2,J}([\alpha,\beta]; \R)} +  \|a_-\|_{V_{2,J}([\alpha,\beta]; \R)} \]
\[\le 2^{3/2} \|a\|^{1/2}_{L^2([\alpha,\beta],\frac{dt}{t})}\|ta'(t)\|^{1/2}_{L^2(t\in [\alpha,\beta],\frac{dt}{t})},\]
which implies \eqref{ftccs2-R} for any real-valued absolutely continuous $a$ by squaring both sides. 

It remains to prove \eqref{ftccs2-pos}. If $a\geq 0$,  we can estimate 
\begin{equation}
    \label{ftca} |a(t_{j+1}) - a(t_j)|^2 \le |a(t_{j+1})^2 - a(t_j)^2| = \int_{t_j}^{t_{j+1}} t (a(t)^2)' \frac{dt}{t} = \int_{t_j}^{t_{j+1}} 2a(t)a'(t)\frac{dt}{t}
\end{equation}
The last equality holds by the fundamental theorem of calculus. 
To see the ineqaulity,  we note that for any two real numbers $0<a<b$, then $a^2<b^2$ and $-2ab<-2a^2$. Thus, $|a-b|^2 = a^2 - 2ab+b^2 \le a^2 - 2a^2 + b^2 = b^2-a^2 = |a^2-b^2|$. 

By the Cauchy-Schwarz inequality in $t$, the display \eqref{ftca} is   estimated
 by 
\[\le 2 \Big(  \int_{t_j}^{t_{j+1}} a(t)^2\frac{dt}{t} \Big)^{1/2}\Big(  \int_{t_j}^{t_{j+1}} (ta'(t))^2\frac{dt}{t} \Big)^{1/2}  \]
\[= 2 \|a\|_{L^2([t_j,t_{j+1}],\frac{dt}{t})} \|ta'(t)\|_{L^2(t\in [t_j,t_{j+1}],\frac{dt}{t})} ,   \]
which shows \eqref{ftccs2-pos}. 
\end{proof}


% \pd{Let's remove this corollary}
% \begin{corollary}\label{lem:ftccs}   \todo{assumtions on $a$}
% \begin{equation}\label{ftccs1}
% \|a\|_{V_{2,J}([2^i, 2^{i+1}]; L^2(\R^d))} \le \|ta'(t)\|_{L^2(t\in [2^i, 2^{i+1}],\frac{dt}{t}; L^2(\R^d))}, \end{equation}
% \begin{equation}\label{ftccs2}
% \|a\|_{V_{2,J}([2^i, 2^{i+1}]; L^2(\R^d))}^2 \le 8 \|a\|_{L^2([2^i, 2^{i+1}],\frac{dt}{t}; L^2(\R^d))}\|ta'(t)\|_{L^2(t\in [2^i, 2^{i+1}],\frac{dt}{t}; L^2(\R^d))}.
% \end{equation}
% \end{corollary}

% \begin{proof} This follows by integrating in $x$. More precisely, 
% we write 
% \[\|a\|^2_{V_{2,J}([2^i, 2^{i+1}]; L^2(\R^d))} = \sup \sum_{j\in[J)} \|a(t_{j+1})- a(t_{j})\|_{L^2}^2     \]
% \[ = \sup  \int_{\R^d} \sum_{j\in[J)} |a(t_{j+1})(x)- a(t_{j})(x)|^2 dx   \]
% where  $\sup$ is  over all $t_0<\cdots<t_J,\, t_j\in [2^i, 2^{i+1}]$ for all $j\in [J+1]$.  

% To see \eqref{ftccs1},  we note  $(\beta-\alpha)/\alpha = (2^{i+1}-2^i)2^{-i} = 1$ and apply \eqref{ftccs1-R} of Lemma \ref{lem:ftccs-R} with $t\mapsto a(t)(x)$ for each fixed $x\in \R^d$. This gives  
% \[\le   \int_{\R^d} \|ta'(t)(x)\|^2_{L^2(t\in [2^i,2^{i+1}], \frac{dt}{t})} dx   =  \|ta'(t)\|^2_{L^2(t\in [2^i, 2^{i+1}],\frac{dt}{t}; L^2(\R^d))} \]
% as desired. 
% %one extra step: Replace the last display by \[ \le   \int_{\R^d} \|ta'(t)(x)\|^2_{L^2(t\in [2^i,2^{i+1}], \frac{dt}{t})} dx   =   \int_{\R^d} \int_{2^i}^{2^{i+1}}|ta'(t)(x)|^2 \frac{dt}{t} dx  \] \[ =  \int_{2^i}^{2^{i+1}} \|ta'(t)\|^2_{L^2}\frac{dt}{t} = \|ta'(t)\|^2_{L^2(t\in [2^i, 2^{i+1}],\frac{dt}{t}; L^2(\R^d))}\]

% To see \eqref{ftccs2}, we apply \eqref{ftccs2-R} of  Lemma \ref{ftccs2} with $t\mapsto a(t)(x)$ for each fixed $x\in \R^d$. This gives 
% \[\|a\|_{V_{2,J}([2^i, 2^{i+1}]; L^2)}^2 \le  8   \int_{\R^d} \|a(t)(x)\|_{L^2(t\in [2^i, 2^{i+1}],\frac{dt}{t})}\|ta'(t)(x)\|_{L^2(t\in [2^i,2^{i+1}],\frac{dt}{t})}  dx  \]
% Applying the Cauchy-Schwarz inequality in $x$ gives
% \[\le 8  \| \|a(t)(x)\|_{L^2(t\in [2^i, 2^{i+1}],\frac{dt}{t})}\|_{L^2(x\in \R^d)} \|\|ta'(t)(x)\|_{L^2(t\in [2^i,2^{i+1}],\frac{dt}{t})} \|_{L^2(x\in \R^d)} \]
% \[=  8   \|a\|_{L^2([2^i, 2^{i+1}],\frac{dt}{t};L^2(\R^d))} \|ta'(t)\|_{L^2(t\in [2^i,2^{i+1}],\frac{dt}{t};L^2(\R^d))}  \]\
% \end{proof}

%\jr{Can use other exponents, but not needed here}

% Before the statement was with L2 replaced by B, but not sure how to prove without splitting into positive and negative parts
% \begin{lemma}[Lemma 13 in \cite{twocommuting}]\label{lem:ftccs} \pd{Check the statement}
% \begin{equation}\label{ftccs1}
% \|a\|_{V_{2,J}([a, b]; B)} \le \|ta'(t)\|_{L^2(t\in [a,b],\frac{dt}{t}; B)}, \end{equation}
% \begin{equation}\label{ftccs2}
% \|a\|_{V_{2,J}([a, b]; B)}^2 \le \|a\|_{L^2([a,b],\frac{dt}{t}; B)}\|ta'(t)\|_{L^2(t\in [a,b],\frac{dt}{t}; B)}.
% \end{equation}
% \end{lemma}
\subsubsection{Bump functions}\label{sec:bump red}

\begin{lemma}\label{lem:ft_deriv_mul}
For a Schwartz function $\phi:\mathbb{R}\to\mathbb{C}$ let
$T\phi(x)=\frac{d}{dx} (x \phi(x))|_x$. Then
for every $m\ge 0$,
\[(\widehat{T\phi})^{(m)}(\xi) = -(m (\widehat{\phi})^{(m)}(\xi) + \xi (\widehat{\phi})^{(m+1)}(\xi)).\]
\end{lemma}

\begin{proof}
We proceed by induction on $m$.
For $m=0$ this follows from the rule for taking the Fourier transform of a derivative and vice versa:
$\widehat{T\phi}(\xi) = -2\pi i \xi\,\widehat{\Phi}(\xi)$,
where $\Phi(x)=x \phi(x)$ and $\widehat{\Phi}(\xi) = (2\pi i)^{-1} \frac{d}{d\xi} \widehat{\phi}(\xi)$.
Suppose the claim holds for $m$. Then by inductive hypothesis and the product rule,
\[ (\widehat{T\phi})^{(m+1)}(\xi) = \frac{d}{d\xi} (-(m (\widehat{\phi})^{(m)}(\xi) + \xi (\widehat{\phi})^{(m+1)}(\xi))) \]
\[ = -((m+1) (\widehat{\phi})^{(m+1)}(\xi) + \xi (\widehat{\phi})^{(m+2)}(\xi)). \]
\end{proof}

\begin{lemma}
\label{lem:bump_diagest}
   For any $u,v\in \R$ and $N\ge 0$, 
   \[\langle u\rangle^{2N}\langle v\rangle^{2N}  \le 2^{2N} \langle u+v\rangle^{N}\langle u-v\rangle^{N}  \] 
\end{lemma}

\ls{This lemma should be removed eventually. If there is still any need for it, we can use Propositions~\ref{orthogonal decay} and \ref{bump triangle} instead. (This concerns the proof of Lemma~\ref{lem:leftbump1_long}).}

\begin{proof}
    We will show 
    \[\langle u+v\rangle^{-1} \langle u-v\rangle^{-1}  \le 2^2 \langle u\rangle^{-2}\langle v\rangle^{-2},\]
    which will imply the desired estimate by raising both sides to the exponent $-N$.
    
   By the triangle inequality, 
    \[\langle u+v\rangle^{-1} \langle u-v\rangle^{-1} \le \langle |u|+|v|\rangle^{-1} \langle |u|+|v|\rangle^{-1}  \]
    \[\le (1 + 2\max(|u|,|v|))^2 \le 2^2(1 + \max(|u|,|v|))^2 \]
    \[\le  2^2(1 + \max(|u|,|v|))^2 (1 + \min(|u|,|v|))^2  =  2^2 \langle u\rangle^{-2}\langle v\rangle^{-2}\]
\end{proof}

\begin{lemma}
\label{lem:transl_pow_v}
    For any $u,v\in \R$ and $N\ge 0$,
    \[\langle u+v\rangle^{N} \le \langle u\rangle^{N}\langle v\rangle^{-N}.\]
\end{lemma}

\ls{Removed $2^N$ from the statement and simplified the proof.}
\begin{proof}
This follows from
    \[1+|u|\le 1+|u+v| +|v| \le (1+|u+v|)(1+|v|),\]
    where we use the triangle inequality for the first inequality.
 %If $|v| \le 1$, this follows from
    %\[1+|u|\le 1+|u+v| +|v|\le 2(1+|u+v|) \le 2(1+|u+v|)(1+|v|),\]
    %where we use the triangle inequality for the first inequality.
    %If $|v|\ge 1$, we estimate
    %\[(1+|u|)|v|^{-1} \le 1 + |u||v|^{-1} \le 1+|u+v||v|^{-1} + |v||v|^{-1} \le 2 (1+|u+v|),\]
     % where we use the triangle inequality for the second inequality.
      %Then we  multiply with $|v|$ on both sides and use $|v|\le 1+|v|$. 
\end{proof}


\begin{lemma} \label{lem:widebump}
Let $k \leq 2$, $t\in [2^{1-k}, 2^{2-k}]$. Then for every $u\in \R$ and $N\ge 0$,
\[\langle u-t\rangle^{N} \le C_{\ref{lem:widebump},N} 2^{-kN}\langle u\rangle^{N},\]
%\[\langle u+v-t\rangle^{N} \le C_{\ref{lem:widebump},N} \langle 2^{k}(u+v)\rangle^{N}, \]
where $C_{\ref{lem:widebump},N}=2^{3N}$.
\end{lemma}
\begin{proof}
This follows from Lemma~\ref{lem:transl_pow_v} and from the fact that
\[
\langle -t \rangle^{-N} \leq (1+2^{2-k})^N \leq 2^{3N} 2^{-kN}.
\]
%We claim 
%\[\langle u+v-t\rangle^{N} \le  2^{3N} \langle t^{-1}2^2(u+v)\rangle^{N}.  \]
%Since  $t^{-1}2^2 > 2^{k}$, this is then 
%\[\le 2^{3N} \langle 2^{k}(u+v)\rangle^{N}  \]
%as desired. The claim follows from 
%\[1+t^{-1}2^2|u+v|  = 1+t^{-1}2^2|u+v- t + t|\]
%\[\le 1+ t^{-1}2^2|u+v- t| + 2^2t^{-1}t \le  2^3 + t^{-1}2^2|u+v- t|  \le 2^3 (1 + |u+v- t|)  \]
%In the last display,  the first inequality follows from the triangle inequality and for the last inequality  we used  $t^{-1}2^2 \le 2^{k-1}2^2 \le 1$  since $k\le -1$. 
\end{proof}

\begin{lemma}\label{L:integral-one-bump}
Let $k\leq -1$. Then for every $N\geq 1$ and $u\in \R$,
\[
\int_1^2\langle u - t2^{-k}\rangle^{N+1}\,dt \leq C_{\ref{L:integral-one-bump},N} 2^{k(1-N)} \langle u \rangle^N,
\]
where $C_{\ref{L:integral-one-bump},N}=\max(2C_{\ref{lem:widebump},N},N^{-1}2^{4N+1})$. 
\end{lemma}

\begin{proof}
By the change of variables, it suffices to prove
\[
\int_{2^{1-k}}^{2^{2-k}}\langle u - t\rangle^{N+1}\,dt \leq C_{\ref{L:integral-one-bump},N} 2^{-kN} \langle u \rangle^N.
\]

Assume that $|u|\geq 2^{3-k}$. Then $\langle u-t \rangle \leq 2^k$ for any $t\in [2^{1-k},2^{2-k}]$, and an application of Lemma~\ref{lem:widebump} yields
\[
\int_{2^{1-k}}^{2^{2-k}}\langle u - t\rangle^{N+1}\,dt \leq
2^k \int_{2^{1-k}}^{2^{2-k}}\langle u - t\rangle^{N}\,dt
\leq 2 C_{\ref{lem:widebump},N} 2^{-kN}\langle u\rangle^{N}.
\]
Conversely, if $|u| <2^{3-k}$ then $1 \leq 2^{4-k} \langle u \rangle$, and thus
\[
\int_{2^{1-k}}^{2^{2-k}}\langle u - t\rangle^{N+1}\,dt \leq
\int_{\R} \langle u - t\rangle^{N+1}\,dt
=\frac{2}{N} \leq \frac{2^{4N+1}}{N} 2^{-kN} \langle u\rangle^{N}.
\]
\end{proof}

\begin{lemma}\label{lem:thetat_offcenter}
Let $k\le -1$. Then for every $N\ge 1$ and $v \in \R^2$,
    \begin{equation}\label{E:two-terms}
    \int_1^2\langle v_0 - t2^{-k}\rangle^{N+1}\langle v_1 - t2^{-k}\rangle^{N+1} \,dt \le C_{\ref{lem:thetat_offcenter},N} 2^{k(1-N)} \sum_{\ell=0}^1 \langle (W_\ell v)_0\rangle^{N} \langle (W_\ell v)_1\rangle^{N}, 
    \end{equation}
    where $C_{\ref{lem:thetat_offcenter},N} = C_{\ref{bump triangle},-2^{-1/2},2^{-1/2},N+1,N+1} C_{\ref{L:integral-one-bump},N} (C_{\ref{orthogonal decay},-2^{-1/2},2^{-1/2},0,1,N,N}+C_{\ref{orthogonal decay},-2^{-1/2},2^{-1/2},1,0,N,N})$.
\end{lemma}

\begin{proof}
By Proposition~\ref{bump triangle}, we estimate the left-hand side of~\eqref{E:two-terms} by
\[
C_{\ref{bump triangle},-2^{-1/2},2^{-1/2},N+1,N+1} \langle (W_1 v)_1\rangle^{N} \sum_{i=0}^1 \int_1^2\langle v_i - t2^{-k}\rangle^{N+1} \,dt.
\]
Let $i\in \{0,1\}$. By Lemma~\ref{L:integral-one-bump}, we have
\[
\int_1^2\langle v_i - t2^{-k}\rangle^{N+1} \,dt
\leq C_{\ref{L:integral-one-bump},N} 2^{k(1-N)} \langle v_i \rangle^N.
\]
Finally, by Proposition~\ref{orthogonal decay},
\[
\sum_{i=0}^1 \langle(W_1 v)_1\rangle^{N} \langle v_i \rangle^N
\]
\[
\leq (C_{\ref{orthogonal decay},-2^{-1/2},2^{-1/2},0,1,N,N}+C_{\ref{orthogonal decay},-2^{-1/2},2^{-1/2},1,0,N,N}) \sum_{\ell=0}^1 \langle (W_\ell v)_0\rangle^{N} \langle (W_\ell v)_1\rangle^{N}.
\]
\end{proof}



\iffalse

\begin{lemma}\label{lem:thetat_offcenter}
For $k\le -1$, $N\ge 2$, 
    \[\int_1^2\langle u - t2^{-k}\rangle^{2(N+1)}\langle v - t2^{-k}\rangle^{2(N+1)} \,dt \]
    \[\le C_{\ref{lem:thetat_offcenter},N} 2^{k} \langle 2^k(u+v)\rangle^{N}  \langle u-v\rangle^{N} \]
    where $C_{\ref{lem:thetat_offcenter},N} = \max{(2^{6N+2}, 2^{2N+3}C_{\ref{lem:widebump},N})}$.
\end{lemma}

\begin{proof}
By  Lemma \ref{lem:bump_diagest} applied with $N+1$ in place of $N$ we estimate the left-hand side by   
\[
2^{2(N+1)} 
\int_1^2
\langle u+v-t2^{1-k}\rangle^{N+1} \langle u-v\rangle^{N+1}
\,dt\]
\begin{equation}
    \label{est_cases}
    \leq 
    2^{2(N+1)}
    2^{k}\langle u-v\rangle^{N+1}  \int_{2^{1-k}}^{2^{2-k}}
\langle u+v-t\rangle^{N+1} 
\,dt
\end{equation}



We estimate the integral. 
If $|u+v|\le 2^{3-k}$, we estimate it by the integral over $\R$, 
\[\int_{\R} 
 \langle u+v-t\rangle^{N+1}
 \,dt  = 2 \int_0^\infty \langle t\rangle^{N+1} = \tfrac{2}{N}.\]
Thus, we can estimate \eqref{est_cases} by
\[
2^{2(N+1)} 
2^{k} \langle u-v\rangle^{N+1}\, \tfrac{2}{N}\]
\[
\leq 2^{2(N+1)} 
2^{k} \langle u-v\rangle^{N} \]
\[\le 2^{6N +2}  2^{k}  \langle 2^k(u+v)\rangle^{N}  \langle u-v\rangle^{N}\]
\[\le \max{(2^{6N+2}, 2^{2N+3}C_{\ref{lem:widebump},N})}  2^{k}  \langle 2^k(u+v)\rangle^{N}  \langle u-v\rangle^{N} \]


If $|u+v| > 2^{3-k} = 2\cdot 2^{2-k} >2t$, we   estimate $1+|u+v-t| > t > 2^{-k}$, and thus 
\[  
\langle u+v-t\rangle^{N+1} 
\le  2^{k} 
\langle u+v-t\rangle^{N} 
 \]
By Lemma \ref{lem:widebump}, this is 
\[\le  2^k C_{\ref{lem:widebump},N} \langle 2^{k}(u+v)\rangle^{N} \]
%  %% This became its own lemma:
% Moreover,  we have 
% \[(1+|u+v-t|)^{-10} \le  2^{30} (1+t^{-1}2^2|u+v|)^{-10} \le 2^{30} (1+2^{k}|u+v|)^{-10}  , \]
% where the second inequality holds since  $t^{-1}2^2 > 2^{k}$, while for the first inequality we use 
% \[1+t^{-1}2^2|u+v|  = 1+t^{-1}2^2|u+v- t + t|\]
% \[\le 1+ t^{-1}2^2|u+v- t| + 2^2t^{-1}t \le  2^3 + t^{-1}2^2|u+v- t|  \le 2^3 (1 + |u+v- t|)  \]
% In the last display,  the first inequality follows from the triangle inequality and for the last inequality  we used  $t^{-1}2^2 \le 2^{k-1}2^2 \le 1$  since $k\le -1$. 
Therefore, if $|u+v|> 2^{3-k}$, the integral in \eqref{est_cases} is estimated by
\[    2^{k} C_{\ref{lem:widebump},N}\int_{2^{1-k}}^{2^{2-k}}
\langle 2^k(u+v)\rangle^{N} 
\,dt  =  2^{k}C_{\ref{lem:widebump},N} 2^{1-k}\langle 2^k(u+v)\rangle^{N}  = 2C_{\ref{lem:widebump},N} \langle 2^k(u+v)\rangle^{N}\]
and thus 
\eqref{est_cases} is estimated by
\[\leq 2^{2N+3}C_{\ref{lem:widebump},N}2^{k}\langle 2^k(u+v)\rangle^{N}  \langle u-v\rangle^{N+1}\]
\[\le \max{(2^{4N}, 2^{2N+3}C_{\ref{lem:widebump},N})}2^{k}\rangle 2^k(u+v)\rangle^{N}  \langle u-v\rangle^{N} \]
 \end{proof}

 \fi

For a function $\phi$ we define 
\[(X\phi)(u) = u\phi(u).\]
\begin{lemma}
    \label{lem:phider}
If  $\phi:\R\to \R$ is  a $C^1$ function, 
    \[-t\partial_t(\phi_{(t)})(u) = \phi_{(t)}(u) + (X\phi')_{(t)}(u) \]
\end{lemma}
\begin{proof}
By the product and chain rules,
\[\partial_t (t^{-1}\phi(t^{-1}u)  ) = -t^{-2}\phi(t^{-1}u)  - t^{-3}u\phi'(t^{-1}u)  \]
    The claim now follows by multiplying with $-t$. 
\end{proof}



\begin{lemma}[Rescaled differences]\todo{too trivial to include?} 
\label{lem:rescaled_diff}
Let $\phi:\R\to\R$ be a Schwartz function such that $\widehat{\phi}$ has support in $B_1(1)$ and equals one on $B_1(\frac12)$.
Then for any integers $k_1< k_2$,
\[ \widehat{\phi_{(2^{k_1})}} - \widehat{\phi_{(2^{k_2})}} \]

 (1) is supported on $\{2^{-k_2-1}\le |\xi|\le 2^{-k_1}\}$

 (2) equals one on $\{2^{-k_2}\le |\xi|\le 2^{-k_1-1}\}$\\

 
\noindent In particular, if $k\ge 1$,
$ \widehat{\phi_{(2^{-k})}} - \widehat{\phi_{(2^k)}} $
is supported on \[\{2^{-k-1}\le |\xi|\le 2^k\} \subset \mathrm{Ann}_1(1, 2^{k+1}),\] and equals one on $\{2^{-k}\le |\xi|\le 2^{k-1}\}\supset \mathrm{Ann}_1(1, 2^{k-1})$.
\end{lemma}


\begin{proof}
This is because $\widehat{\phi_{(2^{k_2})}}$ equals one on $B_1(2^{-k_2-1})$ (and $\widehat{\phi_{(2^{k_1})}}$ equals one on the larger ball $B_1(2^{-k_1-1})$), while
$\widehat{\phi_{(2^{k_1})}}$ is supported on $B_1(2^{-k_1})$ and $\widehat{\phi_{(2^{k_2})}}$ on the smaller small $B_1(2^{-k_2})$.
\end{proof}


   \begin{lemma} 
         \label{lem:int_fct} 
          Let $\psi:\R\to \R$ be a Schwartz function and assume that 
          \begin{equation}
              \label{int_fct_assumption}
              \textup{supp}(\widehat{\psi}) 
\subset  [-1,-2^{-2}]\cup [2^{-2},1] 
          \end{equation}
     Let $\Psi:\R^2\to \R$ be defined as
     \[\Psi(u,v) = \int_1^2 \psi_{(t)}(u) \psi_{(t)}(v) \frac{dt}{t}\]
     and let $\ell\in\Z$.
   Then 
        \begin{enumerate}
    \item  \label{int_fct_symm} For all $(u,v)\in \R^2$,  $\Psi_{(2^\ell)}(u,v) = \Psi_{(2^\ell)}(v,u)$
     \item \label{int_fct_pos}  For all bounded functions $f:\R\to\R$, 
     \begin{equation}
         \int_{\R^2} f(u){f(v)} \Psi_{(2^\ell)}(u,v)\,dudv  
         %= 2^{-k}\int_1^2 \Big| \int_{\R}f(u_0) (\varphi_{4,k})_{(2^{k_j}t)}(u_0) du_0 \Big |^2 \frac{dt}{t} 
         \ge 0 
     \end{equation}
     \item  \label{int_fct_supp}   
     $         \textup{supp}(\widehat{\Psi_{(2^\ell)}}) 
\subset  ([-2^{-\ell},-2^{-3-\ell}]\cup [2^{-3-\ell},2^{-\ell}])^2 \subset \textup{Ann}_1(2^{-\ell},2^3)^2$
     \end{enumerate} 
     \end{lemma}
          \begin{proof}
The symmetry \eqref{int_fct_symm}
 is immediate from the definition. 
 The claim \eqref{int_fct_pos} follows from 
          \[\int_{\R^2} f(u){f(v)} \Psi_\ell(u,v)\,du dv  =  \int_1^2 \Big( \int_{\R}f(u) \psi_{(t)}(u) du \Big )^2 \frac{dt}{t} \ge 0 
          \]
To see \eqref{int_fct_supp}, we use the assumption \eqref{int_fct_assumption} and $1\le t \le 2$. Thus, if  $\widehat{\psi_{(t)}}(\xi) = \widehat{\psi}(t\xi)$ is non-zero, then necessarily $2^{-2}t^{-1}\le |\xi|\le t^{-1} $, so $2^{-3}\le |\xi| \le 1 $ and thus 
\[         \textup{supp}(\widehat{\psi_{(t)}}) 
\subset  [-1,-2^{-3}]\cup [2^{-3},1] \subset \textup{Ann}_1(1,2^3) \]
By rescaling,  %(Lemma \ref{lem:rescaling_supp} with $a=2^{-3},\, b=1, \, r=2^\ell$), 
\[         \textup{supp}(\widehat{\psi_{(2^\ell t)}}) 
\subset  [-2^{-\ell},-2^{-3-\ell}]\cup [2^{-3-\ell},2^{-\ell}] \subset \textup{Ann}_1(2^{-\ell},2^3) \]
Then the bound \eqref{int_fct_supp} 
  follows. 
     \end{proof}

%This is an old version assuming complex values
% \begin{lemma}\label{lem:Phipos}
% Let 
% $\Psi:\mathbb{R}^2\to \mathbb{R}$ be such   that %$\Psi(u_0,u_1)=\Psi(u_1,u_0)$ for all $u=(u_0,u_1)\in \R^2$ and such that 
% for all bounded functions $g:\R\to\C$,
% \begin{equation}\label{nonneg_assumption}
%     \int_{\R^2} g(u_0)\overline{g(u_1)} \Psi(u)\,du \ge 0 
% \end{equation}
% Then for  all $\xi\in \R$, \[\widehat{\Psi}(\xi,-\xi)\geq 0.\]
% \end{lemma}
% \begin{proof}
% We write
%    \[\widehat{\Psi}(\xi,-\xi) = \int \Psi(u) e^{-2\pi i \xi u_0}e^{2\pi i \xi u_1} du \ge 0, \]
%    which is non-negative by the assumption \eqref{nonneg_assumption} applied with  $g(u_0)=e^{-2\pi i \xi u_0}$. 
% \end{proof}


\begin{lemma}\label{lem:Phipos_v2}
Let 
$\Psi:\mathbb{R}^2\to \mathbb{R}$ be such   that $\Psi(u_0,u_1)=\Psi(u_1,u_0)$ for all $(u_0,u_1)\in \R^2$ and such that 
for all bounded functions $g:\R\to\R$,
\begin{equation}\label{nonneg_assumption_v2}
    \int_{\R^2} g(u_0){g(u_1)} \Psi(u)\,du \ge 0 
\end{equation}
Then for all $\xi\in \R$, \[\widehat{\Psi}(\xi,-\xi)\geq 0.\]
\end{lemma}
\begin{proof}
We have
\[\widehat{\Psi}(\xi,-\xi) = \int_{\R^2} \Psi(u) e^{-2\pi i \xi u_0}e^{2\pi i \xi u_1} du =  \int_{\R^2} \Psi(u) \textup{Re}(e^{2\pi i \xi (u_1-u_0)}) du  \]
where we used  $\Psi(u_0,u_1)=\Psi(u_1,u_0)$ for the last equality.
Since  
\[\textup{Re}(e^{2\pi i \xi (u_1-u_0)}) =  \cos(2\pi (u_1-u_0)) = \cos u_0\cos u_1 + \sin u_0\sin u_1,\]
the last display equals
    \[    \int_{\R^2} \Psi(u) \cos u_0\cos u_1 \, du +  \int_{\R^2} \Psi(u) \sin u_0\sin u_1 \, du \ge 0 \]
   and each  integral  is non-negative by the assumption \eqref{nonneg_assumption_v2}.
\end{proof}

%% Re(uv) = Re (x+iy)(a+ib) = xa-yb = Re(u)Re(v) - Im(a)Im(b)
 


\subsubsection{Windows and pairs}\label{sec:windows}
\jr{Clean this up, e.g. should fix $N$ once and for all, see which of these definitions are really needed and minimize}\jr{Also make $C_0$ explicit and name it like any other constant by a subscript reference}
\begin{definition}[$(c,N)$-window]\label{def:cn-window}
Let $c > 0$, $N\in\mathbb{N}$. A {\em $(c,N)$-window} is a Schwartz function $\phi:\R\to\R$ such that
\begin{enumerate}
\item For all $\xi\in \R$, $\widehat{\phi}$ is real and
\[0\le \widehat{\phi}(\xi)\le 1.\]
\item For all $\xi\in \R\setminus [-1,1]$ we have
\begin{equation}\label{ft_window_eq_zero}
\widehat{\phi}(\xi)=0.
\end{equation}
\item For all $\xi\in   [-\frac12,\frac12]$ we have
\begin{equation}\label{ft_window_eq_one}
\widehat{\phi}(\xi)=1.
\end{equation}
\item For all $\xi \in [-1,1]$ and all $m\le N$,
\begin{equation}\label{windowbound}
|\widehat{\phi}^{(m)}(\xi)| \leq c. 
\end{equation}
\end{enumerate}
\end{definition}

\begin{definition}[$(c,N)$-pair]
\label{def:cpair}
Let $c > 0$, $N\in\mathbb{N}$.  A pair $(\phi_0,\phi_1)$ of $(c,N)$-windows is a {\em $(c,N)$-pair} if 
\begin{equation}\label{windowsum}
(\widehat{\phi_0})^2+(1-\widehat{\phi_1})^2=1. 
\end{equation}
\end{definition}
If $(\phi_0,\phi_1)$ is a $(c,N)$-pair, then we write
\[\phi_{i,k}=(\phi_i)_{(2^{k})}\]
for $i\in [2), k\in\mathbb{Z}$

Existence of $(c,N)$-pairs is shown in Lemma \ref{lem:cpair}. 

 \begin{definition}[$N$-universal pair]%\jr{refactor}
 \label{def:unipair}
 Let $N\ge 10$.  %\todo{Placeholder for large number depending on $n$}
    An {\em $N$-universal pair} is a $(C_0,N)$-pair,     where
    \begin{equation}
        \label{def:C0}
        C_0= 1+ 2^{20} \inf c' % The $2^20$ is quite arbitrary
    \end{equation}
  and the infimum is over all $c' > 0$ such that a $(c',N)$-pair exists. Note that $C_0$ depends on $N$ but we do not denote this dependence in notation.
  A function $\phi$ is a {\em window} if $\phi$ is a $(C_0,N)$--window.
\end{definition}

\begin{definition}[Universal pair]
Let 
\[ N_0 = 10^{10^n}. \]
A {\em universal pair} is an $N_0$-universal pair.
\end{definition}

\begin{definition}[Left and right windows]
\label{def:window}
\begin{enumerate}\hfill
   \item   A {\em left window} is a function  $\phi_0:\R\to \R$ such that there exists a function $\phi_1:\R \to \R$ such that $(\phi_0,\phi_1)$ is a universal pair.       \item   A {\em right window} is a function  $\phi_1:\R\to \R$ such that there exists a function $\phi_0:\R \to \R$ such that $(\phi_0,\phi_1)$ is a universal pair.
\end{enumerate}
\end{definition}

\begin{lemma}\label{lem:smoothincr} Let $a,b,r\in\R$ with $a<b$ and  $r>0$. 
There exists a smooth monotone increasing function $\varphi_{a,b,r}:\R\to \R$ with $\varphi_{a,b,r}(x)=0$ for $x\le a$ 
and $\varphi_{a,b,r}(x)=r$ for $x\ge b$. 
\end{lemma}
\begin{proof}
Let $\phi_0(x) = e^{-1/x}\mathbf{1}_{(0,\infty)}(x)$ 
and set
\[\phi(x) = \frac{\phi_0(x)}{\phi_0(x)+\phi_0(1-x)}\]
 Then $\phi$ is smooth since $\phi_0$ is smooth and $\phi_0(x)+\phi_0(1-x)>0$ for all $x$.  Moreover, $\phi(x)=0$ for $x\leq 0$ and  $\phi(x) = 1$ for $x\ge 1$.
The function
\[ \varphi_{a,b,r}(x) =  r\, \phi\Big( \frac{x-a}{b-a}\Big) \]
has the  desired properties. 
\end{proof}

 \begin{lemma}\label{lem:cpair} 
For every $N\in\mathbb{N}$ there exists a $c>0$ and a $(c,N)$-pair as in Definition \ref{def:cpair}. In particular, $C_0<\infty$.
\end{lemma}

\begin{proof}  
In the following, let $\varphi_{a,b,r}$ denote the function from Lemma \ref{lem:smoothincr} for various parameters $a,b,r$.
Let $\psi :[0,\infty) \to \R$ be defined by 
\[\psi(x)=\varphi_{-1,-5/6,1/2}(-x)\mathbf{1}_{[0,\infty)}(x)\]
Then $\psi$ is  a smooth monotone decreasing function  with
\[\psi(x)=1/2 \ \ {\rm for}\ \  x\in [0, 5/6],\] 
\[\psi(x)=0 \ \ {\rm for}\ \  x\in [1, \infty).\]
Let $\rho:[0,\infty) \to \R$ be defined by
\[\rho(x) = \varphi_{1/2, 4/6, \sqrt{3}}(x)\mathbf{1}_{[0,\infty)}(x).\]
Then $\rho$ is  a smooth monotone increasing function with
\[\rho(x)=0 \ \ {\rm for}\ \  x\in [0, 1/2],\] 
\[\rho(x)=\sqrt{3}  \ \ {\rm for}\ \  x\in [4/6, \infty).\]

There exists a smooth even function $\phi_0$ on $\R$
such that its Fourier transform is nonnegative and satisfies on $[0,\infty)$
\begin{equation*}%\label{largex}
(\wh{\phi_0})^2= (4-\rho^2)\psi^2.\end{equation*}
This is because the right-hand side equals $\psi^2$ on $[4/6,\infty)$
and is bounded below by $1/4$ on $[0,5/6]$ and constant one 
on $[0,1/2)$.
There exists a smooth even function $\phi_1$ on $\R$ such that its Fourier transform is nonnegative and fulfills on the interval $[0,\infty)$
\begin{equation*}%\label{smallx}
(1-\wh{\phi_1})^2=1-(4 -\rho^2)\psi^2.\end{equation*}
This is 
because the right-hand side equals $\rho^2/4$ on $[0,5/6]$
and is bounded below by $3/4$ on $[4/6,\infty)$ and constant
on $[0,1/2]$.

The pair $(\phi_0,\phi_1)$ is then   a $c$-pair with
\[c= \max_{0\le m \le N} \max_{j\in [2)} \|{\wh{\phi_j}}^{(m)}\|_\infty <\infty.\]
\end{proof}


\begin{lemma}\label{lem:diagsumone}
Let $(\phi_0,\phi_1)$ be a universal pair.
Let $(k_j)_{j\in [2J)}$ be a strictly increasing sequence of integers. Then
% \begin{equation}\label{eqn:diagsumone}
% (1-\widehat{\phi_{1,k_0}})^2+\sum_{j\in [J)} (\widehat{\phi_{0,k_{2j}}}-\widehat{\phi_{1,k_{2j+1}}})^2+ \sum_{j\in [J-1]} (\widehat{\phi_{0,k_{2j+1}}}-\widehat{\phi_{1,k_{2j+2}}})^2+(\widehat{\phi_{0,k_{2J-1}}})^2
% =1.
% \end{equation}
\begin{equation}\label{eqn:diagsumone}
(1-\widehat{\phi_{1,k_0}})^2+\sum_{j\in [2J-1)} (\widehat{\phi_{0,k_{j}}}-\widehat{\phi_{1,k_{j+1}}})^2+(\widehat{\phi_{0,k_{2J-1}}})^2
=1.
\end{equation}

\end{lemma}

\begin{proof}
Since $(\phi_0, \phi_1)$ is a universal pair,
\begin{equation}\label{phi01-property}
(1-\widehat{{\phi}_{1,k}})^2+(\widehat{{\phi}_{0,k}})^2=1.
\end{equation}
%Observe that $\widehat{\phi_{0,k}}(\xi)\not\in [2)$ implies $\xi\in (2^{-k-1}, 2^{-k})$, and the same holds for $\widehat{\phi_{1,k}}$. 
For every $\xi$, let $k_\xi\in\mathbb{Z}$ be such that $|\xi|\in [2^{-k_\xi-1}, 2^{-k_\xi})$. 
Then for $i=0,1$ we have that $\widehat{\phi_{i,k}}(\xi)=0$ for every $k>k_\xi$ and $\widehat{\phi_{i,k}}(\xi)=1$ for every $k<k_\xi$.
Thus, if $k_\xi$ is not one of the integers $k_j$, then the claim follows (every function occurrence evaluates to $0$ or $1$).
On the other hand, if $k_\xi=k_j$ for some $j\in [J)$, then
the left-hand side is equal to
$(\widehat{\phi_{0,k_\xi}}(\xi))^2 + (1-\widehat{\phi_{1,k_\xi}}(\xi))^2=1$.
\end{proof}


\subsubsection{A smoothing decomposition}
In this section we  decompose the characteristic function $\mathbf{1}_{[0,1)}$ into smoother functions.
Let $\phi_0$ be a window and define  
\begin{equation}\label{eqn:thetadef}
{\theta}={\phi_0}-{(\phi_0)}_{(2)}.
\end{equation}
Define for $k\in\mathbb{Z}$,
\begin{equation}\label{eqn:varphi0kdef}
\varphi_{0,k} = (\mathbf{1}_{[0,1)}\ast\phi_0 - \phi_0)\ast \theta_{(2^{k})}
\end{equation}
\begin{equation}\label{eqn:varphi1kdef}\varphi_{1,k} = \mathbf{1}_{[0,\infty)}\ast\theta_{(2^k)}
\end{equation}
\begin{equation}\label{eqn:varphi2kdef}\varphi_{2,k} = -\mathbf{1}_{[1,\infty)}\ast\theta_{(2^k)}
\end{equation}
\begin{equation}\label{eqn:varphi3def}
\varphi_{3,k} = 2^{k}(\varphi_{0,k})_{(2^{-k})}
\end{equation}
% Set $\gamma=1/2$. \todo{remove $\gamma$ or make arbitrary $<1$} 
Let
\begin{equation}\label{eqn:varphi4kdef}
\varphi_{4,k}(u)=2^k\widetilde{\theta}(u-2^{-k}),
\end{equation}
where $\widetilde{\theta}=\mathbf{1}_{(-\infty,0)}*\theta$ (a primitive of $\theta$).

% In this section let $\chi$ be an even Schwartz function with $\widehat{\chi}(\xi)=0$ for $|\xi|\ge 1$, $\widehat{\chi}(\xi)=1$ for $|\xi|\le \frac12$, and monotone on $[1/2,1]$.
% Define $\theta=\chi-\chi_{(2)}$. \todo{$\chi$ should be a window? Is monotone actually used?}

\begin{lemma}\label{lem:bumpbasic}
Let $\theta$ be as in \eqref{eqn:thetadef}.\\

(i) We have
\[ \mathrm{supp}\;\widehat{\theta} \subset [-1,-2^{-2}]\cup [2^{-2},1] \]

(ii) For every $m\in\mathbb{Z}$
\[ \sum_{\ell=m}^\infty \theta_{(2^\ell)} = (\phi_0)_{(2^m)} \]
and the series converges uniformly.
\end{lemma}

\begin{proof}
(i) This is immediate from the support of windows.
(ii) For  $N\ge m$ and write
\[ S_N := \sum_{\ell=m}^N \theta_{(2^\ell)} =  \sum_{\ell=m}^N (\phi_0)_{(2^\ell)} - (\phi_0)_{(2^{\ell+1})} = (\phi_0)_{(2^m)} - (\phi_0)_{(2^{N+1})} \]
For any $x\in \R$ we  bound
\[|S_N(x)-(\phi_0)_{(2^m)}(x)| = |(\phi_0)_{(2^{N+1})}(x)|  \le 2^{-N-1}\|\phi_0\|_\infty \le 2^{-N} \]
Since $2^{-N}$ converges   to $0$ as $N\to \infty$, the conclusion follows.
\end{proof}

\begin{lemma}\label{lem:chardecomp}
With $\phi_0,\theta$ defined as in \eqref{eqn:thetadef}, the identity
 \[ \mathbf{1}_{[0,1]} = \mathbf{1}_{[0,1]}* \phi_0 +  \sum_{\ell=-\infty}^{-1} \mathbf{1}_{[0,1]} * \theta_{(2^\ell)}  \]   
holds in $L^2$.
\end{lemma}

\begin{rem}
    This identity also holds pointwise a.e., which can be shown using the weak $L^2$ boundedness of the maximally truncated convolution-type singular integrals. However, norm  convergence   suffices for our proof.
\end{rem}


\begin{proof}
For every $N>1$ we have by Plancherel's identity
\begin{equation}
    \label{eq:plan}
    \Big \|\mathbf{1}_{[0,1]} - \mathbf{1}_{[0,1]}*(\phi_0 +   \sum_{\ell=-N}^{-1} \theta_{(2^\ell)} )\Big  \|_2 =\Big \|\widehat{\mathbf{1}_{[0,1]}} - \widehat{\mathbf{1}_{[0,1]}}(\widehat{\phi_0} +   \sum_{\ell=-N}^{-1} \widehat{\theta_{(2^\ell)}} )\Big  \|_2 
\end{equation}
Using the definition of $\theta$ and summing a telescoping sum,  
 \[  \widehat{\mathbf{1}_{[0,1]}}(\xi) - \widehat{\mathbf{1}_{[0,1]}}(\widehat{\phi_0}(\xi) +   \sum_{\ell=-N}^{-1} \widehat{\theta_{(2^\ell)}}(\xi)  ) =   \widehat{\mathbf{1}_{[0,1]}}(\xi)(1-\widehat{(\phi_0)_{2^{-N}}}(\xi) )  \]
 \[=\frac{\sin (\pi \xi)}{\pi \xi}(1 - \widehat{\phi_0}(2^{-N}\xi)) \]
Since  $\widehat{(\phi_0)_{2^{-N}}}(\xi)$ vanishes for  $\xi \in [-2^{N-1},2^{N+1}]$, equals $1$ whenever $|\xi|\geq 2^{N}$, and is bounded in magnitude by $1$,
we have
\[\Big\| \xi\mapsto \frac{\sin (\pi \xi)}{\pi \xi}(1 - \widehat{\chi}(2^{-N}\xi)) \Big\|_2  \leq \|\xi \mapsto \frac{2}{\pi \xi}\mathbf{1}_{[2^{N-1},\infty)}(\xi)\|_2 = 2 \pi^{-1}2^{\frac{-N+1}{2}},\]
where we also used that $\sin$ is bounded by $1$ and we have evaluated the last $L^2$ norm.
Thus, 
\eqref{eq:plan} is dominated by $2 \pi^{-1}2^{\frac{-N+1}{2}}$ and letting $N\to \infty$ shows the  identity   in $L^2$. 
\end{proof}
% \[ \sum_{\ell=m}^{k} \theta_{2^\ell} = \chi_{(2^m)} - \chi_{(2^{k+1})} \]
% \[ \chi_{(2^k)} (x) = 2^{-k} \chi(2^{-k}x)  \]



% \begin{lemma}
% \label{lem:charsplit}
% For each $a\in \R$, the  series  $\sum_{k=-\infty}^{-1} \mathbf{1}_{[a,\infty)}\ast\theta_{(2^k)}$ converges 
% converge pointwise a.e.\todo{obsolete now?}
% \end{lemma}
% \begin{proof}
% Let $\widetilde{\theta}$ be the primitive of $\theta$. 
% We compute
% \[\mathbf{1}_{[a,\infty)}\ast\theta_{(2^k)} = \widetilde{\theta}(2^{-k}(x-a)). \]
% Note that $\widetilde{\theta}$ is a Schwartz function since $\theta$ has mean zero. 
% Thus, there is a constant $C>0$ such that
% \[|\widetilde{\theta}(2^{-k}(x-a))|\le C |x-a|^{-1}2^k\]
% and therefore 
% \[\sum_{k=-\infty}^{-1} |\mathbf{1}_{[a,\infty)}\ast\theta_{(2^k)}| \le C |x-a|^{-1} \sum_{k=-\infty}^{-1}  2^k = C|x-a|^{-1}.\]
% \end{proof}



\begin{lemma}[Smoothing decomposition]\label{lem:smoothingdecomp}
The identity
\begin{equation}\label{eq:seriesexpansion}
\mathbf{1}_{[0,1]} = \phi_0 + \sum_{k=-2}^\infty \varphi_{0,k}  + \sum_{k=-\infty}^{-1} \varphi_{1,k}+ \sum_{k=-\infty}^{-1}\varphi_{2,k}.
\end{equation}
holds in the $L^2$ sense.
\end{lemma}

\begin{proof} 
By Lemma \ref{lem:chardecomp} we have
\begin{equation}
    \label{eq:L2char}
    \mathbf{1}_{[0,1]} =\mathbf{1}_{[0,1]}\ast\phi_0  +  \sum_{k=-\infty}^{-1} \mathbf{1}_{[0,1]}\ast  \theta_{(2^k)}
\end{equation}
where the identity holds in $L^2$. 
Denoting $\varphi=\mathbf{1}_{[0,1]}\ast\phi_0$ and splitting 
\[\mathbf{1}_{[0,1]} = \mathbf{1}_{[0,\infty)} - \mathbf{1}_{[1,\infty)}\]
 we  thus have 
\[\mathbf{1}_{[0,1]} = \varphi +  \sum_{k=-\infty}^{-1} \varphi_{1,k}+ \sum_{k=-\infty}^{-1}\varphi_{2,k},\]
where the sums converge in $L^2$. %a.e. by Lemma \ref{lem:charsplit}. 
It remains to decompose $\varphi$. 
We  write $\varphi= \phi_0 + (\varphi-\phi_0)$
and 
\[
\varphi-\phi_0 = (\varphi - \phi_0)*(\phi_0)_{(2^{-2})},\]
where the last identity holds due to 
\[\widehat{\varphi-\phi_0}(\xi) \widehat{(\phi_0)_{(2^{-2})}}(\xi) = \widehat{\varphi-\phi_0}(\xi) \]
for each $\xi\in \R$. Indeed, the function $\widehat{\varphi-\phi_0}(\xi)=0$ for $|\xi|\ge 1$ and $\widehat{(\phi_0)_{(2^{-2})}}(\xi) = \widehat{\phi_0}(2^{-2}\xi) =1$ for $|\xi|\le 2$.  
By part (ii) of Lemma \ref{lem:bumpbasic} we obtain
\[ (\varphi - \phi_0)*(\phi_0)_{(2^{-2})} =  (\varphi - \phi_0)*\Big( \sum_{k=-2}^\infty \theta_{(2^k)}\Big )  =  \sum_{k=-2}^\infty (\varphi - \phi_0)*\theta_{(2^k)} =   \sum_{k=-2}^\infty \varphi_{0,k}.
\]
The  second equality is justified, for example, by uniform convergence.
\end{proof}


\pd{Revise constants using the new estimates on bumps}

\begin{lemma} \label{lem:theta_decay}
For any $N\ge 1$, 
    \[   |\theta(u)|\le  C_{\ref{lem:theta_decay},N} \langle u\rangle^{N} \]
%     and  for  $(X\theta)(u) = u\theta(u)$,
% \[   |(X\theta)(u)|\le  2^{24}C_0C_{\ref{lem:smoothdecay2}, 23}  (1+|u|)^{-22} \] 
where $C_{\ref{lem:theta_decay},N} = 2^{N+1}C_0C_{\ref{lem:smoothdecay2},N} $
\end{lemma}

\begin{proof}We compute $\widehat{\theta}(\xi) = \widehat{\phi_0}(\xi) - \widehat{\phi_0}(2\xi)$. Thus,   $|\textup{supp}(\widehat{\theta})| =1$ and for any $m\ge 0$, 
\[ \widehat{\theta}^{(m)}(\xi) = \widehat{\phi_0}^{(m)}(\xi) - 2^m\widehat{\phi_0}^{(m)}(2\xi), \]
so
\[ |\widehat{\theta}^{(m)}(\xi)| \le  |\widehat{\phi_0}^{(m)}(\xi)| + 2^m|\widehat{\phi_0}^{(m)}(2\xi)| \le 2^{m+1}C_0, \]
where we used that $\phi_0$ is a window. Applying 
Lemma \ref{lem:smoothdecay2}  with  $\phi = (2^{m+1}C_0)^{-1} \theta $ $N=m$  and  implies the claim. 
\end{proof}

% \subsubsection{Properties of $\varphi_{3,k}$}

\begin{lemma}\label{lem:ft_phi3_eq}
For $k\in\mathbb{Z}$ and $\xi\in\mathbb{R}$,
\[ \widehat{\varphi_{3,k}}(\xi) = 2^k(\widehat{\mathbf{1}_{[0,1)}}-1)(2^{-k}\xi)\widehat{\chi}(2^{-k}\xi)\widehat{\theta}(\xi) \]
\end{lemma}

\begin{proof}
By definition, %$\varphi_{0,k} = (\mathbf{1}_{[0,1)} * \chi - \chi)*\theta_{(2^k)}$ and thus
\[\varphi_{3,k} = 2^k (\mathbf{1}_{[0,1)} * \chi - \chi)_{(2^{-k}) }*\theta.\]
By the convolution theorem,
\[\widehat{\varphi_{3,k}}(\xi) = 2^k(\widehat{\mathbf{1}_{[0,1)}}-1)(2^{-k}\xi)\widehat{\chi}(2^{-k}\xi)\widehat{\theta}(\xi).\]
\end{proof}

\begin{lemma}\label{lem:abs_deriv_ft_phi3_le}
For every $m\ge 0$, $k\in\mathbb{Z}$, $|\xi|\le 1$,
\[ |\widehat{\varphi_{3,k}}^{(m)}(\xi)| \le C_{\ref{lem:abs_deriv_ft_phi3_le}, m}, \]
where 
$C_{\ref{lem:abs_deriv_ft_phi3_le}, m}=C_0^2 2^{6 m+m^2}$.
\end{lemma}

\begin{proof}
By Lemma \ref{lem:ft_phi3_eq},
\[ \widehat{\varphi_{3,k}}(\xi) = 2^k(\widehat{\mathbf{1}_{[0,1)}}-1)(2^{-k}\xi)\widehat{\chi}(2^{-k}\xi)\widehat{\theta}(\xi). \]
By the higher order product rule (two applications of the ordinary product rule), 
\[\widehat{\varphi_{3,k}}^{(m)}(\xi) = 2^k \sum_{\ell=0}^m {\binom{m}{\ell}} \widehat{\theta}^{(m-\ell)}(\xi) ((\widehat{\mathbf{1}_{[0,1)}}-1)(2^{-k}\xi)\widehat{\chi}(2^{-k}\xi))^{(\ell)}  \]
\[= 2^k \sum_{\ell=0}^m {\binom{m}{\ell}} \widehat{\theta}^{(m-\ell)}(\xi)\sum_{\ell'=0}^\ell {\binom{\ell}{\ell'}} (\widehat{\mathbf{1}_{[0,1)}}-1)(2^{-k}\xi)^{(\ell')} (\widehat{\chi}(2^{-k}\xi))^{(\ell-\ell')}  \]
\begin{equation}
    \label{phi3hat_der}
    =2^k \sum_{\ell=0}^m {\binom{m}{\ell}} 2^{-k\ell} \widehat{\theta}^{(m-\ell)}(\xi)\sum_{\ell'=0}^\ell {\binom{\ell}{\ell'}} (\widehat{\mathbf{1}_{[0,1)}}-1)^{(\ell')}(2^{-k}\xi) \widehat{\chi}^{(\ell - \ell')}(2^{-k}\xi) 
\end{equation}
If $\ell=0$, we use 
  Taylor's theorem. Namely, for  $\xi\ge 0$ and some   $\alpha\in [0,\xi]$ we have 
\[\sin \xi = \xi - \frac{\sin \alpha}{2}\xi^2, \]
and thus \[(\widehat{\mathbf{1}_{[0,1)}}-1)(2^{-k}\xi) = \frac{\sin (\pi 2^{-k} \xi)}{\pi 2^{-k} \xi} - 1 = - \frac{\sin \alpha}{2}(\pi 2^{-k}\xi) \]
Therefore,  for each $|\xi|\le 1$, 
\[|(\widehat{\mathbf{1}_{[0,1)}}-1)(2^{-k}\xi)| \le \pi 2^{-k-1} \]
If $\ell\neq 0$, we compute further
\[(\widehat{\mathbf{1}_{[0,1)}}-1)^{(\ell')}(\xi) = (\sin (\pi \xi) (\pi\xi)^{-1})^{(\ell')} \]
\[=\sum_{\ell''=0}^{\ell'} {\binom{\ell'}{\ell''}} \pi^{\ell'-\ell''} \sin^{(\ell'-\ell'')}(\pi\xi) (-1)^{\ell''} (\ell'')! \pi^{-\ell''-1} \xi^{-\ell''-1}\] 
Using that if $\xi$ is  in the support of $\widehat{\varphi_{3,k}}$ implies  then  $|\xi|>2^{-1}$,  we have  for  $\tfrac{1}{2}\le |\xi|\le 1$, 
\[|(\widehat{\mathbf{1}_{[0,1)}}-1)^{(\ell')}(\xi)|  \le \pi^{\ell'} (\ell')! 2^{2\ell'} \le \pi^{m} m^m 2^{2m} \le 2^{4m+m^2} \]
Using this and  $ |\widehat{\chi}^{(\ell-\ell')}(\xi)|\le C_0$, we estimate \eqref{phi3hat_der} and thus $|\widehat{\varphi_{3,k}}(\xi)|$ as  
\[|\widehat{\varphi_{3,k}}(\xi)| \le  2^k  C_0\Big( \pi 2^{-k-1} + \sum_{\ell=1}^m {\binom{m}{\ell}} 2^{-k\ell} \widehat{\theta}^{(m-\ell)}(\xi) \sum_{\ell'=0}^\ell {\binom{\ell}{\ell'}}  \pi^{m} m^m 2^{2m} \Big) \]
Using also $ |\widehat{\theta}^{(m-\ell)}(\xi)|\le C_0$, we obtain 
\begin{equation}
    \label{phi3hat_der_bound}
 |\widehat{\varphi_{3,k}}(\xi)| \le    C_0^2 2^{6m+m^2}
\end{equation}
\end{proof}

\begin{lemma}\label{lem:abs_deriv_ft_Tphi3_le}
For every $m\ge 0$, $k\in\mathbb{Z}$, $|\xi|\le 1$,
\[ |\widehat{T\varphi_{3,k}}^{(m)}(\xi)| \le C_{\ref{lem:abs_deriv_ft_Tphi3_le}, m}, \]
where $C_{\ref{lem:abs_deriv_ft_Tphi3_le}, m}=m C_{\ref{lem:abs_deriv_ft_phi3_le}, m} + C_{\ref{lem:abs_deriv_ft_phi3_le}, m+1}$
\end{lemma}

\begin{proof}
By   Lemma \ref{lem:ft_deriv_mul},
\[\widehat{T\varphi_{3,k}}^{(m)}(\xi) = -(m \widehat{\varphi_{3,k}}^{(m)}(\xi) + \xi \widehat{\varphi_{3,k}}^{(m+1)}(\xi)), \]
so the claim follows by Lemma \ref{lem:abs_deriv_ft_phi3_le}.
\end{proof}

% \subsubsection{Properties of $\varphi_{4,k}$}

\begin{lemma}
\label{lem:theta_prim}
Let $\widetilde{\theta}$ be a primitive of $\theta$, i.e. $\widetilde{\theta}'=\theta$.  Then 
\begin{equation}
    \label{thetaprim_supp}
    \textup{supp} (\widehat{\widetilde{\theta}}) \subset [-1,-2^{-2}]\cup [2^{-2},1] \subset \textup{Ann}_1(1,2^2),
\end{equation}
\begin{equation}
    \label{Tthetaprim_supp}
    \textup{supp} (\widehat{T\widetilde{\theta}}) \subset [-1,-2^{-2}]\cup [2^{-2},1] \subset \textup{Ann}_1(1,2^2)
\end{equation}
and for $N\ge 1$, 
\begin{equation}
    \label{thetaprim_decay}
    |\widetilde{\theta}(u)|\le  C_{\ref{lem:theta_prim},N}  \langle u\rangle^{N}
\end{equation}
\begin{equation}
    \label{Thetaprim_decay}
    |T\widetilde{\theta}(u)|\le  C_{\ref{lem:theta_prim},N} \langle u\rangle^{N}
\end{equation}
where $C_{\ref{lem:theta_prim},N} = 4 C_0 N^{3N}   C_{\ref{lem:smoothdecay2},N}$
\end{lemma}

\begin{proof}
Since $\widehat{\theta}$ is supported in $[-1,-2^{-2}]\cup [2^{-2},1]\subset \textup{Ann}_1(1,2^2)$,    the same holds for 
\[\widehat{\widetilde{\theta}}(\xi) = (2\pi i \xi)^{-1} \widehat{\theta}(\xi)\]   Thus, \eqref{thetaprim_supp} holds.  Since $\widehat{T\widetilde{\theta}}(\xi) = -\xi \widehat{\widetilde{\theta}}(\xi)$, also \eqref{Tthetaprim_supp} holds. 
 

To see \eqref{thetaprim_decay}, 
we use the product rule to estimate for each $m\ge 0$
\[\widehat{\widetilde{\theta}}^{(m)}(\xi) = (2\pi i)^{-1} \sum_{\ell=0}^m {\binom{m}{\ell}}  (-1)^{\ell} \ell! \xi^{-\ell-1} \widehat{\theta}^{(m-\ell)}(\xi) \]
If $\xi$ is in the support of $\widehat{\theta}$, then $|\xi|^{-\ell -1} \le 2^{-2(\ell+1)}$. Also, $\widehat{\theta}(\xi) = \widehat{\phi_0}(\xi) - \widehat{\phi_0}(2\xi)$. 
This gives for each $\xi$  in the support of $\widehat{\theta}$, 
\[|\widehat{\widetilde{\theta}}^{(m)}(\xi)| \le (2\pi)^{-1} \sum_{\ell=0}^m {\binom{m}{\ell} }  (-1)^{\ell} \ell! 2^{-2(\ell+1)} (1+2^{m-\ell}) |\widehat{\phi_0}^{(m-\ell)}(\xi)|  \] 
Since $\phi_0$ is a window, using   Definition \ref{def:window} we bound the last display as
\begin{equation}\label{boundschwartz}
\le (2\pi)^{-1} C_0 \sum_{\ell=0}^m {\binom{m}{\ell}}  (-1)^{\ell} \ell! 2^{-2(\ell+1)} (1+2^{m-\ell})
\end{equation}
We crudely estimate $(2\pi)^{-1}\le 1$. 
$2^{-2(\ell+1)}\le 1$, $1+2^{m-\ell}\le 2^{m+1}$, $\ell!\le m!\le m^m$. 
Therefore, the display \eqref{boundschwartz} can be estimated as 
\[\le C_0   2^{m+1}m^{m} \sum_{\ell=0}^m {\binom{m}{\ell}} = C_02^{2m+1}m^{m} \le 2 C_0 m^{3m} =: c_m
\]
We also compute
\[\widehat{T\widetilde{\theta}}^{(m)}(\xi) = - (\xi \widehat{\widetilde{\theta}'}(\xi))^{(m)} = -\sum_{\ell=0}^m {\binom{m}{\ell}} \xi^{m-\ell} \widehat{\widetilde{\theta}}^{(m+1)}(\xi) \]
and thus, using $|\xi|\le 1$, 
\[|\widehat{T\widetilde{\theta}}^{(m)}(\xi)| \le  |\widehat{\widetilde{\theta}}^{(m+1)}(\xi)| \sum_{\ell=0}^m {\binom{m}{\ell}}  = 2 |\widehat{\widetilde{\theta}}^{(m+1)}(\xi)| \le 2^mC_0 \le c_m  \]
If $m\in \{0,N\}$, $c_m\le C_N$. 
Lemma \ref{lem:smoothdecay2} applied with $c_N^{-1}\widetilde{\theta}$,  together with    \[|\supp(\widehat{\widetilde{\theta}})|=|\supp(\widehat{T\widetilde{\theta}})|\le 2\] now give
\[|\tilde{\theta}(u)|   \le 2 c_N C_{\ref{lem:smoothdecay2},N} \langle u\rangle^{N} \]
and likewise 
\[|T\tilde{\theta}(u)|   \le 2 c_N C_{\ref{lem:smoothdecay2},N} \langle u\rangle^{N} \]
\end{proof}

\begin{lemma}\label{lem:phi4_supp} We have 
(a) $\textup{supp}(\widehat{\varphi_{4,k}}) 
\subset  [-1,-2^{-2}]\cup [2^{-2},1]$\\

(b) $\textup{supp}(\widehat{T\varphi_{4,k}}) 
\subset  [-1,-2^{-2}]\cup [2^{-2},1]$
\end{lemma}
\begin{proof}
(a) By Lemma \ref{lem:theta_prim}, 
\[\textup{supp}(\widehat{\widetilde{\theta}}) \subset [-1,-2^{-2}]\cup [2^{-2},1]\]  
The claim follows since
\[\widehat{\varphi_{4,k}}(\xi)= 2^k e^{2\pi i 2^{-k}}\widehat{\widetilde{\theta}}(\xi).  \]

(b) By (a), $\widehat{\varphi_{4,k}}$ is supported in 
$ [-1,-2^{-2}]\cup [2^{-2},1] $.  
  By Lemma \ref{lem:ft_deriv_mul} applied with $m=1$, 
  \[\widehat{T\varphi_{4,k}}(\xi) = -\xi(\widehat{\varphi_{4,k}})'(\xi) \]
The supports of $\widehat{\varphi_{4,k}}$ and $(\widehat{\varphi_{4,k}})'$ coincide, and  thus the last display is non-zero only if  $|\xi|\le 1$. The claim follows.
\end{proof}

\subsubsection{Miscellany}
Let $\mathbf{f}=(f_0,\dots,f_{n-1})$ be Schwartz functions.


\begin{lemma}[Rescaling]\label{lem:rescaling}Let $\phi\in L^1(\mathbb{R})$.
For all $r,q,\lambda>0$,
\[ \|A_t({\phi_{(\lambda)}})\|_{V_{r,J}(t\in \mathbb{R}; L^q)} =  \|A_t({\phi})\|_{V_{r,J}(t\in \mathbb{R}; L^q)}. \]
\end{lemma}

\begin{proof}
By replacing $t_j \lambda\mapsto t_j$ in the supremum defining the variation seminorm.
\end{proof}


\begin{lemma}\label{lem:norm_A_sum_le_sum}
If $\phi=\sum_{j=1}^\infty \phi_j$ holds in $L^2$ (with $L^2$ functions $\phi, \phi_j$), then
\[ \|A(\phi, \mathbf{f})\|_{L^2(\mathbb{R}^n)} \le \sum_{j=1}^\infty \|A({\phi_j}, \mathbf{f})\|_{L^2(\mathbb{R}^n)}. \]
\end{lemma}

\begin{proof}
Let $E_N=\phi - \sum_{n=1}^N \phi_n$. By assumption, $\|E_N\|_2 \to 0$ as $N\to\infty$. By the triangle inequality, 
\[\|A(\phi, \mathbf{f})\|_{L^2(\mathbb{R}^n)} \le \sum_{n=1}^N\|A({\phi_n}, \mathbf{f})\|_{L^2(\mathbb{R}^n)} + \|A({E_N}, \mathbf{f})\|_{L^2(\mathbb{R}^n)} \]
We apply the Cauchy-Schwarz inequality in $s$ to estimate
\[\|A({E_N}, \mathbf{f})\|_{L^2(\mathbb{R}^n)}  = \Big ( \int_{\R^n}\Big |\int_{\R} \Big(\prod_{i\in [n)} f_i(x+se_i)\Big) E_N(s)\,ds\Big |^2 dx  \Big)^{1/2}\]
\[\le \Big ( \int_{\R^n}\Big (\int_{\R} \Big|\prod_{i\in [n)} f_i(x+se_i)\Big|^2 ds \Big) \Big( \int_{\R}|E_N(s)|^2\,ds\Big )\, dx  \Big)^{1/2} \]
\begin{equation}
    \label{en_L2}
    = \Big( \int_{\R^{n+1}}\Big|\prod_{i\in [n)} f_i(x+se_i)\Big|^2\, d(x,s) \Big)^{1/2} \|E_N\|_{L^2(\R)}
\end{equation}
We claim that this is 
\[\le c_1\|E_N\|_{L^2},\] where $c_1$ is a constant that depends on $f_i$. 
Letting  $N\to \infty$ then finishes the proof. 



To see the claim,  we write $x=(x_0,x')\in \R\times \R^{n-1}$. 
Since the functions $f_i$ are bounded with rapid decay, there exists a constant $C>0$ such that $|f_i|\le C$ for $i\ge 2$, and   
\[|f_1(x+se_1)|\le C \langle x_0\rangle^{10},\]
\[|f_0(x+se_0)| \le C \langle x_0+s\rangle^{5}\langle x'\rangle^{5n} . \]
Using  Lemma \ref{lem:transl_pow_v} (with, say, $N=5$), the last display is 
\[\le 2^5 C \langle x_0\rangle^{-5}\langle s\rangle^{5}\langle x'\rangle^{5n} \]
Thus, we can estimate
\[ \int_{\R^{n+1}}\Big|\prod_{i\in [n)} f_i(x+se_i)\Big|^2 d(x,s)  \le   2^5 C^{2n} \int_{\R^{n+1}}\langle s\rangle^{5} \langle x_0\rangle^{5}\langle x'\rangle^{5n} d(x,s), \]
which is a finite positive constant. 
\end{proof}

\begin{lemma}    \label{lem:form_pos} 
(i) Let $\Psi\in\mathcal{S}'(\mathbb{R}^2)$ be a real-valued distribution such that
\begin{equation}\label{eqn:hyp_form_pos}
\Psi(g^{\otimes 2}) \ge 0
\end{equation}
for all real-valued $g\in \mathcal{S}(\mathbb{R})$. Then for any $\mathbf{F}\in\mathfrak{F}$,  
\[ \Lambda(\Psi)(\F)\ge 0. \]

(ii) In particular, for all real-valued $\psi\in\mathcal{S}'(\mathbb{R})$,
\begin{equation}\label{form_pos_psipsi}
     \Lambda(\psi^{\otimes 2})(\F)\ge 0.
\end{equation}
\end{lemma}


\begin{proof}
(i) By definition, 
%\[ \]
\[\Lambda(\Psi)(\mathbf{F})  =\Psi(\psi_\Lambda(\mathbf{F}))=\Psi\Big(y\mapsto\int_{\R^n} \prod_{h\in [2), i\in [n)} F_{i}(y^h-\Sigma(x), x_{[n)\setminus i})\,dx\Big).  \]

 

Since the integrand is a Schwartz function in $y, x$ (see Proposition \ref{prop:Fx}), we may pull out the integral \cite{Hormander1} so that this equals
\[\int_{\R^n} \Psi((g_x)^{\otimes 2})\, dx \ge 0,\]
where
$g_x(y^0)= \prod_{i\in [n)} F_i(y^0-\Sigma(x), x_{[n)\setminus i})$ for every $(x,y^0)\in\mathbb{R}^n\times\mathbb{R}$.

(ii) This follows by (i) since
\[ (\psi^{\otimes 2})(g^{\otimes 2}) = (\psi(g))^2\ge 0 \]
(since $\psi$ is real-valued).
\end{proof}
 





\begin{lemma} \label{lem:ftc_ATphi}  
Let $x\in \R^2$, $\phi$ a Schwartz function,  $t\in [2^k, 2^{k+1}]$, $k\in \Z$. Let $a(t)=A_t(\phi)(x)$. Then 
\[ \|ta'(t)\|_{L^2(t\in [2^k, 2^{k+1}],\frac{dt}{t})}^2 = \int_{1}^2 |A_{2^k t}(T\phi)(x)|^2\,\frac{dt}{t}. \]
\end{lemma}

\begin{proof} 
Note that
\begin{equation}\label{shortlongftc_reduction_pf3}
t\frac{\partial}{\partial t}(\phi_{(t)}(s)) = -(T\phi)_{(t)}(s),
\end{equation}
where $T\phi(s) = (s \phi(s))'$.
    By this and \eqref{A_def}, 
\[ ta'(t) = t\frac{\partial}{\partial t}\int_{\mathbb{R}} \Big(\prod_{i\in [n)} f_i(x+se_i)\Big)\phi_{(t)}(s)\,ds = -A_t(T\phi)(x). \]
Integrating in $t$ over $[2^k,2^{k+1}]$ and changing variables $t \mapsto 2^k t$ in the $t$-integration on the right-hand side,
the claim follows.
\end{proof}
  




\begin{lemma}
\label{lem:Phij_prop} 
For every integer $\kappa \ge 0$ the following holds.  Let 
    $\Psi:\mathbb{R}^2\to \mathbb{R}$ be a non-zero test function such that for all $u=(u_0,u_1)\in\R^2$, 
\begin{equation}
    \Psi(u_0,u_1) = \Psi(u_1,u_0).
\end{equation} 
Assume also that  
%there is a strictly increasing sequence of integers $(k_j)_{j=1}^J$ such that
\begin{equation}
    \mathrm{supp}(\widehat{\Psi}) \subset \mathrm{Ann}_1(1, 2^{20})^2
\end{equation}
and
for $a\in [3)$ and $\xi \in \R$,
\begin{equation}\label{Phij_psi_der}
    \Big|\tfrac{d^a}{d\xi^{a}} \widehat{\Psi}(\xi,-\xi)\Big| \leq 1.
\end{equation}
% \begin{equation}\label{Phij_psi_decay}
%      |\Psi(u)| \le 2^{-\frac12 {\kappa}} \langle u_0+u_1\rangle_{(2^\kappa)}^{N+2} \langle u_0-u_1\rangle^{N+2} 
% \end{equation}
% for all $u\in\mathbb{R}^2$.
Let $(\phi_0,\phi_1)$ be a universal pair and denote \[\phi_{i,\kappa}=(\phi_i)_{(2^{\kappa})}\]
for $i\in [2), \kappa\in\mathbb{Z}$.
There exists a real even smooth function $\psi:\R\to \R$ such that
\begin{enumerate} 
\item For all $\xi\in\R$,   
    \begin{equation}
    \label{Phij_psi_diag}
        %\widehat{\psi}(\xi)^2 = 2c(\widehat{\phi_{0}}-\widehat{\phi_{1}})(\xi)^2 - \widehat{\Psi}(\xi,-\xi) 
           \widehat{\psi}(\xi)^2 = 2 (\widehat{\phi_{0,2^{-25}}}-\widehat{\phi_{1,2^{25}}})(\xi)^2 - \widehat{\Psi}(\xi,-\xi) 
    \end{equation}
    \item $\widehat{\psi}$ is supported in $\textup{Ann}_1(1,2^{27})$.
    \item  For all $a\in [3)$ and $\xi\in \R$,  
  \begin{equation}
  \label{psi_decay} 
      |\widehat{\psi}^{(a)}(\xi)|\le C_{\ref{lem:Phij_prop}}.
  \end{equation}
    where   $C_{\ref{lem:Phij_prop}}=2^{56}C_0^4$.
    %2^{2^{N+2}}(2^{27N+2} C_0^2)^{2N-1}$
    %C_0^{2} 2^{2^{12}}$.  
    
\end{enumerate}
\end{lemma}
\begin{proof}
% First note that \eqref{Phij_psi_decay} implies
% $\|\Psi\|_1\le 1$. 
% Indeed, we compute   
% \begin{equation}
%     \label{int_eval}
%     \int_{\R} \langle x\rangle^{N+2}dx = \tfrac{2}{N+1},
% \end{equation}
% which gives,  
% \[\|\Psi\|_1 \le 2^{-\frac32 {k}} \int_{\R^2} \langle 2^{-k}(u_0+u_1)\rangle^{N+2} \langle u_0-u_1\rangle^{N+2} du\] 
% \[=2^{-\frac32 {k}} \int_{\R^2}  \langle 2^{-k}u_0\rangle^{N+2} \langle u_0-2u_1\rangle^{N+2}du \]
% \[\le 2^{-\frac32 {k}} \int_{\R}  \langle 2^{-k}u_0\rangle^{N+2}  du_0 = 2^{-\frac32 {k}} \tfrac{2}{(N+1)2^k} \le 1 . \]

       The function $(\widehat{\phi_{0,2^{-25}}}-\widehat{\phi_{1,2^{25}}})^{\otimes 2}$ is supported in 
       $\textup{Ann}_1(1,2^{26})^2$ and constant $1$ on $\textup{Ann}_1(1,2^{24})^2$. 
       Since $|\widehat{\Psi}(\xi,-\xi)| \le 1$
       and   $\widehat{\Psi}$ is supported in $\textup{Ann}_1(1,2^{20})^2$,
       \begin{equation}\label{eqn:lemPhij_proppf}
       |\widehat{\Psi}(\xi,-\xi)|\le (\widehat{\phi_{0,2^{-25}}}-\widehat{\phi_{1,2^{25}}})^{\otimes 2}(\xi,-\xi)  .
       \end{equation}
       Define 
       \begin{equation}
       \label{def_sqrtpsi}
           \widehat{\psi}(\xi) = \Big( 2 (\widehat{\phi_{0,2^{-25}}}-\widehat{\phi_{1,2^{25}}})(\xi)^2 - \widehat{\Psi}(\xi,-\xi)\Big)^{1/2},
       \end{equation}
        This is indeed a smooth function: 
        first, the function under the square root is always non-negative.
        On the complement of $\mathrm{Ann}_1(1, 2^{21})$, $\widehat{\Psi}(\xi,-\xi)=0$, so $\widehat{\psi}(\xi)$ equals 
        \[ 2^{1/2} (\widehat{\phi_{0,2^{-25}}}-\widehat{\phi_{1,2^{25}}})(\xi),\]
        which is a smooth function. 
        Moreover, by \eqref{eqn:lemPhij_proppf} one has $\widehat{\psi}(\xi)^2 \ge 1$
        on the annulus $\mathrm{Ann}_1(1,2^{24})$, so $\psi$ is smooth also on this annulus, which properly contains $\mathrm{Ann}_1(1,2^{21})$ (formally, the interior of $\mathrm{Ann}_1(1,2^{24})$ and the complement of $\mathrm{Ann}_1(1,2^{21})$ form an open cover of $\mathbb{R}$ with two open sets on each of which $\psi$ is smooth, so $\psi$ is smooth on $\R$).
        \begin{enumerate}[leftmargin=*]
        \item This holds by definition.
       \item This follows since $\xi\mapsto \widehat{\Psi}(\xi,-\xi)$ is supported in 
       $\textup{Ann}_1(1,2^{21})$ and $\xi \mapsto (\widehat{\phi_{0,2^{-25}}}-\widehat{\phi_{1,2^{25}}})(\xi)^2$ is supported in $\textup{Ann}_1(1,2^{26})$. 
       \item   
Denote 
\[\vartheta(\xi)=(\widehat{\phi_{0,2^{-25}}}-\widehat{\phi_{1,2^{25}}})(\xi) =  \widehat{\phi_{0}}(2^{-25}\xi)-\widehat{\phi_{1}} (2^{25}\xi)\]
so that 
\begin{equation}
    \label{vartheta_derbd}
    |\vartheta^{(a)}(\xi)| = |2^{-25a}\widehat{\phi_{0}}^{(a)}(2^{-25}\xi) -2^{25a}\widehat{\phi_{1}}^{(a)} (2^{25}\xi)|\le 2^{26 a}C_0
\end{equation}

      If  $\xi\not \in  \textup{Ann}_1(1,2^{21})$, then by \eqref{def_sqrtpsi}, 
      \[\widehat{\psi}(\xi) =2^{1/2} \vartheta(\xi)\]
      and thus for $a\in [3)$, 
            \[|\widehat{\psi}^{(a)}(\xi)| \le  2^{52}C_0 \le C_{\ref{lem:Phij_prop}}. \]  
            
     Let now  $\xi\in  \textup{Ann}_1(1,2^{24})$.  
   We compute  
        \begin{equation}
            \label{varthsq_derbd}
        \Big| \tfrac{d^a}{d\xi^a}(\vartheta(\xi)^2)\Big| = \Big | \sum_{l=0}^a {\binom{a}{l}} \vartheta^{(l)}(\xi) \vartheta^{(a-l)}(\xi) \Big |    \le 
         \sum_{l=0}^a {\binom{a}{l}} 2^{26a}C_0^2 = 2^{27a}  C_0^2  
        \end{equation}
          where we used the bound \eqref{vartheta_derbd}. Let now $f(\xi) =   2\vartheta(\xi)^2-\widehat{\Psi}(\xi,-\xi)$. Then $1\le |f(\xi)|\le 4C_0^2 + 1$, $|f'(\xi)|\le 2^{28}C_0^2+1$, $|f''(\xi)|\le 2^{55}C_0^2+1$        
Thus, with $\widehat{\psi} = f^{1/2}$, 
\[|\widehat{\psi}(\xi)| \le 2C_0+1\]
 \[|\widehat{\psi}'(\xi)| = \Big| \frac{f'(\xi)}{2(f(\xi))^{1/2}}\Big| \le 2^{28}C_0^2\]
 \[|\widehat{\psi}''(\xi)| = \Big|\frac{f''(\xi)}{2(f(\xi))^{1/2}} - \frac{(f'(\xi))^2}{4(f(\xi))^{3/2}}\Big| \le  2^{56}C_0^4\]
%        First we claim that for   $a\in [3)$, 
%        \begin{equation}
%            \label{Psi_derbd}
%            |\widehat{\Psi}^{(a)}(\xi,-\xi)| \le 2^{3N}
%        \end{equation}
%        Indeed, we compute 
%        \[|\widehat{\Psi}^{(a)}(\xi,-\xi)| = \Big | (-2\pi i)^a \int_{\R^2} \Psi(u)(u_0-u_1)^ae^{-2\pi i \xi(u_0-u_1)} du \Big| \]
%        \[\le (2\pi)^a\int_{\R^2} | \Psi(u)(u_0-u_1)^a| du\]
%        Using \eqref{Phij_psi_decay},
%        %and estimating $(1+2^{-k}|u_0+u_1|)^{-10}\le (1+2^{-k}|u_0+u_1|)^{-2}$, 
%        this is
%        \[\le (2\pi)^a 2^{-\tfrac{3}{2} {k}} \int_{\R^2} \langle 2^{-k}(u_0+u_1)\rangle^{N+2} \langle u_0-u_1\rangle^{2}du  \]
%        \begin{equation}
%            \label{int_est}= (2\pi)^a 2^{-\tfrac{3}{2} {k}} \int_{\R^2} \langle 2^{-k}u_0\rangle^{N+2} \langle u_0-2u_1\rangle^{2}du 
%        \end{equation}
%          We have 
%          \[\int_{\R}\langle u_0-2u_1\rangle^{2} du_1  =\frac12 \int_{\R}\langle u_1\rangle^{2} du_1 =  1,   \]
%          where we used that the last integral evaluates to $2$.
%         Therefore, integrating first in $u_1$, \eqref{int_est} equals 
%          \[=  (2\pi)^a 2^{-\tfrac{3}{2} {k}} \int_{\R} \langle 2^{-k}u_0\rangle^{N+2}du_0  \le (2\pi)^a \le 2^{3N}\]
% %         \[=  (2\pi)^a 2^{-\tfrac{1}{2} {k}} \int_{\R} (1+|u_0|)^{-10}du_0  =  (2\pi)^a 2^{-\tfrac{1}{2} {k}} \le   (2\pi)^a \le 2^{26}\]
%           where we used \eqref{int_eval},  $k\ge 0$,  and $a\le N$. 
         
   
%          The estimates \eqref{Psi_derbd} and \eqref{varthsq_derbd} give  
%          \[\Big| \frac{d^a}{d\xi^a}(\widehat{\psi}(\xi)^2)\Big|   \le 2^{27a+1} C_0^2 + 2^{3N}\le B \]
%          with $B= 2^{27N+2} C_0^2$. 
%          We now apply Lemma \ref{derivative of sqrt}
%          for each $\xi \in \textup{Ann}_1(1,2^{24})$
%          with $\phi=B^{-1/2}\widehat{\psi}$, $\lambda=1$, $m=a$,  $c_0=B^{-1/2}$  (since $|\widehat{\psi}(\xi)|\ge 1$ on $\textup{Ann}_1(1,2^{24})$). We obtain (using $a\le N$)
%          \[\Big| \frac{d^a}{d\xi^a}\widehat{\psi}(\xi)\Big| \le C_{\ref{derivative of sqrt},a,B^{-1/2}}\le 2^{2^{N+1}} B^{N} =2^{2^{N+1}}(2^{27N+2} C_0^2)^N \]
         \end{enumerate}
         \end{proof}

 





\begin{lemma}[Bootstrapping]
\label{lem:bootstrap}
Let $\phi:\R\to \R$ be a Schwartz function.
Then 
    \[\|A_{t}(\phi)\|^2_{V_{2,J}(t\in 2^{\Z}; L^2)} \le  4 \cdot \sup \sum_{\ell\in [J)} \|A_{2^{k_{\ell}}}(\phi)\|^2_2     \]
    where the supremum is over all sequences of integers $k_0<\cdots <k_{J-1}$. 
\end{lemma}
{\em Note:} This lemma will only be applied to mean zero bump functions. 

\begin{proof}
Let $k_0<\cdots <k_{J}$ be integers. For $\ell \in [J)$ we write 
\begin{equation}
    \label{expandrhs}
    \|A_{2^{k_{\ell+1}}}(\phi)-A_{2^{k_{\ell}}}(\phi)\|_2^2 = \|A(\phi_{(2^{k_\ell+1})}-\phi_{(2^{k_\ell})})\|_2^2
    \end{equation} 
 \[= \int_{\R^n} A(\phi_{(2^{k_{\ell+1}})}-\phi_{(2^{k_{\ell}})})(x)A(\phi_{(2^{k_{\ell+1}})}-\phi_{(2^{k_{\ell}})})(x)  dx \]
%    \[  = \int_{\R^n} \int_{\mathbb{R}^2} \Big(\prod_{(i,j)\in S_{n-1}} f_i(x+u_j e_i) \Big) K(u)\, du\, dx,\]
    where  $k_0<\ldots < k_J$ is any increasing sequence of integers.  
    %   This can be recognized as $\Lambda_{D_{n-1}, K}(\tau\mathbf{f})$ with 
    % \[ K(u) = (\phi_{(2^{k_{\ell+1}})}-\phi_{(2^{k_{\ell}})})(u_0)(\phi_{(2^{k_{\ell+1}})}-\phi_{(2^{k_{\ell}})})(u_1). \]
% We use the distributive law to write  
% \[K=K_1-K_2\]
% where
% \[K_1 = \phi_{(2^{k_{\ell+1}})}(u_0)(\phi_{(2^{k_{\ell+1}})}-\phi_{(2^{k_{\ell}})})(u_1)  \]
% \[K_2 = \phi_{(2^{k_{\ell}})}(u_0)(\phi_{(2^{k_{\ell+1}})}-\phi_{(2^{k_{\ell}})})(u_1)  \]
% We split the form accordingly, i.e.,    write \eqref{expandrhs} as 
By the distributive law, 
  \[\sum_{\ell\in [J)}  \|A_{2^{k_{\ell+1}}}(\phi)-A_{2^{k_{\ell}}}(\phi)\|_2^2 \]
\[ = \sum_{\ell\in [J)} \int_{\R^n} A(\phi_{(2^{k_{\ell+1}})})(x)A(\phi_{(2^{k_{\ell+1}})}-\phi_{(2^{k_{\ell}})})(x)  dx \]
\[ - \sum_{\ell\in [J)}  \int_{\R^n} A(\phi_{(2^{k_{\ell}})})(x)A(\phi_{(2^{k_{\ell+1}})}-\phi_{(2^{k_{\ell}})})(x)  dx \]
By the Cauchy-Schwarz inequality in $x$ and in the summation,  this is 
\[ \le  \Big ( \sum_{\ell\in [J)}  \|A(\phi_{(2^{k_{\ell+1}})})\|^2_2 \Big)^{1/2} \Big (\sum_{\ell\in [J)} \|A(\phi_{(2^{k_{\ell+1}})}-\phi_{(2^{k_{\ell}})})\|^2_2\Big)^{1/2}\]
\[+ \Big ( \sum_{\ell\in [J)} \|A(\phi_{(2^{k_{\ell}})})\|_2^2\Big)^{1/2} \Big(\sum_{\ell\in [J)} \|A(\phi_{(2^{k_{\ell+1}})}-\phi_{(2^{k_{\ell}})})\|^2_2\Big)^{1/2} \]
And thus, with $V= \|A_{t}(\phi))\|_{V_{2,J}(t\in 2^\mathbb{Z}; L^2)}$, 
\[V^2 \le  V\cdot  \Big( \Big(\sum_{\ell\in [J)} \|A(\phi_{(2^{k_{\ell+1}})})\|^2_2\Big)^{1/2} + \Big( \sum_{\ell\in [J)} \|A(\phi_{(2^{k_{\ell}})})\|^2_2 \Big)^{1/2} \Big)   \]
We may assume $V\neq 0$  as otherwise the claim holds trivially. Dividing by $V$ and estimating the terms on the   right-hand side by the supremum over all increasing sequences of integers gives 
\[V \le  2 \cdot  \Big( \sup \sum_{\ell\in [J)} \|A(\phi_{(2^{k_{\ell}})})\|^2_2\Big)^{1/2}     \]
Squaring both sides gives the desired estimate.  
\end{proof}

 
 

\subsection{On-diagonal from main argument}
% The goal of this section is to prove off-diagonal estimates Propositions \ref{prop:1-off-W} and \ref{lem:1-off-nonW} using Theorem \ref{induct positive terms theorem} for $k=2$.

Throughout this section, let 
$\nu=1$ if $n>2$ and $\nu=2$ if $n=2$.

\begin{lemma}[$\rho$ kernels - reduction variant] \label{lem:rho-kernels-reduction}  Let $a\in A$. For $h\in \N,j\in \Z$ define
\[
\lambda_{h,j}^-=
\begin{cases}
2^{h-1}a(j),&h\geq1,\\
a(j-1),&h=0,
\end{cases}
\qquad
\lambda_{h,j}^+=2^ha(j),
\]
and
\[
\mu_{h,j}^-=2^ha(j-1),
\qquad
\mu_{h,j}^+=2^ha(j).
\]
Let \(\rho_{h,j}\) be the function whose Fourier transform is
\[
\widehat{\rho_{h,j}}(\zeta)
=
\bigl(
\widehat\Phi(\lambda_{h,j}^-\zeta)
-
\widehat\Phi(\lambda_{h,j}^+\zeta)
\bigr)
\bigl(
\g(\mu_{h,j}^-\zeta)
-
\g(\mu_{h,j}^+\zeta)
\bigr)^{-\nu/2}.
\]

Then \(\rho_{h,j}\in W_0(\mathbb R)\). Moreover, if \(h\geq1\), then, with $\lambda=2^ha(j)$,
\begin{equation}
    \label{E:increase-data-rho-mean-estimate}
    |\rho_{h,j}(x+p)-\rho_{h,j}(x)|
    \leq
    C_{\ref{lem:rho-kernels-reduction}}
    \min(1,\lambda^{-1}|p|)
    \left(
    \langle x+p\rangle_{(\lambda)}^2
    +
    \langle x\rangle_{(\lambda)}^2
    \right),
\end{equation}
and if $h\in \N$, then
\begin{equation}
    \label{E:increase-data-rho-zero-estimate}
    |\rho_{h,j}(x)|
    \leq
    C_{\ref{lem:rho-kernels-reduction}}
    \big(
    \langle x\rangle_{(\lambda_{h,j}^-)}^2
    +
    \langle x\rangle_{(\lambda_{h,j}^+)}^2
    \big),
\end{equation}
where $C_{\ref{lem:rho-kernels-reduction}} = 2^{21} (C_{\ref{mean four scale Gaussian kernel},2}+C_{\ref{four scale Gaussian kernel},2})$. 
\end{lemma}

\begin{proof}
If \(h\geq1\), the parameters satisfy
\[
2\mu_{h,j}^-
\leq
2\lambda_{h,j}^-
=
\lambda_{h,j}^+
=
\mu_{h,j}^+,
\]
because \(2a(j-1)\leq a(j)\).  
If $h=0$, we have 
\[
2\mu_{0,j}^-
=
2\lambda_{0,j}^-
\leq
\lambda_{0,j}^+
=
\mu_{0,j}^+.
\]


Now, if $h\ge 1$, we  use Proposition   
\ref{mean four scale Gaussian kernel}   with $N=2$, which  gives   
\[|\rho_{h,j}(x+p)-\rho_{h,j}(x)| \le 2^{18} C_{\ref{mean four scale Gaussian kernel},2} \min(1,2\pi \lambda^{-1}|p|) (\langle x+p\rangle^2_{(\lambda)}+
\langle x\rangle^2_{(\lambda)}), 
\]
which gives \eqref{E:increase-data-rho-mean-estimate}.
Here we also used that by Proposition~\ref{standard bump properties},
\[
\max_{0\leq m\leq2}
\|(\widehat\Phi)^{(m)}\|_\infty
\leq2^{18}.
\]
If $h\in \N$, we use  Proposition~\ref{four scale Gaussian kernel}   with $N=2$, which   gives $\rho_{h,j}\in W_0(\R)$ and
\[ |\rho_{h,j}(x)| \le 2^{18}C_{\ref{four scale Gaussian kernel},2} (\langle x\rangle^2_{(\lambda_{h,j}^-)} + \langle x\rangle^2_{(\lambda_{h,j}^+)}). \]
\end{proof}

\begin{rem}
The value of $\widehat{\rho_{h,j}}$   at the origin is understood through the continuous
extension supplied by Proposition~\ref{four scale Gaussian kernel}.
\end{rem}

\begin{lemma}[affine diagonal cancellation - reduction variant]
\label{lem:affine-diagonal-cancellation-reduction}
Let \(M\in W_0(\mathbb R^2)\) satisfy for every \(\xi\in\mathbb R\)
\[
\widehat M(\xi,-\xi)=0
\]
 Then for every \(x\in\mathbb R\),
\begin{equation}
    \label{E:affine-diagonal-cancellation}
    \int_{\mathbb R}M(x+q,q)\,dq=0
\end{equation}
\end{lemma}

\begin{proof}
    Define
\[
G(x)=\int_{\mathbb R}M(x+q,q)\,dq.
\]
By Proposition~\ref{W_0 fiber integrals}, \(G\in W_0(\mathbb R)\).
By Fubini's theorem and the change of variables
$ u=x+q, \, v=q, $
we have
\[
\widehat G(\xi)
=
\int_{\mathbb R^2}
M(x+q,q)e^{-2\pi ix\xi}\,dq\,dx
=
\int_{\mathbb R^2}
M(u,v)e^{-2\pi i(u-v)\xi}\,du\,dv=
\widehat M(\xi,-\xi)
=
0.
\]
Injectivity of the Fourier transform on \(L^1\) and continuity of $G$ give
\(G=0\) everywhere.
\end{proof}


\begin{proposition}[bracket domination - reduction variant] 
\label{lem:increase-data-bracket-domination}
Let \(a\in A\), and let
$
\mathcal P\subset B_{\dist}(a,1)^2\times[2)$ be such that $\#\mathcal P\leq5$. Assume additionally that
\begin{equation}
    \label{E:increase-data-one-sided-scales}
    t_0(j)\leq a(j),
    \quad
    t_1(j)\leq a(j)
\end{equation}
for every \((t_0,t_1,u)\in\mathcal P\) and \(j\in\mathbb Z\).
Suppose that \(M_j\in W_0(\mathbb R^2)\) satisfies
\[
\widehat M_j(\xi,-\xi)=0
\]
and
\[
|M_j(v)|
\leq
\sum_{(t_0,t_1,u)\in\mathcal P}
\langle(W_uv)_0\rangle_{(t_0(j))}^{3/2}
\langle(W_uv)_1\rangle_{(t_1(j))}^{3/2}.
\]

Let $\lambda_{h,j}^-,\lambda_{h,j}^+$ be as in Lemma \ref{lem:rho-kernels-reduction}.
For $h\in\N$, $j\in \Z$, we define
$ 
\sigma_{h,j} = s(2^{h} a,j) 
$
and
\[
N_{h,j} =
\mathcal F^{-1}\left((\xi,\eta) \mapsto (\widehat{\Phi_{(\lambda_{h,j}^-)}}-\widehat{\Phi_{(\lambda_{h,j}^+)}})(\xi+\eta) \widehat{\sigma_{h,j}}(\xi+\eta)^{-\nu} \widehat{M_j}(\xi,\eta)\right),
\]
% and for $h>0$,
% \[
% N_{h,j} =
% \mathcal F^{-1}\left((\xi,\eta) \mapsto (\widehat{\Phi_{(2^{h-1} a(j))}}-\widehat{\Phi_{(2^{h}a(j))}})(\xi+\eta) \widehat{\sigma_{h,j}}(\xi+\eta)^{-\nu} \widehat{M_j}(\xi,\eta)\right),
% \]
% and
% \[
% N_{0,j} =
% \mathcal F^{-1}\left((\xi,\eta) \mapsto (\widehat{\Phi_{(a(j-1))}}-\widehat{\Phi_{(a(j))}})(\xi+\eta) \widehat{\sigma_{0,j}}(\xi+\eta)^{-\nu} \widehat{M_j}(\xi,\eta)\right).
% \]

Then for every \(h\in\mathbb N\), there exist a finite set
\(\mathcal B_h\), numbers \(u_{h,b}\in[2)\), and pairs
\[
\beta_{h,b}
=
(\beta_{h,b}^0,\beta_{h,b}^1)\in A^2,
\qquad b\in\mathcal B_h,
\]
satisfying
\begin{equation}
    \label{E:increase-data-B-cardinality}
    \#\mathcal B_h
    \leq
    C_{\ref{lem:increase-data-bracket-domination},0},
\end{equation}
\begin{equation}
    \label{E:increase-data-beta-base-distance}
    \dist(a,\beta_{h,b}^\ell)
    \leq
     (1+h)
\end{equation}
for every \(b\in\mathcal B_h\) and \(\ell\in[2)\), and
\begin{equation}
    \label{E:increase-data-bracket-majorant}
    \begin{aligned}
    |N_{h,j}(v)|
    &\leq
    C_{\ref{lem:increase-data-bracket-domination},1}
    2^{-h/3}
    \sum_{b\in\mathcal B_h}
    \langle(W_{u_{h,b}}v)_0
    \rangle_{(\beta_{h,b}^0(j))}^{7/6}
    \langle(W_{u_{h,b}}v)_1
    \rangle_{(\beta_{h,b}^1(j))}^{7/6} 
    \end{aligned}
\end{equation}
where $C_{\ref{lem:increase-data-bracket-domination},0}=2^6$, $C_{\ref{lem:increase-data-bracket-domination},1}
=
2^{10}C_{\ref{lem:rho-kernels-reduction}}$. 
In particular,
\begin{equation}
    \label{E:increase-data-beta-pair-distance}
    \Delta(\beta_{h,b})
    \leq
    2(1+h).
\end{equation}
\end{proposition}


 

\begin{proof}
    See Section \ref{sec:proof-increase-data-bracket-domination}.
\end{proof}

\begin{proposition}[Gaussian domination - reduction variant] 
\label{lem:increase-data-Gaussian-expansion}
Let the notation and conclusions of Proposition~
\ref{lem:increase-data-bracket-domination} hold. For
\(b\in\mathcal B_h\) and
\(m=(m_0,m_1)\in\mathbb N^2\), define
\[
p_{h,b,m}
=
\left(
2^{m_0}\beta_{h,b}^0,
2^{m_1}\beta_{h,b}^1
\right)\in A^2.
\]
Then
\begin{equation}
    \label{E:increase-data-p-distance}
    \Delta(p_{h,b,m})
    \leq
    2(1+h)
    +
    |m|,
\end{equation}
and
\begin{equation}
    \label{E:increase-data-Gaussian-majorant}
    \begin{aligned}
    |N_{h,j}(v)|
    \leq
    C_{\ref{lem:increase-data-Gaussian-expansion}}
    2^{-h/3}
    \sum_{b\in\mathcal B_h}
    \sum_{m\in\mathbb N^2}
    2^{-|m|/6}
    G_{p_{h,b,m}(j),u_{h,b}}(v).
    \end{aligned}
\end{equation}
where $ C_{\ref{lem:increase-data-Gaussian-expansion}}
    =
    2^3C_{\ref{Gaussian domination}}^2C_{\ref{lem:increase-data-bracket-domination},1}.$
The   series on the right-hand side converges for every
\(v\in\mathbb R^2\) and in \(L^1(\mathbb R^2)\).
\end{proposition}

\begin{proof}
    See Section \ref{sec:proof-increase-data-Gaussian-expansion}.
\end{proof}


 


\begin{lemma}
\label{lem:N-reduction} Let $a\in A$ and let $\mathcal P$, $(M_j)_{j\in \Z}$, and $(N_{h,j})_{h\in \N,j\in \Z}$,  be as in Proposition~\ref{lem:increase-data-bracket-domination}. 
Then for each $h\in \N$ and $j\in \Z$,   $N_{h,j}\in W_0(\R^2)$ and 
\begin{equation}
    \label{E:diagonal-band-reduction-N-L1}
    \|N_{h,j}\|_1\leq  C_{\ref{lem:N-reduction}}
\end{equation}
with $C_{\ref{lem:N-reduction}} = 2^9C_{\ref{lem:rho-kernels-reduction}}$
\end{lemma}
\begin{proof}
Let \(\rho_{h,j}\) be as in Lemma \ref{lem:rho-kernels-reduction}. 
Then, 
\[
N_{h,j}=M_j\ast_{(1,1)}\rho_{h,j}.
\]
Indeed, this follows from 
\[
\widehat{M_j\ast_{(1,1)}\rho_{h,j}}(\xi,\eta)
=
\widehat M_j(\xi,\eta)
\widehat{\rho_{h,j}}(\xi+\eta)
=
\widehat{N_{h,j}}(\xi,\eta).
\]
using Proposition~\ref{convolution vector}.
 

By the pointwise bound on \(M_j\), the fact that each
\(W_u\) is measure preserving, and
\[
\int_{\mathbb R}\langle x\rangle_{(t)}^{3/2}\,dx=4,
\]
we have
\begin{equation}
    \label{E:diagonal-band-reduction-M-L1}
    \|M_j\|_1
    \leq
    5\cdot 4^2
    =
    80.
\end{equation}
 To estimate  $\rho_{h,j}$  we   apply  Proposition \ref{lem:rho-kernels-reduction}, 
which gives
\[
|\rho_{h,j}(x)|
\leq
%2^{18}C_{\ref{four scale Gaussian kernel},2}
C_{\ref{lem:rho-kernels-reduction}}
\left(
\langle x\rangle_{(\lambda_{h,j}^-)}^2
+
\langle x\rangle_{(\lambda_{h,j}^+)}^2
\right).
\]
Since
\[
\int_{\mathbb R}\langle x\rangle_{(t)}^2\,dx=2,
\]
it follows that
\[
\|\rho_{h,j}\|_1
\leq
4C_{\ref{lem:rho-kernels-reduction}}
\]
Combining this last estimate with 
\eqref{E:diagonal-band-reduction-M-L1} and the \(L^1\) convolution
inequality proves
\[
\|N_{h,j}\|_1
\leq
80\cdot 4 C_{\ref{lem:rho-kernels-reduction}}
\leq
C_{\ref{lem:N-reduction}}.
\]
\end{proof}



 
\begin{proposition}[increase data - reduction variant]
\label{P:increase-data-reduction}
Let $a\in A$ and let $\mathcal P$, $(M_j)_{j\in \Z}$, and $(N_{h,j})_{h\in\N,j\in\Z}$ be as in Proposition~\ref{lem:increase-data-bracket-domination}.
 For $j\in \Z$ and $y\in (\R^2)^2$ let 
    \[\tilde{M}_{h,j}(y) = |N_{h,j}(y_0)| (\sigma_{h,j}^{\otimes2})(y_1). \]
    
Then with $\mathbf{\tilde{M}}_h=(\tilde{M}_{h,j})_{j\in \Z}$,
    \[\|\mathbf{\tilde{M}}_h\|_{\rm M(2)}\le C_{\ref{P:increase-data-reduction}} 2^{-h/4}\]
    where $C_{\ref{P:increase-data-reduction}} = 2^{22}C_{\ref{induct positive terms theorem}}C_{\ref{lem:increase-data-Gaussian-expansion}}C_{\ref{lem:increase-data-bracket-domination},0}$
    
% Let $\nu=1$ if $n>2$ and $\nu=2$ if $n=2$, and let $a\in A$. Let $\mathcal P \subset B_{\operatorname{dist}}(a,1)^2 \times [2)$ be such that $\# \mathcal P \leq 5$.
% Suppose that $(M_j)_{j\in\Z}\subset W_0(\R^2)$ satisfies for all $\xi \in \R$ and $j\in \Z$ 
% \[
% \widehat{M_j}(\xi,-\xi)=0,
% \]
% and for all $v\in \R^2$ and $j\in \Z$,
% \[ 
% |M_j(v)| \le 
% \sum_{(t_0,t_1,u) \in \mathcal P} \langle (W_uv)_0\rangle_{(t_0(j))}^{3/2} \langle (W_uv)_1\rangle_{(t_1(j))}^{3/2}.
% \]
% For $h\in\N$ and $j\in \Z$, we define
% \[ 
% \sigma_{h,j} = s(2^{h} a,j) 
% \]
% and for $h>0$,
% \[
% N_{h,j} =
% \mathcal F^{-1}\left((\xi,\eta) \mapsto (\widehat{\Phi_{(2^{h-1} a(j))}}-\widehat{\Phi_{(2^{h}a(j))}})(\xi+\eta) \widehat{\sigma_{h,j}}(\xi+\eta)^{-\nu} \widehat{M_j}(\xi,\eta)\right),
% \]
% and
% \[
% N_{0,j} =
% \mathcal F^{-1}\left((\xi,\eta) \mapsto (\widehat{\Phi_{(a(j-1))}}-\widehat{\Phi_{(a(j))}})(\xi+\eta) \widehat{\sigma_{0,j}}(\xi+\eta)^{-\nu} \widehat{M_j}(\xi,\eta)\right).
% \]
\end{proposition} 

\begin{proof}
% \ls{todo}
% Same proof as for Proposition~\ref{induct positive terms imply increase data} (reduce to InductPositiveTerms(2)). Unify the Gaussian domination with that from Section~\ref{sec:main argument}?
By Lemma \ref{lem:N-reduction}, \(N_{h,j}\in W_0(\mathbb R^2)\) and consequently $|N_{h,j}|\in W_0(\mathbb R^2)$. Together with Propositions~\ref{square root Gaussian difference W0}
and~\ref{tensor Wiener}, this gives
\[
\tilde M_{h,j}\in W_0((\mathbb R^2)^2).
\]
Let \(\mathcal B_h\), \(u_{h,b}\), and \(p_{h,b,m}\) be given by
Lemma~\ref{lem:increase-data-Gaussian-expansion}. For
\(b\in\mathcal B_h\) and \(m\in\mathbb N^2\), define
\[
\alpha_h
=
\left(
2^{h-1/2}a,
2^{h-1/2}a
\right)\in A^2
\]
and
\[
\gamma_{h,b,m}
=
\left(
2,
(u_{h,b},0),
(p_{h,b,m},\alpha_h)
\right)\in\Gamma.
\]

Since the second orientation is zero, Definition~
\ref{s multiplier} gives
\[
(s_{\gamma_{h,b,m}})_{1,j}
=
s\Big(
\big( 
(2^{h-1/2}a)^2+
(2^{h-1/2}a)^2
\big )^{1/2},
j
\Big)
=
s(2^ha,j)
=
\sigma_{h,j}.
\]
Consequently,
\begin{equation}
    \label{E:increase-data-positive-term-recognition}
    \bigl(
    \M(
    \gamma_{h,b,m},
    s_{\gamma_{h,b,m}}\otimes s_{\gamma_{h,b,m}},
    1
    )
    \bigr)_j(y) =
    G_{p_{h,b,m}(j),u_{h,b}}(y_0)
    \bigl(\sigma_{h,j}^{\otimes2}\bigr)(y_1).
\end{equation}
Moreover,
$\Delta(\alpha_h)=0$
and therefore
\[ \Delta_{\gamma_{h,b,m}} =
1+\Delta(p_{h,b,m}). \]

Set $
q_n=2-2^{3-n}\in[0,2).$
By Theorem~\ref{induct positive terms theorem}, applied with
\(i=1\),
\[
\left\|
\M(
\gamma_{h,b,m},
s_{\gamma_{h,b,m}}\otimes s_{\gamma_{h,b,m}},
1
)
\right\|_{\mathrm M(2)}
 \leq
C_{\ref{induct positive terms theorem}}
\bigl(
1+\Delta(p_{h,b,m})
\bigr)^{q_n}.
\]

Using Lemma~\ref{lem:increase-data-Gaussian-expansion},
Proposition~\ref{Monotonicity K}, we obtain
\[
\|\tilde{\mathbf M}_h\|_{\mathrm M(2)}
\leq
C_{\ref{lem:increase-data-Gaussian-expansion}}
2^{-h/3}
\sum_{b\in\mathcal B_h}
\sum_{m\in\mathbb N^2}
2^{-|m|/6}
\left\|
\M(
\gamma_{h,b,m},
s_{\gamma_{h,b,m}}\otimes s_{\gamma_{h,b,m}},
1
)
\right\|_{\mathrm M(2)}.
\]
(first sum partial
sums of the Gaussian series, then pass to the full series by Proposition~
\ref{prism BL inequality} and dominated convergence). Hence, together with the estimate on $\#\mathcal{B}$ and $(1+\Delta(p_{h,b,m}))^{q_n}\le 2^{4}(1+h+|m|)^2$, 
\[
\|\tilde{\mathbf M}_h\|_{\mathrm M(2)}
\leq
C_{\ref{induct positive terms theorem}}C_{\ref{lem:increase-data-Gaussian-expansion}}C_{\ref{lem:increase-data-bracket-domination},0}2^4
2^{-h/3}
\sum_{m\in\mathbb N^2}
2^{-|m|/6}
(1+h+|m|)^{2}.
\]
Since 
\[\sup_{h\in \N}2^{-h/12}\sum_{m\in \N^2}2^{-|m|/6}(1+h+|m|)^2\le 2^{18}\]
% Since \(q_n<2\),
% \[
% (1+h+|m|)^{q_n}
% \leq
% (1+h)^2(1+|m|)^2.
% \]
% Thus,
% \[
% \|\tilde{\mathbf M}_h\|_{\mathrm M(2)}
% \leq
% C_{\ref{induct positive terms theorem}}C_{\ref{lem:increase-data-Gaussian-expansion}}C_{\ref{lem:increase-data-bracket-domination},0}2^4
% n^2
% 2^{-h/3}(1+h)^2
% \sum_{m\in\mathbb N^2}
% 2^{-|m|/6}(1+|m|)^2.
% \]
% The sum in \(m\) is $\le 2^{20}$, and
% \[
% 2^{-h/3}(1+h)^2
% =
% 2^{-h/4}
% \left(
% 2^{-h/12}(1+h)^2
% \right)
% \leq
% 2^82^{-h/4}.
% \]
it follows that
\[
\|\tilde{\mathbf M}_h\|_{\mathrm M(2)}
\leq
C_{\ref{P:increase-data-reduction}}2^{-h/4},
\]
as desired.
\end{proof}

 

\subsection{On-diagonal from off-diagonal estimates}
\begin{proposition}[diagonal band - reduction variant]
\label{P:diagonal-band-reduction}
Let $a\in A$ and let $\mathcal P$, $(M_j)_{j\in \Z}$ be as in Proposition~\ref{P:increase-data-reduction}. For $h\in \N_{\geq 1}$ and $j\in \Z$ we define
\[
L_{h,j}= M_j \ast_{(1,1)} (\Phi_{(2^{h-1}a(j))}-\Phi_{(2^{h}a(j))}),
\]
and for $j\in \Z$,
\[
L_{0,j}=M_j \ast_{(1,1)} (\Phi_{(a(j-1))}-\Phi_{(a(j))}).
\]
Let $\mathbf{L}_h=(L_{h,j})_{j\in \Z}$. Then
\[
\sum_{h\in \N} \|\mathbf{L}_h\|_{\rm M(1)} \leq C_{\ref{P:diagonal-band-reduction}}
\]
with  $C_{\ref{P:diagonal-band-reduction}} = 2^4(
    C_{\ref{lem:N-reduction}}
    C_{\ref{P:increase-data-reduction}}
  )^{1/2} + 2^3 C_{\ref{P:increase-data-reduction}}$.  
\end{proposition}

% \ls{todo
% Same proof as for Proposition~\ref{increase data implies diagonal band} (except that the analogous decomposition in Section~\ref{sec:main argument} now has three stick terms corresponding to $h=0$ - should we also have three here for consistency?). Reduce to Proposition~\ref{P:increase-data-reduction}.}


\begin{proof}
    This follows by case distinction from Propositions \ref{P:diagonal-band-reduction case 1} and
     \ref{P:diagonal-band-reduction case 2}.
\end{proof}

\begin{proposition} 
\label{P:diagonal-band-reduction case 1}
 Proposition \ref{P:diagonal-band-reduction} holds in case $n>2$.
\end{proposition}

\begin{proof} Since $n>2$ we have  $\nu=1$.  Let $\sigma_{h,j}$ and $N_{h,j}$ be as  Proposition~\ref{P:increase-data-reduction}.  
We have 
\begin{equation}
    \label{E:diagonal-band-reduction-factorization}
    L_{h,j}
    =
    N_{h,j}\ast_{(1,1)}\sigma_{h,j}.
\end{equation}
Indeed, by Proposition~\ref{convolution vector}, 
\[\widehat{
N_{h,j}\ast_{(1,1)}\sigma_{h,j} 
}(\xi,\eta)=
\widehat{N_{h,j}}(\xi,\eta)
\widehat{\sigma_{h,j}}(\xi+\eta)=
\widehat{L_{h,j}}(\xi,\eta).
\]

Fix
\(J\in\mathbb N_{>0}\) and \(\mathbf F\in\mathfrak F\).
 Proposition 
\ref{Cauchy-Schwarz at k}  with \(k=1\), applied to
\[
\rho_j=C_{\ref{lem:N-reduction}}^{-1}N_{h,j},
\qquad
\varphi_j=\sigma_{h,j},
\]
gives
\[
\Big|
\Lambda_1(\sum_{j\in[J)}L_{h,j})(\mathbf F)
\Big|^2
\leq C_{\ref{lem:N-reduction}} J
\Big|
\Lambda_2(
\sum_{j\in[J)}\tilde M_{h,j}
)(\mathbf F)
\Big|
\]
Here
\[
\tilde M_{h,j}(y)
=
|N_{h,j}(y_0)|
\bigl(\sigma_{h,j}^{\otimes2}\bigr)(y_1),
\]
as in Proposition~\ref{P:increase-data-reduction}. 

 Proposition~\ref{P:increase-data-reduction} and the
definition of the \({\rm M}(2)\)-norm give
\[
\Big |
\Lambda_2 (
\sum_{j\in[J)}\tilde M_{h,j}
)(\mathbf F)
\Big |
\leq
C_{\ref{P:increase-data-reduction}}
2^{-h/4}
J^{1-2^{3-n}}.
\]
It follows that
\[
\Big|
\Lambda_1(\sum_{j\in[J)}L_{h,j})(\mathbf F)
\Big|
\leq
\bigl(
C_{\ref{lem:N-reduction}}
C_{\ref{P:increase-data-reduction}}
\bigr)^{1/2}
2^{-h/8}
J^{1-2^{2-n}},
\]
since 
$
\frac12+\frac12(1-2^{3-n})
=
1-2^{2-n}.
$
Taking the supremum over \(J\) and \(\mathbf F\) therefore gives
\begin{equation}
    \label{E:diagonal-band-reduction-n-greater-2}
    \|\mathbf L_h\|_{\mathrm M(1)}
    \leq
    \bigl(
    C_{\ref{lem:N-reduction}}
    C_{\ref{P:increase-data-reduction}}
    \bigr)^{1/2}
    2^{-h/8}.
\end{equation}
Summing over $h\in \N$ finishes the proof.
 \end{proof}

 


 



\begin{proposition} 
\label{P:diagonal-band-reduction case 2}
 Proposition \ref{P:diagonal-band-reduction} holds in case $n=2$.
\end{proposition}

\begin{proof}
    Now $\nu=2$ and 
    \[
L_{h,j}
=
N_{h,j}\ast_{(1,1)}
(\sigma_{h,j}\ast\sigma_{h,j})
\]
where    $\sigma_{h,j}$   and $N_{h,j}$ are  as   in Proposition~\ref{P:increase-data-reduction}.   
Proposition 
\ref{Cauchy-Schwarz at n-1} gives
\[
\Big|
\Lambda_1(\sum_{j\in[J)}L_{h,j})(\mathbf F)
\Big|
\leq
\sup_{\widetilde{\mathbf F}\in\mathfrak F}
\Big|
\Lambda_2(
\sum_{j\in[J)}\tilde M_{h,j}
)(\tilde{\mathbf F})
\Big|
\]
with $\tilde M_{h,j}$ as in Proposition~\ref{P:increase-data-reduction}.  
Proposition \ref{P:increase-data-reduction} now gives
\begin{equation}
    \label{E:diagonal-band-reduction-n-equals-2}
    \|\mathbf L_h\|_{\mathrm M(1)}
    \leq
    C_{\ref{P:increase-data-reduction}}2^{-h/4}.
\end{equation}
Summing in $h\in \N$ finishes the proof. 
\end{proof}
 
 


\begin{lemma}\label{lem:L1-reduction}
% Let $a\in A$ and let $\mathcal P$, $(M_j)_{j\in \Z}$ be as in Proposition~\ref{P:increase-data-reduction}. 
Let $a\in A$ and suppose that $(M_j)_{j\in\Z}\subset W_0(\R^2)$ satisfies for all $\xi \in \R$ and $j\in \Z$ 
\[
\widehat{M_j}(\xi,-\xi)=0.
\]
Then 
\[    \lim_{t\to\infty}
    \| M_j\ast_{(1,1)}\Phi_{(t)} \|_1
    =0
    \]
\end{lemma}

\begin{proof}
%     Define
% \[
% G_j(x)=\int_{\mathbb R}M_j(x+q,q)\,dq.
% \]
% The function \(G_j\) is well defined and belongs to \(W_0(\mathbb R)\)
% by Proposition~\ref{W_0 fiber integrals}. Moreover, by Fubini's
% theorem and the change of variables \(u=x+q\), \(v=q\),
% \[
% \widehat{G_j}(\xi)
% =
% \int_{\mathbb R^2}
% M_j(x+q,q)e^{-2\pi i x\xi}\,dq\,dx
% =
% \widehat{M_j}(\xi,-\xi)
% =
% 0.
% \]
By Lemma \ref{lem:affine-diagonal-cancellation-reduction}, 
\begin{equation}
    \label{E:vanishing-diagonal-reduction-line-cancellation}
    \int_{\mathbb R}M_j(x+q,q)\,dq=0
\end{equation}
for every \(x\in\mathbb R\).
Let
\[
R_{t,j}=M_j\ast_{(1,1)}\Phi_{(t)}.
\]
Using \eqref{E:vanishing-diagonal-reduction-line-cancellation} and
changing variables \(q=y-p\), we obtain
\[
\begin{aligned}
R_{t,j}(x+y,y)
&=
\int_{\mathbb R}
\Phi_{(t)}(y-q)M_j(x+q,q)\,dq
\\
&=
\int_{\mathbb R}
\bigl(\Phi_{(t)}(y-q)-\Phi_{(t)}(y)\bigr)
M_j(x+q,q)\,dq.
\end{aligned}
\]
It follows from the triangle inequality and Fubini's theorem that
    \[
    \|R_{t,j}\|_1
    \le
    \int_{\mathbb R^2}
    \omega(q/t)\,
    |M_j(x+q,q)|\,dx\,dq,
    \]
    where
    \[
    \omega(s)
    =
    \int_{\mathbb R}|\Phi(z-s)-\Phi(z)|\,dz.
    \]
    By continuity of translations and the $L^1$ norm,
    \(\omega(q/t)\to0\) as \(t\to\infty\) for   \(q\in\mathbb R\).
    Also,
    \[
    0\le \omega(q/t)\le 2\|\Phi\|_1.
    \]
The claim now follows by the 
    dominated convergence theorem.
\end{proof}
 

\begin{proposition}[vanishing diagonal - reduction variant]
\label{P:vanishing-diagonal-reduction}
Let $a\in A$ and let $\mathcal P$, $(M_j)_{j\in \Z}$ be as in Proposition~\ref{P:increase-data-reduction}. In addition, assume that $\widehat{M_j}(\xi,\eta)$ is supported in the set where $|\xi+\eta| \leq 2^{-1} (a(j-1))^{-1}$. 
%\ls{Maybe replace $2^{-1}$ by some other number.} 
Let $\mathbf{M}=(M_j)_{j\in \Z}$. Then 
\[
\|\mathbf{M}\|_{\rm M(1)} \leq C_{\ref{P:vanishing-diagonal-reduction}}
\]
with $C_{\ref{P:vanishing-diagonal-reduction}} = C_{\ref{P:diagonal-band-reduction}}$.
\end{proposition}

\begin{proof}
%  \ls{todo; 
% Same proof as for Proposition~\ref{diagonal band implies vanishing diagonal} (reduce to Proposition~\ref{P:diagonal-band-reduction}). The terms with $h<0$ vanish by the Fourier support condition.}
Since    $|\xi+\eta|\le 2^{-1} (a(j-1))^{-1}$ and   $\widehat{\Phi}=1$ on $[-1/2,1/2]$, 
\[M_j*_{(1,1)} \Phi_{a(j-1)} = M_j\]
Let $L_{h,j}$ be as in Proposition \ref{P:increase-data-reduction}. For $H\in \N$ we telescope
\[\sum_{h=0}^H L_{h,j} = M_j*_{(1,1)} \Big( \Phi_{(a(j-1))}-\Phi_{(a(j))} + \sum_{h=1}^{H}
( \Phi_{(2^{h-1}a(j))} - \Phi_{(2^ha(j))} ) \Big)\]
\[= M_j*_{(1,1)} ( \Phi_{(a(j-1))} - \Phi_{(2^Ha(j))} )\]
\[= M_j-M_j*_{(1,1)}\Phi_{(2^Ha(j))} =: M_j-R_{H,j}.\]
Let $J\ge 1$ and $\F\in \mathfrak{F}$. By linearity and the triangle inequality, 
%\[\sum_{j\in [J)} M_j = \sum_{h=0}^H \sum_{j\in [J)}L_{h,j} + \]
\[|\Lambda_1 (\sum_{j\in [J)} M_j )(\F)| \le  \sum_{h=0}^H |\Lambda_1( \sum_{j\in [J)}L_{h,j})(\F)| + |\Lambda_1( \sum_{j\in [J)}R_{H,j})(\F)| \]
By Propositions  \ref{prism BL inequality} and \ref{M to K}, 
\[|\Lambda_1( \sum_{j\in [J)}R_{H,j})(\F)| \le \|\sum_{j\in [J)}R_{H,j}\|_1 \le \sum_{j\in [J)}\|R_{H,j}\|_1,\]
which tends to zero as $H\to \infty$ by Proposition \ref{lem:L1-reduction}.

Therefore, setting $c_J = \min(1,J^{-1+2^{2-n}})$, 
\[c_J |\Lambda_1 (\sum_{j\in [J)} M_j )(\F)| \le  \sum_{h\in\N} c_J |\Lambda_1( \sum_{j\in [J)}L_{h,j})(\F)|\le \sum_{h\in \N} \|\mathbf{L}_h\|_{\rm M(1)} \le C_{\ref{P:diagonal-band-reduction}}\]
where the last inequality is by Proposition 
\ref{P:diagonal-band-reduction}. Taking the supremum over
$J$ and $\mathbf F\in\mathfrak F$ yields the claim.
\end{proof}


Fix a universal pair $(\phi_0,\phi_1)$. For $\lambda>0$ and $i\in[2)$ write
\[ \phi_{i,\lambda} = (\phi_i)_{(\lambda)}. \]

\begin{proposition}\label{P:one-scale-estimate-window}
Let $a\in A$. For $j\in \Z$, we set
\[
M_j=(\phi_{0,a(j-1)})^{\otimes 2} - (\phi_{0,a(j)})^{\otimes 2}
\]
and $\mathbf{M}=(M_j)_{j\in \Z}$. Then
\begin{equation}\label{E:telescoping-estimate}
\|\mathbf{M}\|_{\rm M(1)} \leq C_{\ref{P:one-scale-estimate-window}},
\end{equation}
where $C_{\ref{P:one-scale-estimate-window}}=8$.
%Let $\lambda >0$. Then for every $\mathbf{F} \in \mathfrak{F}$, 
%\[
%|\Lambda_1(\phi_{0,\lambda}^{\otimes 2})(\mathbf{F})| \leq C_{\ref{L:one-scale-estimate-window}},
%\]
%where $C_{\ref{L:one-scale-estimate-window}}=4$.
\end{proposition}

\begin{proof}
For a given $J\in \mathbb N_{>0}$, 
\begin{equation}\label{E:telescoping}
\sum_{j\in [J-1)} M_j=(\phi_{0,a(-1)})^{\otimes 2}-(\phi_{0,a(J-1)})^{\otimes 2}.
\end{equation}
Let $\lambda \in \{a(-1), a(J-1)\}$.
A combination of Propositions~\ref{prism BL inequality} and~\ref{M to K} yields for any $\mathbf{F} \in \mathfrak{F}$,
\begin{equation}\label{E:est-l1-norm}
|\Lambda_1 ((\phi_{0,\lambda})^{\otimes 2})(\mathbf{F})| \leq
\|K_1((\phi_{0,\lambda})^{\otimes 2})\|_1 \leq\|(\phi_{0,\lambda})^{\otimes 2}\|_1
=\|\phi_0\|_1^2.
\end{equation}
We then use the Cauchy-Schwarz inequality, Plancherel's identity and properties \textup{(i)} and \textup{(ii)} in Definition~\ref{def:cn-window} to obtain
\begin{equation}\label{E:est-of-l1-norm}
\|\phi_0\|_1 \leq \sqrt{2} \|\phi_0\|_2 =\sqrt{2} \|\widehat{\phi_0}\|_2 \leq 2.
\end{equation}
Combining~\eqref{E:telescoping}, \eqref{E:est-l1-norm} and \eqref{E:est-of-l1-norm} yields~\eqref{E:telescoping-estimate}.
\end{proof}

\begin{lemma}\label{L:fourier-transform-window}
Let $a\in A$. Then for any $j\in \Z$,  
\[
\mathcal F({\phi}_{0,a(j-1)}-{\phi}_{1,a(j)})^2=
(\widehat{{\phi}_{0,a(j-1)}})^2-(\widehat{{\phi}_{0,a(j)}})^2.
\]
\end{lemma}

\begin{proof}
We have
\begin{equation}\label{E:square-fourier-transform}
\mathcal F({\phi}_{0,a(j-1)}-{\phi}_{1,a(j)})^2=
(\widehat{{\phi}_{0,a(j-1)}})^2 - 2 \widehat{{\phi}_{0,a(j-1)}} \widehat{{\phi}_{1,a(j)}} +(\widehat{{\phi}_{1,a(j)}})^2.
\end{equation}
Using~\eqref{ft_window_eq_zero},~\eqref{ft_window_eq_one} and the fact that $2a(j-1) \leq a(j)$, we see that $\widehat{{\phi}_{0,a(j-1)}}(\xi)=1$ whenever $\widehat{{\phi}_{1,a(j)}}(\xi) \neq 0$. Thus, the right-hand side of~\eqref{E:square-fourier-transform} can be rewritten as
\[
(\widehat{{\phi}_{0,a(j-1)}})^2 +(1-\widehat{{\phi}_{1,a(j)}})^2-1.
\]
The conclusion follows by applying~\eqref{windowsum}.
\end{proof}

\begin{lemma} \label{lem:scaleest}
    For any $x\in \R$ and $\lambda, s>0$,
    \[\langle x \rangle^2_{(\lambda s)} \le \max\{\lambda,\lambda^{-1}\}\langle x \rangle^2_{(s)}\]
\end{lemma}
\begin{proof}
    Let $r=|x|/s$. Then
    \[\frac{\langle x \rangle^2_{(\lambda s)} }{\langle x \rangle^2_{(s)}} = \lambda^{-1}\Big(\frac{1+r}{1+\frac{r}{\lambda}} \Big)^2\]
    If $\lambda\ge 1$, then
    \[1+r\le \lambda\Big(1+\frac{r}{\lambda} \Big),\]
    so the quotient is at most $\lambda$. If $0<\lambda <1$, then
    \[1+\frac{r}{\lambda}\ge 1+ r,\]
    so the quotient is at most $\lambda^{-1}$.
\end{proof}

\begin{proposition}[induct positive terms - reduction variant, non-Whitney]
\label{P:induct-positive-terms-reduction-non-whitney}
Let $a\in A$ and let
\[M_j=({\phi}_{0,a(j-1)}-{\phi}_{1,a(j)})^{\otimes 2}.\]
Then with $\M=(M_j)_{j\in \Z}$,  \[\|\M\|_{\rm M(1)}\le C_{\ref{P:induct-positive-terms-reduction-non-whitney}},\]
where $C_{\ref{P:induct-positive-terms-reduction-non-whitney}}=C_{\ref{P:one-scale-estimate-window}}+ 64 \max(1,(c/2\pi)^4) C_{\ref{lem:smoothdecay2},2}^2C_{\ref{P:vanishing-diagonal-reduction}}$.
\end{proposition}

\begin{proof}
For $j\in \Z$, we write $M_j=M_{0,j}+M_{1,j}$, where
\[
M_{0,j}= ({\phi}_{0,a(j-1)})^{\otimes 2} -({\phi}_{0,a(j)})^{\otimes 2}
\]
and
\[
M_{1,j}=({\phi}_{0,a(j-1)}-{\phi}_{1,a(j)})^{\otimes 2} + ({\phi}_{0,a(j)})^{\otimes 2} -({\phi}_{0,a(j-1)})^{\otimes 2}.
\]
For $i\in [2)$, we also set $\mathbf{M}_i=(M_{i,j})_{j\in \Z}$. By Proposition~\ref{P:one-scale-estimate-window},
\[
\|\mathbf{M}_0\|_{\rm M(1)} \leq C_{\ref{P:one-scale-estimate-window}}.
\]

We proceed by estimating $\mathbf{M}_1$. 
Using Lemma~\ref{L:fourier-transform-window} and the fact that $\phi_0$ is an even function, we deduce that for every $\xi \in \R$,
\[
\widehat{M_{1,j}}(\xi,-\xi)=0.
\]
Let $b=a/4$, then $b\in A$ by Proposition~\ref{Operations on spaced sequences}.
If $(\xi,\eta) \in \R^2$ is such that $\widehat{M_{1,j}}(\xi,\eta) \neq 0$ then $|\xi+\eta| \leq |\xi|+|\eta| \leq 2(a(j-1))^{-1}=2^{-1} (b(j-1))^{-1}$.

By Lemma~\ref{lem:smoothdecay2}, we get for $i\in [2)$ and $x\in \R$,
\[
|\phi_i(x)| \leq 2 C_{\ref{lem:smoothdecay2},2} \max(1,(c/2\pi)^2) \langle x \rangle^2.
\]
Thus, for $v\in \R^2$,
\[
|M_{1,j}(v)|
\leq 4 C_{\ref{lem:smoothdecay2},2}^2 \max(1,(c/2\pi)^4)\Big[
(\langle v_0 \rangle^2_{(a(j-1))}+\langle v_0 \rangle^2_{(a(j))})
(\langle v_1 \rangle^2_{(a(j-1))}+\langle v_1 \rangle^2_{(a(j))})
\]
\[
+\langle v_0 \rangle^2_{a(j)} \langle v_1 \rangle^2_{(a(j))}+\langle v_0 \rangle^2_{(a(j-1))} \langle v_1 \rangle^2_{(a(j-1))} \Big].
\]
Applying Lemma~\ref{lem:scaleest} and lowering all exponents to $3/2$, this is bounded by
\[
64 C_{\ref{lem:smoothdecay2},2}^2 \max(1,(c/2\pi)^4)
\sum_{(t_0,t_1,u) \in \mathcal P} \langle (W_u v)_0 \rangle^{3/2}_{(t_0(j))} \langle (W_u v)_1 \rangle^{3/2}_{(t_1(j))},
\]
where $\mathcal P$ consists of all triples of the form $(t_0,t_1,0)$, where for each $i\in [2)$ either $t_i(j)=b(j-1)$ for all $j\in \Z$, or $t_i(j)=b(j)$ for all $j\in \Z$. In particular, $\# \mathcal P=4$. Applying Proposition~\ref{P:vanishing-diagonal-reduction}, we thus get
\[
\|\mathbf{M}_1\|_{\rm M(1)} \leq 64 \max(1,(c/2\pi)^4) C_{\ref{lem:smoothdecay2},2}^2C_{\ref{P:vanishing-diagonal-reduction}}.
\]
\end{proof}

 

\begin{proposition}[induct positive terms - reduction variant, non-Whitney, skip terms]
\label{P:induct-positive-terms-reduction-non-whitney-skip}
Let $a\in A$ and let
\[\tilde{M}_j=({\phi}_{0,a(2j)}-{\phi}_{1,a(2j+1)})^{\otimes 2}.\]
Then with $\tilde{\M}=(\tilde{M}_j)_{j\in \Z}$, 
\begin{equation}\label{E:est-skipped-terms}
\|\tilde{\M}\|_{\rm M(1)}\le C_{\ref{P:induct-positive-terms-reduction-non-whitney-skip}},
\end{equation}
where $C_{\ref{P:induct-positive-terms-reduction-non-whitney-skip}}= 2^{1-2^{2-n}} C_{\ref{P:induct-positive-terms-reduction-non-whitney}}$.
\end{proposition}



\begin{proof}
Let $\gamma=(1,0,(a,a))$, then $\gamma \in \Gamma$. For $j\in \Z$, we set
\[
X_{0,j}=({\phi}_{0,a(2j-1)}-{\phi}_{1,a(2j)})^{\otimes 2}.
\]
Then $X=(X_{i,j})_{i\in [1),j\in \Z} \in \mathcal X_1$. By Proposition~\ref{positive terms}, we deduce that for any $j\in \Z$ and $\mathbf{F} \in \mathfrak{F}$,
\begin{equation}\label{E:intermediate-terms-nonnegative}
\Lambda_1(X_{0,j})(\mathbf{F}) \geq 0.
\end{equation}
Let $(M_j)_{j\in \Z}$ be as in Proposition~\ref{P:induct-positive-terms-reduction-non-whitney} and let $J \in \N_{>0}$. By~\eqref{E:intermediate-terms-nonnegative} and Proposition~\ref{P:induct-positive-terms-reduction-non-whitney}, we get that for every $\mathbf{F} \in \mathfrak{F}$,
\[
\Lambda_1(\sum_{j\in [J)} \tilde{M}_j)(\mathbf{F})
\leq \Lambda_1(\sum_{j\in [2J)} M_j)(\mathbf{F})
\leq C_{\ref{P:induct-positive-terms-reduction-non-whitney}} 2^{1-2^{2-n}} J^{1-2^{2-n}}.
\]
This yields~\eqref{E:est-skipped-terms}.
\end{proof}
 
 
\begin{proposition}[induct positive terms - reduction variant, Whitney, with gap]
\label{P:induct-positive-terms-reduction-whitney-gap}
Let $a\in {\rm A}$ satisfy   $a(j)\geq 2^{51}a(j-1)$
for every $j\in\mathbb Z$. Let $M_j:\mathbb{R}^2\to \mathbb{R}$ be of the form
\[ M_j=\Psi_{(a(j))} \]
for every $j\in\Z$, where $\Psi:\mathbb{R}^2\to \mathbb{R}$ is a test function satisfying $\Psi(v_0,v_1)=\Psi(v_1,v_0)$ for all $v=(v_0,v_1)\in\mathbb{R}^2$,
\begin{equation}
    \label{p1-on-W_pos}
    \int_{\mathbb{R}^2} g(v_0){g(v_1)}\Psi(u)\,du\ge 0
\end{equation}
for all bounded measurable $g:\mathbb{R}\to\mathbb{R}$,
\begin{equation}
    \label{p1-on-W_supp}
    \mathrm{supp}(\widehat{\Psi}) \subset \mathrm{Ann}_1(1, 2^{20})^2,
\end{equation}
for $m\in [3)$ and $\xi \in \R$,
\begin{equation}
    \label{p1-on-W_der}
    \Big|\tfrac{d^m}{d\xi^{m}} \widehat{\Psi}(\xi,-\xi)\Big| \leq 1,
\end{equation}
and for all $v\in\mathbb{R}^2$,
\begin{equation}
    \label{p1-on-W_decay}
     |\Psi(v)| \le \sum_{u=0}^1 \langle(W_u v)_0\rangle^{3/2} \langle (W_u v)_1\rangle^{3/2}.
\end{equation}
Then with $\M=(M_j)_{j\in\Z}$,
\[\|\M\|_{\rm M(1)} \le C_{\ref{P:induct-positive-terms-reduction-whitney-gap}}.\]
where $C_{\ref{P:induct-positive-terms-reduction-whitney-gap}} = 2C_{\ref{prop:1-on-nonW-gap}} +2^{60}(C^2_0 + C^2_{\ref{psi_decay}}) C_{\ref{P:vanishing-diagonal-reduction}}$
\end{proposition}

\begin{proof}
 We use Lemma \ref{lem:Phij_prop} 
to obtain a function $\psi$ satisfying \eqref{Phij_psi_diag} and \eqref{psi_decay}. Define the sequence  $\alpha$  as $\alpha(2j)=2^{-25}a(j)$ and $\alpha(2j+1)=2^{25}a(j)$. Then $\alpha \in A$. 
  By  \eqref{Phij_psi_diag}, for each $j\in \Z$, 
\begin{equation}
    \label{psi_diag}
    \widehat{\psi_{(a(j))}}(\xi)^2 = 2 (\widehat{\phi_{0,\alpha(2j)}}-\widehat{\phi_{1,\alpha(2j+1)}})(\xi)^2 - \widehat{\Psi_{(a(j))}}(\xi,-\xi). 
\end{equation}
Decompose  $M_j=M_{0,j}-M_{1,j}$ with
\[M_{0,j} =   \Psi_{(a(j))} + (\psi_{(a(j))})^{\otimes 2}  , \]
\[M_{1,j} =   (\psi_{(a(j))})^{\otimes 2} .  \]
  Let $J\ge 1,\, M = \sum_{j\in[J)}M_j$, and $\F\in \mathfrak{F}$. 
 The condition \eqref{p1-on-W_pos}, Lemma \ref{Lambda1 to Lambda}, and Lemma  \ref{lem:form_pos}  applied term by term  give  that     \[\Lambda_1(M)(\F)\ge 0,\quad \Lambda_1(\sum_{j\in[J)}M_{1,j})(\F) \ge 0\]
Therefore, 
  \begin{equation}
    \label{K0_split}
    |\Lambda_1(M)(\F)| = \Lambda_1(M)(\F) \end{equation}
    \[=\Lambda_1(\sum_{j\in[J)}M_{0,j})(\F)-\Lambda_1(\sum_{j\in[J)}M_{1,j})(\F)\le \Lambda_1(\sum_{j\in[J)}M_{0,j})(\F).\]

We now write $M_{0,j} = M_{2,j} + M_{3,j}$,   
where \[M_{2,j} = 2(\phi_{0,\alpha(2j)} -\phi_{1,\alpha(2j+1)} )^{\otimes 2}\] 
and   $M_{3,j} = M_{0,j} - M_{2,j}$.
Proposition \ref{P:induct-positive-terms-reduction-non-whitney}   gives  
\[\|(M_{2,j})_{j\in\Z}\|_{\rm M(1)}\le 2C_{\ref{prop:1-on-nonW-gap}}.\]
By multilinearity of $\Lambda_1$,  it thus remains to estimate  $\|(M_{3,j})_{j\in\Z}\|_{\rm M(1)}$, which will be done by an application of  Proposition~\ref{P:vanishing-diagonal-reduction}.

 We now  verify the  assumptions of Proposition~\ref{P:vanishing-diagonal-reduction}.
Evaluating the Fourier transform at $(\xi,-\xi)$, 
\[   \widehat{M_{3,j}}(\xi,-\xi)  
=  \widehat{\Psi_{(a(j))}}(\xi,-\xi) + \widehat{\psi_{(a(j))}}(\xi)^2- 2 (\widehat{\phi_{0,\alpha(2j)}} -\widehat{\phi_{1,\alpha(2j+1)}} )(\xi)^2  = 0\]
where we used \eqref{psi_diag} and the fact that $\psi$ and windows are even. 

Next we verify that if $(\xi,\eta)$ is in the support of $\widehat{M_{3,j}}$,  then  
\begin{equation}
    \label{xi_plus_eta}
    |\xi+\eta|\le 2^{29}a(j)^{-1} \le 2^{-1}a(j-1)^{-1}
\end{equation}
It suffices to prove the first inequality; the second   follows from     $a(j)\ge 2^{51}a(j-1)$. 
% Recall that we have
% \[\textup{supp}(\widehat{\Psi}) \subset \textup{Ann}_1(1,2^{20})^2 
% %\subset \textup{Ann}_2(1,2^{21}) 
% \]
Since $\textup{supp}(\widehat{\phi_{0,2^{-25}}})\subset [-2^{25},2^{25}]$ and $\widehat{\phi_{0,2^{-25}}} =1$ on $[-2^{24},2^{24}]$ and 
 $\textup{supp}(\widehat{\phi_{1,2^{25}}})\subset [-2^{-25},2^{-25}]$ and $\widehat{\phi_{1,2^{25}}} =1$ on $[-2^{-26},2^{-26}]$, 
 % \[\textup{supp}(\widehat{\phi_{0,2^{-25}}} -\widehat{\phi_{1,2^{25}}} ) \subset \textup{Ann}_1(1,2^{26})\]
 % and thus 
\[\textup{supp}((\widehat{\phi_{0,2^{-25}}} -\widehat{\phi_{1,2^{25}}} )^{\otimes 2} ) \subset \textup{Ann}_1(1,2^{26})^2  
%\subset  \textup{Ann}_2(1,2^{27})
\]
By Lemma \ref{lem:Phij_prop}, also 
\[\textup{supp}(\widehat{\psi}^{\otimes 2})\subset \textup{Ann}_1(1,2^{27})^2 
%\subset \textup{Ann}_2(1,2^{28})
\]
Therefore, together with \eqref{p1-on-W_supp} we obtain that  the function \[ {\Psi} + \psi^{\otimes 2} - 2({\phi_{0,2^{-25}}} -{\phi_{1,2^{25}}} )^{\otimes 2}  \] has Fourier support in $\textup{Ann}_1(1,2^{28})^2$. 
Rescaling by  $a(j)$   we obtain  \eqref{xi_plus_eta}.

Now we examine the decay of  $M_{3,j}$.    Using \eqref{windowbound} and  Lemma \ref{lem:smoothdecay2} with $N=2$ 
gives for $j=0,1$, $|\phi_j(u_0)|  \le  2^3 C_0 \langle u_0 \rangle^{2}$ and thus by distributivity and lemma \ref{lem:scaleest}, 
\[|(\phi_{0,2^{-25}}-\phi_{1,2^{25}})^{\otimes 2}(v)| \le 2^6C^2_0 \sum_{\substack{t_0,t_1\in \{2^{-25}, 2^{25}\}}}\langle v_0\rangle^{2}_{(t_0)} \langle v_1\rangle^{2}_{(t_1)} \le 2^{58}C^2_0  \langle v_0\rangle^{2}  \langle v_1\rangle^{2}   \]
By Lemma \ref{lem:smoothdecay2} with $N=2$, using the derivative bound \eqref{psi_decay} and   $|\mathrm{supp}\;\widehat{\psi}|\le 2^{28}$,   
%and thus
%\[|(\phi_0-\phi_1)^{\otimes 2}(v)| \le 2^8C^2_0 \langle v_0\rangle^{2} \langle v_1\rangle^{2} \]
%By Lemma \ref{lem:smoothdecay2} again  with $N=2$,  
\begin{equation}
    \label{decay_2}
  |\psi^{\otimes 2}(v)| \le 2^{60} C^2_{\ref{psi_decay}} \langle v_0\rangle^{2} \langle v_1\rangle^{2} 
  \end{equation}
  Together with the decay assumption \eqref{p1-on-W_decay}, rescaling, and $\langle\cdot\rangle^2 \le \langle\cdot\rangle^{3/2}$ we obtain
 \[c^{-1}|M_{3,j}(v)| \le   \langle (W_1 v)_0\rangle_{a(j)}^{3/2} \langle (W_1 v)_1\rangle_{a(j)}^{3/2}+\langle v_0\rangle_{a(j)}^{3/2}  \langle v_1\rangle_{a(j)}^{3/2}   .\]
with $c =  2^{60}(C^2_0 + C^2_{\ref{psi_decay}})$ (one can in fact  take $1 + 2^{59}C_0^2 + 2^{60}C^2_{\ref{psi_decay}}$).
Applying Proposition~\ref{P:vanishing-diagonal-reduction}  to $c^{-1}M_{3,j}$ yields
\[\|(M_{3,j})_{j\in\Z}\|_{\rm M(1)} \le   2^{60}(C^2_0 + C^2_{\ref{psi_decay}}) C_{\ref{P:vanishing-diagonal-reduction}},\]
which completes the proof.
\end{proof}

% \pd{\[|\phi(x)|\leq 2^{N} \max(\|\widehat{\phi}\|_\infty, (2\pi )^{-N}\|\widehat{\phi}^{(N)}\|_\infty)|\supp (\widehat{\phi})|   \langle x\rangle^{N}
% \]
% with $C_{\ref{lem:smoothdecay2}, N}=2^{N} $. }
 

\begin{proposition}[induct positive terms - reduction variant, Whitney]
\label{P:induct-positive-terms-reduction-whitney}
Let $a\in {\rm A}$. Let $M_j:\mathbb{R}^2\to \mathbb{R}$ be of the form
\[ M_j=\Psi_{(a(j))} \]
for every $j\in\Z$, where $\Psi:\mathbb{R}^2\to \mathbb{R}$ is a test function satisfying $\Psi(u_0,u_1)=\Psi(u_1,u_0)$ for all $u=(u_0,u_1)\in\mathbb{R}^2$, and \eqref{p1-on-W_pos}, \eqref{p1-on-W_supp}, \eqref{p1-on-W_der}, \eqref{p1-on-W_decay}
Then with $\M=(M_j)_{j\in\Z}$,
\[\|\M\|_{\rm M(1)} \le C_{\ref{P:induct-positive-terms-reduction-whitney}},\]
where $C_{\ref{P:induct-positive-terms-reduction-whitney}} = 2^6C_{\ref{P:induct-positive-terms-reduction-whitney-gap}}$.
\end{proposition}

\begin{proof}
    For \(s\in[51)\), define
\[
a^{(s)}(k):=a(51k+s),
\qquad
M^{(s)}_k:=M_{51k+s},
\qquad k\in\mathbb Z.
\]
Since \(a\in A\),
\[
a^{(s)}(k)
=a(51k+s)
\geq 2^{51}a(51(k-1)+s)
=2^{51}a^{(s)}(k-1),
\]
and we have
\[
M^{(s)}_k=\Psi_{(a^{(s)}(k))}.
\]
Fix \(J\geq1\) and \(\F\in\mathfrak F\).  For
\(s\in[51)\), let
\[
J_s
:=
\#\{j\in[J):j\equiv s\pmod {51}\}.
\]
Then
\[
\sum_{\substack{j\in[J)\\j\equiv s\;({\rm mod}\;51)}}M_j
=
\sum_{k\in[J_s)}M^{(s)}_k.
\]
By linearity in
the kernel and the triangle inequality,  
\[
|
\Lambda_1(\sum_{j\in[J)}M_j)(\F)
|
\leq
\sum_{s\in[51)}
|
\Lambda_1(
\sum_{k\in[J_s)}M^{(s)}_k
)(\F)
|\]
Proposition \ref{P:induct-positive-terms-reduction-whitney-gap} applied for each $s\in [51)$ and finally summing in $s$ give
\[\|\mathbf{M}\|_{\rm M(1)}\leq 51C_{\ref{P:induct-positive-terms-reduction-whitney-gap}}
\leq C_{\ref{P:induct-positive-terms-reduction-whitney}}.
\]
% Taking the supremum over $J$ and $\F$ then gives the claim. 
\end{proof}

\begin{proposition}[induct positive terms - reduction variant, Whitney, product-type]
\label{P:induct-positive-terms-reduction-whitney-product}
Let $a\in A$ and for $j\in \Z$ let $M_j:\mathbb{R}^2\to \mathbb{R}$  be given by
\[ M_j=   (\psi_{(a(j))})^{\otimes 2}, \]
 where   $\psi:\mathbb{R}\to \mathbb{R}$ satisfies
%there is a strictly increasing sequence of integers $(k_j)_{j\in [J)}$ such that for all $j\in [J)$
\begin{equation}\label{1-on-W_supp}
\mathrm{supp}(\widehat{\psi}) \subset \mathrm{Ann}_1(1, 2^{20})
\end{equation}
and for $m\in [3)$,
\begin{equation}\label{E:derivative-assumption}
\|\widehat{\psi}^{(m)}\|_\infty \leq 1.
\end{equation}
Then for $\M=(M_j)_{j\in \Z}$,
\[ \|\M\|_{\rm M(1)}\le C_{\ref{P:induct-positive-terms-reduction-whitney-product}}.  \]
\end{proposition}

\begin{proof}
We apply Proposition \ref{P:induct-positive-terms-reduction-whitney} with $\Psi = 2^{-46}\psi\otimes \psi$.  Then 
\[\Psi(v_0,v_1) = 2^{-46}\psi(v_0)\psi(v_1) = \Psi(v_1,v_0),\] 
the positivity assumption \eqref{p1-on-W_pos} is satisfied by  \eqref{form_pos_psipsi} from   Lemma \ref{lem:form_pos}, and 
the support containment \eqref{p1-on-W_supp} follows immediately from \eqref{1-on-W_supp}. 
The derivative bound \eqref{E:derivative-assumption} implies
\[\tfrac{d^m}{d\xi^m}|\widehat{\Psi}(\xi,-\xi)| \le 1,  \]
so \eqref{p1-on-W_der} holds.
  By Lemma~\ref{lem:smoothdecay2} applies with $N=2$, 
$\psi\le 2^{23}\langle \cdot \rangle^2$ and thus
\[\Psi(v_0,v_1) = 2^{-46}\psi(v_0)\psi(v_1) \le \langle v_0 \rangle^2 \langle v_1\rangle^2\]
which gives \eqref{p1-on-W_decay}. 
\end{proof}



 

% \begin{lemma}
% Let $\kappa \in \R$, $\kappa \geq 0$ and $x\in \R$. Then
% \[
% 2^{-\kappa/2} \langle x\rangle^{2}_{(2^{\kappa})} \leq \langle x\rangle^{3/2}.
% \]
% \end{lemma}

% \ls{This lemma is no longer needed due to the reformulation of bump estimates in Section~\ref{sec:bump red}.}

% \begin{proof}
% Using $2^\kappa \ge 1$ we estimate
% \[(2^\kappa+|x|)^2=(2^\kappa+|x|)^{1/2}(2^\kappa+|x|)^{3/2} \ge 2^{\kappa/2}(1+|x|)^{3/2},\]
% This yields
% \[2^{-\kappa/2} \langle x\rangle^{2}_{(2^{\kappa})}  = 2^{\kappa/2}(2^\kappa+|x|)^{-2} \le 2^{\kappa/2}2^{-\kappa/2}(1+|x|)^{-3/2} = \langle x\rangle^{3/2}\]
% \end{proof}
% \pd{TODO: use this in the next section}




\subsection{Final reduction: proof of main theorem}\label{sec:pfmainthm}
%\jr{refactor: e.g. $\Lambda_n$ should now be $\Lambda_1$ etc.} 
In this section we prove Theorem \ref{thm:nct main real} using Propositions \ref{prop:1-on-W}, \ref{prop:1-on-W-cor} and \ref{prop:1-on-nonW}.
Let $n\ge 2$. 
Let $\mathbf{f}=(f_0,\dots,f_{n-1})$ be Schwartz functions. % $\Lambda_1$ as in \eqref{prism Brascamp--Lieb} and $A$ as in \eqref{A_def}.
Throughout this section we abbreviate
\begin{equation}
A(\chi) = A(\chi,\mathbf{f}),
\end{equation}
since $\mathbf{f}$ will not change.
We will also write
\begin{equation}
    \label{avg_rescale}
    A_t(\chi) = A(\chi_{(t)}),
\end{equation}


Let $(\phi_0, \phi_1)$ be a universal pair
and let $\theta, \varphi_{0,k}, \varphi_{1,k}, \varphi_{2,k}, \varphi_{3,k}, \varphi_{4,k}$ be defined as in \eqref{eqn:thetadef}-\eqref{eqn:varphi4kdef}.
In what follows we denote 
\[\alpha(n) = 1-2^{-n+2}.\]

\begin{lemma}\label{lem:main_aux1}
Suppose that $\psi:\R\to\R$ is a Schwartz function with
\begin{equation}\label{main_aux1_supp}
{\rm supp}(\widehat{\psi})\subset \mathrm{Ann}_1(1, 2^{19}),
\end{equation}
\begin{equation}
\label{main_aux1_decay}
|\psi (u)| \le  \langle u\rangle^{2N_{\ref{prop:1-on-W}}}.
\end{equation}
Then for every $t\in [1,2]$, $J\ge 1$, and every increasing sequence of integers $(k_j)_{j\in [J)}$,
\begin{equation}\label{main_aux1} 
\sum_{j \in [J)} \|A_{2^{k_j}t}(\psi)\|_2^2 \le C_{\ref{lem:main_aux1}} J^{\alpha(n)},
\end{equation}
where $C_{\ref{lem:main_aux1}}=C_{\ref{prop:1-on-W}}$.
\end{lemma}
% \begin{rem}
% If $t=1$, the constant $2^{40}$ can be omitted.
% \end{rem}

\begin{proof}
By Lemma \ref{A to Lambda} and linearity the left-hand side equals
\[ \Lambda_1(M_t)(\F), \]
with $\F$   as in Lemma \ref{A to Lambda} and 
\[ M_t = \sum_{j\in [J)} (\psi_{(2^{k_j} t)})^{\otimes 2} \]
The claim follows from Proposition \ref{prop:1-on-W-cor} applied with the function $2^{-2N_{\ref{prop:1-on-W}}} \psi_{(t)}$.  
Then \eqref{1-on-W_supp} follows from \eqref{main_aux1_supp} using $t\in [1,2]$. 
We verify \eqref{1-on-W_decay} using \eqref{main_aux1_decay} and $t\in[1,2]$:
\[ 2^{-2N_{\ref{prop:1-on-W}}} |\psi_{(t)}(u)| \le 2^{-2N_{\ref{prop:1-on-W}}} \langle u\rangle_{(t)}^{2N_{\ref{prop:1-on-W}}} \le \langle u \rangle^{2N_{\ref{prop:1-on-W}}}. \]
\end{proof}

%\jr{For now only $n=3$}

% \begin{lemma}[Long and short variation OLD]\label{lem:shortlongftc_reduction}
% Let $J\ge 1$ and $\phi$ a Schwartz function and $\mathbf{f}=(f_0,\dots,f_{n-1})$ Schwartz functions.
% Suppose that $A\in (0,\infty)$ is such that
% for every increasing sequence of integers $(k_i)_{i=0}^J$,
% % \[ \sum_{i=0}^J \int_1^2 \| A_{2^{k_i}t}(T\phi) \|_2^2 \frac{dt}{t} \le A \]
% \[ \int_1^2 \|A_{2^{k_i}t}(T\phi)\|_{\ell^2(i\in [J+1]; L^2)}^2\,\frac{dt}{t} \le A \]
% where $T\phi(x) = (x \phi(x))'$ and
% \[ \|A_{t}(\phi))\|^2_{V_{2,J}(t\in 2^\mathbb{Z}; L^2)} \le A \]
% % \[ \sum_{i=0}^J \| A_{2^{k_i}}(\phi) - A_{2^{k_{i-1}}}(\phi) \|_2^2 \le A. \]
% Then
% \[ \|A_{t}(\phi))\|_{V_{2,J}(t\in \mathbb{R}; L^2)}^2 \le 10 A. \]
% \end{lemma}

% \begin{proof}
% This is by Lemma \ref{lem:shortlongjumps} and Lemma \ref{lem:ftccs} 
% \todo{insert details}
% \end{proof}


\begin{lemma}[Long and short variation]\label{lem:shortlongftc_reduction}
Let $J\ge 1$ and $\phi$ a Schwartz function.
Assume that the function $T\phi(s) = (s \phi(s))'$ satisfies \eqref{main_aux1_supp} and \eqref{main_aux1_decay}.
Suppose that $A\in (0,\infty)$ is such that
\begin{equation}\label{shortlongftc_reduction1} \|A_{t}(\phi)\|^2_{V_{2,J}(t\in 2^\mathbb{Z}; L^2)} \le A 
\end{equation} 
% \[ \sum_{i=0}^J \| A_{2^{k_i}}(\phi) - A_{2^{k_{i-1}}}(\phi) \|_2^2 \le A. \]
Then
\begin{equation}\label{shortlongftc_reduction2}
\|A_{t}(\phi)\|_{V_{2,J}(t\in \mathbb{R}; L^2)}^2 \le 2^4 C_{\ref{lem:main_aux1}} J^{\alpha(n)} + 2A.
\end{equation}
\end{lemma}

\begin{proof}
Let $(t_j)_{j\in [J+1)}$ be an increasing sequence of real numbers.
By Lemma \ref{lem:shortlongjumps} (with $r=2$, $a(t) = A_t(\phi)$ and $B=L^2$),
\begin{equation}\label{shortlongftc_reduction_pf1} \Big(\sum_{j\in [J)} \|A_{t_{j+1}}(\phi)- A_{t_j}(\phi)\|_2^2\Big)^{1/2} \le 2 \Big(\sum_{k\in \kappa} \|A_t(\phi)\|_{V_{2,J}(t\in [2^k, 2^{k+1}]; L^2)}^2\Big)^{1/2}
\end{equation}
\begin{equation}\label{shortlongftc_reduction_pf2}
+\|A_t(\phi)\|_{V_{2,J}(t\in 2^\mathbb{Z}; L^2)}.
\end{equation}
where $\kappa$ is the set of $k\in \mathbb{Z}$ such that $2^{k}\le t_j<2^{k+1}$ for some $j\in [J+1)$.
 Fix $x\in\mathbb{R}^n$ and $k\in\mathbb{Z}$. Set $a(t)=A_t(\phi)(x)$ and use Lemma \ref{lem:ftccs-R} (specifically, \eqref{ftccs1-R}) to estimate
 \[ \|A_t(\phi)(x)\|_{V_{2,J}(t\in [2^k, 2^{k+1}])}^2 \le 
 \|ta'(t)\|_{L^2(t\in [2^k, 2^{k+1}],\frac{dt}{t})}^2.  \]
% Note that
% \begin{equation}\label{shortlongftc_reduction_pf3}
% t\frac{\partial}{\partial t}(\phi_{(t)}(s)) = -(T\phi)_{(t)}(s),
% \end{equation}
% where $T\phi(s) = (s \phi(s))'$.
% For each $k\in\mathbb{Z}$ we estimate
% \[ \|A_t(\phi)\|_{V_{2,J}(t\in [2^k, 2^{k+1}]; L^2)}^2 \le \int_{\mathbb{R}^n}  \|A_t(\phi)(x)\|_{V_{2,J}(t\in [2^k, 2^{k+1}])}^2\, dx.\]
% Fix $x\in\mathbb{R}^n$ and $k\in\mathbb{Z}$. Set $a(t)=A_t(\phi)(x)$ and use Lemma \ref{lem:ftccs-R} (specifically, \eqref{ftccs1-R}) to estimate
% \[ \|A_t(\phi)(x)\|_{V_{2,J}(t\in [2^k, 2^{k+1}])}^2 \le 
% \|ta'(t)\|_{L^2(t\in [2^k, 2^{k+1}],\frac{dt}{t})}^2.  \]
% By \eqref{A_def} and \eqref{shortlongftc_reduction_pf3},
% \[ ta'(t) = t\frac{\partial}{\partial t}\int_{\mathbb{R}} \Big(\prod_{i\in [n)} f_i(x+se_i)\Big)\phi_{(t)}(s)\,ds = -A_t({T\phi})(x). \]
% From this, and changing variables $t \mapsto 2^k t$ in the $t$-integration,
By Lemma~\ref{lem:ftc_ATphi}, 
\[ \|ta'(t)\|_{L^2(t\in [2^k, 2^{k+1}],\frac{dt}{t})}^2 = \int_{1}^2 |A_{2^k t}(T\phi)(x)|^2\,\frac{dt}{t}. \]
Thus for every $x\in\mathbb{R}^2$ and $k\in\mathbb{Z}$, 
\[ \|A_t(\phi)(x)\|_{V_{2,J}(t\in [2^k, 2^{k+1}])}^2 \le\int_{1}^2 |A_{2^k t}(T\phi)(x)|^2\,\frac{dt}{t}.  \]
Integrating over $x\in\mathbb{R}^2$ and interchanging the integrals over $x$ and $t$ gives
\[ \|A_t(\phi)\|_{V_{2,J}(t\in [2^k, 2^{k+1}]; L^2)}^2 \le \int_1^2 \|A_{2^k t}(T\phi)\|_2^2 \frac{dt}{t}. \]
Using this to estimate the sum on the right in \eqref{shortlongftc_reduction_pf1} and interchanging the summation over $k$ with the integral over $t$,
\[
\sum_{k\in \kappa} \|A_t(\phi)\|_{V_{2,J}(t\in [2^k, 2^{k+1}]; L^2)}^2 \le \int_1^2 \sum_{k\in\kappa} \|A_{2^k t}(T\phi)\|_2^2\,\frac{dt}{t}.
\]
For each fixed $t\in [1,2]$ we use Lemma \ref{lem:main_aux1} (as $T\phi$ satisfies \eqref{main_aux1_supp} and \eqref{main_aux1_decay} by assumption) and that $\#\kappa\le J+1$ to obtain
\[ \sum_{k\in\kappa} \|A_{2^k t}(T\phi)\|_2^2 \le 2 C_{\ref{lem:main_aux1}} J^{\alpha(n)}. \]
(here we have crudely estimated $\#\kappa^\frac12 \le (J+1)^{\alpha(n)} \le 2 J^{\alpha(n)}$).
Now the claim follows from
\[ \sum_{j\in [J)} \|A_{t_{j+1}}(\phi)- A_{t_j}(\phi)\|_2^2 \le 8 \sum_{k\in \kappa} \|A_t(\phi)\|_{V_{2,J}(t\in [2^k, 2^{k+1}]; L^2)}^2 + 2 \|A_t(\phi)\|_{V_{2,J}(t\in 2^\mathbb{Z}; L^2)}^2 \]
\[ \le 2^4 C_{\ref{lem:main_aux1}} J^{\alpha(n)} + 2A.  \]
\end{proof}



% \begin{lemma}\label{lem:mainbump1_short}[NOW OBSOLETE]
% \begin{equation}\label{mainbump1_short} 
% \|A_{2^{k_i}t}(T\chi)\|_{\ell^2(i\in [J+1]; L^2(t\in [1,2],\frac{dt}{t}; L^2))}^2 \lesssim J^{\alpha(n)} ,
% \end{equation}
% \end{lemma}

% \begin{proof}
% The left-hand side equals
% \[ \int_1^2 \|A_{2^{k_i}t}(T\chi)\|_{\ell^2(i\in [J+1]; L^2)}^2 \frac{dt}t \]
% Apply Lemma \ref{lem:main_aux1} to $\phi=T\chi$.

% \end{proof}


\begin{lemma}\label{lem:mainbump1_long1}
Then for every $J\ge 1$ and every strictly increasing sequence of integers $(k_j)_{j\in[J)}$, 
\begin{equation}\label{mainbump1_long1} 
\sum_{j \in [J)} \|A_{2^{k_j}}(\phi_0) - A_{2^{k_{j}}}(\phi_1)\|_{2}^2  \le C_{\ref{lem:mainbump1_long1}} J^{\alpha(n)}
\end{equation}
where $C_{\ref{lem:mainbump1_long1}} = (2C_0C_{\ref{lem:smoothdecay2}, 2N_{\ref{prop:1-on-W}}})^2 C_{\ref{lem:main_aux1}}$.
\end{lemma}

\begin{proof}
This follows by 
Lemma \ref{lem:main_aux1} (with $t=1$ and $c_1^{-1}(\phi_0 - \phi_1)$ in place of $\phi$, where $c_1=2C_0C_{\ref{lem:smoothdecay2},2N_{\ref{prop:1-on-W}}}$), after rewriting the left-hand side  as
\[\sum_{j\in [J)} \|A_{2^{k_{j}}}(\phi_0-\phi_1)\|_{2}^2 = c_1^{2}\sum_{j\in[J)} \|A_{2^{k_j}}(c_1^{-1}(\phi_0-\phi_1))\|_{2}^2  \] 
and verifying  
 \eqref{main_aux1_supp} and \eqref{main_aux1_decay}, i.e.
\[{\rm supp}\big(c_1^{-1}(\widehat{\phi_0}- \widehat{\phi_1})\big)\subset \mathrm{Ann}_1(1, 2^{19}),\]
\begin{equation}
    \label{mainbump1_long_decay}
    c_1^{-1}\big|\phi_0(u) - \phi_1(u)\big| \le  \langle u \rangle^{2N_{\ref{prop:1-on-W}}}.
\end{equation}
Since $\widehat{\phi_1}$ and $\widehat{\phi_0}$ are both supported in $[-1,1]$ and constant one on $[-2^{-1},2^{-1}]$, the function $\widehat{\phi_0}- \widehat{\phi_1}$ is supported in $\mathrm{Ann}_1(1,2)\subset \mathrm{Ann}_1(1,2^{19})$ and the same holds for any constant multiple of it. To verify the decay assumption, we use  
\eqref{windowbound}, which implies
\[\|\widehat{\phi_0}^{(2N_{\ref{prop:1-on-W}})} - \widehat{\phi_1}^{(2N_{\ref{prop:1-on-W}})}\|_\infty \leq 2C_0  \ .\]
We also have $\|\widehat{\phi_0} - \widehat{\phi_1}\|_\infty \le 2 \le 2C_0$ and $\textup{supp}(\widehat{\phi_0}-\widehat{\phi_1})\subset [-1,1]$. 
 Lemma \ref{lem:smoothdecay2} applied with $\phi$ replaced by $(2C_0)^{-1}(\phi_0-\phi_1)$ and $M=2N_{\ref{prop:1-on-W}}$ and using that  $|\supp(\widehat{\phi_1})|=2$   gives 
\[\big|\phi_0(u) - \phi_1(u)\big| \le   4 C_0 C_{\ref{lem:smoothdecay2},2N_{\ref{prop:1-on-W}}} \langle u \rangle^{2N_{\ref{prop:1-on-W}}}, \]
which implies  \eqref{mainbump1_long_decay}. 
\end{proof}

\begin{lemma}\label{lem:mainbump1_long2}
Then for every $J\ge 1$ and every strictly increasing sequence of integers $(m_j)_{j\in[J+1]}$,
\begin{equation}\label{mainbump1_long2} 
\sum_{j\in[J)} \|
A_{2^{m_{j}}}(\phi_0)-A_{2^{m_{j+1}}}(\phi_1) \|_{2}^2  \le C_{\ref{lem:mainbump1_long2}} J^{\alpha(n)},\end{equation}
where $C_{\ref{lem:mainbump1_long2}}=2 C_{\ref{prop:1-on-nonW}}$.
\end{lemma}

\begin{proof}
By Lemma \ref{A to Lambda}applied to the function $\phi_{0,m_{j+1}}-\phi_{1,m_j}$ and linearity, 
the left-hand side equals $\Lambda_1(M)(\F)$ with $\F$ as in Lemma \ref{A to Lambda}, where 
\[ M = \sum_{j\in[J)} (\phi_{0,m_{j}}-\phi_{1,m_{j+1}})^{\otimes 2}. \]
We may assume without loss of generality that $J$ is even (by adding one more jump if needed).
For $r\in [2)$, $j\in [J/2)$ set $k_{j}^r = m_{j+r}.$
Then $M=M^0+M^1$, where
\[ M^r = \sum_{j\in[J/2)} (\phi_{0,k_{2j}^r}-\phi_{1,k_{2j+1}^r})^{\otimes 2}. \]
For each $r\in [2)$ we apply Proposition \ref{prop:1-on-nonW} to obtain
\[ |\Lambda_1(M^r)(\F)| \le C_{\ref{prop:1-on-nonW}} (J/2)^{\alpha(n)}. \]
Then by the triangle inequality,
\[ |\Lambda_1(M)(\F)| \le |\Lambda_1(M^0)(\F)| + |\Lambda_1 (M^1)(\F)| \le 2 C_{\ref{prop:1-on-nonW}} J^{\alpha(n)}. \]
\end{proof}

\begin{lemma}\label{lem:mainbump1_long}
Let $J\ge 1$. Then
\begin{equation}\label{mainbump1_long} 
\|A_{t}(\phi_0)\|^2_{V_{2,J}(t\in 2^\mathbb{Z}; L^2)}  \le C_{\ref{lem:mainbump1_long}} J^{\alpha(n)},
\end{equation}
where $C_{\ref{lem:mainbump1_long}}=2(C_{\ref{lem:mainbump1_long2}} + C_{\ref{lem:mainbump1_long1}})$.
\end{lemma}

\begin{proof}
Let $(k_j)_{j\in [J+1)}$ be a strictly increasing sequence of integers.
By the triangle inequality,
\[ \Big(\sum_{j\in [J)} \|A_{2^{k_{j+1}}}(\phi_0)- A_{2^{k_j}}(\phi_0)\|_2^2\Big)^{1/2} \]
\[ \le \Big(\sum_{j\in [J)} \|A_{2^{k_{j}}}(\phi_0)- A_{2^{k_{j+1}}}(\phi_1)\|_2^2\Big)^{1/2} 
+ \Big(\sum_{j\in [J)} \|A_{2^{k_{j+1}}}(\phi_0)- A_{2^{k_{j+1}}}(\phi_1)\|_2^2\Big)^{1/2}.
\]
By Lemma \ref{lem:mainbump1_long2} and Lemma \ref{lem:mainbump1_long1} this is 
\[ \le C_{\ref{lem:mainbump1_long2}}^{\frac12} J^{\frac{\alpha(n)}{2}} + C_{\ref{lem:mainbump1_long1}}^{\frac12} J^{\frac{\alpha(n)}{2}} \le C_{\ref{lem:mainbump1_long}}^{\frac12} J^{\frac{\alpha(n)}{2}}.  \]
\end{proof}

\begin{lemma}\label{lem:mainbump1}
Let $J\ge 1$. Then
\begin{equation}\label{mainbump1}
\|A_{t}(\phi_0)\|_{V_{2,J}(t\in \mathbb{R}; L^2)}^2 \le C_{\ref{lem:mainbump1}} J^{\alpha(n)},
\end{equation}
where $C_{\ref{lem:mainbump1}}=2^6 C_0^2 C_{\ref{lem:smoothdecay2},20}^2 C_{\ref{lem:main_aux1}} + 2C_{\ref{lem:mainbump1_long}}$. 
\end{lemma}

\begin{proof}
Let $c_1=2C_0 C_{\ref{lem:smoothdecay2},20}$.
We apply Lemma \ref{lem:shortlongftc_reduction} with $\phi=c_1^{-1}\phi_0$ and $A=c_1^{-2}C_{\ref{lem:mainbump1_long}} J^{\alpha(n)}$.
The assumption \eqref{shortlongftc_reduction1} follows from Lemma \ref{lem:mainbump1_long} after multiplying by $c_1^{-2}$ on both sides.
Then \eqref{shortlongftc_reduction2} gives
\[\|A_{t}(\phi_0)\|_{V_{2,J}(t\in \mathbb{R}; L^2)}^2 \le c_1^2( 2^4 C_{\ref{lem:main_aux1}} J^{\alpha(n)} + 2A) = C_{\ref{lem:mainbump1}}J^{\alpha(n)}.\]
It remains to verify that $T(c_1^{-1}\phi_0)=c_1^{-1} (s \phi_0(s))'$ satisfies the assumptions \eqref{main_aux1_supp}, \eqref{main_aux1_decay}.

The support assumption \eqref{main_aux1_supp} follows since $\widehat{T\phi}$ has the same support as $(\widehat{\phi})'$ (by Lemma~\ref{lem:ft_deriv_mul}) and by using the assumptions \eqref{ft_window_eq_one}, \eqref{ft_window_eq_zero} on $\phi_0$.

The decay assumption \eqref{main_aux1_decay} 
follows by Lemma \ref{lem:smoothdecay2} (with $N=2N_{\ref{prop:1-on-W}}$ and $\phi=c_1^{-1}\phi_0$ and using \eqref{windowbound} with $m=2N_{\ref{prop:1-on-W}}$ and \eqref{ft_window_eq_zero} which implies $|\mathrm{supp}\,\widehat{\phi_0}|\le 2$).
\end{proof}




\begin{lemma}\label{lem:main_aux2}
Let $\psi$ satisfy \eqref{main_aux1_supp}, \eqref{main_aux1_decay},   
% \begin{equation}%\label{main_aux1_supp}
% {\rm supp}(\widehat{\psi})\subset \mathrm{Ann}_1(0, 19),
% \end{equation}
and  
% \begin{equation}
% \label{main_aux2_decay}
% |\psi (u)| + |u\psi'(u)| \le  (1+|u|)^{-20}.
% \end{equation}
\begin{equation}
\label{main_aux2_decay}
|T\psi (u)|  \le  \langle u \rangle^{2N_{\ref{prop:1-on-W}}}
\end{equation}
for all $u\in \R$.
Then
\begin{equation}\label{main_aux2} 
\|A_{t}(\psi)\|_{V_{2,J}(t\in \mathbb{R}; L^2)}^2 \le C_{\ref{lem:main_aux2}} J^{\alpha(n)},
\end{equation}
where $C_{\ref{lem:main_aux2}} = 2^5 C_{\ref{lem:main_aux1}}$. 
\end{lemma}

\begin{proof}
We  apply Lemma \ref{lem:shortlongftc_reduction} with $\phi=\psi$, so we   verify its   assumptions. Since by assumption $\psi$ satisfies \eqref{main_aux1_supp}, then also $T\psi(s) = (s \psi(s))'$  satisfies \eqref{main_aux1_supp} as $\widehat{T\psi}(\xi) = -\xi\widehat{\psi}'(\xi)$. The decay assumption on $T\psi$ of Lemma \ref{lem:shortlongftc_reduction} is satisfied by \eqref{main_aux2_decay}. 
%That $T\psi$ satisfies  \eqref{main_aux1_decay} follows from 
%$T\psi(s) = \psi(s)  + s \psi'(s)$
%and \eqref{main_aux2_decay}. 
We will show that \eqref{shortlongftc_reduction1} holds with $A=2C_{\ref{lem:main_aux1}}J^{1/2}$, i.e. 
 \begin{equation}\label{eq:main_aux2_long}
\|A_{t}(\psi)\|^2_{V_{2,J}(t\in 2^\mathbb{Z}; L^2)} \le   2^2 C_{\ref{lem:main_aux1}}J^{1/2}. 
\end{equation}
 Lemma \ref{lem:shortlongftc_reduction}   then  gives
\[\|A_{t}(\psi)\|_{V_{2,J}(t\in \mathbb{R}; L^2)}^2 \le  2^4 C_{\ref{lem:main_aux1}} J^{\alpha(n)} + 2^3 C_{\ref{lem:main_aux1}}J^{1/2} \le 2^5 C_{\ref{lem:main_aux1}}J^{1/2} \]
 as desired. 
By Lemma \ref{lem:bootstrap}, 
    \[\|A_{t}(\psi)\|^2_{V_{2,J}(t\in 2^{\Z}; L^2)} \le  4 \cdot \sup \sum_{\ell\in [J)} \|A_{2^{k_{\ell}}}(\psi)\|^2_2     \]
By assumption,   \eqref{main_aux1_supp} and \eqref{main_aux1_decay}  are satisfied for $\psi$.     Lemma \ref{lem:main_aux1} applied with $\psi$ and $t=1$  gives  \eqref{eq:main_aux2_long}.
\end{proof}                  

\begin{lemma}\label{lem:mainbump2}
For all $k\ge -2$,
\begin{equation}\label{mainbump2}
\|A_{t}(\varphi_{0,k}) \|_{V_{2,J}(t\in \mathbb{R}; L^2)}^2\le C_{\ref{lem:mainbump2}} 2^{-2k} J^{\alpha(n)}.
\end{equation}
where $C_{\ref{lem:mainbump2}}= C_{\ref{lem:main_aux2}} C_{\ref{lem:abs_deriv_ft_Tphi3_le}, 20}^2C_{\ref{lem:smoothdecay2},20}^2$.
\end{lemma}

\begin{proof}  
By Lemma \ref{lem:rescaling} (rescaling) and multiplying with $2^k$ on both sides, it suffices to show
\begin{equation}      \label{mainbump2_1}
\|A_{t}(\varphi_{3,k})\|_{V_{2,J}(t\in \mathbb{R}; L^2)}^2\le C_{\ref{lem:mainbump2}} J^{\alpha(n)},
\end{equation}
This follows from Lemma \ref{lem:main_aux2} applied with $\psi$ being $c_1^{-1}\varphi_{3,k}$, where $c_1= C_{\ref{lem:abs_deriv_ft_Tphi3_le}, 20}C_{\ref{lem:smoothdecay2},20}.$
To apply this lemma we need to verify that
\begin{equation}
    \label{mainbump2_supp}
    \textup{supp}(\widehat{\varphi_{3,k}}) \subset \textup{Ann}_1(1,2^{19}),
\end{equation}
\begin{equation}
 \label{mainbump2_dec1}
    c_1^{-1}|\varphi_{3,k}(u)| \le \langle u \rangle^{20},
\end{equation}
\begin{equation}
 \label{mainbump2_dec2}
    |T(c_1^{-1}\varphi_{3,k})(u)| = c_1^{-1}|T\varphi_{3,k}(u)| \le \langle u \rangle^{20}.
\end{equation}
We begin with  \eqref{mainbump2_supp}. 
By Lemma \ref{lem:ft_phi3_eq},
\[\widehat{\varphi_{3,k}}(\xi) = 2^k(\widehat{\mathbf{1}_{[0,1)}}-1)(2^{-k}\xi)\widehat{\chi}(2^{-k}\xi)\widehat{\theta}(\xi).\]
Since $\widehat{\theta}$ is supported in $[-1,-2^{-2}]\cup [2^{-2},1]\subset \textup{Ann}_1(1,2^{19})$, also $\widehat{\varphi_{3,k}}$
is supported in $\textup{Ann}_1(1,2^{19})$. 

To see \eqref{mainbump2_dec1}, first use \eqref{mainbump2_supp} and Lemma \ref{lem:abs_deriv_ft_phi3_le} with $m=2N_{\ref{prop:1-on-W}}$:
\[ |\widehat{\varphi_{3,k}}^{(20)}(\xi)| \le C_{\ref{lem:abs_deriv_ft_phi3_le}, 2N_{\ref{prop:1-on-W}}} \]
Then apply Lemma \ref{lem:smoothdecay2} (with $N=2N_{\ref{prop:1-on-W}}$) together with $|\textup{supp} (\widehat{\varphi_{3,k}})|\le 1$ to see
\[|\varphi_{3,k}(u)| \le C_{\ref{lem:abs_deriv_ft_phi3_le}, 2N_{\ref{prop:1-on-W}}} C_{\ref{lem:smoothdecay2},2N_{\ref{prop:1-on-W}}} \langle u \rangle^{2N_{\ref{prop:1-on-W}}},\]
which implies \eqref{mainbump2_dec1} (using also $C_{\ref{lem:abs_deriv_ft_phi3_le}, 2N_{\ref{prop:1-on-W}}}\le C_{\ref{lem:abs_deriv_ft_Tphi3_le}, 2N_{\ref{prop:1-on-W}}}$).

To prove \eqref{mainbump2_dec2} use Lemma \ref{lem:abs_deriv_ft_Tphi3_le} with $m=2N_{\ref{prop:1-on-W}}$:
\[ 
|\widehat{T\varphi_{3,k}}^{(2N_{\ref{prop:1-on-W}})}(\xi)| \le C_{\ref{lem:abs_deriv_ft_Tphi3_le}, 2N_{\ref{prop:1-on-W}}}. \]
Here we have also used that the support of $\widehat{T\varphi_{3,k}}$ is contained in, say $\{|\xi|\le \frac12\}$ (by  \eqref{mainbump2_supp} and Lemma \ref{lem:ft_deriv_mul}).
This also implies $|\textup{supp} (\widehat{T\varphi_{3,k}})|\le 1$.
Then apply Lemma \ref{lem:smoothdecay2} (with $N=2N_{\ref{prop:1-on-W}}$ and $|\textup{supp} (\widehat{T\varphi_{3,k}})|\le 1$) to obtain
\[ |T\varphi_{3,k}(u)| \le C_{\ref{lem:abs_deriv_ft_Tphi3_le}, 2N_{\ref{prop:1-on-W}}} C_{\ref{lem:smoothdecay2},2N_{\ref{prop:1-on-W}}} \langle u \rangle^{2N_{\ref{prop:1-on-W}}}, \]
which implies \eqref{mainbump2_dec2}.
\end{proof}


\begin{lemma}\label{lem:leftbump}
For every $k\le -1$,
\begin{equation}\label{leftbump}
\|A_{t}(\varphi_{1,k})\|_{V_{2,J}(t\in \mathbb{R}; L^2)}^2  \le C_{\ref{lem:leftbump}} 2^{2k} J^{\alpha(n)}, \end{equation}
where $C_{\ref{lem:leftbump}} = C_{\ref{lem:main_aux2}}C_{\ref{lem:theta_prim},2N_{\ref{prop:1-on-W}}}^2$.
\end{lemma}

\begin{proof}
We have \[\varphi_{1,k} = \mathbf{1}_{(-\infty,0)}*\theta_{(2^k)} = 2^k  \widetilde{\theta}_{(2^k)},\]
where $\widetilde{\theta}=\mathbf{1}_{(-\infty,0)}*\theta$ is a primitive of $\theta$.
By Lemma \ref{lem:rescaling} (rescaling) it suffices to show
    \[  \|A_{t}(\widetilde{\theta})\|_{V_{2,J}(t\in \mathbb{R}; L^2)}^2  \le C_{\ref{lem:leftbump}}  J^{\alpha(n)}.\]
This will follow from Lemma \ref{lem:main_aux2} applied with $\psi=C_{\ref{lem:theta_prim},2N_{\ref{prop:1-on-W}}}^{-1}\widetilde{\theta}$  
after we verify the necessary assumptions, %\eqref{main_aux1_supp} and \eqref{main_aux1_decay}, 
i.e. 
\[{\rm supp}(C_{\ref{lem:theta_prim},2N_{\ref{prop:1-on-W}}}^{-1}(\widehat{\widetilde{\theta}}))\subset \mathrm{Ann}_1(1, 2^{19}),\]
\begin{equation}
    C_{\ref{lem:theta_prim},2N_{\ref{prop:1-on-W}}}^{-1}|\widetilde{\theta}(u)| \le  \langle u \rangle^{2N_{\ref{prop:1-on-W}}}, 
\end{equation}
and 
\begin{equation}
    |T(C_{\ref{lem:theta_prim},2N_{\ref{prop:1-on-W}}}^{-1}\widetilde{\theta})(u)| =  C_{\ref{lem:theta_prim},2N_{\ref{prop:1-on-W}}}^{-1}|T\widetilde{\theta}(u)| \le  \langle u \rangle^{2N_{\ref{prop:1-on-W}}}.
\end{equation}
This follows   from Lemma \ref{lem:theta_prim} applied  with $N=2N_{\ref{prop:1-on-W}}$. 
\end{proof}

\begin{lemma}\label{lem:leftbump1_short1}
Let $\gamma=\frac12$. For every $k\le -1$  and every increasing sequence of integers $(k_j)_{j\in [J)}$,
\begin{equation}
\sum_{j\in [J)} 2^{-(1+\gamma)k}\int_{1}^2 \|A_{2^{k_j}t}(\varphi_{4,k})\|_2^2 \,\frac{dt}{t}  \le C_{\ref{lem:leftbump1_short1}} J^{\alpha(n)}
   \end{equation}
   where $C_{\ref{lem:leftbump1_short1}}=C_{\ref{prop:1-on-W}} C^2_{\ref{lem:theta_prim},2N_{\ref{prop:1-on-W}}+2} C_{\ref{lem:thetat_offcenter},N_{\ref{prop:1-on-W}}} $.  
\end{lemma}

\begin{proof}
By Lemma \ref{A to Lambda} applied to the function $(\varphi_{4,k})_{(2^{k_j} t)}$ the left-hand side equals
\[ \sum_{j\in [J)} 2^{-(1+\gamma)k}\int_{1}^2 \Lambda_1((\varphi_{4,k})_{(2^{k_j} t)}^{\otimes 2})(\F)     \,\frac{dt}{t} = \Lambda_1(M)(\F), \]
(with $\F$ as in Lemma \ref{A to Lambda}), 
where
\[ M = \sum_{j\in [J)} 2^{-(1+\gamma)k}\int_{1}^2 (\varphi_{4,k})_{(2^{k_j} t)}^{\otimes 2}  \,\frac{dt}{t} = \sum_{j\in [J)} \Psi_{(2^{k_j})} \]
where  
\[ \Psi = 2^{-(1+\gamma)k} \int_{1}^2 (\varphi_{4,k})_{(2^{k_j} t)}^{\otimes 2}  \,\frac{dt}{t} \]


We will apply Proposition \ref{prop:1-on-W} to the sequence of functions $c_1^{-1}\Psi_{(2^{k_j})}$ where $c_1 = C^2_{\ref{lem:theta_prim},2N_{\ref{prop:1-on-W}}+2} C_{\ref{lem:thetat_offcenter},N_{\ref{prop:1-on-W}}} $, so we verify the  assumptions of this proposition.
The assumptions 
\eqref{p-1-on-W_pos} and \eqref{p-1-on-W_supp}
follow from  Lemma \ref{lem:int_fct} applied with $\ell=k_j$ and 
\[\psi=2^{-(1+\gamma)k/2}c_1^{-1/2}\varphi_{4,k}.\]
The assumption \eqref{int_fct_assumption} of Lemma \ref{lem:int_fct}  
is satisfied due to Lemma \ref{lem:phi4_supp}.
 

 






It remains to verify the decay assumption \eqref{p-1-on-W_decay} of Proposition \ref{prop:1-on-W}. 
We will show 
\[|\Psi(u,v)| \le c_1 2^{(2-\gamma)k} \langle 2^k(u+v)\rangle ^{N_{\ref{prop:1-on-W}}}  \langle u-v\rangle^{N_{\ref{prop:1-on-W}}} \]
Then \eqref{p-1-on-W_decay} follows by multiplying with $c_1^{-1}$ on both sides. 

To see this estimate, by the triangle inequality and $t^{-1}\le 1$,
\begin{equation}%\label{centeredest}
|\Psi(u,v)| \le  2^{-(1+\gamma)k}\int_1^2 |\varphi_{4,k}(t^{-1}u)\varphi_{4,k}(t^{-1}v)|\,{dt} 
\end{equation}
\[=  2^{(1-{\gamma})k} \int_1^2|\widetilde{\theta}(t^{-1}(u-t2^{-k}))\widetilde{\theta}(t^{-1}(v-t2^{-k}))|\,dt\]
\[\le C^2_{\ref{lem:theta_prim},2N_{\ref{prop:1-on-W}}+2}  
2^{(1-{\gamma})k} \int_1^2\langle u - t2^{-k}\rangle^{2N_{\ref{prop:1-on-W}}+2}\langle v - t2^{-k}\rangle^{2N_{\ref{prop:1-on-W}}+2} \,dt \]
where we used \eqref{thetaprim_decay} of Lemma \ref{lem:theta_prim} with $N=2N_{\ref{prop:1-on-W}}+2$.  By Lemma \ref{lem:thetat_offcenter} with $N=N_{\ref{prop:1-on-W}}$, this is 
\[\le C^2_{\ref{lem:theta_prim},2N_{\ref{prop:1-on-W}}+2} C_{\ref{lem:thetat_offcenter}, N_{\ref{prop:1-on-W}}} 2^{(2-\gamma)k} \langle 2^k(u+v)\rangle^{N_{\ref{prop:1-on-W}}}  \langle u-v\rangle^{N_{\ref{prop:1-on-W}}} \]
\end{proof}


\begin{lemma}\label{lem:leftbump1_short2}
Let $\gamma=\frac12$. For every $k\le -1$ and every increasing sequence of integers $(k_j)_{j\in [J)}$,
\begin{equation}
\sum_{j\in [J)} 2^{(1-\gamma)k}\int_{1}^2 \|A_{2^{k_j}t}(T\varphi_{4,k}) \|_2^2 \,\frac{dt}{t} \le C_{\ref{lem:leftbump1_short2}} J^{\alpha(n)}
\end{equation}
where $C_{\ref{lem:leftbump1_short2}}= 2^3 C_{\ref{prop:1-on-W}}  C_{\ref{lem:thetat_offcenter}, N_{\ref{prop:1-on-W}}} \max ( C_{\ref{lem:theta_prim},2N_{\ref{prop:1-on-W}}+2}, C_{\ref{lem:theta_decay},2N_{\ref{prop:1-on-W}}+2}, C_{\ref{lem:theta_decay},2N_{\ref{prop:1-on-W}}+3} )^2$.
\end{lemma}

\begin{proof}
By Lemma \ref{A to Lambda} applied to the function $(T\varphi_{4,k})_{(2^{k_j} t)}$, the left-hand side equals
\[ \sum_{j\in [J)} 2^{(1-\gamma)k}\int_{1}^2 \Lambda_1( (T\varphi_{4,k})_{(2^{k_j} t)}^{\otimes 2}  )(\F)   \,\frac{dt}{t} = \Lambda_1(M)(\F), \]
(with $\F$ as in Lemma \ref{A to Lambda}),  where
\[ M = \sum_{j\in [J)} 2^{(1-\gamma)k}\int_{1}^2 (T\varphi_{4,k})_{(2^{k_j} t)}^{\otimes 2}   \,\frac{dt}{t} =\sum_{j\in [J)} \Psi_{(2^{k_j})} \]
where 
\[ \Psi   = 2^{(1-\gamma)k}\int_{1}^2 (T\varphi_{4,k})_{(t)}^{\otimes 2}  \,\frac{dt}{t} \] 

We will apply Proposition \ref{prop:1-on-W} to the function $c_1^{-1}\Psi$ where 
\[c_1 = 2^3 \max ( C_{\ref{lem:theta_prim},2N_{\ref{prop:1-on-W}}+2}, C_{\ref{lem:theta_decay},2N_{\ref{prop:1-on-W}}+2}, C_{\ref{lem:theta_decay},2N_{\ref{prop:1-on-W}}+3} )^2,\]
so we verify the  assumptions of this proposition. 
The assumptions 
\eqref{p-1-on-W_pos} and \eqref{p-1-on-W_supp}
follow from  Lemma \ref{lem:int_fct} applied with $\ell=k_j$ and 
\[\psi=2^{(1-\gamma)k/2}c_1^{-1/2}T\varphi_{4,k}.\] 
(The assumption \eqref{int_fct_assumption} of Lemma \ref{lem:int_fct}  
is satisfied due to Lemma \ref{lem:phi4_supp}, since the support does not change after multiplying with a constant.) 

It remains to verify the decay assumption \eqref{p-1-on-W_decay} of Proposition \ref{prop:1-on-W}. We will show
\begin{equation}
\label{bump_1short1_decay}
    |\Psi(u,v)| \le c_1 2^{(2-\gamma)k} \langle 2^k(u+v)\rangle^{N_{\ref{prop:1-on-W}}}  \langle u-v\rangle^{N_{\ref{prop:1-on-W}}}
\end{equation}
Then \eqref{p-1-on-W_decay} follows  by multiplying with $c_1^{-1}$ on both sides. 

By the definition of $T\varphi_{4,k}$, 
\[\Psi(u,v) = 2^{(3-\gamma)k}\int_{1}^2 t\partial_t(t^{-1}\widetilde{\theta}(t^{-1}u-2^{-k}))(u) t\partial_t(t^{-1}\widetilde{\theta}(t^{-1}u-2^{-k}))(v)  \,\frac{dt}{t}\]
To compute the derivatives we   apply  Lemma \ref{lem:phider}  with   $\phi(u) = \widetilde{\theta}(u-2^{-k})$, giving
\begin{equation}
    \label{thetatilde_tder}
    -t\partial_t(t^{-1}\widetilde{\theta}(t^{-1}u-2^{-k})) = t^{-1}\widetilde{\theta}(t^{-1}u-2^{-k}) + (X(\widetilde{\theta}' (\cdot - 2^{-k}))_{(t)}(u)
\end{equation}
where  we compute
\[ (X(\widetilde{\theta}' (\cdot - 2^{-k}))_{(t)}(u)  = t^{-2}u {\theta} (t^{-1}u-2^{-k}) \]
\[ = t^{-1} (X\theta)(t^{-1}u-2^{-k})  + t^{-2} 2^{-k} {\theta} (t^{-1}u-2^{-k}),  \]

For two functions $\phi_1,\phi_2$ let us denote 
\[I_k(\phi_1,\phi_2) = \int_{1}^2  | \phi_1(t^{-1}(u-t2^{-k})) \phi_2(t^{-1}(v-t2^{-k}))|  \,dt \]
 Using \eqref{thetatilde_tder}, together with $t^{-1}\le 1$,   $t^{-1}u-2^{-k} = t^{-1}(u-t2^{-k})$, and the triangle inequality, 
we estimate 
\[|\Psi(u,v)| \le  2^{(3-\gamma)k}I_k(\widetilde{\theta},\widetilde{\theta})\]
\[+ 2^{(3-\gamma)k} I_k(\widetilde{\theta},X{\theta}) + 2^{(2-\gamma)k} I_k(\widetilde{\theta},{\theta})\]
\[+ 2^{(3-\gamma)k} I_k(X{\theta},\widetilde{\theta})+ 2^{(2-\gamma)k} I_k({\theta},\widetilde{\theta})\]
\[+ 2^{(3-\gamma)k} I_k(X{\theta},X{\theta}) + 2^{(1-\gamma)k} I_k({\theta},{\theta})\]
% \[ 2^{(3-\gamma)k}\int_{1}^2 |\widetilde{\theta}(t^{-1}(u-t2^{-k})) \widetilde{\theta}(t^{-1}(v-t2^{-k}))|  \,dt \]
% \[ + 2^{(3-\gamma)k}\int_{1}^2 |\widetilde{\theta}(t^{-1}(u-t2^{-k}))   (X\theta) (t^{-1}(v-t2^{-k}))|  \,dt \] 
% \[ + 2^{(2-\gamma)k}\int_{1}^2 |\widetilde{\theta}(t^{-1}(u-t2^{-k}))   \theta (t^{-1}(v-t2^{-k}))|  \,dt \] 
% \[ + 2^{(3-\gamma)k}\int_{1}^2  | (X\theta)(t^{-1}(u-t2^{-k})) \widetilde{\theta}(t^{-1}(v-t2^{-k}))|  \,dt \]
% \[ + 2^{(2-\gamma)k}\int_{1}^2  | \theta(t^{-1}(u-t2^{-k})) \widetilde{\theta}(t^{-1}(v-t2^{-k}))|  \,dt \]
% \[ + 2^{(3-\gamma)k}\int_{1}^2 |(X\theta)  (t^{-1}(u-t2^{-k}))  (X\theta) (t^{-1}(v-t2^{-k}))|   \,dt \]
% \[ + 2^{(1-\gamma)k}\int_{1}^2 |\theta (t^{-1}(u-t2^{-k}))  \theta (t^{-1}(v-t2^{-k}))|   \,dt \]

We now state the estimates on the bump functions that will be used. 
 Lemma \ref{lem:theta_decay} with $N=2N_{\ref{prop:1-on-W}}+2$ gives the estimate
\[   |\theta(u)|\le  C_{\ref{lem:theta_decay},2N_{\ref{prop:1-on-W}}+2} \langle u\rangle^{2N_{\ref{prop:1-on-W}}+2} \]
Lemma \ref{lem:theta_decay} with $N=2N_{\ref{prop:1-on-W}}+3$ give the estimate 
\[|(X\theta)(u)| = |u\theta(u)| \le \langle u\rangle|\theta(u)|\]
\[ \le C_{\ref{lem:theta_decay},23}  \langle u\rangle \langle u\rangle^{2N_{\ref{prop:1-on-W}}+3} = C_{\ref{lem:theta_decay},23}\langle u\rangle^{2N_{\ref{prop:1-on-W}}+2} \]
%\[   |(X\theta)(u)|\le  C_{\ref{lem:theta_decay},23}  (1+|u|)^{-22} \] 
The estimate  \eqref{thetaprim_decay}  from Lemma \ref{lem:theta_prim} with $N=2N_{\ref{prop:1-on-W}}+2$ gives 
\[   |\widetilde{\theta}(u)|\le  C_{\ref{lem:theta_prim},2N_{\ref{prop:1-on-W}}+2}  \langle u\rangle^{2N_{\ref{prop:1-on-W}}+2} \]
 


Using these bounds to estimate each $I_k$ and crudely estimating $2^{(3-\gamma)k}\le 2^{(1-\gamma)k}$, $2^{(2-\gamma)k}\le 2^{(1-\gamma)k} $, gives 
\[|\Psi(u,v)|  \le  c_2 2^{(1-{\gamma})k} \int_1^2\langle u - t2^{-k}\rangle^{2N_{\ref{prop:1-on-W}}+2}\langle v - t2^{-k}\rangle^{2N_{\ref{prop:1-on-W}}+2} \,dt  \]
with $c_2=2^3 \max ( C_{\ref{lem:theta_prim},2N_{\ref{prop:1-on-W}}+2}, C_{\ref{lem:theta_decay},2N_{\ref{prop:1-on-W}}+2}, C_{\ref{lem:theta_decay},2N_{\ref{prop:1-on-W}}+3} )^2 $. 
Applying Lemma \ref{lem:thetat_offcenter} with $N = N_{\ref{prop:1-on-W}}$, this is 
\[\le c_2  C_{\ref{lem:thetat_offcenter},N_{\ref{prop:1-on-W}}} 2^{(2-\gamma)k} \langle 2^k(u+v)\rangle^{N_{\ref{prop:1-on-W}}}  \langle u-v\rangle^{N_{\ref{prop:1-on-W}}} \]
\end{proof}

\begin{lemma}\label{lem:leftbump1_long}
Let $\gamma=\frac12$. For every $k\le -1$,
\begin{equation}
\|A_{t}(\varphi_{4,k})\|^2_{V_{2,J}(t\in 2^\mathbb{Z}; L^2)} \le C_{\ref{lem:leftbump1_long}} 2^{\gamma k} J^{\alpha(n)}
\end{equation}
where $C_{\ref{lem:leftbump1_long}} = C_{\ref{prop:1-on-W}} C^2_{\ref{lem:theta_prim},2N_{\ref{prop:1-on-W}}} C_{\ref{lem:widebump},M}2^{2N_{\ref{prop:1-on-W}}}$.  
\end{lemma}


\begin{proof}
By Lemma \ref{lem:bootstrap}, 
        \[\|A_{t}(\varphi_{4,k})\|^2_{V_{2,J}(t\in 2^{\Z}; L^2)} \le  4 \cdot \sup \sum_{j\in [J)} \|A_{2^{k_{j}}}(\varphi_{4,k})\|^2_2     \]
        By Lemma \ref{A to Lambda} applied to the function $(\varphi_{4,k})_{(2^{k_j})}$ and using linearity,  the sum on the right-hand side equals
 $\Lambda_1(M)(\F)$ (with $\F$ as in Lemma \ref{A to Lambda}), where
\[ M = \sum_{j\in [J)} (\varphi_{4,k})_{(2^{k_j})}\otimes (\varphi_{4,k})_{(2^{k_j})}  \]
We will apply Proposition \ref{prop:1-on-W} to the functions $c_1^{-1}\Psi$, where   $c_1= C^2_{\ref{lem:theta_prim},2N_{\ref{prop:1-on-W}}}C_{\ref{lem:widebump},N_{\ref{prop:1-on-W}}}2^{2N_{\ref{prop:1-on-W}}}$  and 
\[\Psi = 2^{-\gamma k} \varphi_{4,k}\otimes \varphi_{4,k},\]  after verifying the necessary assumptions. The support property \eqref{p-1-on-W_supp} follows from Lemma \ref{lem:phi4_supp} after taking tensor products and multiplying by a constant. The symmetry property is immediate. 
The positivity property \eqref{p-1-on-W_pos} follows from \eqref{form_pos_psipsi} of Lemma \ref{lem:form_pos} applied with $\psi=(\varphi_{4,k})_{(2^{k_j})}$. 
% \[ \int_{\R^2} f(u)\overline{f(v)} (\varphi_{4,k})_{(2^{k_j})}(u) (\varphi_{4,k})_{(2^{k_j})}(v) \,dudv  =  \Big| \int_{\R} f(u) (\varphi_{4,k})_{(2^{k_j})}(u)  \,du\Big|^2 \]
 To see the decay \eqref{p-1-on-W_decay}, we   write
 \[ |\Psi(u,v)| = 2^{-\gamma k}|(\varphi_{4,k})(u) (\varphi_{4,k})(v)| \]
 \[= 2^{k(2-\gamma)}  |\widetilde{\theta}(u-2^{-k})\widetilde{\theta}(v-2^{-k})|\]
 \[\le  2^{k(2-\gamma)} C^2_{\ref{lem:theta_prim},2N_{\ref{prop:1-on-W}}} \langle u-2^{-k}\rangle^{2N_{\ref{prop:1-on-W}}}\langle v-2^{-k}\rangle^{2N_{\ref{prop:1-on-W}}} \]
where  we used \eqref{thetaprim_decay} of Lemma \ref{lem:theta_prim} with $N=2N_{\ref{prop:1-on-W}}$. 
By Lemma \ref{lem:bump_diagest}, this is 
\[\le 2^{k(2-\gamma)}  C^2_{\ref{lem:theta_prim},20} 2^{2N_{\ref{prop:1-on-W}}} \langle u+v-2^{1-k}\rangle^{N_{\ref{prop:1-on-W}}}\langle u-v\rangle^{N_{\ref{prop:1-on-W}}} \]
By Lemma \ref{lem:widebump} applied with $t=2^{1-k}$ and $N = N_{\ref{prop:1-on-W}}$, this is
\[\le c_1 2^{k(2-\gamma)} \langle 2^k(u+v)\rangle^{N_{\ref{prop:1-on-W}}}\langle u-v\rangle^{N_{\ref{prop:1-on-W}}}. \]
\end{proof}



\begin{lemma}\label{lem:leftbump1}
% For every $\gamma<1$ there exists a constant\todo{specify constant} $C<\infty$ such that for every $k\le -1$,
Let $\gamma=\frac12$. For every $k\le -1$,
\begin{equation}\label{leftbump1}
\|A_{t}(\varphi_{2,k})\|_{V_{2,J}(t\in \mathbb{R}; L^2)}^2  \le C_{\ref{leftbump1}} 2^{\gamma k} J^{\alpha(n)}, \end{equation}
where  $C_{\ref{lem:leftbump1}} = 2^3(C_{\ref{lem:leftbump1_short1}}^{1/2} C_{\ref{lem:leftbump1_short2}}^{1/2}  + C_{\ref{lem:leftbump1_long}})$. 
\end{lemma}

\begin{proof}
We have 
\[\varphi_{2,k}(u)=(2^k\widetilde{\theta}(u-2^{-k}))_{(2^k)} = (\varphi_{4,k})_{(2^k)}(u),\]
so by Lemma \ref{lem:rescaling} (rescaling) it suffices to show
\begin{equation}
\|A_{t}({\varphi_{4,k}}) \|_{V_{2,J}(t\in \mathbb{R}; L^2)}^2\le C_{\ref{leftbump1}} 2^{\gamma k} J^{\alpha(n)}
\end{equation}

Let $(t_j)_{j\in [J+1)}$ be an increasing sequence of real numbers.
By Lemma \ref{lem:shortlongjumps} (with $r=2$, $a(t) = A_t({\varphi_{4,k}}) $ and $B=L^2$),
\begin{equation}  \Big(\sum_{j\in [J)} \|A_{t_{j+1}}({\varphi_{4,k}}) - A_{t_j}({\varphi_{4,k}}) \|_2^2\Big)^{1/2} \le 2 \Big(\sum_{\ell \in L} \|A_t({\varphi_{4,k}}) \|_{V_{2,J}(t\in [2^\ell, 2^{\ell+1}]; L^2)}^2\Big)^{1/2}
\end{equation}
\begin{equation} 
+\|A_t({\varphi_{4,k}}) \|_{V_{2,J}(t\in 2^\mathbb{Z}; L^2)}.
\end{equation}
where $L$ is the set of $\ell\in \mathbb{Z}$ such that $2^{\ell}\le t_j<2^{\ell+1}$ for some $j\in [J+1)$.
Squaring and using $(a+b)^2\le 2(a^2+b^2)$ we obtain
\begin{equation}\label{shortlongftc_reduction_pf12} \sum_{j\in [J)} \|A_{t_{j+1}}(\varphi_{4,k})- A_{t_j}(\varphi_{4,k})\|_2^2  \le 2^3 \sum_{\ell \in L} \|A_t(\varphi_{4,k})\|_{V_{2,J}(t\in [2^\ell, 2^{\ell+1}]; L^2)}^2
\end{equation}
\begin{equation}\label{shortlongftc_reduction_pf22}
+ 2 \|A_t(\varphi_{4,k})\|^2_{V_{2,J}(t\in 2^\mathbb{Z}; L^2)}.
\end{equation}

 Now we use Lemma \ref{lem:leftbump1_long} to estimate 
 \begin{equation}
     \label{leftbump_long_est}
     \|A_{t}(\varphi_{4,k})\|^2_{V_{2,J}(t\in 2^\mathbb{Z}; L^2)} \le C_{\ref{lem:leftbump1_long}} 2^{\gamma k}J^{\alpha(n)}
 \end{equation}
We will also show
\begin{equation}
    \label{leftbump_short_est}
    \sum_{\ell \in L} \|A_t(\varphi_{4,k})\|_{V_{2,J}(t\in [2^\ell, 2^{\ell+1}]; L^2)}^2 \le C_{\ref{lem:leftbump1_short1}}^{1/2} C_{\ref{lem:leftbump1_short2}}^{1/2}\ 2^{\gamma k} J^{\alpha(n)}
\end{equation}
The estimates \eqref{leftbump_short_est}, \eqref{leftbump_long_est}, and \eqref{shortlongftc_reduction_pf22} together give
\[\|A_{t}(\varphi_{4,k})\|_{V_{2,J}(t\in \mathbb{R}; L^2)}^2  \le (2^3C_{\ref{lem:leftbump1_short1}}^{1/2} C_{\ref{lem:leftbump1_short2}}^{1/2}  + 2C_{\ref{lem:leftbump1_long}}) 2^{\gamma k} J^{\alpha(n)} = C_{\ref{lem:leftbump1}} 2^{\gamma k} J^{\alpha(n)}  \]
as desired. 
To see \eqref{leftbump_short_est}, 
for each $\ell\in\mathbb{Z}$ we estimate
\[ \|A_t(\varphi_{4,k})\|_{V_{2,J}(t\in [2^\ell, 2^{\ell+1}]; L^2)}^2 \le \int_{\mathbb{R}^n}  \|A_t(\varphi_{4,k})(x)\|_{V_{2,J}(t\in [2^\ell, 2^{\ell+1}])}^2\, dx.\]
Fix $x\in\mathbb{R}^n$ and $\ell\in\mathbb{Z}$. Set $a(t)=A_t(\varphi_{4,k})(x)$ and use Lemma \ref{lem:ftccs-R} (specifically, \eqref{ftccs2-R}) to estimate for each $x$
\[ \|A_t(\varphi_{4,k})(x)\|_{V_{2,J}(t\in [2^\ell, 2^{\ell+1}])}^2 \le  2^3
\|a(t)\|_{L^2(t\in [2^\ell, 2^{\ell+1}],\frac{dt}{t})} \|ta'(t)\|_{L^2(t\in [2^\ell, 2^{\ell+1}],\frac{dt}{t})}.  \]
By Lemma \ref{lem:ftc_ATphi}, 
\[ \|ta'(t)\|_{L^2(t\in [2^\ell, 2^{\ell+1}],\frac{dt}{t})}^2 = \int_{1}^2 |A_{2^\ell t}(T\varphi_{4,k})(x)|^2\,\frac{dt}{t}. \]
Thus for every $x\in\mathbb{R}^2$ and $\ell\in\mathbb{Z}$,  
\[ \|A_t(\varphi_{4,k})(x)\|_{V_{2,J}(t\in [2^\ell, 2^{\ell+1}])}^2\]
\[ \le 2^3 \Big( 2^{-k} \int_{1}^2 |A_{2^\ell t}(\varphi_{4,k})(x)|^2\,\frac{dt}{t} \Big)^{1/2} \Big( 2^{k}  \int_{1}^2 |A_{2^\ell t}(T{\varphi_{4,k}})(x)|^2\,\frac{dt}{t} \Big)^{1/2}.  \]
We integrate over $x\in\mathbb{R}^2$ and sum in $\ell\in L$.  Then we apply the Cauchy-Schwarz inequality in $x$ and in the summation.   We also interchange the integrals over $x$ with the integrals $t$ and the summation. This gives
\[ \sum_{\ell \in L}\|A_t(\varphi_{4,k})\|_{V_{2,J}(t\in [2^\ell, 2^{\ell+1}]; L^2)}^2 \]
\[\le \Big( 2^{-k} \int_1^2 \sum_{\ell\in L}\|A_{2^\ell t}(\varphi_{4,k})\|_2^2 \frac{dt}{t} \Big)^{1/2}\Big( 2^k \int_1^2 \sum_{\ell \in L}\|A_{2^\ell t}(T\varphi_{4,k})\|_2^2 \frac{dt}{t} \Big)^{1/2}   . \]
For each fixed $t\in [1,2]$ we use Lemma \ref{lem:leftbump1_short1} and Lemma \ref{lem:leftbump1_short2} and that $L\le J$.
%(we estimate $\#L^\frac12 \le (J+1)^\frac12 \le 2 J^\frac12$). 
This bounds the last display by 
$C_{\ref{lem:leftbump1_short1}}^{1/2} C_{\ref{lem:leftbump1_short2}}^{1/2} 2^{\gamma k} J^{\alpha(n)},$
which shows \eqref{leftbump_short_est}.
\end{proof}
 
%\begin{equation}\label{remainbump1}
%\sum_{j=1}^{J} \| M^{\varphi_{\infty,k}}_{t_j}(f_0,f_1,f_2) - M^{\varphi_{2,\infty}}_{t_{j-1}}(f_0,f_1,f_2)\|_{2}^2  \lesssim 2^{ k}J \end{equation}

\subsubsection{\texorpdfstring{Proof of Theorem \ref{thm:nct main real}}{Proof of main theorem}}
Fix $j\in [J)$.
By Lemma \ref{lem:smoothingdecomp},
$\|A_{t_{j + 1}}(\mathbf{1}_{[0,1]}) - A_{t_{j}}(\mathbf{1}_{[0,1]})\|_{{L}^2(\R^n)}$
is no larger than 
\[ \|A_{t_{j+1}}(\chi) - A_{t_{j}}(\chi)\|_{L^2(\R^n)}  + \sum_{k=-2}^\infty\|A_{t_{j+1}}(\varphi_{0,k}) - A_{t_{j}}(\varphi_{0,k})\|_{L^2(\R^n)}  + \]
\[ +\sum_{k=-\infty}^{-1}\|A_{t_{j+1}}(\varphi_{1,k}) - A_{t_{j}}(\varphi_{1,k})\|_{L^2(\R^n)}  + \sum_{k=-\infty}^{-1}\|A_{t_{j+1}}({\varphi_{2,k}}) - A_{t_{j}}({\varphi_{2,k}})\|_{L^2(\R^n)}. \]
Here we have also applied Lemma \ref{lem:norm_A_sum_le_sum} to interchange series and integrals.
After squaring this expression and summing over $j\in [J)$ and applying the triangle inequality in $\ell^2([J))$ (the series are now over non-negative terms, so no further justification is needed) we apply
Lemma \ref{lem:mainbump1}, Lemma \ref{lem:mainbump2}, Lemma \ref{lem:leftbump}, and Lemma \ref{lem:leftbump1} to obtain
\[ \sum_{j\in [J)} \|A_{t_{j + 1}}(\mathbf{1}_{[0,1]}) - A_{t_{j}}(\mathbf{1}_{[0,1]})\|_{{L}^2(\R^n)}^2 \le C_{\ref{lem:mainbump1}} J^{\alpha(n)} + \sum_{k=-2}^\infty C_{\ref{lem:mainbump2}} 2^{-2k} J^{\alpha(n)} + \] 
\[ + \sum_{k=-\infty}^{-1} C_{\ref{lem:leftbump}} 2^{2k} J^{\alpha(n)} + \sum_{k=-\infty}^{-1} C_{\ref{lem:leftbump1}} 2^{\frac12 k} J^{\alpha(n)} \le C J^{\alpha(n)},  \]
where 
$C = C_{\ref{lem:mainbump1}} + \tfrac{2^5}{3} C_{\ref{lem:mainbump2}} + \tfrac13 C_{\ref{lem:leftbump}} + (1 + 2^{\frac12}) C_{\ref{lem:leftbump1}}$.

\bibliographystyle{plain}
\bibliography{References}

\appendix
 
\section{Supplements: Automatically generated proofs}\label{sec:machine generated}
\jr{Needs to be regenerated in blueprint style -- will do this in the end}

In this section we collect all arguments that were machine-generated.

\subsection{Comments on automation}
These comments are more general than the present issue; ours is a good sample use case.

Consider Gaussian domination, Proposition \ref{Gaussian domination combined}.
Above we prove in detail only the near case $h>0$, $u(i)=0$ in Proposition \ref{Gaussian domination combined}. The other cases are similar with the obvious modifications. I convinced myself of this on paper, but I think there is no value in typing the full details by hand, because LLMs can do such tasks very well. Partly because of this possibility, I also don't think there is any value in including these details in a traditional math paper. They naturally belong into the blueprint. Similarly, all the other proofs in the blueprint can be generated by LLM when provided with the appropriate hints where needed. We will have to still make the decisions on division and formulation of statements.

The subsections below are 100\% generated by GPT 5.5, each by a version of the prompt (not that the prompt matters a lot):
\begin{lstlisting}[breaklines]
Attached is a draft. Your task is to complete the Gaussian domination estimate (Proposition Gaussian domination combined) for the case h>0, u_i = 1. This is very similar to the case that is already written. You should very closely stick to the style that already exists, avoid introducing any new notation and provide the complete proof in LateX. You should ideally not use any other theorems than the ones stated in the draft.
\end{lstlisting}
A human does not necessarily have to read the sections below in detail, because:

1. We are already independently convinced in the ``traditional'' sense that the request above is well-posed (i.e. it is actually true that this is a straightforward modification of the existing and already trusted argument, and does not involve additional difficulties).

2. At the line-by-line level, GPT 5.5 has checked its own output for correctness in a separate request. Anecdotally, it is great at catching local correctness errors, especially when arguments are very detailed. What it is {\em not} good at is having an actual understanding of underlying ideas.

3. I expect an LLM will also do the eventual formalization of the proofs based on this as input (but not the statements).\\

In summary, we never actually have to read, or even {\em trust} the text below.

\subsection{Remaining cases of Proposition \ref{Gaussian domination combined}}

\subsubsection{Near case $h>0$ and $u(i)=1$}
\begin{proof}[Proof of Proposition \ref{Gaussian domination combined} for the near case, $\iota=(h,0)$, $h>0$, and $u_i=1$]
Here we consider the case $\iota=(h,0)\in\mathcal{I}_{\gamma,i}$ with $h>0$ and $u_i=1$.
As before we write $d_i=\Delta(a_i)$. Let $\lambda$ denote the sequence
\[
    \lambda(j)=2^h a_{i,2}(j+d_i).
\]
For fixed $j\in\Z$, let $\rho=\rho_j$ be the unique Schwartz function with
\[
\widehat{\rho}(\tau)=
\frac{\widehat{\psi_{(2^{h-1}a_{i,2}(j+d_i))}}(\tau)
-\widehat{\psi_{(2^{h}a_{i,2}(j+d_i))}}(\tau)}
{\widehat{s(2^h a_{i,2}(\cdot+d_i),j)}(\tau)^\nu}\, .
\]
By Proposition \ref{four scale Gaussian kernel}, for every $N\ge 1$ and $p\in\R$,
\[
    |\rho(p)|\lesssim_N \langle p\rangle^N_{(\lambda(j))},
\]
and
\[
    |\rho'(p)|\lesssim_N \lambda(j)^{-1}\langle p\rangle^N_{(\lambda(j))}.
\]
In what follows we fix $j$ and omit the argument $j$ from the scales when there is no danger of confusion.

Let
\[
    F=(N_{\gamma,\iota})_{i,j},\qquad H=(H_\gamma)_{i,j}.
\]
Then, by the definition of $N_{\gamma,\iota}$,
\[
    F(v)=\int_\R \rho(p)H(v_1-p,v_2-p)\,dp .
\]
We use the cancellation of $H$ in the diagonal direction.  Write
\[
    v_1=w_1-w_2,\qquad v_2=w_1+w_2 .
\]
Changing variables in the integral and subtracting zero, using Proposition \ref{H vanishing}, gives
\[
F(v)=\int_\R \big(\rho(w_1+p)-\rho(w_1)\big)
H(-w_2-p,w_2-p)\,dp .
\]
By Proposition \ref{mean value bump estimate},
\[
|\rho(w_1+p)-\rho(w_1)|
\lesssim
\min(1,\lambda^{-1}|p|)
\big(\langle w_1+p\rangle^N_{(\lambda)}
+\langle w_1\rangle^N_{(\lambda)}\big).
\]

Let
\[
    t_{1,\ell}=a_{i,\ell}(j-1),\qquad
    t_{2,\ell}=a_{i,\ell}(j),
    \qquad \ell=1,2.
\]
Since $u_i=1$, we have
\[
    (s_\gamma)_{i,j}=\Sigma(\sqrt{2}a_{i,2},j).
\]
Thus Proposition \ref{diagonal square root} gives
\[
    |(s_\gamma)_{i,j}(x)|
    \lesssim
    \langle x\rangle^8_{(t_{1,2})}
    +\langle x\rangle^2_{(t_{2,2})}.
\]
Moreover, by the definition of $G_{q,1}$,
\[
    |(G_\gamma)_{i,j-1}(v)|+|(G_\gamma)_{i,j}(v)|
    \lesssim
    \sum_{m=1}^2
    \left\langle \frac{v_1+v_2}{\sqrt{2}}\right\rangle^8_{(t_{m,1})}
    \left\langle \frac{-v_1+v_2}{\sqrt{2}}\right\rangle^8_{(t_{m,2})}.
\]
Consequently,
\[
\begin{aligned}
|H(v)|
&\lesssim
\big(\langle v_1\rangle^8_{(t_{1,2})}
+\langle v_1\rangle^2_{(t_{2,2})}\big)
\big(\langle v_2\rangle^8_{(t_{1,2})}
+\langle v_2\rangle^2_{(t_{2,2})}\big)  \\
&\quad
+\sum_{m=1}^2
    \left\langle \frac{v_1+v_2}{\sqrt{2}}\right\rangle^8_{(t_{m,1})}
    \left\langle \frac{-v_1+v_2}{\sqrt{2}}\right\rangle^8_{(t_{m,2})}.
\end{aligned}
\]
Substituting $v=(-w_2-p,w_2-p)$ and absorbing harmless constants coming from $\sqrt{2}$, we obtain
\[
\begin{aligned}
|H(-w_2-p,w_2-p)|
&\lesssim
\big(\langle -w_2-p\rangle^8_{(t_{1,2})}
+\langle -w_2-p\rangle^2_{(t_{2,2})}\big) \\
&\quad\quad\times
\big(\langle w_2-p\rangle^8_{(t_{1,2})}
+\langle w_2-p\rangle^2_{(t_{2,2})}\big)  \\
&\quad
+\sum_{m=1}^2
\langle p\rangle^8_{(t_{m,1})}
\langle w_2\rangle^8_{(t_{m,2})}.
\end{aligned}
\]

Let
\[
    \mathcal{T}_0=\{t_{1,1},t_{1,2},t_{2,1},t_{2,2}\}.
\]
For every $r\in\mathcal{T}_0$, the definition of $d_i=\Delta(a_i)$ gives
\[
    \frac{r}{\lambda}\le 2^{-h}.
\]
Indeed, each of the four scales in $\mathcal{T}_0$ is bounded by $a_{i,2}(j+d_i)$ up to an absolute shift allowed by $d_i$, while
$\lambda=2^h a_{i,2}(j+d_i)$.

We first estimate the part containing the two factors depending on $w_2+p$ and $w_2-p$.  Using
\[
    |p|\le \frac12|w_2+p|+\frac12|w_2-p|
\]
and, for $r\in\mathcal{T}_0$,
\[
\min(1,\lambda^{-1}|p\pm w_2|)
\le
(\lambda^{-1}|p\pm w_2|)^{1/2}
\le
2^{-h/2}(1+r^{-1}|p\pm w_2|)^{1/2},
\]
we estimate the corresponding integrand by a finite sum of terms of the form
\[
C2^{-h/2}
\langle w_1+\alpha p\rangle^N_{(\lambda)}
\langle -w_2-p\rangle^{n_-}_{(r_-)}
\langle w_2-p\rangle^{n_+}_{(r_+)}
\]
with $\alpha\in[2)$, $r_\pm\in\mathcal{T}_0$, and $n_\pm\ge 3/2$.

If $\alpha=0$, Proposition \ref{two bump estimate} gives an upper bound by a constant times
\[
    2^{-h/2}
    \langle w_1\rangle^N_{(\lambda)}
    \langle w_2\rangle^{\min(n_-,n_+)}_{(\max(r_-,r_+))}.
\]
If $\alpha=1$, we use Proposition \ref{bump triangle} to separate the two narrow bumps, then Proposition \ref{two bump estimate}, and finally Proposition \ref{orthogonal decay}.  This gives an upper bound by a constant times
\[
\begin{aligned}
2^{-h/2}\big(
&\langle w_1\rangle^{n_0}_{(\lambda)}
 \langle w_2\rangle^{n_0}_{(r_-)}
+\langle w_1+w_2\rangle^{n_0}_{(\lambda)}
 \langle w_1-w_2\rangle^{n_0}_{(r_-)}  \\
&+\langle w_1\rangle^{n_0}_{(\lambda)}
 \langle w_2\rangle^{n_0}_{(r_+)}
+\langle w_1-w_2\rangle^{n_0}_{(\lambda)}
 \langle w_1+w_2\rangle^{n_0}_{(r_+)}
\big),
\end{aligned}
\]
where $n_0=\min(n_-,n_+)\ge 3/2$.

It remains to estimate the terms coming from the last line in the bound for
$|H(-w_2-p,w_2-p)|$.  These terms are bounded by expressions of the form
\[
\min(1,\lambda^{-1}|p|)
\big(\langle w_1+p\rangle^N_{(\lambda)}
+\langle w_1\rangle^N_{(\lambda)}\big)
\langle p\rangle^8_{(r_1)}
\langle w_2\rangle^8_{(r_2)}
\]
with $r_1,r_2\in\mathcal{T}_0$.  Since $r_1/\lambda\le 2^{-h}$,
\[
\min(1,\lambda^{-1}|p|)
\le
(\lambda^{-1}|p|)^{1/2}
\le
2^{-h/2}(1+r_1^{-1}|p|)^{1/2}.
\]
Thus, after lowering the exponent $8$ by $1/2$, Proposition \ref{two bump estimate} gives
\[
\int_\R
\min(1,\lambda^{-1}|p|)
\big(\langle w_1+p\rangle^N_{(\lambda)}
+\langle w_1\rangle^N_{(\lambda)}\big)
\langle p\rangle^8_{(r_1)}
\langle w_2\rangle^8_{(r_2)}
\,dp
\lesssim
2^{-h/2}
\langle w_1\rangle^{3/2}_{(\lambda)}
\langle w_2\rangle^{3/2}_{(r_2)}.
\]

Combining the previous estimates and lowering all decay exponents to $3/2$, we have bounded $|F(v)|$ by a finite sum of terms of the form
\[
    C2^{-h/2}
    \langle (W_u v)_1\rangle^{3/2}_{(r_1)}
    \langle (W_u v)_2\rangle^{3/2}_{(r_2)},
\]
where $u\in[2)$ and one of the two scales $r_1,r_2$ is $\lambda$ while the other belongs to
\[
    \mathcal{T}_1=\{\max(r,r'):\ r,r'\in\mathcal{T}_0\}.
\]
Here we used that $w_1=2^{-1/2}(W_1v)_1$, $w_2=2^{-1/2}(W_1v)_2$, while
$w_1-w_2=v_1$ and $w_1+w_2=v_2$.

Let $\mathcal{P}\subset{\rm A}^2$ be the finite collection of all pairs $(r_1,r_2)$ such that one component is $\lambda$ and the other component is in $\mathcal{T}_1$.  Define
\[
    \mathcal{B}=[2)\times\mathcal{P}.
\]
Then $\#\mathcal{B}$ is bounded by an absolute constant.  For $b=(u,p)\in\mathcal{B}$, set $u_b=u$.  The preceding estimate becomes
\[
    |F(v)|
    \lesssim
    2^{-h/2}
    \sum_{b=(u,p)\in\mathcal{B}}
    \langle (W_u v)_1\rangle^{3/2}_{(p_1(j))}
    \langle (W_u v)_2\rangle^{3/2}_{(p_2(j))}.
\]
By Proposition \ref{Gaussian domination}, each summand is bounded by
\[
C
\sum_{m_1,m_2\in\N}
2^{-m_1/2}2^{-m_2/2}
g_{(2^{m_1}p_1(j))}((W_u v)_1)
g_{(2^{m_2}p_2(j))}((W_u v)_2).
\]
For $m=(m_1,m_2)\in\N^2$ define
\[
    p_{b,m}(j)=\big(2^{m_1}p_1(j),2^{m_2}p_2(j)\big).
\]
Then
\[
    |F(v)|
    \le
    C2^{-h/2}
    \sum_{m\in\N^2}2^{-|m|/2}
    \sum_{b\in\mathcal{B}}
    G_{p_{b,m}(j),u_b}(v).
\]
Since here $|\iota_1|=h$, this is the desired pointwise domination.

It remains only to verify the distance bound.  Every element of $\mathcal{T}_1$ is obtained from
$a_i^1(\cdot-1)$, $a_i^1$, $a_{i,2}(\cdot-1)$, and $a_{i,2}$ by taking a maximum of two such sequences.  Hence, by Propositions \ref{Operations on spaced sequences} and \ref{Properties of distance of sequences},
\[
    \dist(r,a_{i,2})\lesssim 1+d_i
    \qquad\text{for all } r\in\mathcal{T}_1.
\]
On the other hand,
\[
    \lambda=2^h a_{i,2}(\cdot+d_i),
\]
so again by Proposition \ref{Properties of distance of sequences},
\[
    \dist(\lambda,a_{i,2})\lesssim h+d_i.
\]
Thus, for every eligible pair $p=(p_1,p_2)\in\mathcal{P}$,
\[
    \dist(p_1,p_2)\lesssim 1+d_i+h.
\]
After multiplying the two components by $2^{m_1}$ and $2^{m_2}$, Proposition \ref{Properties of distance of sequences} gives
\[
\Delta(p_{b,m})
\le
C_1(1+d_i+h+m_1+m_2)
=
C_1(1+\Delta(a_i)+|\iota_1|+|m|).
\]
This completes the proof of the case $h>0$ and $u_i=1$.
\end{proof}


\subsubsection{Far case $h<0$}
\begin{proof}[Proof of Proposition \ref{Gaussian domination combined} for the far case, $\iota=(h,0)$, $h<0$]
Here we consider the case $\iota=(h,0)\in\mathcal{I}_{\gamma,i}$ with $h<0$.
As before we write $d_i=\Delta(a_i)$ and fix $j\in\mathbb Z$. We often omit the argument $j$ where clear from context. Let
\[
    \lambda=\lambda(j):=2^h a_{i,2}(j-d_i-1).
\]
Let $\rho$ be the unique Schwartz function with
\[
\widehat{\rho}(\tau)=
\frac{\widehat{\psi_{(2^h a_{i,2}(j-d_i-1))}}(\tau)
      -\widehat{\psi_{(2^{h+1}a_{i,2}(j-d_i-1))}}(\tau)}
     {\widehat{s(2^h a_{i,2}(\cdot-d_i),j)}(\tau)^{\nu}} .
\]
Indeed, in Proposition \ref{four scale Gaussian kernel} we take
\[
 \lambda_-=2^h a_{i,2}(j-d_i-1),\qquad
 \lambda_+=2^{h+1}a_{i,2}(j-d_i-1),
\]
and
\[
 \mu_-=2^h a_{i,2}(j-d_i-1),\qquad
 \mu_+=2^h a_{i,2}(j-d_i).
\]
Then
\[
    \mu_-\le \lambda_- \ll_2 \lambda_+ \le \mu_+,
\]
where the last inequality follows from the multiplicative spacing of $a_{i,2}$.
Thus, by Proposition \ref{four scale Gaussian kernel}, for every sufficiently large $N$ and every $p\in\mathbb R$,
\begin{equation}\label{eqn:Gaussianbumppfneg-rho}
    |\rho(p)|\lesssim_N \langle p\rangle^N_{(\lambda)}.
\end{equation}
The numerator in the definition of $\widehat{\rho}$ vanishes in a neighborhood of the origin. Hence
\[
    \widehat{\rho}(0)=0,
\]
and therefore
\begin{equation}\label{eqn:Gaussianbumppfneg-mean-zero}
    \int_{\mathbb R}\rho(p)\,dp=0.
\end{equation}
From \eqref{eqn:Gaussianbumppfneg-rho}, choosing $N$ large enough, we also have
\begin{equation}\label{eqn:Gaussianbumppfneg-moment}
    \int_{\mathbb R}|p|\,|\rho(p)|\,dp\lesssim \lambda.
\end{equation}

Let
\[
    F=(N_{\gamma,\iota})_{i,j},\qquad H=(H_\gamma)_{i,j}.
\]
By the definition of $N_{\gamma,\iota}$,
\[
    F(v)=\int_{\mathbb R}\rho(p)H(v_1-p,v_2-p)\,dp.
\]
Using \eqref{eqn:Gaussianbumppfneg-mean-zero}, we subtract the value at $p=0$ and obtain
\begin{equation}\label{eqn:Gaussianbumppfneg-F}
    F(v)=\int_{\mathbb R}\rho(p)\big(H(v_1-p,v_2-p)-H(v_1,v_2)\big)\,dp.
\end{equation}
This is the point where the case $h<0$ differs from the case $h>0$: here the decay in $h$ comes from the small scale $\lambda$ and the cancellation of $\rho$.

We shall use the following elementary consequence of \eqref{eqn:Gaussianbumppfneg-rho}. If $r_l\ge \lambda$, $|c_l|\le \sqrt2$, and $L\le 2$, then
\begin{equation}\label{eqn:Gaussianbumppfneg-tail}
    \int_{\mathbb R}|p|\,|\rho(p)|\prod_{l=1}^{L}
    \langle x_l-c_l p\rangle_{(r_l)}^2\,dp
    \lesssim
    \lambda\prod_{l=1}^{L}\langle x_l\rangle_{(r_l)}^{3/2}.
\end{equation}
Indeed,
\[
    \langle x_l-c_l p\rangle_{(r_l)}^2
    \lesssim
    (1+|p|/r_l)^2\langle x_l\rangle_{(r_l)}^{3/2},
\]
and since $r_l\ge \lambda$, the remaining integral is bounded by a constant times $\lambda$ by \eqref{eqn:Gaussianbumppfneg-rho}, provided $N$ is chosen large.

We also use the differentiated form of the one-dimensional estimate following from Proposition \ref{diagonal square root}: for $a\in{\rm A}$,
\begin{equation}\label{eqn:Gaussianbumppfneg-sprime}
    |s(a,j)'(x)|
    \lesssim
    a(j-1)^{-1}\langle x\rangle_{(a(j-1))}^{8}
    +a(j)^{-1}\langle x\rangle_{(a(j))}^{2}.
\end{equation}
This follows by differentiating the representation used in the proof of Proposition \ref{diagonal square root}. We also use the elementary Gaussian estimate
\[
    |(g_{(r)})'(x)|\lesssim r^{-1}\langle x\rangle_{(r)}^{8}.
\]

{\em Case $u(i)=0$.}
Let
\[
    t_{1,\ell}(j)=a_{i,\ell}(j-1),\qquad
    t_{2,\ell}(j)=a_{i,\ell}(j),\qquad \ell=1,2,
\]
and put
\[
    t_m^+=\max(t_{m,1},t_{m,2}),\qquad m=1,2.
\]
Let
\[
    \mathcal{T}_0=\{t_{1,1},t_{1,2},t_{2,1},t_{2,2},t_1^+,t_2^+\}.
\]
By the definition of $d_i=\Delta(a_i)$ and by the multiplicative spacing,
\begin{equation}\label{eqn:Gaussianbumppfneg-scale0}
    \lambda\le 2^h r
\end{equation}
for every $r\in\mathcal{T}_0$.

With $s=(s_\gamma)_{i,j}$, using $u(i)=0$,
\[
    |H(v)|
    \le |s(v_1)s(v_2)|
      +\sum_{m=1}^{2}g_{(t_{m,1})}(v_1)g_{(t_{m,2})}(v_2).
\]
Combining this estimate with \eqref{eqn:Gaussianbumppfneg-sprime} and the Gaussian derivative estimate gives
\begin{equation}\label{eqn:Gaussianbumppfneg-DH0}
    |(\partial_1+\partial_2)H(v)|
    \lesssim
    \sum_{r_-,r_+\in\mathcal{T}_0}
    (r_-^{-1}+r_+^{-1})
    \langle v_1\rangle_{(r_-)}^2
    \langle v_2\rangle_{(r_+)}^2.
\end{equation}
Using the mean value theorem in the diagonal direction in \eqref{eqn:Gaussianbumppfneg-F}, then using
\eqref{eqn:Gaussianbumppfneg-DH0}, \eqref{eqn:Gaussianbumppfneg-tail}, and \eqref{eqn:Gaussianbumppfneg-scale0}, we obtain
\[
\begin{aligned}
    |F(v)|
    &\lesssim
    \sum_{r_-,r_+\in\mathcal{T}_0}
    (r_-^{-1}+r_+^{-1})
    \int_{\mathbb R}|p|\,|\rho(p)|
    \int_0^1
    \langle v_1-\theta p\rangle_{(r_-)}^2
    \langle v_2-\theta p\rangle_{(r_+)}^2
    d\theta\,dp                                                        \\
    &\lesssim
    2^h
    \sum_{r_-,r_+\in\mathcal{T}_0}
    \langle v_1\rangle_{(r_-)}^{3/2}
    \langle v_2\rangle_{(r_+)}^{3/2}.
\end{aligned}
\]
Since $h<0$, we have $2^h\le 2^{h/2}=2^{-|h|/2}$. Hence
\begin{equation}\label{eqn:Gaussianbumppfneg-final0}
    |F(v)|
    \lesssim
    2^{-|h|/2}
    \sum_{r_-,r_+\in\mathcal{T}_0}
    \langle (W_0v)_1\rangle_{(r_-)}^{3/2}
    \langle (W_0v)_2\rangle_{(r_+)}^{3/2}.
\end{equation}

{\em Case $u(i)=1$.}
Let again
\[
    t_{1,\ell}(j)=a_{i,\ell}(j-1),\qquad
    t_{2,\ell}(j)=a_{i,\ell}(j),\qquad \ell=1,2.
\]
Then, by the definition of $d_i=\Delta(a_i)$ and by multiplicative spacing,
\begin{equation}\label{eqn:Gaussianbumppfneg-scale1}
    \lambda\le 2^h t_{m,\ell}
\end{equation}
for $m=1,2$ and $\ell=1,2$. Let $z=W_1v$. Since
\[
    W_1(1,1)=(\sqrt2,0),
\]
a diagonal translation of $v$ translates only the first $W_1$-coordinate.

With $s=(s_\gamma)_{i,j}$, using $u(i)=1$,
\[
    |H(v)|
    \le |s(v_1)s(v_2)|
      +\sum_{m=1}^{2}
      g_{(t_{m,1})}(z_1)g_{(t_{m,2})}(z_2).
\]
Here $s=\Sigma(\sqrt2 a_{i,2},j)$, and the harmless factor $\sqrt2$ is absorbed into the constants. Therefore \eqref{eqn:Gaussianbumppfneg-sprime} gives
\[
\begin{aligned}
    |(\partial_1+\partial_2)(s(v_1)s(v_2))|
    \lesssim
    \sum_{r_-,r_+\in\{a_{i,2}(\cdot-1),a_{i,2}\}}
    (r_-^{-1}+r_+^{-1})
    \langle v_1\rangle_{(r_-)}^2
    \langle v_2\rangle_{(r_+)}^2.
\end{aligned}
\]
For the Gaussian terms, since the diagonal derivative differentiates only the first $W_1$-coordinate,
\[
    |(\partial_1+\partial_2)
    \big(g_{(t_{m,1})}(z_1)g_{(t_{m,2})}(z_2)\big)|
    \lesssim
    t_{m,1}^{-1}
    \langle z_1\rangle_{(t_{m,1})}^2
    \langle z_2\rangle_{(t_{m,2})}^2.
\]
Using these two estimates in \eqref{eqn:Gaussianbumppfneg-F}, and then using
\eqref{eqn:Gaussianbumppfneg-tail} together with \eqref{eqn:Gaussianbumppfneg-scale1}, gives
\[
\begin{aligned}
    |F(v)|
    \lesssim 2^h
    \sum_{r_-,r_+\in\{a_{i,2}(\cdot-1),a_{i,2}\}}
    \langle v_1\rangle_{(r_-)}^{3/2}
    \langle v_2\rangle_{(r_+)}^{3/2}                                      \\
    +2^h\sum_{m=1}^{2}
    \langle (W_1v)_1\rangle_{(t_{m,1})}^{3/2}
    \langle (W_1v)_2\rangle_{(t_{m,2})}^{3/2}.
\end{aligned}
\]
Again $2^h\le 2^{h/2}=2^{-|h|/2}$, and hence
\begin{equation}\label{eqn:Gaussianbumppfneg-final1}
\begin{aligned}
    |F(v)|
    \lesssim
    2^{-|h|/2}
    \sum_{r_-,r_+\in\{a_{i,2}(\cdot-1),a_{i,2}\}}
    \langle (W_0v)_1\rangle_{(r_-)}^{3/2}
    \langle (W_0v)_2\rangle_{(r_+)}^{3/2}                                  \\
    +2^{-|h|/2}
    \sum_{m=1}^{2}
    \langle (W_1v)_1\rangle_{(t_{m,1})}^{3/2}
    \langle (W_1v)_2\rangle_{(t_{m,2})}^{3/2}.
\end{aligned}
\end{equation}

We now finish both cases at once. Let $\mathcal{B}$ be the finite set of all pairs $(u',p)$ which occur in
\eqref{eqn:Gaussianbumppfneg-final0} in the case $u(i)=0$, and in
\eqref{eqn:Gaussianbumppfneg-final1} in the case $u(i)=1$.
Thus $u'\in[2)$ and $p=(p_1,p_2)\in{\rm A}^2$.
The cardinality of $\mathcal{B}$ is bounded by an absolute constant.

By Proposition \ref{Gaussian domination}, applied with exponent $3/2$, each term in
\eqref{eqn:Gaussianbumppfneg-final0} and \eqref{eqn:Gaussianbumppfneg-final1} is bounded by
\[
    C
    \sum_{m_1,m_2\in\mathbb N}
    2^{-m_1/2}2^{-m_2/2}
    g_{(2^{m_1}p_1(j))}((W_{u'}v)_1)
    g_{(2^{m_2}p_2(j))}((W_{u'}v)_2).
\]
For $b=(u',p)\in\mathcal{B}$ put $u_b=u'$, and for
$m=(m_1,m_2)\in\mathbb N^2$ define
\[
    p_{b,m}(j):=(2^{m_1}p_1(j),2^{m_2}p_2(j)).
\]
Then $p_{b,m}\in{\rm A}^2$ by Proposition \ref{Operations on spaced sequences}, and the last display gives
\[
    |F(v)|
    \le
    C2^{-|h|/2}
    \sum_{m\in\mathbb N^2}2^{-|m|/2}
    \sum_{b\in\mathcal{B}}
    G_{p_{b,m}(j),u_b}(v).
\]
Since here $\iota=(h,0)$, this is exactly the required pointwise estimate.

It remains to check the distance bound. Before the final dyadic dilations by $2^{m_1}$ and $2^{m_2}$, every sequence appearing in one of the pairs $p=(p_1,p_2)$ is obtained from $a_i^1$ and $a_{i,2}$ by at most one shift and, in the case $u(i)=0$, possibly by taking the maximum of two such shifted sequences. By Propositions \ref{Operations on spaced sequences} and \ref{Properties of distance of sequences}, and directly from the definition of $\dist$ for the maximum, we have
\[
    \Delta(p)\lesssim 1+\Delta(a_i).
\]
After multiplying the two components by $2^{m_1}$ and $2^{m_2}$, Proposition \ref{Properties of distance of sequences} gives
\[
    \Delta(p_{b,m})
    \lesssim
    1+\Delta(a_i)+m_1+m_2
    \le
    C_1(1+\Delta(a_i)+|h|+|m|).
\]
This proves the required estimate for the far case $h<0$, in both cases $u(i)=0$ and $u(i)=1$.
\end{proof}

\subsubsection{Intermediate case $h=0$}
\begin{proof}[Proof of Proposition \ref{Gaussian domination combined} for the intermediate case, $\iota=(0,l)$, $|l|\le d_i$]
Here $\iota=(0,l)$ and $d_i=\Delta(a_i)$. Since $\iota_1=0$, the factor
$2^{-|\iota_1|/2}$ in the conclusion is equal to $1$. Let
\[
    \nu=1 \quad \text{if } k<n-1,\qquad 
    \nu=2 \quad \text{if } k=n-1.
\]
Define the shifted sequences
\[
    \lambda_0(j):=a_{i,2}(j+l-1),\qquad 
    \lambda_1(j):=a_{i,2}(j+l).
\]
Fix $j\in\mathbb Z$ and write $\lambda_\epsilon=\lambda_\epsilon(j)$,
$\epsilon\in[2)$. Let $\rho$ be the unique Schwartz function such that
\[
    \widehat{\rho}(\tau)
    =
    \frac{
    \widehat{\psi_{(\lambda_0)}}(\tau)
    -
    \widehat{\psi_{(\lambda_1)}}(\tau)}
    {\widehat{s(a_{i,2}(\cdot+l),j)}(\tau)^\nu}.
\]
Equivalently, by the definition of $s$,
\[
    \widehat{\rho}(\tau)
    =
    \bigl(\widehat{\psi_{(\lambda_0)}}(\tau)
    -
    \widehat{\psi_{(\lambda_1)}}(\tau)\bigr)
    \bigl(g(\lambda_0\tau)-g(\lambda_1\tau)\bigr)^{-\nu/2}.
\]
Since $a_{i,2}$ is multiplicatively spaced, $2\lambda_0\le \lambda_1$.
Hence Proposition \ref{four scale Gaussian kernel}, applied with
\[
    \mu_-=\lambda_-=\lambda_0,\qquad 
    \lambda_+=\mu_+=\lambda_1,
\]
gives
\begin{equation}\label{gdc-int-rho-size}
    |\rho(p)|
    \lesssim
    \langle p\rangle^{3/2}_{(\lambda_0)}
    +
    \langle p\rangle^{3/2}_{(\lambda_1)} .
\end{equation}

Let
\[
    F=(N_{\gamma,(0,l)})_{i,j},\qquad 
    H=(H_\gamma)_{i,j}.
\]
By the definition of $N_{\gamma,(0,l)}$,
\[
    \widehat F(\xi,\eta)=\widehat{\rho}(\xi+\eta)\widehat H(\xi,\eta).
\]
Thus, for $v=(v_1,v_2)\in\mathbb R^2$,
\[
    F(v)=\int_{\mathbb R}\rho(p)H(v_1-p,v_2-p)\,dp.
\]
We use the same coordinates as in the near case,
\[
    v_1=w_1-w_2,\qquad v_2=w_1+w_2.
\]
Changing variables in the last integral gives
\begin{equation}\label{gdc-int-direct-convolution}
    F(v)=
    \int_{\mathbb R}
    \rho(w_1+p)H(-w_2-p,w_2-p)\,dp .
\end{equation}
No cancellation is used in this case.

We now introduce the finite set of scales
\[
\begin{aligned}
\mathcal T_0=\{&
\lambda_0,\lambda_1,
a_i^1(\cdot-1),a_{i,2}(\cdot-1),
a_i^1,a_{i,2},\\
&\max(a_i^1(\cdot-1),a_{i,2}(\cdot-1)),
\max(a_i^1,a_{i,2})\}.
\end{aligned}
\]
By Proposition \ref{Operations on spaced sequences}, all elements of
$\mathcal T_0$ are in ${\rm A}$. Let $\mathcal T$ be the set of all pointwise
maxima of two elements of $\mathcal T_0$, enlarged, if necessary, by the fixed
dilates arising from the harmless factors $\sqrt 2$ in the coordinates
$W_1v=\sqrt 2(w_1,w_2)$. This remains a finite set of cardinality bounded by
an absolute constant.

We first assume $u_i=0$. Put
\[
    t_{1,r}(j):=a_{i,r}(j-1),\qquad 
    t_{2,r}(j):=a_{i,r}(j),
    \qquad r=1,2,
\]
and
\[
    t_m^+(j):=\max(t_{m,1}(j),t_{m,2}(j)),
    \qquad m=1,2.
\]
We omit the argument $j$ in the following estimates. By Proposition
\ref{diagonal square root} and the definition of $H_\gamma$,
\[
\begin{aligned}
    |H(x,y)|
    \lesssim&
    \bigl(\langle x\rangle^8_{(t_1^+)}
          +
          \langle x\rangle^2_{(t_2^+)}\bigr)
    \bigl(\langle y\rangle^8_{(t_1^+)}
          +
          \langle y\rangle^2_{(t_2^+)}\bigr)\\
    &+
    \sum_{m=1}^2
    \langle x\rangle^8_{(t_{m,1})}
    \langle y\rangle^8_{(t_{m,2})}.
\end{aligned}
\]
Decreasing all powers to $3/2$, every term obtained from
\eqref{gdc-int-rho-size} and \eqref{gdc-int-direct-convolution} is bounded by
a constant times
\begin{equation}\label{gdc-int-direct-term}
    \int_{\mathbb R}
    \langle w_1+p\rangle^{3/2}_{(\lambda_\epsilon)}
    \langle -w_2-p\rangle^{3/2}_{(r_-)}
    \langle w_2-p\rangle^{3/2}_{(r_+)}
    \,dp,
\end{equation}
where $\epsilon\in[2)$ and $r_-,r_+\in\mathcal T_0$.

By Proposition \ref{bump triangle},
\[
\begin{aligned}
    &\langle -w_2-p\rangle^{3/2}_{(r_-)}
    \langle w_2-p\rangle^{3/2}_{(r_+)}\\
    &\qquad\lesssim
    \langle w_2\rangle^{3/2}_{(r_-)}
    \langle w_2-p\rangle^{3/2}_{(r_+)}
    +
    \langle -w_2-p\rangle^{3/2}_{(r_-)}
    \langle w_2\rangle^{3/2}_{(r_+)} .
\end{aligned}
\]
Applying Proposition \ref{two bump estimate} to the remaining $p$-integrals,
we obtain
\[
\begin{aligned}
    \eqref{gdc-int-direct-term}
    \lesssim&
    \langle w_1+w_2\rangle^{3/2}_{(\max(\lambda_\epsilon,r_+))}
    \langle w_2\rangle^{3/2}_{(r_-)}\\
    &+
    \langle w_1-w_2\rangle^{3/2}_{(\max(\lambda_\epsilon,r_-))}
    \langle w_2\rangle^{3/2}_{(r_+)} .
\end{aligned}
\]
Proposition \ref{orthogonal decay} turns each of these products into a
finite sum of products of two bumps either in the coordinates $(w_1,w_2)$ or
in the coordinates $(w_1-w_2,w_1+w_2)$. Since
\[
    W_0v=(w_1-w_2,w_1+w_2),
    \qquad
    W_1v=\sqrt 2(w_1,w_2),
\]
we have shown that, when $u_i=0$,
\begin{equation}\label{gdc-int-direct-bracket-u0}
    |F(v)|
    \lesssim
    \sum_{u_0\in[2)}
    \sum_{r_1,r_2\in\mathcal T}
    \langle (W_{u_0}v)_1\rangle^{3/2}_{(r_1(j))}
    \langle (W_{u_0}v)_2\rangle^{3/2}_{(r_2(j))}.
\end{equation}

It remains to check the case $u_i=1$. In this case
$s_\gamma=\Sigma(\sqrt 2\,a_{i,2},j)$, and Proposition
\ref{diagonal square root} gives, after suppressing harmless absolute factors
in the scales,
\[
    |s_\gamma(x)|
    \lesssim
    \langle x\rangle^8_{(a_{i,2}(j-1))}
    +
    \langle x\rangle^2_{(a_{i,2}(j))}.
\]
Also,
\[
    (G_\gamma)_{i,j'}(x,y)
    =
    g_{(a_i^1(j'))}((W_1(x,y))_1)
    g_{(a_{i,2}(j'))}((W_1(x,y))_2)
\]
for $j'=j-1,j$. Hence
\[
\begin{aligned}
    |H(x,y)|
    \lesssim&
    \bigl(\langle x\rangle^8_{(a_{i,2}(j-1))}
          +
          \langle x\rangle^2_{(a_{i,2}(j))}\bigr)
    \bigl(\langle y\rangle^8_{(a_{i,2}(j-1))}
          +
          \langle y\rangle^2_{(a_{i,2}(j))}\bigr)\\
    &+
    \sum_{m=1}^2
    \langle (W_1(x,y))_1\rangle^8_{(a_i^1(j+m-2))}
    \langle (W_1(x,y))_2\rangle^8_{(a_{i,2}(j+m-2))}.
\end{aligned}
\]
The first line is of the same form as the estimate already treated in the
case $u_i=0$, and therefore gives the same contribution as in
\eqref{gdc-int-direct-bracket-u0}.

For the second line, substitute $(x,y)=(-w_2-p,w_2-p)$. Then, up to harmless
absolute constants,
\[
    (W_1(x,y))_1=-p,\qquad (W_1(x,y))_2=w_2.
\]
Thus the terms to estimate are bounded by
\begin{equation}\label{gdc-int-direct-u1-term}
    \int_{\mathbb R}
    \langle w_1+p\rangle^{3/2}_{(\lambda_\epsilon)}
    \langle p\rangle^{3/2}_{(r_-)}
    \langle w_2\rangle^{3/2}_{(r_+)}
    \,dp,
\end{equation}
with $\epsilon\in[2)$ and $r_-,r_+\in\mathcal T_0$. By Proposition
\ref{two bump estimate},
\[
    \eqref{gdc-int-direct-u1-term}
    \lesssim
    \langle w_1\rangle^{3/2}_{(\max(\lambda_\epsilon,r_-))}
    \langle w_2\rangle^{3/2}_{(r_+)}.
\]
This is already a product of two bumps in the coordinates of $W_1v$. Hence
\eqref{gdc-int-direct-bracket-u0} also holds when $u_i=1$.

We have proved in both cases $u_i=0$ and $u_i=1$ that
\begin{equation}\label{gdc-int-direct-bracket}
    |F(v)|
    \lesssim
    \sum_{u_0\in[2)}
    \sum_{r_1,r_2\in\mathcal T}
    \langle (W_{u_0}v)_1\rangle^{3/2}_{(r_1(j))}
    \langle (W_{u_0}v)_2\rangle^{3/2}_{(r_2(j))}.
\end{equation}
By Proposition \ref{Gaussian domination}, each product on the right-hand side
is bounded by
\[
    C
    \sum_{m_1,m_2\in\mathbb N}
    2^{-m_1/2}2^{-m_2/2}
    g_{(2^{m_1}r_1(j))}((W_{u_0}v)_1)
    g_{(2^{m_2}r_2(j))}((W_{u_0}v)_2).
\]
Define
\[
    \mathcal B:=[2)\times\mathcal T\times\mathcal T .
\]
This set has cardinality bounded by an absolute constant. For
$b=(u_0,r_1,r_2)\in\mathcal B$, set
\[
    u_b:=u_0,
\]
and, for $m=(m_1,m_2)\in\mathbb N^2$, define
\[
    p_{b,m}:=(2^{m_1}r_1,2^{m_2}r_2)\in{\rm A}^2.
\]
Then \eqref{gdc-int-direct-bracket} implies
\[
    |(N_{\gamma,(0,l)})_{i,j}(v)|
    \le
    C
    \sum_{m\in\mathbb N^2}
    2^{-|m|/2}
    \sum_{b\in\mathcal B}
    G_{p_{b,m}(j),u_b}(v).
\]
This is the desired pointwise estimate, since $|\iota_1|=0$.

It remains to verify the distance bound. Every element of $\mathcal T_0$ has
distance at most $d_i+O(1)$ from $a_{i,2}$. This is immediate for
$a_{i,2}$ and its shift, follows from $\dist(a_i^1,a_{i,2})=d_i$ for
$a_i^1$ and its shift, and follows for
$\lambda_0=a_{i,2}(\cdot+l-1)$ and
$\lambda_1=a_{i,2}(\cdot+l)$ because $|l|\le d_i$. Taking pointwise maxima,
and also taking the harmless fixed dilates used above, changes this bound by
at most an absolute constant. Therefore, for all $r\in\mathcal T$,
\[
    \dist(r,a_{i,2})\lesssim 1+d_i.
\]
Hence, for $b=(u_0,r_1,r_2)$ and $m=(m_1,m_2)$,
\[
\begin{aligned}
    \Delta(p_{b,m})
    &=
    \dist(2^{m_1}r_1,2^{m_2}r_2)\\
    &\le
    \dist(2^{m_1}r_1,r_1)
    +
    \dist(r_1,a_{i,2})
    +
    \dist(a_{i,2},r_2)
    +
    \dist(r_2,2^{m_2}r_2)\\
    &\lesssim
    m_1+m_2+1+d_i.
\end{aligned}
\]
Since $d_i=\Delta(a_i)$ and $|\iota_1|=0$, this gives
\[
    \Delta(p_{b,m})
    \le C_1(1+\Delta(a_i)+|\iota_1|+|m|).
\]
This completes the proof of the intermediate case.
\end{proof}

\subsection{Proof of Proposition \ref{Gaussian domination reduction variant}}\label{Gaussian domination reduction proof}

\begin{proof}
Fix $h\in\Z$ and $j\in\Z$. Let $\rho$ be the unique Schwartz function with
\[
    \widehat{\rho}(\tau)
    =
    \bigl(
    \widehat{\psi_{(2^h a(j))}}(\tau)
    -
    \widehat{\psi_{(2^{h+1}a(j))}}(\tau)
    \bigr)
    \widehat{\sigma_{j,h}}(\tau)^{-1}.
\]
The quotient is interpreted at the origin by its removable value. Indeed, the
numerator vanishes in a neighborhood of the origin, while
$\widehat{\sigma_{j,h}}$ can vanish only at the origin.

Set
\[
    \lambda=2^h a(j).
\]
By Proposition \ref{four scale Gaussian kernel}, applied with
\[
    \lambda_-=2^h a(j),\qquad
    \lambda_+=2^{h+1}a(j),
\]
and
\[
    \mu_-=2^{h+1}a(j-1),\qquad
    \mu_+=2^{h+1}a(j),
\]
we have
\[
    |\rho(p)|\lesssim \langle p\rangle^2_{(\lambda)}
\]
and
\[
    |\rho'(p)|
    \lesssim
    \lambda^{-1}\langle p\rangle^2_{(\lambda)}.
\]
Here the hypotheses of Proposition \ref{four scale Gaussian kernel} hold
because
\[
    2^{h+1}a(j-1)\le 2^h a(j)
    \le 2^{h+1}a(j)
    =
    2^{h+1}a(j).
\]

Let $F=N_{j,h}$. Then
\[
    F(v)
    =
    \int_{\R}\rho(p)M_j(v_0-p,v_1-p)\,dp .
\]
Write
\[
    v_0=w_0-w_1,\qquad v_1=w_0+w_1 .
\]
Changing variables gives
\[
    F(v)
    =
    \int_{\R}\rho(w_0+p)M_j(-w_1-p,w_1-p)\,dp .
\]
Using the pointwise hypothesis on $M_j$, and absorbing harmless absolute
constants in the scales, we have
\[
    |M_j(-w_1-p,w_1-p)|
    \lesssim
    \sum_{w\in[2)}
    2^{-w\kappa/2}
    \langle p\rangle^2_{(2^{w\kappa}a(j))}
    \langle w_1\rangle^2_{(a(j))}.
\]
It is enough to estimate the contribution of each fixed $w\in[2)$.

First assume
\[
    h\le w\kappa .
\]
Then $2^h a(j)\le 2^{w\kappa}a(j)$. Hence by Proposition
\ref{two bump estimate},
\[
\begin{aligned}
&\int_{\R}
    |\rho(w_0+p)|
    \langle p\rangle^2_{(2^{w\kappa}a(j))}
    \,dp  \\
&\qquad\lesssim
    \int_{\R}
    \langle w_0+p\rangle^2_{(2^h a(j))}
    \langle p\rangle^2_{(2^{w\kappa}a(j))}
    \,dp                                      \\
&\qquad\lesssim
    \langle w_0\rangle^2_{(2^{w\kappa}a(j))}.
\end{aligned}
\]
Therefore this part is bounded by
\[
    C2^{-w\kappa/2}
    \langle w_0\rangle^2_{(2^{w\kappa}a(j))}
    \langle w_1\rangle^2_{(a(j))}.
\]
Since in the present case
\[
    \max(w\kappa,h_+)=w\kappa,
\]
this is bounded by
\[
    C2^{-\max(w\kappa,h_+)/2}
    \langle w_0\rangle^{3/2}_{(2^{\max(w\kappa,h_+)}a(j))}
    \langle w_1\rangle^{3/2}_{(a(j))}.
\]

It remains to consider the case
\[
    h>w\kappa .
\]
Then in particular $h>0$. The diagonal cancellation assumption gives
\[
    \int_{\R}M_j(-w_1-p,w_1-p)\,dp=0.
\]
Indeed, the Fourier transform in the variable $w_1$ of the left-hand side is
the restriction of $\widehat{M_j}$ to the line $\eta=-\xi$.

Thus we may subtract zero:
\[
    F(v)
    =
    \int_{\R}
    \bigl(\rho(w_0+p)-\rho(w_0)\bigr)
    M_j(-w_1-p,w_1-p)\,dp .
\]
By Proposition \ref{mean value bump estimate},
\[
\begin{aligned}
|\rho(w_0+p)-\rho(w_0)|
&\lesssim
\min(1,2^{-h}a(j)^{-1}|p|)  \\
&\qquad\times
\bigl(
\langle w_0+p\rangle^2_{(2^h a(j))}
+
\langle w_0\rangle^2_{(2^h a(j))}
\bigr).
\end{aligned}
\]
Since $h>w\kappa$,
\[
\begin{aligned}
\min(1,2^{-h}a(j)^{-1}|p|)
&\le
(2^{-h}a(j)^{-1}|p|)^{1/2}       \\
&\le
2^{-(h-w\kappa)/2}
\bigl(1+2^{-w\kappa}a(j)^{-1}|p|\bigr)^{1/2}.
\end{aligned}
\]
Hence
\[
\begin{aligned}
&2^{-w\kappa/2}
\min(1,2^{-h}a(j)^{-1}|p|)
\langle p\rangle^2_{(2^{w\kappa}a(j))}  \\
&\qquad\lesssim
2^{-h/2}
\langle p\rangle^{3/2}_{(2^{w\kappa}a(j))}.
\end{aligned}
\]
Using Proposition \ref{two bump estimate}, and the fact that
$2^h a(j)>2^{w\kappa}a(j)$, we obtain
\[
\begin{aligned}
&\int_{\R}
\bigl(
\langle w_0+p\rangle^2_{(2^h a(j))}
+
\langle w_0\rangle^2_{(2^h a(j))}
\bigr)
\langle p\rangle^{3/2}_{(2^{w\kappa}a(j))}
\,dp                                      \\
&\qquad\lesssim
\langle w_0\rangle^{3/2}_{(2^h a(j))}.
\end{aligned}
\]
Therefore this part is bounded by
\[
    C2^{-h/2}
    \langle w_0\rangle^{3/2}_{(2^h a(j))}
    \langle w_1\rangle^{3/2}_{(a(j))}.
\]
Since in the present case
\[
    \max(w\kappa,h_+)=h,
\]
this is again bounded by
\[
    C2^{-\max(w\kappa,h_+)/2}
    \langle w_0\rangle^{3/2}_{(2^{\max(w\kappa,h_+)}a(j))}
    \langle w_1\rangle^{3/2}_{(a(j))}.
\]

Combining the two cases, we have shown that
\[
    |F(v)|
    \lesssim
    \sum_{w\in[2)}
    2^{-\max(w\kappa,h_+)/2}
    \langle w_0\rangle^{3/2}_{(2^{\max(w\kappa,h_+)}a(j))}
    \langle w_1\rangle^{3/2}_{(a(j))}.
\]
Since
\[
    W_1v=\sqrt2\,(w_0,w_1),
\]
the harmless factors of $\sqrt2$ may be absorbed into the constants and into
fixed absolute dilates of the scales. By Proposition \ref{Gaussian domination},
each product of bracket bumps is bounded by
\[
    C
    \sum_{m\in\N^2}
    2^{-|m|/2}
    g_{(2^{m_0}2^{\max(w\kappa,h_+)}a(j))}((W_1v)_0)
    g_{(2^{m_1}a(j))}((W_1v)_1).
\]

Now take
\[
    \mathcal B_h=[2).
\]
For $b\in\mathcal B_h$, set
\[
    u_b=1,\qquad w_b=b,
\]
and for $m=(m_0,m_1)\in\N^2$ define
\[
    p_{b,m}(j)
    =
    \bigl(
    2^{m_0}2^{\max(w_b\kappa,h_+)}a(j),
    2^{m_1}a(j)
    \bigr).
\]
Then $p_{b,m}\in{\rm A}^2$ by Proposition
\ref{Operations on spaced sequences}, and the previous estimate gives
\[
 |N_{j,h}(v)|
 \le
 C
 \sum_{m\in\N^2}2^{-|m|/2}
 \sum_{b\in\mathcal B_h}
 2^{-\max(w_b\kappa,h_+)/2}
 G_{p_{b,m}(j),u_b}(v).
\]

It remains to check the distance estimate. By Proposition
\ref{Properties of distance of sequences},
\[
    \dist(a,2^{m_1}a)\le m_1.
\]
Also, directly from the definition of $\dist$, for every $t\ge0$,
\[
    \dist(a,2^t a)\le \lceil t\rceil .
\]
Therefore
\[
\begin{aligned}
\Delta(p_{b,m})
&=
\dist\bigl(
2^{m_0}2^{\max(w_b\kappa,h_+)}a,
2^{m_1}a
\bigr)                                      \\
&\le
m_0+\lceil \max(w_b\kappa,h_+)\rceil+m_1  \\
&\le
C_1(1+w_b\kappa+|h|+|m|).
\end{aligned}
\]
Finally,
\[
    \#\mathcal B_h=2,
\]
so the cardinality bound holds after increasing $C_0$. The non-negative
series on the right-hand side converges pointwise and has finite $L^1$ norm
because each two-dimensional Gaussian has $L^1$ norm $1$ and
\[
    \sum_{m\in\N^2}2^{-|m|/2}<\infty .
\]
This completes the proof.
\end{proof}


\subsection{Proofs of properties of multiplicatively spaced sequences}\label{sec:auto multiplicative sequences}.

\begin{proof}[Proof of Proposition \ref{Extension of sequences}]
Define $b:\Z\to\R_{>0}$ by
\[
 b(j)=
 \begin{cases}
  2^j a(0), & j<0,\\
  a(j), & j\in [J),\\
  2^{j-J+1}a(J-1), & j\ge J.
 \end{cases}
\]
Then clearly $b(j)=a(j)$ for every $j\in[J)$, while the desired formulas for
$j<0$ and $j\ge J$ hold by definition.

It remains to check that $b\in {\rm A}$. Let $j\in\Z$. If $j<-1$, then
\[
 b(j+1)=2^{j+1}a(0)=2b(j).
\]
If $j=-1$, then
\[
 b(0)=a(0)=2\cdot 2^{-1}a(0)=2b(-1).
\]
If $j\in [J-1)$, then by assumption
\[
 b(j+1)=a(j+1)\ge 2a(j)=2b(j).
\]
If $j=J-1$, then
\[
 b(J)=2a(J-1)=2b(J-1).
\]
Finally, if $j\ge J$, then
\[
 b(j+1)=2^{j+1-J+1}a(J-1)=2b(j).
\]
Thus for every $j\in\Z$ we have $2b(j)\le b(j+1)$, so $b\in{\rm A}$.

The conditions in the statement prescribe the value of $b(j)$ for every
$j\in\Z$, since the three ranges $j<0$, $j\in[J)$, and $j\ge J$ cover
$\Z$ and are disjoint. Hence such a $b$ is unique.
\end{proof}

\begin{proof}[Proof of Proposition \ref{Operations on spaced sequences}]
For (i), let $j\in\Z$. If $\max(a,b)(j)=a(j)$, then
\[
 2\max(a,b)(j)=2a(j)\le a(j+1)\le \max(a,b)(j+1).
\]
The case $\max(a,b)(j)=b(j)$ is identical. Hence $\max(a,b)\in{\rm A}$.

For (ii), for every $j\in\Z$,
\[
 2(t a)(j)=t(2a(j))\le t a(j+1)=(t a)(j+1),
\]
so $t\cdot a\in{\rm A}$.

For (iii), for every $j\in\Z$,
\[
 2c(j)=2a(j+n)\le a(j+n+1)=c(j+1).
\]
Thus $c\in{\rm A}$.
\end{proof}

\begin{proof}[Proof of Proposition \ref{Properties of distance of sequences}]
For (i), if $\dist(a,b)=0$, then for every $j\in\Z$,
\[
 a(j)\le b(j)\le a(j),
\]
and hence $a(j)=b(j)$.

For (ii), suppose first that $\dist(a,b)\le k<\infty$. Then for every $j\in\Z$,
\[
 a(j-k)\le b(j)\le a(j+k).
\]
Applying this with $j-k$ in place of $j$ gives
\[
 b(j-k)\le a(j),
\]
and applying it with $j+k$ in place of $j$ gives
\[
 a(j)\le b(j+k).
\]
Therefore
\[
 b(j-k)\le a(j)\le b(j+k)
\]
for every $j\in\Z$, so $\dist(b,a)\le k$. By symmetry, $\dist(a,b)=\dist(b,a)$.

For (iii), if either $\dist(a,b)$ or $\dist(b,c)$ is infinite, there is nothing to prove. Otherwise let
\[
 \dist(a,b)\le k,\qquad \dist(b,c)\le l.
\]
Then for every $j\in\Z$,
\[
 b(j-l)\le c(j)\le b(j+l).
\]
Using the inequalities for $a$ and $b$ on the two outside terms gives
\[
 a(j-l-k)\le b(j-l)\le c(j)\le b(j+l)\le a(j+l+k).
\]
Hence
\[
 \dist(a,c)\le k+l.
\]
Taking the minimum over admissible $k$ and $l$ gives
\[
 \dist(a,c)\le \dist(a,b)+\dist(b,c).
\]

For (iv), if $\dist(a,b)$ is infinite, there is nothing to prove. Otherwise let $\dist(a,b)\le k$. Since $c(j)=b(j+1)$, for every $j\in\Z$,
\[
 a(j+1-k)\le c(j)\le a(j+1+k).
\]
Because $a\in{\rm A}$ is increasing, we have
\[
 a(j-k-1)\le a(j+1-k)
\]
and
\[
 a(j+1+k)=a(j+k+1).
\]
Therefore
\[
 a(j-(k+1))\le c(j)\le a(j+(k+1)),
\]
so $\dist(a,c)\le k+1$. Taking the minimum over admissible $k$ gives
\[
 \dist(a,c)\le \dist(a,b)+1.
\]

For (v), for every $k\in\Z_{\ge 0}$, the inequalities
\[
 a(j-k)\le b(j)\le a(j+k)
\]
hold for every $j\in\Z$ if and only if
\[
 t a(j-k)\le t b(j)\le t a(j+k)
\]
hold for every $j\in\Z$. Hence
\[
 \dist(t\cdot a,t\cdot b)=\dist(a,b).
\]

For (vi), first suppose $h\ge 0$. Iterating the defining inequality for $a\in{\rm A}$ gives
\[
 2^h a(j)\le a(j+h)
\]
for every $j\in\Z$. Also, since $a$ is increasing,
\[
 a(j-h)\le a(j)\le 2^h a(j).
\]
Thus
\[
 a(j-h)\le 2^h a(j)\le a(j+h),
\]
so $\dist(a,2^h a)\le h=|h|$.

If $h<0$, write $-h>0$. Iterating the defining inequality gives
\[
 2^{-h}a(j+h)\le a(j),
\]
and hence
\[
 a(j+h)\le 2^h a(j).
\]
Also, since $a$ is increasing,
\[
 2^h a(j)\le a(j)\le a(j-h).
\]
Therefore
\[
 a(j-(-h))\le 2^h a(j)\le a(j+(-h)),
\]
so $\dist(a,2^h a)\le -h=|h|$.

(iv) Since $a,b\in{\rm A}$, for every $j\in\Z$,
\[
2c(j)=\sqrt{(2a(j))^2+(2b(j))^2}
\le \sqrt{a(j+1)^2+b(j+1)^2}=c(j+1).
\]
Thus $c\in{\rm A}$.
\end{proof}

\subsection{Proofs of properties of Gaussians}
\label{sec:auto gaussian properties}

% \begin{proof}[Proof of Proposition \ref{square root one minus Gaussian}]
% For every $x\in\R$ we have $0<\g(x)\le 1$, and hence
% $1-\g(x)\ge 0$. Thus $f$ is well defined using the nonnegative square
% root. Since $\g$ is continuous and the square root is continuous on
% $[0,\infty)$, the function $f$ is continuous.

% Let $|x|\le 1/2$ and put $u=\pi x^2$. Then $0\le u\le \pi/4<1$. Using
% $e^{-u}\le 1-u+u^2/2$, we get
% \[
% 1-e^{-u}\ge u-\frac{u^2}{2}\ge \frac u2 .
% \]
% Therefore
% \[
% f(x)^2=1-e^{-\pi x^2}\ge \frac{\pi x^2}{2}\ge \frac{x^2}{4},
% \]
% and hence $f(x)\ge \frac12 |x|$.

% Finally, if $|x|\ge 1/2$, then $0\le 1-\g(x)\le 1$. For every
% $a\in[0,1]$ we have $a\le \sqrt a\le 1$, so with $a=1-\g(x)$ this gives
% \[
% 1-\g(x)\le f(x)\le 1 .
% \]
% \end{proof}

\begin{proof}[Proof of Proposition \ref{Gaussian bump decay}]
We first prove that, for every $m\in\N$ and $x\in\R$,
\[
|\g^{(m)}(x)|
\le
\bigl(50(m+1)\bigr)^{m/2}e^{-\pi x^2/2}.
\]
For $m=0$ this is immediate. For $m\ge 1$, repeated differentiation gives
\[
\g^{(m)}(x)=P_m(x)e^{-\pi x^2},
\]
where
\[
P_m(x)
=
m!\sum_{j=0}^{\lfloor m/2\rfloor}
\frac{(-1)^{m+j}2^{m-2j}\pi^{m-j}x^{m-2j}}
{j!(m-2j)!}.
\]
Indeed, one proves this using induction and
$P_{m+1}(x)=P_m'(x)-2\pi xP_m(x).$
Set $d=m-2j$. We use
\[
|x|^d e^{-\pi x^2/2}
\le
\left(\frac{d+1}{\pi e}\right)^{d/2}
\]
with the convention that the right hand side is $1$ when $d=0$.
Also,
\[
\frac{m!}{j!d!}
=
\binom{m}{d}\frac{(2j)!}{j!}
\le
2^m m^j .
\]
Hence each summand in $|P_m(x)|e^{-\pi x^2/2}$ is at most
\[
2^m(\pi m)^j
\left(\frac{4\pi(d+1)}{e}\right)^{d/2}
\le
\bigl(19(m+1)\bigr)^{m/2}.
\]
Since there are at most $2^{m/2}$ summands, we get
\[
|P_m(x)|e^{-\pi x^2/2}
\le
\bigl(38(m+1)\bigr)^{m/2}
\le
\bigl(50(m+1)\bigr)^{m/2}.
\]
This proves the displayed Gaussian derivative bound.

Next, for every $N\in\N$,
\[
(1+|x|)^N e^{-\pi x^2/2}
\le
\bigl(10(N+1)\bigr)^{N/2}.
\]
Equivalently,
\[
e^{-\pi x^2/2}
\le
\bigl(10(N+1)\bigr)^{N/2}\langle x\rangle^N .
\]
Therefore
\[
|\g^{(m)}(x)|
\le
\bigl(50(m+1)\bigr)^{m/2}
\bigl(10(N+1)\bigr)^{N/2}
\langle x\rangle^N.
\]
Finally,
\[
\bigl(50(m+1)\bigr)^{m/2}
\bigl(10(N+1)\bigr)^{N/2}
\le
\bigl(100(m+N+1)\bigr)^{(m+N)/2}.
\]
This proves the claim.
\end{proof}

\begin{proof}[Proof of Proposition \ref{Elementary Gaussian properties}]
By Proposition \ref{Gaussian bump decay}, every derivative of $\g$ decays
faster than any polynomial. Hence $\g\in \mathcal S(\R)$, and therefore
$\g\in W_0(\R)$.

We next prove the Fourier transform identity. With the convention
\[
\widehat h(\xi)=\int_\R h(x)e^{-2\pi i x\xi}\,dx,
\]
set
\[
I(\xi):=\int_\R e^{-\pi x^2}e^{-2\pi i x\xi}\,dx .
\]
Completing the square gives
\[
e^{-\pi x^2}e^{-2\pi i x\xi}
=
e^{-\pi (x+i\xi)^2}e^{-\pi \xi^2}.
\]
Thus, by shifting the contour from $\R$ to $\R+i\xi$,
\[
I(\xi)
=
e^{-\pi \xi^2}\int_\R e^{-\pi x^2}\,dx
=
e^{-\pi \xi^2}.
\]
Here we used the standard Gaussian integral
$\int_\R e^{-\pi x^2}\,dx=1.$
Hence $\widehat{\g}=\g$.

Let $\lambda>0$ and write
\[
\g_{(\lambda)}(x)=\lambda^{-1}\g(x/\lambda).
\]
Then, by the change of variables $x=\lambda y$,
\[
\widehat{\g_{(\lambda)}}(\xi)
=
\int_\R \lambda^{-1}\g(x/\lambda)e^{-2\pi i x\xi}\,dx
=
\int_\R \g(y)e^{-2\pi i y(\lambda \xi)}\,dy
=
\g(\lambda \xi).
\]
This proves \eqref{eq:scaleFT}.

Finally, using \eqref{eq:scaleFT} and $\widehat{\g}=\g$, we have
\[
\widehat{\g_{(\lambda)}*\g_{(\mu)}}(\xi)
=
\widehat{\g_{(\lambda)}}(\xi)\widehat{\g_{(\mu)}}(\xi)
=
\g(\lambda \xi)\g(\mu \xi)
=
e^{-\pi(\lambda^2+\mu^2)\xi^2}.
\]
Since
\[
e^{-\pi(\lambda^2+\mu^2)\xi^2}
=
\g(\sqrt{\lambda^2+\mu^2}\,\xi)
=
\widehat{\g_{(\sqrt{\lambda^2+\mu^2})}}(\xi),
\]
Fourier inversion gives
\[
\g_{(\lambda)}*\g_{(\mu)}
=
\g_{(\sqrt{\lambda^2+\mu^2})}.
\]
\end{proof}


\begin{proof}[Proof of Proposition \ref{auxiliary function B}]
Set
\[
    \alpha:=\sqrt{\pi},\qquad
    E(\xi):=e^{-\alpha|\xi|},\qquad
    s(\xi):=\sqrt{1-\g(\xi)}.
\]
Then \(B=1-s-E\). The function \(B\) is smooth on \(\R\setminus\{0\}\), even, and
continuous on \(\R\), since \(s(0)=0\) and \(E(0)=1\).

We first control the origin. For \(\xi\ge0\), put \(u=\alpha \xi\). Since
\[
    1-e^{-u^2}=u^2\int_0^1 e^{-tu^2}\,dt,
\]
we can write
\[
    s(\xi)=uR(u),\qquad
    R(u):=\left(\int_0^1 e^{-tu^2}\,dt\right)^{1/2}.
\]
For \(0\le u\le1\), elementary differentiation under the integral sign gives
\[
    \frac12\le R(u)\le1,\qquad |R'(u)|\le1,\qquad |R''(u)|\le5.
\]
Thus, with
\[
    F(u):=1-uR(u)-e^{-u},
\]
we have \(B(\xi)=F(\alpha\xi)\) for \(\xi\ge0\), \(F(0)=F'(0)=0\), and
\[
    |F''(u)|
    =
    |-2R'(u)-uR''(u)-e^{-u}|
    \le 8
    \qquad (0\le u\le1).
\]
Since \(\alpha^{-1}>1/2\), it follows that
\[
    |B''(\xi)|\le 8\alpha^2=8\pi\le26
    \qquad (0<|\xi|\le1/2),
\]
and \(B'\) extends continuously to \(0\) by setting \(B'(0)=0\).

It remains to estimate the tails. If \(|\xi|\ge1/2\), then \(s(\xi)\ge1/2\). For \(\xi>0\),
\[
    B'(\xi)
    =
    -\frac{\pi\xi\g(\xi)}{s(\xi)}
    +
    \alpha e^{-\alpha\xi},
\]
and
\[
    B''(\xi)
    =
    \frac{-\pi\g(\xi)+2\pi^2\xi^2\g(\xi)}{s(\xi)}
    +
    \frac{\pi^2\xi^2\g(\xi)^2}{s(\xi)^3}
    -
    \pi e^{-\alpha\xi}.
\]
By evenness, the same absolute value bounds hold on the negative half-line. Hence, for
\(|\xi|\ge1/2\),
\[
    |B(\xi)|\le \g(\xi)+E(\xi),
\]
\[
    |B'(\xi)|\le 2\pi|\xi|\g(\xi)+\alpha E(\xi),
\]
and
\[
    |B''(\xi)|
    \le
    2\pi\g(\xi)+12\pi^2\xi^2\g(\xi)+\pi E(\xi).
\]

We now integrate. On \(|\xi|\le1/2\), the trivial bound \(|B|\le3\) gives contribution at
most \(3\). On the complement,
\[
    \int_\R(\g+E)
    =
    1+\frac{2}{\alpha}
    \le3.
\]
Thus \(\|B\|_1\le8\).

For \(B'\), the local bound \(|B'(\xi)|\le26|\xi|\) on \(|\xi|\le1/2\) gives local
contribution at most \(7\). The tail contributes at most
\[
    2\pi\int_\R |\xi|\g(\xi)\,d\xi
    +
    \alpha\int_\R E(\xi)\,d\xi
    =
    2+2=4.
\]
Thus \(\|B'\|_1\le20\).

Finally, for \(B''\), the local contribution is at most \(26\). The tail contributes at most
\[
\begin{aligned}
    \int_\R
    \bigl(2\pi\g(\xi)+12\pi^2\xi^2\g(\xi)+\pi E(\xi)\bigr)\,d\xi
    &=
    2\pi+12\pi^2\frac{1}{2\pi}+\pi\frac{2}{\alpha}  \\
    &=
    8\pi+2\sqrt{\pi}
    \le30.
\end{aligned}
\]
Therefore
\[
    \|B''\|_1\le 26+30=56.
\]
This proves the claimed integrability and the explicit estimates. The stated Borel measurability
of the extension of \(B''\) is immediate from continuity on each half-line and the definition at
the single point \(0\).
\end{proof}

\begin{proof}[Proof of Proposition \ref{square root of Gaussian decay}]
Set
\[
    m(\xi):=1-\sqrt{1-\g(\xi)}.
\]
We first construct \(\rho\). By Proposition \ref{auxiliary function B},
\[
    m(\xi)=B(\xi)+e^{-\sqrt{\pi}|\xi|}.
\]
Both terms on the right-hand side are integrable, so \(m\in L^1(\R)\). Define
\[
    \rho:=\widecheck m.
\]
We will prove that \(\rho\in W_0(\R)\), that it satisfies the claimed decay estimate, and that
\(\widehat\rho=m\).

Let \(\alpha=\sqrt{\pi}\) and \(E_\alpha(\xi)=e^{-\alpha|\xi|}\). By Proposition
\ref{poisson to abel},
\[
    E_\alpha(\xi)=\widehat p(\alpha \xi).
\]
By the scaling rule for the Fourier transform,
\[
    \widecheck{E_\alpha}(x)
    =
    \alpha^{-1}p(x/\alpha)
    =
    \frac{2\alpha}{\alpha^2+(2\pi x)^2}
    =
    \frac{2/\sqrt{\pi}}{1+4\pi x^2}.
\]
In particular,
\begin{equation}\label{eq:abel-decay}
    0\le \widecheck{E_\alpha}(x)\le 2\langle x\rangle^{-2}.
\end{equation}

It remains to estimate \(\widecheck B\). By Proposition \ref{auxiliary function B},
\[
    \|B\|_1\le 8,
    \qquad
    \|B''\|_1\le 100.
\]
Therefore
\[
    |\widecheck B(x)|\le 8
    \qquad (x\in\R).
\]
For \(x\neq 0\), integration by parts twice gives
\[
    |\widecheck B(x)|
    \le
    \frac{\|B''\|_1}{4\pi^2x^2}
    \le
    \frac{100}{4\pi^2x^2}.
\]
Equivalently, for all \(x\in\R\),
\begin{equation}\label{eq:Bcheck-decay}
    |\widecheck B(x)|
    \le
    \left(8+\frac{100}{4\pi^2}\right)\langle x\rangle^{-2}.
\end{equation}
Combining \eqref{eq:abel-decay} and \eqref{eq:Bcheck-decay}, we obtain
\[
    |\rho(x)|
    \le
    \left(10+\frac{100}{4\pi^2}\right)\langle x\rangle^{-2}
    \le
    100\langle x\rangle^{-2}.
\]
The function \(\rho\) is continuous because \(m\in L^1\). The preceding decay estimate implies
\(\rho\in W_0(\R)\): indeed, for \(|x|\le 2\) the local \(W_0\)-supremum is bounded, while for
\(|x|>2\) and \(|x-y|\le 1\) we have \(|y|\ge |x|/2\), so the local supremum is bounded by a
constant multiple of \(\langle x\rangle^{-2}\), which is integrable on \(\R\).
It remains to prove non-negativity. For \(0\le t\le 1\), the binomial series gives
\[
    1-\sqrt{1-t}
    =
    \sum_{k=1}^\infty a_k t^k,
    \qquad
    a_k:=(-1)^{k+1}\binom{1/2}{k}>0.
\]
Moreover \(\sum_{k=1}^\infty a_k=1\). Hence, for every \(\xi\in\R\),
\[
    m(\xi)
    =
    \sum_{k=1}^\infty a_k \g(\xi)^k
    =
    \sum_{k=1}^\infty a_k \g(\sqrt{k}\xi).
\]
By Proposition \ref{Elementary Gaussian properties},
\[
    \g(\sqrt{k}\xi)=\widehat{\g_{(\sqrt{k})}}(\xi).
\]
The series
\[
    \sum_{k=1}^\infty a_k \g_{(\sqrt{k})}(x)
\]
converges in \(L^1(\R)\), because \(\|\g_{(\sqrt{k})}\|_1=1\) and
\(\sum_k a_k=1\). Therefore termwise Fourier transformation is justified, and uniqueness of the
Fourier transform on \(L^1\) gives
\[
    \rho(x)
    =
    \sum_{k=1}^\infty a_k \g_{(\sqrt{k})}(x).
\]
Every term on the right-hand side is non-negative, so \(\rho(x)\ge 0\) for every \(x\in\R\).
\end{proof}

\subsection{Proof of Proposition \ref{derivative of diagonal square root}}\label{sec:derivative of diagonal square root}

\begin{proof}
By algebra, we obtain
\[
\widehat{s}(\xi)
=
\g(t_02^{-\frac12}\xi)\sqrt{1-\g(r\xi)},
\]
where $r=(t_1^2-t_0^2)^{\frac12}$.
Let $\rho$ be the function from Proposition
\ref{square root of Gaussian decay} satisfying
\[
\widehat{\rho}(\xi)=1-\sqrt{1-\g(\xi)}.
\]
Then
\[
\widehat{s}(\xi)
=
\g(\sigma\xi)(1-\widehat{\rho}(r\xi)),
\]
where $\sigma=t_02^{-\frac12}$.

The power series identity
\[
1-\sqrt{1-z}
=
\sum_{k=1}^{\infty}
\frac{1}{2k-1}\binom{2k}{k}4^{-k}z^k
\]
has nonnegative coefficients whose sum is $1$. Since
\[
\|\g(\sqrt{k}\,\cdot)\|_1=k^{-\frac12},
\]
Proposition \ref{Elementary Gaussian properties} gives
\[
\rho(x)
=
\sum_{k=1}^{\infty}
\frac{1}{2k-1}\binom{2k}{k}4^{-k}
\g_{(\sqrt{k})}(x).
\]
The series and its differentiated series converge uniformly by
Proposition \ref{Gaussian bump decay}. Hence
\[
\rho'(x)
=
\sum_{k=1}^{\infty}
\frac{1}{2k-1}\binom{2k}{k}4^{-k}
\g_{(\sqrt{k})}'(x).
\]
By Proposition \ref{Gaussian bump decay},
\[
\left|\g_{(\sqrt{k})}'(x)\right|
\le
C_{\ref{Gaussian bump decay},1,2}
k^{-\frac12}\langle x\rangle_{(\sqrt{k})}^2
=
C_{\ref{Gaussian bump decay},1,2}
(\sqrt{k}+|x|)^{-2}
\le
C_{\ref{Gaussian bump decay},1,2}
\langle x\rangle^2.
\]
It follows that
\[
|\rho'(x)|
\le
C_{\ref{Gaussian bump decay},1,2}
\langle x\rangle^2.
\]
Consequently,
\[
|(\rho_{(r)})'(x)|
\le
C_{\ref{Gaussian bump decay},1,2}
r^{-1}\langle x\rangle_{(r)}^2.
\]

We have
\[
s=\g_{(\sigma)}-\g_{(\sigma)}*\rho_{(r)}
\]
and therefore
\[
s'=\g_{(\sigma)}'
-\g_{(\sigma)}*(\rho_{(r)})'.
\]
We use Proposition \ref{Gaussian bump decay} to obtain
\[
|\g_{(\sigma)}'(x)|
\le
C_{\ref{Gaussian bump decay},1,N}
\sigma^{-1}\langle x\rangle_{(\sigma)}^N
\]
and Proposition \ref{two bump estimate}, using $r>\sigma$ by
the assumptions on $t_i$, $i=0,1$, to obtain
\[
\begin{aligned}
|(\g_{(\sigma)}*(\rho_{(r)})')(x)|
\le{}&
C_{\ref{Gaussian bump decay},0,2}
C_{\ref{Gaussian bump decay},1,2}
C_{\ref{two bump estimate},2,2}
\\
&\times
r^{-1}\langle x\rangle_{(r)}^2.
\end{aligned}
\]
This gives
\[
\begin{aligned}
|s'(x)|
\le{}&
\max\bigl(
C_{\ref{Gaussian bump decay},1,N},
C_{\ref{Gaussian bump decay},0,2}
C_{\ref{Gaussian bump decay},1,2}
C_{\ref{two bump estimate},2,2}
\bigr)
\\
&\times
\bigl(
\sigma^{-1}\langle x\rangle_{(\sigma)}^N
+
r^{-1}\langle x\rangle_{(r)}^2
\bigr).
\end{aligned}
\]
Since $\sigma=t_02^{-\frac12}$ and
$\frac{\sqrt{3}}{2}t_1\le r\le t_1$, we have
\[
\sigma^{-1}\langle x\rangle_{(\sigma)}^N
\le
2t_0^{-1}\langle x\rangle_{(t_0)}^N
\]
and
\[
r^{-1}\langle x\rangle_{(r)}^2
\le
2t_1^{-1}\langle x\rangle_{(t_1)}^2.
\]
The desired estimate follows.
\end{proof}


\subsection{Proof of Proposition \ref{lem:increase-data-bracket-domination}}
\label{sec:proof-increase-data-bracket-domination}
\begin{proof}
Set
\[
r=\frac76,
\qquad
C_\rho=C_{\ref{lem:rho-kernels-reduction}},
\]
and, for \(s>0\), define
\[
\omega_s(x)=\langle x\rangle_{(s)}^r.
\]


We first record several elementary estimates that will be used
throughout the proof. Since
\[
r=\frac76<\frac32<2,
\]
we have
\begin{equation}
    \label{E:increase-data-weaken-exponents}
    \langle x\rangle_{(s)}^2
    \leq
    \langle x\rangle_{(s)}^{3/2}
    \leq
    \omega_s(x).
\end{equation}
Moreover,
\begin{equation}
    \label{E:increase-data-r-bump-L1}
    \|\omega_s\|_1
    =
    2\int_0^\infty(1+x)^{-7/6}\,dx
    =
    12
    \leq 2^4.
\end{equation}
Directly from the definition of the bracket bumps,
\begin{equation}
    \label{E:increase-data-fixed-scaling}
    \omega_s(x/\sqrt2)\leq2\omega_s(x),
    \qquad
    \omega_s(\sqrt2x)\leq\omega_s(x).
\end{equation}

Proposition~\ref{two bump estimate}, applied with both exponents
equal to \(r\), gives
\begin{equation}
    \label{E:increase-data-two-bump-corollary}
    \int_{\mathbb R}
    \omega_s(x-p)\omega_t(y-p)\,dp
    \leq
    2^6
    \omega_{\max(s,t)}(x-y).
\end{equation}
Indeed,
\[
C_{\ref{two bump estimate},r,r}
=
2^{1+r}\left(1+(r-1)^{-1}\right)
=
7\cdot2^{13/6}
<
2^6.
\]

We will also use the following consequence of Propositions~
\ref{bump triangle} and~\ref{orthogonal decay}. Suppose that
\[
v_0=w_0-w_1,
\qquad
v_1=w_0+w_1.
\]
Then
\[
W_1v=(\sqrt2w_0,\sqrt2w_1).
\]
Let
\[
e_0=(1,0),
\qquad
e_1=(0,1),
\]
and
\[
\theta_0=2^{-1/2}(1,1),
\qquad
\theta_1=2^{-1/2}(-1,1).
\]
Thus
\[
v\cdot\theta_0=(W_1v)_0,
\qquad
v\cdot\theta_1=(W_1v)_1,
\qquad
w_1=2^{-1/2}v\cdot\theta_1.
\]
For either \(i=0\) or \(i=1\), the constant in Proposition~
\ref{orthogonal decay}, applied to \(e_i\) and \(\theta_1\), is
\[
(2|e_i\cdot\theta_0|^{-1})^r
=
(2\sqrt2)^{7/6}
=
2^{7/4}
<
2^2.
\]
Combining this with
\eqref{E:increase-data-fixed-scaling}, we obtain
\begin{equation}
    \label{E:increase-data-mixed-zero}
    \begin{aligned}
    \omega_s(v_0)\omega_t(w_1)
    \leq 2^3\Big(
    &
    \omega_s(v_0)\omega_t(v_1)
    \\
    &+
    \omega_s((W_1v)_0)
    \omega_t((W_1v)_1)
    \Big)
    \end{aligned}
\end{equation}
and
\begin{equation}
    \label{E:increase-data-mixed-one}
    \begin{aligned}
    \omega_s(v_1)\omega_t(w_1)
    \leq 2^3\Big(
    &
    \omega_t(v_0)\omega_s(v_1)
    \\
    &+
    \omega_s((W_1v)_0)
    \omega_t((W_1v)_1)
    \Big).
    \end{aligned}
\end{equation}

We next express \(N_{h,j}\) as a convolution. Let
\(\rho_{h,j}\) be as in Lemma~
\ref{lem:rho-kernels-reduction}. By the definition of
\(\sigma_{h,j}=s(2^ha,j)\),
\[
\widehat{\sigma_{h,j}}(\zeta)
=
\left(
\g(2^ha(j-1)\zeta)
-
\g(2^ha(j)\zeta)
\right)^{1/2}.
\]
It follows from Proposition~\ref{convolution vector} that
\begin{equation}
    \label{E:increase-data-N-rho-convolution}
    N_{h,j}
    =
    M_j\ast_{(1,1)}\rho_{h,j}.
\end{equation}

Fix \(h\in\mathbb N\), \(j\in\mathbb Z\), and \(v\in\mathbb R^2\),
and write
\[
v_0=w_0-w_1,
\qquad
v_1=w_0+w_1.
\]
Changing variables in
\eqref{E:increase-data-N-rho-convolution}, we obtain
\begin{equation}
    \label{E:increase-data-convolution-coordinates}
    N_{h,j}(v)
    =
    \int_{\mathbb R}
    \rho_{h,j}(w_0+p)
    M_j(-w_1-p,w_1-p)\,dp.
\end{equation}

We first treat the case \(h\geq1\). Set
\[
\lambda=2^ha(j).
\]
By Lemma~\ref{lem:affine-diagonal-cancellation-reduction},
applied with \(x=-2w_1\),
\[
\int_{\mathbb R}
M_j(-w_1-p,w_1-p)\,dp=0.
\]
Consequently,
\begin{equation}
    \label{E:increase-data-cancelled-convolution}
    \begin{aligned}
    N_{h,j}(v)
    =
    \int_{\mathbb R}
    \bigl(
    \rho_{h,j}(w_0+p)-\rho_{h,j}(w_0)
    \bigr)
    M_j(-w_1-p,w_1-p)\,dp.
    \end{aligned}
\end{equation}
Lemma~\ref{lem:rho-kernels-reduction}, together with
\eqref{E:increase-data-weaken-exponents}, gives
\begin{equation}
    \label{E:increase-data-cancelled-majorant}
    \begin{aligned}
    |N_{h,j}(v)|
    \leq
    C_\rho
    \sum_{\alpha\in[2)}
    \int_{\mathbb R}
    &\min(1,\lambda^{-1}|p|)
    \omega_\lambda(w_0+\alpha p)
    \\
    &\times
    |M_j(-w_1-p,w_1-p)|\,dp.
    \end{aligned}
\end{equation}

Fix
\[
(t_0,t_1,u)\in\mathcal P,
\]
and suppress the argument \(j\) in \(t_0(j)\) and \(t_1(j)\).
The additional one-sided assumption \eqref{E:increase-data-one-sided-scales} implies
\begin{equation}
    \label{E:increase-data-scale-comparison}
    t_\ell
    \leq
    a(j)
    \leq
    2^ha(j)
    =
    \lambda,
    \qquad
    \ell\in[2).
\end{equation}

Suppose first that \(u=0\). Then
\[
\begin{aligned}
&
\langle(-w_1-p)\rangle_{(t_0)}^{3/2}
\langle(w_1-p)\rangle_{(t_1)}^{3/2}
\\
&\qquad=
\langle w_1+p\rangle_{(t_0)}^{3/2}
\langle w_1-p\rangle_{(t_1)}^{3/2}.
\end{aligned}
\]
Since
\[
2|p|
\leq
|w_1+p|+|w_1-p|
\]
and \(\min(1,s)\leq s^{1/3}\), we have
\[
\begin{aligned}
\min(1,\lambda^{-1}|p|)
&\leq
2^{-1/3}\lambda^{-1/3}
\left(
|w_1+p|^{1/3}
+
|w_1-p|^{1/3}
\right)
\\
&\leq
2^{-h/3}
\Big(
(1+t_0^{-1}|w_1+p|)^{1/3}
\\
&\hspace{34mm}
+
(1+t_1^{-1}|w_1-p|)^{1/3}
\Big),
\end{aligned}
\]
where we used
\[
t_\ell/\lambda\leq2^{-h}.
\]
Multiplying by the two \(3/2\)-bumps and using
\[
\frac32-\frac13=\frac76=r,
\]
we obtain
\begin{equation}
    \label{E:increase-data-loss-u-zero}
    \begin{aligned}
    &\min(1,\lambda^{-1}|p|)
    \langle w_1+p\rangle_{(t_0)}^{3/2}
    \langle w_1-p\rangle_{(t_1)}^{3/2}
    \\
    &\qquad\leq
    2 \cdot 2^{-h/3}
    \omega_{t_0}(w_1+p)
    \omega_{t_1}(w_1-p).
    \end{aligned}
\end{equation}

Suppose now that \(u=1\). Since
\[
W_1(-w_1-p,w_1-p)
=
(-\sqrt2p,\sqrt2w_1),
\]
we have
\[
\begin{aligned}
\min(1,\lambda^{-1}|p|)
&\leq
(\lambda^{-1}|p|)^{1/3}
\\
&=
2^{-1/6}
(t_0/\lambda)^{1/3}
(t_0^{-1}|\sqrt2p|)^{1/3}
\\
&\leq
 2^{-h/3}
(1+t_0^{-1}|\sqrt2p|)^{1/3}.
\end{aligned}
\]
It follows that
\begin{equation}
    \label{E:increase-data-loss-u-one}
    \begin{aligned}
    &\min(1,\lambda^{-1}|p|)
    \langle\sqrt2p\rangle_{(t_0)}^{3/2}
    \langle\sqrt2w_1\rangle_{(t_1)}^{3/2}
    \\
    &\qquad\leq
     2\cdot 2^{-h/3}
    \omega_{t_0}(\sqrt2p)
    \omega_{t_1}(\sqrt2w_1).
    \end{aligned}
\end{equation}

We now integrate the expressions resulting from
\eqref{E:increase-data-loss-u-zero} and
\eqref{E:increase-data-loss-u-one}.

Assume first that \(u=0\). For \(\alpha=0\), let
\[
\tau=\max(t_0,t_1).
\]
By \eqref{E:increase-data-two-bump-corollary},
\[
\begin{aligned}
&
\omega_\lambda(w_0)
\int_{\mathbb R}
\omega_{t_0}(w_1+p)
\omega_{t_1}(w_1-p)\,dp
\\
&\qquad\leq
2^6
\omega_\lambda(w_0)
\omega_\tau(2w_1).
\end{aligned}
\]
Using
\[
(W_1v)_0=\sqrt2w_0,
\qquad
(W_1v)_1=\sqrt2w_1
\]
and \eqref{E:increase-data-fixed-scaling}, this is at most
\begin{equation}
    \label{E:increase-data-u0-alpha0}
    2^7
    \omega_\lambda((W_1v)_0)
    \omega_\tau((W_1v)_1).
\end{equation}

For \(\alpha=1\), Proposition~\ref{bump triangle}, applied with
\[
w_1
=
\frac12(w_1+p)+\frac12(w_1-p),
\]
has constant one and gives
\begin{equation}
    \label{E:increase-data-narrow-triangle}
    \begin{aligned}
    &
    \omega_{t_0}(w_1+p)
    \omega_{t_1}(w_1-p)
    \\
    &\qquad\leq
    \omega_{t_0}(w_1+p)\omega_{t_1}(w_1)
    +
    \omega_{t_0}(w_1)\omega_{t_1}(w_1-p).
    \end{aligned}
\end{equation}
Since \(\lambda\geq t_0,t_1\), Proposition~
\ref{two bump estimate} gives
\[
\begin{aligned}
&
\int_{\mathbb R}
\omega_\lambda(w_0+p)
\omega_{t_0}(w_1+p)
\omega_{t_1}(w_1-p)\,dp
\\
&\qquad\leq
2^6\Big(
\omega_\lambda(v_0)\omega_{t_1}(w_1)
+
\omega_\lambda(v_1)\omega_{t_0}(w_1)
\Big).
\end{aligned}
\]
Applying \eqref{E:increase-data-mixed-zero} and
\eqref{E:increase-data-mixed-one}, the last expression is at most
\begin{equation}
    \label{E:increase-data-u0-alpha1}
    \begin{aligned}
    2^9\Big(
    &
    \omega_\lambda(v_0)\omega_{t_1}(v_1)
    +
    \omega_\lambda((W_1v)_0)
    \omega_{t_1}((W_1v)_1)
    \\
    &+
    \omega_{t_0}(v_0)\omega_\lambda(v_1)
    +
    \omega_\lambda((W_1v)_0)
    \omega_{t_0}((W_1v)_1)
    \Big).
    \end{aligned}
\end{equation}

Suppose next that \(u=1\). For \(\alpha=0\),
\eqref{E:increase-data-r-bump-L1} gives
\[
\begin{aligned}
&
\omega_\lambda(w_0)
\int_{\mathbb R}\omega_{t_0}(\sqrt2p)\,dp
\,
\omega_{t_1}(\sqrt2w_1)
\\
&\qquad\leq
2^4
\omega_\lambda(w_0)
\omega_{t_1}(\sqrt2w_1)
\\
&\qquad\leq
2^5
\omega_\lambda((W_1v)_0)
\omega_{t_1}((W_1v)_1).
\end{aligned}
\]
For \(\alpha=1\), we use
\[
\omega_{t_0}(\sqrt2p)\leq\omega_{t_0}(p)
\]
and Proposition~\ref{two bump estimate}. Since
\(\lambda\geq t_0\), we obtain
\[
\begin{aligned}
&
\int_{\mathbb R}
\omega_\lambda(w_0+p)
\omega_{t_0}(\sqrt2p)\,dp
\,
\omega_{t_1}(\sqrt2w_1)
\\
&\qquad\leq
2^6
\omega_\lambda(w_0)
\omega_{t_1}(\sqrt2w_1)
\\
&\qquad\leq
2^7
\omega_\lambda((W_1v)_0)
\omega_{t_1}((W_1v)_1).
\end{aligned}
\]
Thus, for both \(\alpha\in[2)\),
\begin{equation}
    \label{E:increase-data-u1-alpha}
    \begin{aligned}
    &
    \int_{\mathbb R}
    \omega_\lambda(w_0+\alpha p)
    \omega_{t_0}(\sqrt2p)
    \omega_{t_1}(\sqrt2w_1)\,dp
    \\
    &\qquad\leq
    2^7
    \omega_\lambda((W_1v)_0)
    \omega_{t_1}((W_1v)_1).
    \end{aligned}
\end{equation}

Combining
\eqref{E:increase-data-cancelled-majorant},
\eqref{E:increase-data-loss-u-zero},
\eqref{E:increase-data-loss-u-one},
\eqref{E:increase-data-u0-alpha0},
\eqref{E:increase-data-u0-alpha1}, and
\eqref{E:increase-data-u1-alpha}, we conclude that, when
\(h\geq1\), every resulting product of two \(r\)-bumps occurs with
coefficient at most
\[
2\cdot2^9 C_\rho\,2^{-h/3}
=
2^{10}C_\rho\,2^{-h/3}.
\]
Every such product is in either the \(W_0\)- or the
\(W_1\)-coordinates. Its two scale sequences are among
\[
2^ha,\qquad
t_0,\qquad
t_1,\qquad
\max(t_0,t_1).
\]

It remains to treat \(h=0\). Define the two sequences
\[
d_0(j)=a(j-1),
\qquad
d_1(j)=a(j).
\]
By Lemma~\ref{lem:rho-kernels-reduction},
\eqref{E:increase-data-convolution-coordinates}, and
\eqref{E:increase-data-weaken-exponents},
\begin{equation}
    \label{E:increase-data-h-zero-initial}
    \begin{aligned}
    |N_{0,j}(v)|
    \leq
    C_\rho
    \sum_{\epsilon\in[2)}
    \int_{\mathbb R}
    \omega_{d_\epsilon(j)}(w_0+p)
    |M_j(-w_1-p,w_1-p)|\,dp.
    \end{aligned}
\end{equation}

Fix again \((t_0,t_1,u)\in\mathcal P\), and define
\[
r_{\epsilon,\ell}
=
\max(d_\epsilon,t_\ell),
\qquad
\epsilon,\ell\in[2).
\]
If \(u=0\), then Proposition~\ref{bump triangle}, followed by
Proposition~\ref{two bump estimate}, gives
\[
\begin{aligned}
&
\int_{\mathbb R}
\omega_{d_\epsilon(j)}(w_0+p)
\omega_{t_0(j)}(w_1+p)
\omega_{t_1(j)}(w_1-p)\,dp
\\
&\qquad\leq
2^6\Big(
\omega_{r_{\epsilon,0}(j)}(v_0)
\omega_{t_1(j)}(w_1)
+
\omega_{r_{\epsilon,1}(j)}(v_1)
\omega_{t_0(j)}(w_1)
\Big).
\end{aligned}
\]
Using
\eqref{E:increase-data-mixed-zero} and
\eqref{E:increase-data-mixed-one}, this is bounded by
\begin{equation}
    \label{E:increase-data-h0-u0}
    \begin{aligned}
    2^9\Big(
    &
    \omega_{r_{\epsilon,0}(j)}(v_0)
    \omega_{t_1(j)}(v_1)
    \\
    &+
    \omega_{r_{\epsilon,0}(j)}((W_1v)_0)
    \omega_{t_1(j)}((W_1v)_1)
    \\
    &+
    \omega_{t_0(j)}(v_0)
    \omega_{r_{\epsilon,1}(j)}(v_1)
    \\
    &+
    \omega_{r_{\epsilon,1}(j)}((W_1v)_0)
    \omega_{t_0(j)}((W_1v)_1)
    \Big).
    \end{aligned}
\end{equation}

If \(u=1\), then
\[
\omega_{t_0(j)}(\sqrt2p)
\leq
\omega_{t_0(j)}(p).
\]
Proposition~\ref{two bump estimate} and
\eqref{E:increase-data-fixed-scaling} therefore give
\begin{equation}
    \label{E:increase-data-h0-u1}
    \begin{aligned}
    &
    \int_{\mathbb R}
    \omega_{d_\epsilon(j)}(w_0+p)
    \omega_{t_0(j)}(\sqrt2p)
    \omega_{t_1(j)}(\sqrt2w_1)\,dp
    \\
    &\qquad\leq
    2^7
    \omega_{r_{\epsilon,0}(j)}((W_1v)_0)
    \omega_{t_1(j)}((W_1v)_1).
    \end{aligned}
\end{equation}
Thus every product produced in the case \(h=0\) has coefficient at
most
\[
2^9C_\rho
\leq
2^{10}C_\rho.
\]

We now collect all products obtained above, retaining
multiplicities. Each product has the form
\[
\omega_{\beta^0(j)}((W_uv)_0)
\omega_{\beta^1(j)}((W_uv)_1)
\]
for some \(u\in[2)\) and some
\(\beta=(\beta^0,\beta^1)\in A^2\).
For every such product introduce one label \(b\), and denote the
resulting index set by \(\mathcal B_h\). Set
\[
u_{h,b}=u,
\qquad
\beta_{h,b}=\beta.
\]
Since
\[
2^{10}C_\rho
=
C_{\ref{lem:increase-data-bracket-domination},1},
\]
the preceding estimates prove
\[
\begin{aligned}
|N_{h,j}(v)|
&\leq
C_{\ref{lem:increase-data-bracket-domination},1}
2^{-h/3}
\\
&\quad\times
\sum_{b\in\mathcal B_h}
\langle(W_{u_{h,b}}v)_0
\rangle_{(\beta_{h,b}^0(j))}^{7/6}
\langle(W_{u_{h,b}}v)_1
\rangle_{(\beta_{h,b}^1(j))}^{7/6}.
\end{aligned}
\]

We verify the cardinality bound. When \(h\geq1\), a summand of the
pointwise majorant for \(M_j\) with \(u=0\) produces one term from
\(\alpha=0\) and four terms from \(\alpha=1\), hence at most five
terms. A summand with \(u=1\) produces at most two terms. Therefore,
\[
\#\mathcal B_h
\leq
5\#\mathcal P
\leq
25
\]
when \(h\geq1\). When \(h=0\), a summand with \(u=0\) produces four
terms for each \(\epsilon\in[2)\), hence at most eight terms, while
a summand with \(u=1\) produces at most two terms. Thus
\[
\#\mathcal B_0
\leq
8\#\mathcal P
\leq
40.
\]
Consequently,
\[
\#\mathcal B_h
\leq
40
\leq
2^6
=
C_{\ref{lem:increase-data-bracket-domination},0}.
\]

Finally, we verify the distance bound. If \(h\geq1\), every scale
sequence occurring above is one of
\[
2^ha,\qquad
t_0,\qquad
t_1,\qquad
\max(t_0,t_1).
\]
By Proposition~\ref{Properties of distance of sequences},
\[
\dist(a,2^ha)\leq h.
\]
By hypothesis,
\[
\dist(a,t_\ell)\leq1.
\]
Moreover, since
\[
a(j-1)\leq t_\ell(j)\leq a(j+1),
\]
we also have
\[
a(j-1)
\leq
\max(t_0(j),t_1(j))
\leq
a(j+1),
\]
and hence
\[
\dist(a,\max(t_0,t_1))\leq1.
\]

When \(h=0\), every scale sequence occurring above is one of
\[
t_0,\qquad
t_1,\qquad
\max(d_\epsilon,t_\ell).
\]
The sequence \(d_0=a(\,\cdot-1)\) satisfies
\[
\dist(a,d_0)\leq1,
\]
while \(d_1=a\). Since both \(d_\epsilon\) and \(t_\ell\) belong to
\(B_{\dist}(a,1)\), their pointwise maximum also belongs to
\(B_{\dist}(a,1)\). All these sequences belong to \(A\) by
Proposition~\ref{Operations on spaced sequences}.

It follows in all cases that
\[
\dist(a,\beta_{h,b}^\ell)
\leq
1+h
\]
Finally, by symmetry and the triangle inequality for \(\dist\),
\[
\begin{aligned}
\Delta(\beta_{h,b})
&=
\dist(\beta_{h,b}^0,\beta_{h,b}^1)
\\
&\leq
\dist(\beta_{h,b}^0,a)
+
\dist(a,\beta_{h,b}^1)
\\
&\leq
2(1+h),
\end{aligned}
\]
which is
\eqref{E:increase-data-beta-pair-distance}.
\end{proof}


\subsection{Proof of Proposition \ref{lem:increase-data-Gaussian-expansion} }
\label{sec:proof-increase-data-Gaussian-expansion}


\begin{proof}
Set
\[
C_G=C_{\ref{Gaussian domination}}
\qquad\text{and}\qquad
C_B=C_{\ref{lem:increase-data-bracket-domination},1}.
\]
We prove the proposition with
\begin{equation}
    \label{E:increase-data-Gaussian-expansion-constant}
    C_{\ref{lem:increase-data-Gaussian-expansion}}
    =
    2^3C_G^2C_B.
\end{equation}

Fix \(h\in\mathbb N\), \(b\in\mathcal B_h\), and
\(m=(m_0,m_1)\in\mathbb N^2\). Since
\(\beta_{h,b}^\ell\in A\) for \(\ell\in[2)\), Proposition~
\ref{Operations on spaced sequences} implies that
\[
2^{m_\ell}\beta_{h,b}^\ell\in A.
\]
Consequently,
\[
p_{h,b,m}
=
\left(
2^{m_0}\beta_{h,b}^0,
2^{m_1}\beta_{h,b}^1
\right)
\in A^2.
\]

We first prove the distance estimate. By the triangle inequality
for \(\dist\), followed by Proposition~
\ref{Properties of distance of sequences},
\[
\begin{aligned}
\Delta(p_{h,b,m})
&=
\dist\left(
2^{m_0}\beta_{h,b}^0,
2^{m_1}\beta_{h,b}^1
\right)
\\
&\leq
\dist\left(
2^{m_0}\beta_{h,b}^0,
2^{m_0}\beta_{h,b}^1
\right)
\\
&\quad+
\dist\left(
2^{m_0}\beta_{h,b}^1,
2^{m_1}\beta_{h,b}^1
\right)
\\
&=
\dist\left(
\beta_{h,b}^0,
\beta_{h,b}^1
\right)
+
\dist\left(
\beta_{h,b}^1,
2^{m_1-m_0}\beta_{h,b}^1
\right)
\\
&\leq
\Delta(\beta_{h,b})
+
|m_1-m_0|
\\
&\leq
\Delta(\beta_{h,b})
+
m_0+m_1.
\end{aligned}
\]
Here the equality in the fourth line follows by multiplying both
sequences in the first distance by \(2^{-m_0}\), and the subsequent
inequality follows from
\[
\dist(c,2^qc)\leq |q|,
\qquad q\in\mathbb Z.
\]
Since \(|m|=m_0+m_1\), estimate
\eqref{E:increase-data-beta-pair-distance} now gives
\[
\Delta(p_{h,b,m})
\leq
2(1+h)
+
|m|,
\]
which is \eqref{E:increase-data-p-distance}.

We next prove the Gaussian majorant. Fix \(j\in\mathbb Z\) and
\(v\in\mathbb R^2\). Proposition~\ref{Gaussian domination}, applied
with
\[
N=\frac76,
\qquad
s=\beta_{h,b}^\ell(j),
\qquad
x=(W_{u_{h,b}}v)_\ell,
\]
gives, for each \(\ell\in[2)\),
\begin{equation}
    \label{E:increase-data-one-coordinate-Gaussian}
    \begin{aligned}
    &
    \left\langle
    (W_{u_{h,b}}v)_\ell
    \right\rangle_{(\beta_{h,b}^\ell(j))}^{7/6}
    \\
    &\qquad\leq
    C_G2^{7/6}
    \sum_{m_\ell\in\mathbb N}
    2^{-m_\ell/6}
    \g_{(2^{m_\ell}\beta_{h,b}^\ell(j))}
    \left((W_{u_{h,b}}v)_\ell\right).
    \end{aligned}
\end{equation}
All terms in these series are nonnegative. Multiplying
\eqref{E:increase-data-one-coordinate-Gaussian} for
\(\ell=0\) and \(\ell=1\), we obtain
\[
\begin{aligned}
&
\left\langle
(W_{u_{h,b}}v)_0
\right\rangle_{(\beta_{h,b}^0(j))}^{7/6}
\left\langle
(W_{u_{h,b}}v)_1
\right\rangle_{(\beta_{h,b}^1(j))}^{7/6}
\\
&\quad\leq
C_G^2 2^{7/3}
\sum_{m\in\mathbb N^2}
2^{-|m|/6}
\\
&\qquad\qquad\times
\g_{(2^{m_0}\beta_{h,b}^0(j))}
\left((W_{u_{h,b}}v)_0\right)
\g_{(2^{m_1}\beta_{h,b}^1(j))}
\left((W_{u_{h,b}}v)_1\right).
\end{aligned}
\]
By Definition~\ref{2D Gaussians}, the product of Gaussians on the
right-hand side is
\[
G_{p_{h,b,m}(j),u_{h,b}}(v).
\]
Since
\[
2^{7/3}\leq2^3,
\]
it follows that
\begin{equation}
    \label{E:increase-data-two-coordinate-Gaussian}
    \begin{aligned}
    &
    \left\langle
    (W_{u_{h,b}}v)_0
    \right\rangle_{(\beta_{h,b}^0(j))}^{7/6}
    \left\langle
    (W_{u_{h,b}}v)_1
    \right\rangle_{(\beta_{h,b}^1(j))}^{7/6}
    \\
    &\qquad\leq
    2^3C_G^2
    \sum_{m\in\mathbb N^2}
    2^{-|m|/6}
    G_{p_{h,b,m}(j),u_{h,b}}(v).
    \end{aligned}
\end{equation}

Applying \eqref{E:increase-data-two-coordinate-Gaussian} to every
summand in
\eqref{E:increase-data-bracket-majorant}, we conclude that
\[
\begin{aligned}
|N_{h,j}(v)|
&\leq
2^3C_G^2C_B
2^{-h/3}
\\
&\quad\times
\sum_{b\in\mathcal B_h}
\sum_{m\in\mathbb N^2}
2^{-|m|/6}
G_{p_{h,b,m}(j),u_{h,b}}(v).
\end{aligned}
\]
In view of
\eqref{E:increase-data-Gaussian-expansion-constant}, this is
\eqref{E:increase-data-Gaussian-majorant}.

It remains to verify the two convergence assertions. Fix
\(h\in\mathbb N\) and \(j\in\mathbb Z\). Since
\[
0\leq\g_{(s)}(x)
=
s^{-1}\g(s^{-1}x)
\leq s^{-1},
\]
we have, for every \(b\in\mathcal B_h\), \(m\in\mathbb N^2\),
and \(v\in\mathbb R^2\),
\[
\begin{aligned}
G_{p_{h,b,m}(j),u_{h,b}}(v)
&\leq
\frac{1}{
2^{m_0}\beta_{h,b}^0(j)
}
\frac{1}{
2^{m_1}\beta_{h,b}^1(j)
}
\\
&=
\frac{2^{-|m|}}{
\beta_{h,b}^0(j)\beta_{h,b}^1(j)
}.
\end{aligned}
\]
Therefore,
\[
\begin{aligned}
&
\sum_{m\in\mathbb N^2}
2^{-|m|/6}
G_{p_{h,b,m}(j),u_{h,b}}(v)
\\
&\qquad\leq
\frac{1}{
\beta_{h,b}^0(j)\beta_{h,b}^1(j)
}
\sum_{m\in\mathbb N^2}
2^{-7|m|/6}
\\
&\qquad=
\frac{1}{
\beta_{h,b}^0(j)\beta_{h,b}^1(j)
}
\left(
\sum_{q\in\mathbb N}2^{-7q/6}
\right)^2
\\
&\qquad\leq
\frac{2^2}{
\beta_{h,b}^0(j)\beta_{h,b}^1(j)
},
\end{aligned}
\]
where in the final line we used
\[
2^{-7q/6}\leq2^{-q}
\qquad\text{and}\qquad
\sum_{q\in\mathbb N}2^{-q}=2.
\]
Since \(\mathcal B_h\) is finite, the full nonnegative series
converges for every \(v\in\mathbb R^2\).

Finally, since \(W_{u_{h,b}}\) is unitary and every rescaled
one-dimensional Gaussian has integral one,
\[
\left\|
G_{p_{h,b,m}(j),u_{h,b}}
\right\|_1
=
1.
\]
Consequently,
\[
\begin{aligned}
&
\sum_{b\in\mathcal B_h}
\sum_{m\in\mathbb N^2}
\left\|
2^{-|m|/6}
G_{p_{h,b,m}(j),u_{h,b}}
\right\|_1
\\
&\qquad=
\#\mathcal B_h
\left(
\sum_{q\in\mathbb N}2^{-q/6}
\right)^2.
\end{aligned}
\]
Writing \(q=6d+r\), where \(d\in\mathbb N\) and \(r\in[6)\), gives
\[
\begin{aligned}
\sum_{q\in\mathbb N}2^{-q/6}
&=
\sum_{r\in[6)}
2^{-r/6}
\sum_{d\in\mathbb N}2^{-d}
\\
&\leq
6\cdot2
=
12
\leq2^4.
\end{aligned}
\]
Thus
\[
\begin{aligned}
&
\sum_{b\in\mathcal B_h}
\sum_{m\in\mathbb N^2}
\left\|
2^{-|m|/6}
G_{p_{h,b,m}(j),u_{h,b}}
\right\|_1
\\
&\qquad\leq
2^8\#\mathcal B_h
\leq
2^8
C_{\ref{lem:increase-data-bracket-domination},0}
<
\infty.
\end{aligned}
\]
The series therefore converges absolutely in
\(L^1(\mathbb R^2)\), completing the proof.
\end{proof}









\iffalse


\newpage
\section{Old subsections 6.2 and 6.3 - will eventually be deleted}

\subsection{\texorpdfstring{Off-diagonal from main argument}{Off-diagonal from main argument}}\label{sec:off-diag}
The goal of this section is to prove off-diagonal estimates Propositions \ref{prop:1-off-W} and \ref{lem:1-off-nonW} using Theorem \ref{induct positive terms theorem} for $k=2$.

Fix a universal pair $(\phi_0,\phi_1)$. For $\lambda>0$ and $i\in[2)$ write
\[ \phi_{i,\lambda} = (\phi_i)_{(\lambda)} \]

% \jr{i.e. in some way redoing the final step of the induction. This seems like the most straightforward way of doing it, but it would be nice if this was not necessary. But this may require heavily reworking the final reduction, which is probably not worth doing.}
% Throughout this section, let 
% \[\beta(n) = \left \{\begin{array}{ll}
%    2  & \text{if }n=2 \\
%    1  & \text{if }n>2
% \end{array} \right .\]

\pd{$N$ in   props}\jr{can set this to $8$, for example}


\begin{lemma}\label{all pieces vanishing}
Let $a\in \rm A$ and $f:\R\to\C$ continuous and compactly supported. 
Suppose $J\ge 1$ and $\sum_{j\in [J)} f(a(j) \xi)=0$ for all $\xi\in\R$. Then we must have $f(\xi)=0$ for all $\xi\in\R$.
\end{lemma}

\begin{proof}
Suppose not. Then since $f$ has compact support we may choose $\xi_0\not=0$ so that $f(\xi_0)\not=0$ and $f(\xi)=0$ for all $\xi$ with $|\xi|\ge 2 |\xi_0|$. 
Then set $\xi_1 = \xi_0 / a(0)$. Then
\[ 0 = \sum_{j\in [J)} f(a(j)\xi_1) = f(a(0)\xi_1) + \sum_{j\ge 1} f(a(j)\xi_1) \]
But since $a(j)\ge 2a(0)$ for $j\ge 1$ we have $|a(j)\xi_1|\ge 2 a(0) |\xi_1| = 2 |\xi_0|$, so $f(a(j)\xi_1)=0$ for all $j\ge 1$.
Thus $0=f(a(0)\xi_1)=f(\xi_0)$, contradiction.
\end{proof}


 \jr{This is a very technical variant with essentially an identical proof. Do we want to unify the Gaussian dominations?}\pd{We could try to unify at least with "reduction variant 2" below, which is $h=0$ but sometimes $a(j+1)$ in place of $a(j)$}
\begin{proposition}[Gaussian domination, reduction variant]\label{Gaussian domination reduction variant}   Let $\beta\in \{1,2\}$. 
There exist constants $C_0,C_1,C_2\in (0,\infty)$ such that the following holds:
Suppose $(M_j)_{j\in\Z}\subset W_0(\R^2)$and $\widehat{M_j}(\xi,-\xi)=0$ for all $\xi\in\R, j\in\Z$ 
and assume $a\in\rm A$ and $\kappa\ge 0$, $N\in\R, N\ge 2$ such that
\[ |M_j(v)| \le 2^{-\kappa/2}\langle v_0+v_1\rangle_{(2^{\kappa} a(j))}^N \langle v_0-v_1\rangle_{(a(j))}^N + \langle v_0+v_1\rangle_{(a(j))}^N \langle v_0-v_1\rangle_{(a(j))}^N\]
Let $h\in\Z$ and let $N_{j,h}=N_{j,h,\beta}$ be defined by
\[\widehat{N_{j,h}}(\xi,\eta) = (\widehat{\psi_{(2^h a(j))}}-\widehat{\psi_{(2^{h+1}a(j))}})(\xi+\eta) \widehat{\sigma_{j,h}}(\xi+\eta)^{-\beta} \widehat{M_j}(\xi,\eta),\]
\[ \widehat{\sigma_{j,h}} = \widehat{s(2^{h+1} a,j)}. \]

Then there exists a finite set $\mathcal{B}=\mathcal{B}_h$ with $\#\mathcal{B}\le C_0$ and for every $b\in \mathcal{B}$ a $u_b\in [2)$ and for every $(b,m)\in \mathcal{B}\times \N^2$ a $p_{b,m}\in {\rm A}^2$ and a $w_b\in [2)$ such that
\[ \Delta(p_{b,m})\le C_1 (1 + w_b\kappa+ |h|+|m|). \]
and for every $j,h\in\Z$, $v\in\R^2$,
\[ |N_{j,h}(v)| \le C_2 \sum_{m\in\mathbb{N}^2}  2^{-|m|/2}\sum_{b\in \mathcal{B}_{h}} 2^{-\max(w_b \kappa,h_+)/2} G_{p_{b,m}(j),u_b}(v). \]
\end{proposition}

The proof follows the same steps as in the proof of Proposition \ref{Gaussian domination combined}. The  machine-generated details can be found in \S \ref{Gaussian domination reduction proof}. 



 

\begin{lemma} 
\label{reduction apply Gaussian}  
Let $\beta\in \{1,2\}$. 
There exist a constants $C_{\ref{reduction apply Gaussian}}\in (0,\infty)$ such that the following holds. 
Suppose $(M_j)_{j\in\Z}\subset W_0(\R^2)$and $\widehat{M_j}(\xi,-\xi)=0$ for all $\xi\in\R, j\in\Z$ 
and assume $a\in\rm A$ and $\kappa\ge 0$, $N\in\R, N\ge 2$ such that
\[ |M_j(v)| \le 2^{-\kappa/2}\langle v_0+v_1\rangle_{(2^{\kappa} a(j))}^N \langle v_0-v_1\rangle_{(a(j))}^N + \langle v_0+v_1\rangle_{(a(j))}^N \langle v_0-v_1\rangle_{(a(j))}^N\]
Let $h\in\Z$ and let $N_{j,h} = N_{j,h,\beta}$ be defined by
\[\widehat{N_{j,h}}(\xi,\eta) = (\widehat{\psi_{(2^h a(j))}}-\widehat{\psi_{(2^{h+1}a(j))}})(\xi+\eta) \widehat{\sigma_{j,h}}(\xi+\eta)^{-\beta} \widehat{M_j}(\xi,\eta),\]
\[ \widehat{\sigma_{j,h}} = \widehat{s(2^{h+1} a,j)}. \]
 For $j\in \Z$ let 
    \[\tilde{M}_j(y) = |N_{j,h}(y_0)| (\sigma_{j,h}\otimes \sigma_{j,h})(y_1). \]
    Then with $\mathbf{\tilde{M}}=(\tilde{M})_{j\in \Z}$,
    \[\|\mathbf{\tilde{M}}\|_{\rm M(2)}\le C_{\ref{reduction apply Gaussian}}\]
\end{lemma}


\begin{proof}
 Let $J\ge 1$ and $\F\in \mathfrak{F}$.    By Proposition \ref{Gaussian domination reduction variant}, Proposition \ref{Monotonicity K}, and the dominated convergence theorem we estimate 
\[ |\Lambda_2(\sum_{j\in [J)}\tilde{M}_j)(\F)| \le  \sum_{m\in\mathbb{N}^2} 2^{-|m|/2} \sum_{b\in \mathcal{B}_{h}} 2^{-\max(w_b \kappa,h_+)/2}
|\Lambda_2(\sum_{j\in [J)} \M(\gamma_{b,m}, s_{\gamma_{b,m}}^{\otimes 2}, 1)_j)(\F)|, \]
where for every $(b,m)\in \mathcal{B}_h\times\N^2$, the data $\gamma_{b,m}=(2, u_{b,m}, a_{b,m})\in\Gamma$ is given by
\[ u_{b,m}=(u_b, 0)\in [2)^2, \]
\[ (a_{b,m})_{i'} = \left\{\begin{array}{ll}
p_{b,m} & \text{if}\;i'=0\\
\alpha & \text{if}\;i'=1,
\end{array}\right.
\] %Note $a_{b,m}$ is an element of $(\rm A^2)^2$.
where $\alpha=(2^{h+1/2} a, 2^{h+1/2} a)\in {\rm A}^2$
and $\mathcal{B}_h, p_{b,m}, u_b, w_b$ are as given by Proposition \ref{Gaussian domination reduction variant} with $\Delta(p_{b,m})\le C(1+w_b\kappa+|h|+|m|)$. 
Note $(s_{\gamma_{b,m}})_{1,j}=\sigma_{j,h}$.
We also have $\Delta(\alpha)=0$, so
\[ \Delta_{\gamma_{b,m}} \le C(1 +w_b\kappa+ |h| + |m|). \]
Thus by Theorem \ref{induct positive terms theorem} with $k=2$ and $i=1$, for every $h\in\Z$ and every $(b,m)\in\mathcal{B}_h\times \N^2$
\[ \|\M(\gamma_{b,m}, s_{\gamma_{b,m}}^{\otimes 2}, 1)\|_{\rm M(2)} \lesssim (1+w_b\kappa+|h|+|m|)^{2-2^{3-n}}. \]
Summing over $h\in\Z_{\ge -100}$ and $(b,m)\in \mathcal{B}_h\times \N^2$ now gives the claim:
\[ \sum_{h\in\Z_{\ge -100}} \sum_{m\in\mathbb{N}^2} 2^{-|m|/2} \sum_{b\in \mathcal{B}_{h}} 2^{-\max(w_b \kappa,h_+)/2} (1+w_b\kappa+|h|+|m|)^{2-2^{3-n}} \lesssim 1, \]
uniformly in $\kappa$.
\end{proof}




\begin{proposition}[Prism, off-diagonal, Whitney] \pd{done, now part of Prop \ref{P:diagonal-band-reduction}}
\label{prop:1-off-W}
 There exist $N=N_{\ref{prop:1-off-W}}\ge 1$ such that the following holds. Let $\kappa\ge 0$ and $a\in {\rm A}$. Let $M_j:\mathbb{R}^2\to \mathbb{R}$ be of the form
\[ M_j=\Phi_{(a(j))} \]
for every $j\in\Z$, where $\Phi:\mathbb{R}^2\to \mathbb{R}$ is a test function satisfying 
\[\textup{supp}(\widehat{\Phi}) \subset \textup{Ann}_2(1,2^{30})\]
Assume further that  for all $(u,v)\in \R^2$, 
\begin{equation}
|\Phi(u,v)| \leq  2^{-\frac12 \kappa} \langle u+v\rangle_{(2^\kappa)}^N
\langle u-v\rangle^N 
+\langle u+v\rangle^N \langle u-v\rangle^N.
\end{equation}
Assume that $\sum_{j\in [J)} \widehat{M_{j}}  (\xi,-\xi) =0$ for all $\xi\in\R$.
  Then with $\M=(M_j)_{j\in\Z}$,  \[\|{\M}\|_{{\rm M}(1)}\le C_{\ref{prop:1-off-W}}\]
  with $C_{\ref{prop:1-off-W}}=...$ 
\end{proposition}

\begin{proof}
    This follows by case distinction form Propositions \ref{prop:1-off-W case 1} and
     \ref{prop:1-off-W case 2}.
\end{proof}

\begin{proposition} \label{prop:1-off-W case 1}\pd{done, now part of Prop \ref{P:diagonal-band-reduction}}
    Proposition \ref{prop:1-off-W} holds in case $n>2$.
\end{proposition}

\begin{proof} 
First observe that the assumption self-improves to $\widehat{M_j}(\xi,-\xi)=0$ for all $\xi\in\R$ and all $j\in\Z$ (see Lemma \ref{all pieces vanishing}).

Write for $j\in\Z$
\[ M_j =  \sum_{h\in\mathbb{Z}}  N_{j,h} *_{1,1} \sigma_{j,h} \]
where for $h\in\Z$, $(\xi,\eta)\in\R^2,$
\[ \widehat{N_{j,h}}(\xi,\eta) = (\widehat{\psi_{(2^h a(j))}}-\widehat{\psi_{(2^{h+1}a(j))}})(\xi+\eta) \widehat{\sigma_{j,h}}(\xi+\eta)^{-1} \widehat{M_j}(\xi,\eta), \]
\[ \widehat{\sigma_{j,h}} = \widehat{s(2^{h+1} a,j)}. \]
Note that $N_{j,h}=0$ for $h<-100$ (say).
Fix $h\in\Z_{\ge -100}$. We apply Proposition \ref{Cauchy-Schwarz at k} for $k=1$   with $\rho_j=N_{j,h}$ and $\varphi_j=\sigma_{j,h}$. Then for all $J\ge 1$, $\F\in\mathfrak{F}$, 
\[ |\Lambda_1(\sum_{j\in [J)}M_j)(\F)| \le |\Lambda_2(\sum_{j\in [J)} \tilde{M}_j)(\F)|^{1/2} J^{1/2} \]
where for $y\in (\R^2)^2$,
% \[ \tilde{M}(y) = \sum_{j\in [J)} \tilde{M}_j(y), \]
\[ \tilde{M}_j(y) = |N_{j,h}(y_0)| (\sigma_{j,h}\otimes \sigma_{j,h})(y_1). \]
Proposition \ref{reduction apply Gaussian}   yields
\[\|(\tilde{M}_j)_{j\in \Z}\|_{\rm M(2)} \le C_{\ref{reduction apply Gaussian}}\]
which finishes the proof.
\end{proof}


\begin{proposition} \label{prop:1-off-W case 1} \pd{done, now part of Prop \ref{P:diagonal-band-reduction}}
    Proposition \ref{prop:1-off-W} holds in case $n=2$.
\end{proposition}

\begin{proof}
We again observe that the assumption self-improves to $\widehat{M_j}(\xi,-\xi)=0$ for all $\xi\in\R$ and all $j\in\Z$ (see Lemma \ref{all pieces vanishing}).

Write for $j\in\Z$
\[ M_j =  \sum_{h\in\mathbb{Z}}  N_{j,h} *_{1,1} (\sigma_{j,h}*\sigma_{j,h}) \]
where for $h\in\Z$, $(\xi,\eta)\in\R^2,$
\[ \widehat{N_{j,h}}(\xi,\eta) = (\widehat{\psi_{(2^h a(j))}}-\widehat{\psi_{(2^{h+1}a(j))}})(\xi+\eta) \widehat{\sigma_{j,h}}(\xi+\eta)^{-2} \widehat{M_j}(\xi,\eta), \]
\[ \widehat{\sigma_{j,h}} = \widehat{s(2^{h+1} a,j)}. \]
Note that $N_{j,h}=0$ for $h<-100$ (say).
Fix $h\in\Z_{\ge -100}$. We apply Proposition \ref{Cauchy-Schwarz at n-1} for $n=2$   with $\rho_j=N_{j,h}$ and $\varphi_j=\sigma_{j,h}$. Then for all $\F\in\mathfrak{F}$
\[ |\Lambda_1(\sum_{j\in [J)}M_j)(\F)| \le \sup_{\mathbf{\tilde{F}}\in\mathfrak{F}}|\Lambda_{2}(\sum_{j\in [J)} \tilde{M}_j)(\mathbf{\tilde{F}})|. \]
where for $y\in (\R^2)^2$,
% \[ \tilde{M}(y) = \sum_{j\in [J)} \tilde{M}_j(y), \]
\[ \tilde{M}_j(y) = |N_{j,h}(y_0)| (\sigma_{j,h}\otimes \sigma_{j,h})(y_1). \]
Proposition \ref{reduction apply Gaussian}  yields
\[\|(\tilde{M}_j)_{j\in \Z}\|_{\rm M(2)} \le C_{\ref{reduction apply Gaussian}},\]
which finishes the proof. 
\end{proof}


We now turn our attention to an on-diagonal non-Whitney proposition which has an additional gap assumption. This proposition will be used in the proofs of Propositions  \ref{prop:1-on-nonW} and \ref{prop:1-on-W}. Then, Proposition \ref{prop:1-on-W} will also be used in the proof of Proposition \ref{prop:1-on-nonW}.




The proof will use the following lemmas. 
 

\begin{lemma}
    \label{lem:Phi_bumpest}
    Let $N\ge 1$, $a\in A$, and    $b_1(j),b_2(j)\in \{a(j),a(j+1)\}$, $j\in \Z$. Then 
    \[\int_{\R} \langle u_0+ p\rangle^{3N+2}_{(b_1(j))} \langle u_1+ p\rangle^{N}_{(a(j+1))} \langle p\rangle^{2N+2}_{(b_2(j))}  dp \]
\[\le C_{\ref{lem:Phi_bumpest},N} ( \langle u_0\rangle^N_{(a(j))} \langle u_1\rangle^N_{(a(j+1))} + \langle u_0\rangle^N_{(a(j+1))} \langle u_1\rangle^N_{(a(j+1))} )\]
with $C_{\ref{lem:Phi_bumpest},N} =2^{8N}$. 
\end{lemma}
\begin{proof}
In the following estimates we will repeatedly use 
 \[\langle u_i+ p\rangle^{\nu}_{(b_i(j))}  \le 2^{\nu} \langle u_i\rangle^{\nu}_{(b_i(j))} \langle b_i(j)^{-1} p\rangle ^{-\nu}\]
valid for any $\nu\ge 0$, which follows from Lemma \ref{lem:transl_pow_v} applied to  $b_i(j)^{-1} u_i$ and $b_i(j)^{-1} p$.
 
Case 1:  $b_1(j)=b_2(j)=a(j+1)$. We estimate the   integral  by 
\[  2^{2N}\langle u_0\rangle^N_{(a(j+1))} \langle u_1\rangle^N_{(a(j+1))}\int_{\R}\langle a(j+1)^{-1} p\rangle ^{-2N} \langle p\rangle^{2N+2}_{(a(j+1))} \le 2^{2N+1} \langle u_0\rangle^N_{(a(j+1))} \langle u_1\rangle^N_{(a(j+1))},\]
 where we integrated in $p$ in the final step (the integral of $(1+|p|)^{-2}$ equals $2$).



Case 2: $b_1(j)=a(j),\, b_2(j)=a(j)$.
The integral is estimated by
\begin{equation}
    \label{case2_int}
    2^{2N}\langle u_0\rangle^{N}_{(a(j))} \langle u_1\rangle^{N}_{(a(j+1))} \int_{\R}   \langle a(j)^{-1} p\rangle ^{-N} \langle a(j+1)^{-1} p\rangle ^{-N} \langle  p\rangle_{a(j)} ^{2N+2}  dp.
\end{equation}
Observe
\[\langle a(j)^{-1} p\rangle ^{-N} \langle a(j+1)^{-1} p\rangle ^{-N} \langle  p\rangle_{a(j)} ^{2N+2}   =  \langle a(j+1)^{-1} p\rangle ^{-N}  \langle  p\rangle_{a(j)} ^{N+2} . \]
If  $a(j+1)^{-1} p \le 1$,  we estimate  $\langle a(j+1)^{-1} p\rangle ^{-N}\le 1$ and thus
\[\langle a(j+1)^{-1} p\rangle ^{-N}  \langle  p\rangle_{a(j)}^{N+2} \le \langle p\rangle_{a(j)}^2\]
 If   $a(j+1)^{-1} p \ge 1$, we  use   $a(j)/a(j+1)\le 1$, giving
 \[\langle a(j+1)^{-1} p\rangle ^{-N}  \langle  p\rangle_{a(j)} ^{N+2} \le 2^N ({a(j)}{a(j+1)^{-1}}   )^{N} \langle p\rangle_{a(j)}^{2} \le 2^N\langle p\rangle_{a(j)}^{2} \]
  Integrating in $p$,  \eqref{case2_int} is bounded 
by 
\[2^{3N+1}\langle u_0\rangle^N_{(a(j))} \langle u_1\rangle^N_{(a(j+1))}.\]

Case 3: $b_1(j)=a(j+1),\, b_2(j)=a(j)$.
We estimate the integral by
\begin{equation}
    \label{case3_int}
    2^{2N}\langle u_0\rangle^{N}_{(a(j+1))} \langle u_1\rangle^{N}_{(a(j+1))} \int_{\R}    \langle a(j+1)^{-1} p\rangle ^{-2N} \langle  p\rangle_{a(j)} ^{2N+2}  dp\end{equation}
    \[ \le 2^{4N+1}\langle u\rangle^N_{(a(j+1))} \langle v\rangle^N_{(a(j+1))}.\]

The last inequality follow as in Case 2, with $2N$ in place of $N$.
 

   
Case 4: $b_1(j)=a(j),\, b_2(j)=a(j+1)$.
 We first do a change of variables $p\to p-u_0$ which gives for the integral in question
 \[\int_{\R} \langle  p\rangle^{3N+2}_{(a(j))} \langle u_1-u_0+ p\rangle^{N}_{(a(j+1))} \langle p-u_0\rangle^{2N}_{(a(j+1))}  dp,\]
which is then estimated by 
\[ 2^{3N}\langle u_1-u_0\rangle^{N}_{(a(j+1))} \langle u_0\rangle^{2N}_{(a(j+1))} \int_{\R}   \langle a(j+1)^{-1} p\rangle ^{-3N}   \langle  p\rangle_{a(j)} ^{3N+2}  dp\]
 \[\le 2^{6N+1} \langle u_1-u_0\rangle^{N}_{(a(j+1))} \langle u_0\rangle^{2N}_{(a(j+1))} ,\]
estimating the $p$-integral as in Case 2.
But this  can now be   estimated by
\[\le 2^{7N+1} \langle u_0\rangle^{N}_{(a(j+1))} \langle u_1 \rangle^{N}_{(a(j+1))},\] 
as desired.
\end{proof}

 In the following we  write $\varphi_{i,a(j)} = (\varphi_i)_{a(j)}$. 
\begin{lemma}[Gaussian domination, reduction variant 2]
\label{lem:gauss-off-nonW}Let $\beta\in \{1,2\}$. 
For $0\le i \le 3$ let $\varphi_i$ be a test function satisfying satisfying for $N\in \R$,  $N\ge 8$, 
\begin{equation}
    \label{Phi_varphi_decay}
    |\varphi_i(x)|\le \langle x \rangle^{N}
\end{equation}
Assume further that $\widehat{\varphi_{3,2^{2}}}$ is  supported in $[-1,1]$ and constant $1$ in $[-1/2,1/2]$. 
 Let   $a\in A$  and  for $j\in \Z$   let   $N_j = N_{j,\beta}$ be defined by 
 \[\widehat{N_{j}}(\xi,\eta) = ((\widehat{\varphi_{0,a(j)}}-\widehat{\varphi_{1,a(j+1)}})\otimes \widehat{\varphi_{2,a(j+1)} })(\xi,\eta) (\widehat{\varphi_{3,a(j)}}-\widehat{\varphi_{3,a(j+1)}})(\xi+\eta)\widehat{\sigma_j}(\xi+\eta)^{-\beta},\]
  % \[\Phi_j = (({\phi_{0,a(j)}}-{\phi_{1,2^{-4}a(j+1)}}) \otimes   {\phi_{1,a(j+1)} }) *_{(1,1)}  {\rho_j}\]
\[\widehat{\sigma_j} =  \widehat{s(2^{-2}a,j+1)}  \]
Then for $b\in [2)$ and $m\in \N^2$ there exist  $p_{b,m} \in A^2$ such that  
\[\Delta(p_{b,m})\le 1+ |m|\]
and for all $u\in \R^2$, 
\[|N_{j}(u)| \le  C_{\ref{lem:gauss-off-nonW},N}\sum_{(m_1,m_2)\in {\N}^2}  2^{-|m|} \sum_{b\in [2)} G_{p_{b,m}(j),0}(u).
\]
with $C_{\ref{lem:gauss-off-nonW},N} = 8C_{\ref{four scale Gaussian kernel},0,N}C^2_{\ref{Gaussian domination}}C_{\ref{lem:Phi_bumpest}, N}$
\end{lemma}

\begin{proof}    Let $N$ be as in \eqref{Phi_varphi_decay}. 
Let 
\[\widehat{\rho_j}  = (\widehat{\varphi_{3,a(j)}}-\widehat{\varphi_{3,a(j+1)}})\cdot \widehat{\sigma_j}^{-\beta}\]
Using Lemma~\ref{four scale Gaussian kernel} with  $(\lambda_-,\lambda_+,\mu_-,\mu_+)=2^{-2}(a(j),a(j+1),a(j),a(j+1))$, $\nu=-\beta/2$, $m=0$, and $\psi=\varphi_{3,2^2}$ we estimate
\[ |\rho_j(x)| \le C_{\ref{four scale Gaussian kernel},0,N} (\langle x\rangle^{N}_{(2^{-2}a(j))} + \langle x\rangle^{N}_{(2^{-2}a(j+1))}) \le C_{\ref{four scale Gaussian kernel},0,N} (\langle x\rangle^{N}_{(a(j))} + \langle x\rangle^{N}_{(a(j+1))}) . \]
\ls{The assumptions of Proposition~\ref{four scale Gaussian kernel} are currently stated in terms of derivative bounds for $\widehat{\varphi_3}$, rather than as in~\eqref{Phi_varphi_decay}. However, maybe in applications, \eqref{Phi_varphi_decay} is proved by estimating derivatives of $\widehat{\varphi_3}$ anyway?}
Passing to the spatial side, 
\[N_j(u) =  \int_{\R}{\varphi_{0,a(j)}}-{\varphi_{1,a(j+1)}}(u_0+p)  {\varphi_{2,a(j+1)} } (u_1+p)  {\rho_j}(p) dp\]
By the triangle inequality, the bound on $\rho_j$, and the assumptions \eqref{Phi_varphi_decay}, 
\[|N_j(u)| \le   \int_{\R}|{\varphi_{0,a(j)}}-{\varphi_{1,a(j+1)}}|(u_0+p)  |{\varphi_{2,a(j+1)} }| (u_1+p)  |{\rho_j}|(p) dp \]
\[\le C_{\ref{four scale Gaussian kernel},0,N} \sum_{\substack{b_1(j),b_2(j)\in  \\ \{a(j),a(j+1)\}}} \int_{\R} \langle u_0+ p\rangle^{N}_{(b_1(j))} \langle u_1+ p\rangle^{N}_{(a(j+1))} \langle p\rangle^{N}_{(b_2(j))}  dp .\]
Lemma \ref{lem:Phi_bumpest} gives, using $N\ge 8$ (we apply the lemma with $(N-2)/3\ge 2$ in place of $N$ and estimate $C_{\ref{lem:Phi_bumpest}, (N-2)/3}\le C_{\ref{lem:Phi_bumpest}, N}$), 
\[|N_j(u)| \le 4C_{\ref{four scale Gaussian kernel},0,N} C_{\ref{lem:Phi_bumpest}, N} ( \langle u_0\rangle^2_{(a(j))} \langle u_1\rangle^2_{(a(j+1))} + \langle u_0\rangle^2_{(a(j+1))} \langle u_1\rangle^2_{(a(j+1))} )\]
\ls{Is this why we need to assume $N\geq 8$? Can we reduce to $N=2$ by mimicking the proof of Proposition~\ref{Gaussian domination combined}, case $h=0$? (Maybe this is what we were planning to do anyway by combining the various Gaussian dominations.)} \pd{Yes this is one place, it uses $N\ge 8$ but it is almost certainly inefficient, this was written before we decided to optimize $N$.}
Applying Gaussian domination from Lemma \ref{Gaussian domination}, once with  $s=a(j)$ and once with $s=a(j+1)$, gives 
\[|N_j(u)| \le 8C_{\ref{four scale Gaussian kernel},0,N}C^2_{\ref{Gaussian domination}}C_{\ref{lem:Phi_bumpest}, N} \sum_{m\in {\N}^2}  2^{-|m|} (G_{p_{0,m},0}+G_{p_{1,m},0}),\, \]
where   
\[p_{0,m}=(2^{m_0}a, 2^{m_1}a(\cdot +1)),\]
\[p_{1,m}=(2^{m_0}a(\cdot +1), 2^{m_1}a(\cdot +1)).\]
Then   $\Delta(p_{b,m}) \le 1+  |m|$. 
\end{proof}


\begin{lemma}
\label{reduction apply Gaussian 2}  
Let $\beta\in \{1,2\}$. 
For $0\le i \le 3$ let $\varphi_i$ be a test function satisfying satisfying for $N\in \R$,  $N\ge 8$, 
\begin{equation}
    \label{Phi_varphi_decay}
    |\varphi_i(x)|\le \langle x \rangle^{N}
\end{equation}
Assume further that $\widehat{\varphi_{3,2^{2}}}$ is  supported in $[-1,1]$ and constant $1$ in $[-1/2,1/2]$. 
 Let   $a\in A$  and  for $j\in \Z$   let   $N_j = N_{j,\beta}$ be defined by 
 \[\widehat{N_{j}}(\xi,\eta) = ((\widehat{\varphi_{0,a(j)}}-\widehat{\varphi_{1,a(j+1)}})\otimes \widehat{\varphi_{2,a(j+1)} })(\xi,\eta) (\widehat{\varphi_{3,a(j)}}-\widehat{\varphi_{3,a(j+1)}})(\xi+\eta)\widehat{\sigma_j}(\xi+\eta)^{-\beta},\]
  % \[\Phi_j = (({\phi_{0,a(j)}}-{\phi_{1,2^{-4}a(j+1)}}) \otimes   {\phi_{1,a(j+1)} }) *_{(1,1)}  {\rho_j}\]
\[\widehat{\sigma_j} =  \widehat{s(2^{-2}a,j+1)}  \]
 For $j\in \Z$ let 
    \[\tilde{M}_j(y) = |N_{j}(y_0)| \sigma_{j}^{\otimes 2}(y_1). \]
    Then with $\mathbf{\tilde{M}}=(\tilde{M})_{j\in \Z}$,
    \[\|\mathbf{\tilde{M}}\|_{\rm M(2)}\le C_{\ref{reduction apply Gaussian 2},N,n}\]
    with $C_{\ref{reduction apply Gaussian 2},N,n} = C_{\ref{lem:gauss-off-nonW},N}\,C_{\ref{induct positive terms theorem}}\,n^2$.
\end{lemma}



\begin{proof}
    We   estimate $N_j$ using the Gaussian domination in Lemma \ref{lem:gauss-off-nonW}, which gives $p_{0,m},p_{1,m}\in A^2$,  $\Delta(p_{b,m}) \le 1+  |m|$, such that 
\[|N_j| \le C_{\ref{lem:gauss-off-nonW},N}\sum_{m\in {\N}^2}  2^{-|m|}\sum_{b\in [2)}  G_{p_{b,m},0} \, \]
Therefore, 
\[\tilde{M}_j \le C_{\ref{lem:gauss-off-nonW},N}  \sum_{m\in {\N}^2}  2^{-|m|} \sum_{b\in [2)}  \tilde{M}_{b,j}  \]
where
\[\tilde{M}_{b,j}(y)=G_{p_{0,m},0}(y_0)\sigma_j^{\otimes 2} (y_1). \]

Let $J\ge 1$ and $\F\in \mathfrak{F}$.
By monotonicity (Lemma \ref{Monotonicity K}) and multilinearity of $\Lambda_2$, 
\[\Lambda_2(\sum_{j\in [J)}\tilde{M}_j)(\F) \le C_{\ref{lem:gauss-off-nonW},N} \sum_{m\in {\N}^2}2^{-|m|}  \sum_{b\in [2)} \Lambda_2(\sum_{j\in [J)}\tilde{M}_{b,j})(\F)  \]

 
We will now   apply Theorem \ref{induct positive terms theorem} for $k=2$, for which we need to verify that the expressions take the required form. 
For $b\in [2)$, the tuple  $\tilde{\M}_b = (\tilde{M}_{b,j})_{j\in \Z}$ can be recognized as $\M(\gamma_b, s_{\gamma_b}\otimes s_{\gamma_b}, 1)$ with $\gamma_b=(2,0,(p_{b,m},q_{b}))$, where    
\[q_{0} = (2^{-2}a(\cdot -1 ), 2^{-2}a(\cdot -1 ))\]
\[q_{1} = (2^{-2}a(\cdot -1 ), 2^{-2}a(\cdot -1 ))\]
Note $\Delta(q_{b})=0 $.

Applying Theorem \ref{induct positive terms theorem} for $k=2$  gives for $b\in [2)$
\[\|\tilde{\M}_b\|_{\rm M(2)} \le C_{\ref{induct positive terms theorem}}n^2 (1+|m|)^{2-2^{3-n}},\] 
This in turn gives 
\[\|\mathbf{\tilde{M}}\|_{\rm M(1)} \le C_{\ref{lem:gauss-off-nonW},N}C_{\ref{induct positive terms theorem}}n^2   \sum_{|m|\in {\N}^2}  2^{-|m|} (1+ |m| )^{2-2^{3-n}} \le C_{\ref{lem:gauss-off-nonW},N}C_{\ref{induct positive terms theorem}}n^2 \]
and completes the proof.
\end{proof}

 

\begin{lemma}[Prism, off-diagonal, non-Whitney] \pd{done, now part of Prop \ref{P:diagonal-band-reduction}}
\label{lem:1-off-nonW}
 Let    $a\in \rm A$ be such that $a(j+1)\ge 2^{10} a(j)$ $j\in \Z$. Let \[M_j=(\phi_{0,a(j)}-\phi_{1,2^{-4}a(j+1)})\otimes  \phi_{1,a(j+1)}.\]
Then with $\M=(M_j)_{j\in \Z}$, \[\|\M\|_{\rm M(1)}\le C_{\ref{lem:1-off-nonW},n}\]
where $C_{\ref{lem:1-off-nonW},n} = (2^{N_0+6}C_0)^3 C_{\ref{reduction apply Gaussian 2},N_0,n}$ .
\end{lemma}


\begin{proof}
This follows by case distinction from Lemma \ref{lem:1-off-nonW case 1} and Lemma \ref{lem:1-off-nonW case 2}
\end{proof}

\begin{lemma}\label{lem:1-off-nonW case 1}\pd{done, now part of Prop \ref{P:diagonal-band-reduction}}
    Lemma \label{lem:1-off-nonW} holds in case $n>2$.
\end{lemma}
 
\begin{proof} 
If $(\xi,\eta)$ is in the support of  \[\widehat{M}_j=(\widehat{\phi_{0,a(j)}}-\widehat{\phi_{1,2^{-4}a(j+1)}})  \otimes  \widehat{\phi_{1,a(j+1)} },\] then 
\begin{equation}
    \label{k1region}
    2^3a(j+1)^{-1} - a(j+1)^{-1}\le |\xi+\eta|\le  a(j)^{-1} + a(j+1)^{-1}.
\end{equation}
Note that  $\widehat{(\phi_0)_{2^{-2}{a}(j)}}(\xi+\eta)$ is constant one and $\widehat{(\phi_0)_{2^{-2}{a}(j+1)}}(\xi+\eta)$ is constant zero whenever \eqref{k1region} holds.
 So we can write   
\[\widehat{M_j}(\xi,\eta) =   \widehat{N_j}(\xi,\eta) \widehat{\sigma_j}(\xi+\eta)\]
with  
\[
\widehat{N_j}(\xi,\eta)=\widehat{M_j} (\xi,\eta) (\widehat{(\phi_0)_{2^{-2}a(j)}}-\widehat{(\phi_0)_{2^{-2}a(j+1)}})(\xi+\eta) \widehat{\sigma_j}(\xi+\eta)^{-1},  \]
\[\widehat{\sigma_j} = \widehat{s(2^{-2}a,j+1)} \]
Passing to the spatial side using Lemma \ref{convolution vector} we have
\[ M_j =  \int_{\R} (N_j*_{(1,1)} \sigma_j)(p) dp\]
We now apply the  Cauchy-Schwarz Proposition \ref{Cauchy-Schwarz at k} with $\rho_j=N_j$ and $\varphi_j = \sigma_j$:  for any $J\ge 1$,  $\F\in \mathfrak{F}$,  
\begin{equation}
|\Lambda_1(\sum_{j\in [J)}M_j)(\F)| \le |\Lambda_2(\sum_{j\in [J)}\tilde{M}_j)(\F)|^{\frac12} J^{\frac12},
\end{equation} 
where for $y=(y_0,y_1)\in (\R^2)^2$,
\begin{equation}
\tilde{M}_j(y)= |N_j(y_0)| \sigma_j^{\otimes 2} (y_1).
\end{equation}
We will apply Proposition \ref{reduction apply Gaussian 2}   with
$\varphi_0=\phi_0$, $\varphi_1 = \phi_{1,2^{-4}}$, $\varphi_2=\phi_1$, $\varphi_3 = \phi_{0,2^{-2}}$. 
For these $\varphi_i$ we have the estimates $|\varphi_i(x)|\le 2^{N_0+6}C_0\langle x \rangle^{N_0}$ which follow from Lemma \ref{lem:smoothdecay2} and the derivative bound  on windows (where $N_0$ is as in  the definition of the universal pair)
% \[|\phi_0(x)| \le 2^{N+1}C_0\langle x \rangle \]
% \[|\phi_{0,2^{-2}}(x)| \le  2^{N+3}C_0\langle x \rangle\]
% \[|\phi_1(x)|\le 2^{N+1}C_0\langle x \rangle\]
% \[|\phi_{1,2^{-4}}(x)| \le 2^{N+5}C_0\langle x \rangle\]


Proposition \ref{reduction apply Gaussian 2} now gives 
\[\|(\tilde{M}_j)_{j\in \Z}\|_{\rm M(2)}\le (2^{N_0+6}C_0)^3 C_{\ref{reduction apply Gaussian 2},N_0,n},\] which implies the desired estimate.
\end{proof}
 
\begin{lemma}\label{lem:1-off-nonW case 2}\pd{done, now part of Prop \ref{P:diagonal-band-reduction}}
    Lemma \label{lem:1-off-nonW} holds in case $n=2$.
\end{lemma}
 
\begin{proof}
If $(\xi,\eta)$ is in the support of  \[\widehat{M}_j=(\widehat{\phi_{0,a(j)}}-\widehat{\phi_{1,2^{-4}a(j+1)}})  \otimes  \widehat{\phi_{1,a(j+1)} },\] then 
\begin{equation}
    \label{k1region_2}
    2^3a(j+1)^{-1} - a(j+1)^{-1}\le |\xi+\eta|\le  a(j)^{-1} + a(j+1)^{-1}.
\end{equation}
Note that  $\widehat{(\phi_0)_{2^{-2}{a}(j)}}(\xi+\eta)$ is constant one and $\widehat{(\phi_0)_{2^{-2}{a}(j+1)}}(\xi+\eta)$ is constant zero whenever \eqref{k1region_2} holds.
 So we can write   
\[\widehat{M_j}(\xi,\eta) =   \widehat{N_j}(\xi,\eta) \widehat{\sigma_j}(\xi+\eta)^{2}\]
with  
\[
\widehat{N_j}(\xi,\eta)=\widehat{M_j} (\xi,\eta) (\widehat{(\phi_0)_{2^{-2}a(j)}}-\widehat{(\phi_0)_{2^{-2}a(j+1)}})(\xi+\eta) \widehat{\sigma_j}(\xi+\eta)^{-2},  \]
\[\widehat{\sigma_j} = \widehat{s(2^{-2}a,j+1)} \]
Passing to the spatial side using Lemma \ref{convolution vector} we have
\[ M_j =  \int_{\R} N_j*_{(1,1)} (\sigma_j*\sigma_j)(p) dp\]
We now apply the  Cauchy-Schwarz lemma \ref{Cauchy-Schwarz at n-1} with $\rho_j=N_j$ and $\varphi_j = \sigma_j$:  for any $J\ge 1$,  $\F\in \mathfrak{F}$,  
\begin{equation}
|\Lambda_1(\sum_{j\in [J)}M_j)(\F)| \le \sup_{\mathbf{\tilde{F}}\in\mathfrak{F}}|\Lambda_{2}(\sum_{j\in [J)}\tilde{M}_j)(\mathbf{\tilde{F}})|.
\end{equation} 
where for $y=(y_0,y_1)\in (\R^2)^2$,
\begin{equation}
\tilde{M}_j(y)= |N_j(y_0)| \sigma_j^{\otimes 2} (y_1).
\end{equation}
Proposition \ref{reduction apply Gaussian 2}  now gives $\|(\tilde{M}_j)_{j\in \Z}\|_{\rm M(2)}\le C_{\ref{lem:1-off-nonW}, n}$, which implies the desired estimate.

% We now estimate $N_j$ using the Gaussian domination in Lemma \ref{lem:gauss-off-nonW}, which we apply with $\varphi_0 = \phi_0$, $\varphi_1 = \phi_{1,2^{-4}}$, $\varphi_2=\phi_1$, $\varphi_3 = \phi_{0,2^{-2}}$. 
% Lemma \ref{lem:gauss-off-nonW} gives $p_{0,m},p_{1,m}\in A^2$,  $\Delta(p_{b,m}) \le 1+  |m|$, such that 
% \[|N_j| \le C\sum_{m\in {\N}^2}  2^{-|m|}\sum_{b\in [2)}  G_{p_{b,m},0} \, \]
% Therefore, 
% \[\tilde{M}_j \le  \sum_{m\in {\N}^2}  2^{-|m|} \sum_{b\in [2)}  \tilde{M}_{b,j}  \]
% where
% \[\tilde{M}_{b,j}(y)=G_{p_{0,m},0}(y_0)\sigma_j^{\otimes 2} (y_1). \]
% Thus, by monotonicity (Lemma \ref{Monotonicity K}) and multilinearity of the form $\Lambda_2$, for any $\mathbf{\tilde{F}}\in \mathfrak{F}$, 
% \[\Lambda_2(\sum_{j\in [J)}\tilde{M}_j)(\mathbf{\tilde{F}}) \le C  \sum_{m\in {\N}^2}2^{-|m|}  \sum_{b\in [2)} \Lambda_2(\sum_{j\in [J)}\tilde{M}_{b,j})(\mathbf{\tilde{F}})  \]
 
 

% Our goal is now to apply Theorem \ref{induct positive terms theorem} for $k=2$, for which we need to verify that the expressions take the required form. 
% For $b\in [2)$, the tuple  $\tilde{\M}_b = (\tilde{M}_{b,j})_{j\in \Z}$ can be recognized as $\M(\gamma_b, s_{\gamma_b}\otimes s_{\gamma_b}, 1)$ with $\gamma_b=(2,0,(p_{b,m},q_{b}))$, where    
% \[q_{0} = (2^{-2}a(\cdot -1 ), 2^{-2}a(\cdot -1 ))\]
% \[q_{1} = (2^{-2}a(\cdot -1 ), 2^{-2}a(\cdot -1 ))\]
% Note $\Delta(q_{b})=0 $.

% Applying Theorem \ref{induct positive terms theorem} for $k=2$  gives for $b\in [2)$
% \[\|\tilde{\M}_b\|_{\rm M(2)} \le C (1+|m|)^{2-2^{3-n}},\] 
% and thus  
% \[\|\M\|_{\rm M(1)} \le C    \sum_{|m|\in {\N}^2}  2^{-|m|} (1+ |m| )^{2-2^{3-n}} \le C\]
% This completes the proof.
\end{proof}




\subsection{On-diagonal from off-diagonal estimates}\label{sec:on-diag}


In this section we prove Propositions \ref{prop:1-on-W}, \ref{prop:1-on-W-cor}, and \ref{prop:1-on-nonW}  
 using Propositions \ref{prop:1-off-W} and \ref{lem:1-off-nonW}. 
 

 

\begin{prop}[$k=1$ prism, on-diagonal, non-Whitney, with gap]\label{prop:1-on-nonW-gap} \todo{Prop \ref{P:induct-positive-terms-reduction-non-whitney}}
Let    $a\in \rm A$ with $a(j+1)\ge 2^{10} a(j)$ for each $j\in \Z$.  For $j\in \Z$ let 
\[M_j=({\phi}_{0,a(2j)}-{\phi}_{1,a(2j+1)})^{\otimes 2}.\]
Then with $\M=(M_j)_{j\in \Z}$, \[\|\M\|_{\rm M(1)}\le C_{\ref{prop:1-on-nonW-gap},n}\]
where $C_{\ref{prop:1-on-nonW-gap},n} =   2^4(C_{\ref{lem:smoothdecay2},2}C_0)^2 + 2C_{\ref{lem:1-off-nonW},n} + 2^{6N_{\ref{prop:1-off-W}}+10} C_0^4 C_{\ref{lem:smoothdecay2}, 2N_{\ref{prop:1-off-W}}}  C_{\ref{prop:1-off-W},N_{\ref{prop:1-off-W}}}$
% $2^2(2C_{\ref{lem:smoothdecay2},2}C_0)^2 + 2C_{\ref{lem:1-off-nonW},N,n}+  2^6 2^{2N_{\ref{prop:1-off-W}}} (2^{2N_{\ref{prop:1-off-W}}+2}C_0^2)^2C_{\ref{lem:smoothdecay2}, 2N_{\ref{prop:1-off-W}}}  C_{\ref{prop:1-off-W},N_{\ref{prop:1-off-W}}} $.
\end{prop}

\ls{Maybe this can be proved without using distributions? The original induction statement in n-transformations.tex essentially included this as the case $k=n$ (the meaning of $k$ has changed since then, so now it is $k=1$). See Proposition 1.3 there.}

\ls{More specifically, if $\phi_j:=\phi_{0,a(j)}-\phi_{1,a(j+1)}$ then $(\widehat{\phi_j})^2=(\widehat{\phi_{0,a(j)}})^2-(\widehat{\phi_{0,a(j+1)}})^2$. So, the proposition becomes (after adding positive terms corresponding to the missing intermediate scales) InductPositiveTerms(1), up to a replacement of the Gaussians under the square root by windows. The proof of the main argument should then apply essentially verbatim. In a way, it would make sense to call this proposition InductPositiveTerms(1) because the Gaussian variant is never needed for $k=1$. On the other hand, I understand that we probably do not want to deal with windows in Section~\ref{sec:main argument}.}

\begin{proof}[Proof of Proposition \ref{prop:1-on-nonW-gap}]  
Let $J\ge 1$ be an integer, let 
${M} = \sum_{j\in [J)} M_j $, 
and let $\tilde{M}_1$ denote the distribution \[
\tilde{M}_1 =
(\delta_0-\phi_{1,a(0)})^{\otimes 2}+\sum_{j\in [2J-1)} (\phi_{0,a(j)}-\phi_{1,a(j+1)})^{\otimes 2} + 
(\phi_{0,a(2J-1)})^{\otimes 2},\]
where $\delta_0$ is the Dirac $\delta$ at $0$.

Let $\F\in \mathfrak{F}$. By Lemma \ref{lem:form_pos} (ii) applied to each summand, 
\[ \Lambda(\tilde{M}_1-{M})(\F) = \Lambda((\delta_0-\phi_{1,a(0)})^{\otimes 2})(\F) \] 
\[ + \sum_{j\in [J-1)} \Lambda((\phi_{0,a(2j+1)}-\phi_{1,a(2j+2)})^{\otimes 2})(\F) + \Lambda(\phi_{0,a(2J-1)}^{\otimes 2})(\F)\ge 0 \]
and also $\Lambda({M})(\F)\ge 0$.
By Lemma \ref{Lambda1 to Lambda}, $ \Lambda_1({M})(\F)= \Lambda({M})(\F)$. 
Therefore, 
\[ |\Lambda_1({M})(\F)| = |\Lambda({M})(\F)| =  \Lambda({M})(\F) \]
\[= \Lambda(\tilde{M}_1)(\F) - \Lambda(\tilde{M}_1-{M})(\F) \le  \Lambda(\tilde{M}_1)(\F), \]
so that it now suffices to estimate $\Lambda(\tilde{M}_1)(\F)$.
%and thus in turn $\|N_{\ref{prop:1-on-W}}\|_{\rm M(1)}$.


It will be important that 
\[\widehat{\tilde{M}_1}(\xi,\xi)=1\] for all $\xi\in\mathbb{R}$ by Lemma 
\ref{lem:diagsumone}.


 
Using  multilinearity we write 
\begin{equation}\label{M1split}
\Lambda(\tilde{M}_1)(\F) = \Lambda(\tilde{M}_2)(\F) + \Lambda(\tilde{M}_3)(\F) + \Lambda(\tilde{M}_4)(\F)
\end{equation}
with
\begin{equation}\label{eqn:ondiagK2def}
\tilde{M}_{2} = 
\delta_0^{\otimes 2}  -(\delta_0 - \phi_{1,2^{-4}a(0)})\otimes \phi_{1,a(0)} -  \phi_{1,a(0)}\otimes (\delta_0 - \phi_{1,2^{-4}a(0)}) 
\end{equation}

 \[\tilde{M}_3 = \sum_{j\in[2J-1)}M_{3,j} ,\qquad \tilde{M}_4 = \sum_{j\in[2J)}M_{4,j}\]
\begin{equation}\label{eqn:ondiagK3def}
M_{3,j} = - \bigl(\phi_{0,a(j)}-\phi_{1,2^{-4}a(j+1)}\bigr)\otimes \phi_{1,a(j+1)}
- \phi_{1,a(j+1)}\otimes \bigl(\phi_{0,a(j)}-\phi_{1,2^{-4}a(j+1)}\bigr)
\end{equation}
\begin{equation}\label{eqn:ondiagK4def}
M_{4,j}=-\bigl(\phi_{1,2^{-4}a(j)}-\phi_{1,a(j)}\bigr)\otimes \phi_{1,a(j)}
-\phi_{1,a(j)}\otimes \bigl(\phi_{1,2^{-4}a(j)}-\phi_{1,a(j)}\bigr)
\end{equation}
\begin{equation}\label{eq:ondiagK4def2}
+\phi_{0,a(j)}^{\otimes 2}
-\phi_{1,a(j)}^{\otimes 2}.
\end{equation}

Evaluated on the frequency diagonal $\xi=-\eta$, $\widehat{\tilde{M}_2}$ equals one. Indeed, $\widehat{\delta_0^{\otimes 2}} = 1$, 
\[\textup{supp} ( (1-\widehat{(\phi_{1,2^{-4}a(0)})} \otimes \widehat{\phi_{1,a(0)}} ) \subset [2^{3}a(0)^{-1},\infty)\times [-a(0)^{-1},a(0)^{-1}] \]
 which is disjoint from $\xi=-\eta$,   and by symmetry   the frequency support of the third term in $\tilde{M}_2$ is disjoint from the diagonal as well.


On the other hand, $\widehat{\tilde{M}_3}$ is equal to zero on that diagonal since 
\[\textup{supp} ( \bigl(\widehat{\phi_{0,a(j)}}-\widehat{\phi_{1,2^{-4}a(j+1)}}\bigr)\otimes \widehat{\phi_{1,a(j+1)})} \]
\[\subset ([-a(j)^{-1},-2^{3}a(j+1)^{-1}] \cup  [2^{3}a(j+1)^{-1},2^{-k_j}]) \times [-a(j+1)^{-1},a(j+1)^{-1}],
\]
(use Lemma \ref{lem:rescaled_diff}) which is disjoint from $\xi=-\eta$ for each $j$. 
Thus, since $\widehat{\tilde{M}_1}$ is equal to one on the diagonal, we must have that $\widehat{\tilde{M}_4}$ equals zero on the diagonal.

Using distributivity and applying 
 Lemma \ref{special prism BL inequality}  gives 
\begin{equation}\label{eqn:ondiagM2triv}
|\Lambda(\tilde{M}_2)(\F)|\le C_{\ref{eqn:ondiagM2triv}} 
\end{equation}
with $C_{\ref{eqn:ondiagM2triv}} = 2^2(2C_{\ref{lem:smoothdecay2},2}C_0)^2$,
  where  we used that $\|\phi_{1,2^{-4}a(0)}\|_1 =\|\phi_{1,a(0)}\|_1=\|\phi_{1}\|_1\le 2C_0C_{\ref{lem:smoothdecay2},2}$, which follows from Lemma \ref{lem:smoothdecay2} with $N=2$ and integrating $\langle x\rangle^{2}$.
 
% with $C_{\ref{eqn:ondiagK2triv}}= 2^3(C_0 C_{\ref{lem:smoothdecay2},2})^2 $,


 
Lemma \ref{lem:1-off-nonW} together with Lemma \ref{Lambda1 to Lambda} give
\begin{equation}
    \label{eqn:ondiagM3}
    |\Lambda(\tilde{M_3})(\F)| = |\Lambda_1(\tilde{M_3})(\F)|\le 2C_{\ref{lem:1-off-nonW},n }J^{\alpha(n)}
\end{equation}
where
\[\alpha(n) = 1-2^{-n+2}.\]



To bound $\Lambda(\tilde{M_4})$ we will use Proposition \ref{prop:1-off-W}. 
We verify its assumptions. First, 
\[\textup{supp}((\widehat{\phi_{1,2^{-4}a(j)}} -\widehat{\phi_{1,a(j)}})\otimes \widehat{\phi_{1,a(j+1)}})\]\[ \subset ([-2^{4}a(j)^{-1},-2^{-1}a(j)^{-1}]\cup [2^{-1}a(j)^{-1},2^{4}a(j)^{-1}] )\times [-a(j)^{-1},a(j)^{-1}]\]
and 
\[ \textup{supp}( \widehat{\phi_{0,a(j)}^{\otimes 2}}
-\widehat{\phi_{1,a(j)}^{\otimes 2}}) \subset [-a(j)^{-1},a(j)^{-1}]^2\setminus [-2^{-1}a(j)^{-1},2^{-1}a(j)^{-1}]^2\]
Thus, the Fourier support of $M_{4,j}$ is contained in 
\[\textup{Ann}_2(-a(j),2^{30}).\]
 By Lemma \ref{lem:smoothdecay2} applied with $N=2N_{\ref{prop:1-off-W}}$, Lemma \ref{lem:bump_diagest} with $N=N_{\ref{prop:1-off-W}}$,  and using the bound on the derivatives of windows,  
 \begin{equation}
     \label{ondgnw_const}
     |(\phi_{1,2^{-4}}-\phi_{1})\otimes \phi_{1}|(u,v)   \le {C_{\ref{ondgnw_const}}} \langle u+v \rangle^{N_{\ref{prop:1-off-W}}} \langle u-v\rangle^{N_{\ref{prop:1-off-W}}}
 \end{equation}
\[| \phi_{0}^{\otimes 2}
-\phi_{1}^{\otimes 2}|(u,v)  \le {C_{\ref{ondgnw_const}}}  \langle u+v\rangle^{N_{\ref{prop:1-off-W}}} \langle u-v\rangle^{N_{\ref{prop:1-off-W}}}\]
with ${C_{\ref{ondgnw_const}}}=2^4 2^{2N_{\ref{prop:1-off-W}}} (2^{2N_{\ref{prop:1-off-W}}+2}C_0^2)^2C_{\ref{lem:smoothdecay2}, 2N_{\ref{prop:1-off-W}}}$
Therefore,    
\[|M_{4,j}|\le 3{C_{\ref{ondgnw_const}}} \langle u+v\rangle^{N_{\ref{prop:1-off-W}}} \langle u-v\rangle^{N_{\ref{prop:1-off-W}}}.\]
By Proposition \ref{prop:1-off-W} applied with $\kappa=0$ together with Lemma \ref{Lambda1 to Lambda}, 
\begin{equation}
    \label{eqn:ondiagM4}
     |\Lambda(\tilde{M}_4)(\F)| =  |\Lambda_1(\tilde{M}_4)(\F)|  \le  3{C_{\ref{ondgnw_const}}}C_{\ref{prop:1-off-W},N_{\ref{prop:1-off-W}}} J^{\alpha(n)}.
    \end{equation}
 

Using \eqref{M1split} and estimates \eqref{eqn:ondiagM2triv}, \eqref{eqn:ondiagM3}, \eqref{eqn:ondiagM4} gives   $|\Lambda_1(M)(\F)|\le C_{\ref{prop:1-on-nonW-gap},n}J^{\alpha(n)}$. Since $J\ge 1 $ was arbitrary, we obtain $\|\M\|_{\rm M(1)}\le C_{\ref{prop:1-on-nonW-gap},n}$. 
\end{proof}

 

 

\begin{prop}[$k=1$ prism, on-diagonal, Whitney, with gap]\label{prop:1-on-W-gap} \pd{done In Prop \ref{P:induct-positive-terms-reduction-whitney-gap}}
There exists %$N=N_{\ref{prop:1-on-W-gap}}\ge 1$ and 
$C_{\ref{prop:1-on-W-gap}}\in (0,\infty)$ such that the following holds: Let $\kappa\ge 0$. Let   $a\in {\rm A}$   satisfy $a(j+1)\ge 2^{100} a(j)$. Let $M_j:\mathbb{R}^2\to \mathbb{R}$ be of the form
\[ M_j=\Psi_{(a(j))} \]
for every $j\in\Z$, where $\Psi:\mathbb{R}^2\to \mathbb{R}$ is a test function satisfying $\Psi(u_0,u_1)=\Psi(u_1,u_0)$ for all $u=(u_0,u_1)\in\mathbb{R}^2$,
\begin{equation}
    \label{p-1-on-W_pos_gap}
    \int_{\mathbb{R}^2} g(u_0){g(u_1)}\Psi(u)\,du\ge 0
\end{equation}
for all bounded measurable $g:\mathbb{R}\to\mathbb{R}$,
\begin{equation}
    \label{p-1-on-W_supp_gap}
    \mathrm{supp}(\widehat{\Psi}) \subset \mathrm{Ann}_1(1, 2^{20})^2,
\end{equation}
and
\begin{equation}
    \label{p-1-on-W_decay_gap}
     |\Psi(u)| \le 2^{-\kappa/2} 
     \langle u_0+u_1\rangle_{(2^\kappa)}^{N_{\ref{prop:1-off-W}}}
     \langle u_0-u_1\rangle^{N_{\ref{prop:1-off-W}}}
\end{equation}
for all $u\in\mathbb{R}^2$.
Then with $\M=(M_j)_{j\in\Z}$
\[\|\M\|_{\rm M(1)} \le C_{\ref{prop:1-on-W-gap}, n}\]
where $C_{\ref{prop:1-on-W-gap}, n} = 2C_{\ref{prop:1-on-nonW-gap}} + C_{\ref{prop:1-off-W},n}2^{56+2N_{\ref{prop:1-off-W}}}C^2_{\ref{lem:smoothdecay2},2N_{\ref{prop:1-off-W}}}(C_0^2 + C^2_{\ref{lem:Phij_prop},2N_{\ref{prop:1-off-W}}})$
\end{prop}


\begin{proof}[Proof of Proposition \ref{prop:1-on-W-gap}]
 We use Lemma \ref{lem:Phij_prop} with   $N=2N_{\ref{prop:1-off-W}}$
to obtain a function $\psi$ satisfying \eqref{Phij_psi_diag} and \eqref{psi_decay} for all $0\le m \le 2N_{\ref{prop:1-off-W}}$. 
\ls{So, the constant $N$ in~\eqref{p-1-on-W_decay_gap} seems to be $2N_{\ref{prop:1-off-W}}+2$.}
Let $\alpha\in A$ be defined as $\alpha(2j)=2^{-25}a(j)$ and $\alpha(2j+1)=2^{25}a(j)$ 
  Then by \eqref{Phij_psi_diag}, for  each $j\in \Z$, 
\begin{equation}
    \label{psi_diag}
    \widehat{\psi_{(a(j))}}(\xi)^2 = 2 (\widehat{\phi_{0,\alpha(2j)}}-\widehat{\phi_{1,\alpha(2j+1)}})(\xi)^2 - \widehat{\Psi_{(a(j))}}(\xi,-\xi). 
\end{equation}


Decompose  $M_j=M_{0,j}-M_{1,j}$ with
\[M_{0,j} =   \Psi_{(a(j))} + (\psi_{(a(j))})^{\otimes 2}  , \]
\[M_{1,j} =   (\psi_{(a(j))})^{\otimes 2} .  \]
  Let $J\ge 1,\, M = \sum_{j\in[J)}M_j$, and $\F\in \mathfrak{F}$.  The condition \ref{p-1-on-W_pos_gap} and Lemma \ref{lem:form_pos}  applied term by term  gives that     \[\Lambda_1(M)(\F)\ge 0,\quad \Lambda_1(\sum_{j\in[J)}M_{1,j})(\F) \ge 0\]
Therefore, 
  \begin{equation}
    \label{K0_split}
    |\Lambda_1(M)(\F)| = \Lambda_1(M)(\F) \end{equation}
    \[=\Lambda_1(\sum_{j\in[J)}M_{0,j})(\F)-\Lambda_1(\sum_{j\in[J)}M_{1,j})(\F)\le \Lambda_1(\sum_{j\in[J)}M_{0,j})(\F).\]


We now write $M_{0,j} = M_{2,j} + M_{3,j}$,   
where \[M_{2,j} = 2(\phi_{0,\alpha(2j)} -\phi_{1,\alpha(2j+1)} )^{\otimes 2}\] 
and   $M_{3,j} = M_{0,j} - M_{2,j}$.
% so that 
% \[ \Lambda_1(\sum_{j\in[J)}M_{0,j})) = \Lambda_1(\sum_{j\in[J)}M_{3,j})) + \Lambda_1(\sum_{j\in[J)}M_{2,j})). \]
Proposition \ref{prop:1-on-nonW-gap}   gives
\[\|(M_{2,j})_{j\in\Z}\|_{\rm M(1)}\le 2C_{\ref{prop:1-on-nonW-gap}}.\]
By multilinearity of $\Lambda_1$,  it thus remains to estimate  $\|(M_{3,j})_{j\in\Z}\|_{\rm M(1)}$, which will be done by an application of  Proposition~\ref{prop:1-off-W}.

 
 We now  verify the  assumptions of Proposition~\ref{prop:1-off-W}.
Evaluating the Fourier transform at $(\xi,-\xi)$, 
\[ \sum_{j\in [J)} \widehat{M_{3,j}}(\xi,-\xi) %= \sum_{j\in [J)} \widehat{\Psi_{j}}(\xi,-\xi) + \widehat{\psi_j}(\xi)\widehat{\psi_j}(-\xi) - c_\Psi  (\widehat{\phi_{0,j}} -\widehat{\phi_{1,j}} )(\xi)(\widehat{\phi_{0,j}} -\widehat{\phi_{1,j}} )(-\xi)
= \sum_{j\in [J)} \widehat{\Psi_{(a(j))}}(\xi,-\xi) + \widehat{\psi_{(a(j))}}(\xi)^2- 2 (\widehat{\phi_{0,\alpha(2j)}} -\widehat{\phi_{1,\alpha(2j+1)}} )(\xi)^2  = 0\]
where we used \eqref{psi_diag} and the fact that $\psi$ and windows are even. 
We have
\[\textup{supp}(\widehat{\Psi}) \subset \textup{Ann}_1(1,2^{20})^2 \subset \textup{Ann}_2(1,2^{21}) \]
Since $\textup{supp}(\widehat{\phi_{0,2^{-25}}})\subset [-2^{25},2^{25}]$ and $\widehat{\phi_{0,2^{-25}}} =1$ on $[-2^{24},2^{24}]$ and 
 $\textup{supp}(\widehat{\phi_{1,2^{25}}})\subset [-2^{-25},2^{-25}]$ and $\widehat{\phi_{1,2^{25}}} =1$ on $[-2^{-24},2^{-24}]$, we have
 \[\textup{supp}(\widehat{\phi_{0,2^{-25}}} -\widehat{\phi_{1,2^{25}}} ) \subset \textup{Ann}_1(1,2^{26})\]
 and thus 
\[\textup{supp}((\widehat{\phi_{0,2^{-25}}} -\widehat{\phi_{1,2^{25}}} )^{\otimes 2} ) \subset \textup{Ann}_1(1,2^{26})^2  \subset  \textup{Ann}_2(1,2^{27})\]
By Lemma \ref{lem:Phij_prop}, also 
\[\textup{supp}(\widehat{\psi}^{\otimes 2})\subset \textup{Ann}_1(1,2^{27})^2 \subset \textup{Ann}_2(1,2^{28})\]
Therefore, the function 
\[\widetilde{\Psi} = \Psi + \psi^{\otimes 2} - 2(\widehat{\phi_{0,2^{-25}}} -\widehat{\phi_{1,2^{25}}} )^{\otimes 2}  \]
satisfies 
\[\textup{supp}(\widehat{\widetilde{\Psi}})\subset \textup{Ann}_2(1,2^{30})\]

  
Let $N=N_{\ref{prop:1-off-W}}$. Using \eqref{windowbound} and  Lemma \ref{lem:smoothdecay2} with $2N$  applied to the functions  $C_{0}^{-1}\phi_0$ and $C_{0}^{-1}\phi_1$,
\[|\phi_0-\phi_1|(u)\le |\phi_0(u)|+|\phi_1(u)| \le  2^2C_{\ref{lem:smoothdecay2},2N}C_0 \langle u\rangle^{2N}\] and thus
\[|(\phi_0-\phi_1)^{\otimes 2}|(u,v) \le 2^4C^2_{\ref{lem:smoothdecay2},2N}C^2_0 \langle u\rangle^{2N} \langle v\rangle^{2N} \]
By Lemma \ref{lem:bump_diagest}, this is 
\begin{equation}
\label{decay_1}
\le 2^{4+2N}C^2_{\ref{lem:smoothdecay2},2N}C^2_0 \langle u+v \rangle^{N} \langle u-v \rangle^{N}
\end{equation}
Similarly, by Lemma \ref{lem:smoothdecay2} applied with $2N$ and $C_{\ref{lem:Phij_prop},2N}^{-1}\psi$   (using that $C_{\ref{lem:Phij_prop},2N}^{-1}\psi$ satisfies \eqref{psi_decay}) and using $|\mathrm{supp}\;\widehat{\psi}|\le 2^{28}$,
\begin{equation}
    \label{decay_2}
  |\psi^{\otimes 2}|(u,v) \le 2^{56}C^2_{\ref{lem:smoothdecay2},2N} C^2_{\ref{lem:Phij_prop},2N} \langle u\rangle^{2N} \langle v\rangle^{2N} 
  \end{equation}
\[  \le 2^{56+2N}C^2_{\ref{lem:smoothdecay2},2N} C^2_{\ref{lem:Phij_prop},2N} \langle u+v \rangle^{N} \langle u-v \rangle^{N}\]
where for the last inequality we used Lemma \ref{lem:bump_diagest}. 
The estimates \eqref{decay_1} and \eqref{decay_2},  together with the assumption \eqref{p-1-on-W_decay_gap}, give
\[|\widetilde{\Psi}(u,v)| \le c_1 (2^{-\frac{1}{2}\kappa}\langle u+v\rangle_{(2^\kappa)}^{N}\langle u-v\rangle^{N} + \langle u+v \rangle^{N} \langle u-v \rangle^{N} )\]
with $c_1=2^{56+2N}C^2_{\ref{lem:smoothdecay2},2N}(C_0^2 + C^2_{\ref{lem:Phij_prop},2N})$
Applying Proposition~\ref{prop:1-off-W} yields
\[\|(M_{3,j})_{j\in\Z}\|_{\rm M(1)} \le c_1 C_{\ref{prop:1-off-W},n},\]
which completes the proof.
%Together with \eqref{K0_split} and \eqref{K3_split} this gives
%\[|\Lambda_n(K)| \le (C_{\ref{prop:1-off-W}}c_1  + 2C_{\ref{prop:1-on-nonW-gap}}) J^{\frac12} \le   C_{\ref{prop:1-on-W-gap}}J^{\frac12}. \]
\end{proof}




\begin{prop}[$k=1$ prism, on-diagonal, Whitney]\label{prop:1-on-W} \pd{done In Prop \ref{P:induct-positive-terms-reduction-whitney}}
Let $\kappa\ge 0$ and $a\in {\rm A}$. Let $M_j:\mathbb{R}^2\to \mathbb{R}$ be of the form
\[ M_j=\Psi_{(a(j))} \]
for every $j\in\Z$, where $\Psi:\mathbb{R}^2\to \mathbb{R}$ is a test function satisfying $\Psi(u_0,u_1)=\Psi(u_1,u_0)$ for all $u=(u_0,u_1)\in\mathbb{R}^2$,
\begin{equation}
    \label{p-1-on-W_pos}
    \int_{\mathbb{R}^2} g(u_0){g(u_1)}\Psi(u)\,du\ge 0
\end{equation}
for all bounded measurable $g:\mathbb{R}\to\mathbb{R}$,
\begin{equation}
    \label{p-1-on-W_supp}
    \mathrm{supp}(\widehat{\Psi}) \subset \mathrm{Ann}_1(1, 2^{20})^2,
\end{equation}
and 
\begin{equation}
    \label{p-1-on-W_decay}
     |\Psi(u)| \le 2^{-\kappa/2} 
     \langle u_0+u_1\rangle_{(2^\kappa)}^{N_{\ref{prop:1-on-W}}}
     \langle u_0-u_1\rangle^{N_{\ref{prop:1-on-W}}}
\end{equation}
for all $u\in\mathbb{R}^2$, where $N_{\ref{prop:1-on-W}}=N_{\ref{prop:1-off-W}}$.
Then with $\M=(M_j)_{j\in\Z}$
\[\|\M\|_{\rm M(1)} \le C_{\ref{prop:1-on-W},n}.\]
where $C_{\ref{prop:1-on-W},n} = 100C_{\ref{prop:1-on-W-gap}, n}$.
\end{prop}

\begin{proof}
%Let $N=N_{\ref{prop:1-on-W-gap}}$.
Let $J\ge 1$ and write $[J) = \cup_{s \in [100)} \mathcal{I}_s$
where 
\[\mathcal{I}_s = \{j\in [J): j  \equiv s \; \textup{mod} \; 100\}.\]
Let $M=\sum_{j\in [J)} M_j$ and $\F\in\mathfrak{F}$.
We split using multilinearity
\[\Lambda_1(M)(\F) = \sum_{s\in [100)} \Lambda_1(M)(\F). \]
The claim now follows from Proposition \ref{prop:1-on-W-gap} applied with $\kappa=0$    for   $s\in [100)$, and   summing in $s$.
\end{proof}



\begin{proposition}
\label{prop:1-on-W-cor} \pd{done In Prop \ref{P:induct-positive-terms-reduction-whitney-product}}
Let $a\in A$ and for $j\in \Z$ let $M_j:\mathbb{R}^2\to \mathbb{R}$  be given by
\[ M_j=   (\psi_{(a(j))})^{\otimes 2}, \]
 where   $\psi:\mathbb{R}\to \mathbb{R}$ satisfies
%there is a strictly increasing sequence of integers $(k_j)_{j\in [J)}$ such that for all $j\in [J)$
\begin{equation}\label{1-on-W_supp}
\mathrm{supp}(\widehat{\psi}) \subset \mathrm{Ann}_1(1, 2^{20})
\end{equation}
and 
\begin{equation}\label{1-on-W_decay}
|\psi (u)| \le \langle u\rangle^{2N_{\ref{prop:1-off-W}}}
\end{equation}
for all $u\in\mathbb{R}$.
Then for $\M=(M_j)_{j\in \Z}$
\[ \|\M\|_{\rm M(1)}\le C_{\ref{prop:1-on-W-cor},n}  \]
 with $C_{\ref{prop:1-on-W-cor},n} = C_{\ref{prop:1-on-W},n}$. 
\end{proposition}

\begin{proof}
%Let $N=2N_{\ref{prop:1-on-W}}$. 
We apply Proposition \ref{prop:1-on-W} with $\Psi = \psi\otimes \psi$.  Then
\[\supp(\widehat{\Psi})\subset \textup{Ann}_1(1,2^{20})^2\]
and for all $u\in \R^2$, 
\[|\Psi(u)|\leq \langle u_0\rangle ^{2N_{\ref{prop:1-off-W}}} \langle u_1 \rangle^{2N_{\ref{prop:1-off-W}}} \leq \langle u_0+u_1\rangle^{N_{\ref{prop:1-off-W}}}\langle u_0-u_1\rangle^{N_{\ref{prop:1-off-W}}}\]
Moreover,
\[\Psi(u_0,u_1) = \psi(u_0)\psi(u_1) = \Psi(u_1,u_0)\]
and the positivity assumption \eqref{p-1-on-W_pos} is satisfied by \eqref{form_pos_psipsi} in Lemma \ref{lem:form_pos}. 
\end{proof}




%The following Proposition \ref{prop:1-on-nonW} will be proved using Propositions  \ref{prop:1-on-W} and \ref{prop:1-on-nonW-gap}. 
\begin{prop}[$k=1$ prism, on-diagonal, non-Whitney]\label{prop:1-on-nonW} \todo{Prop \ref{P:induct-positive-terms-reduction-non-whitney}}
There exists $C_{\ref{prop:1-on-nonW}}\in (0,\infty)$ such that the following holds.

Let  $(k_j)_{j\in\Z}$ be a strictly increasing sequence of integers and let $a\in \rm A$ be given by $a(j)=2^{k_j}$ for $j\in \Z$, and  let
\[M_j=({\phi}_{0,a(2j)}-{\phi}_{1,a(2j+1)})^{\otimes 2}.\]
Then with $\M=(M_j)_{j\in \Z}$,  \[\|\M\|_{\rm M(1)}\le C_{\ref{prop:1-on-nonW},n}.\]
where $C_{\ref{prop:1-on-nonW},n} =C_{ \ref{prop:1-on-nonW-gap},n } + 10C_{\ref{prop:1-on-W-cor},n}$
\end{prop}


 
\begin{proof}[Proof of Proposition \ref{prop:1-on-nonW}] 
For  $J\ge 1$ we   write 
\[\sum_{j\in [J)} M_j =\Big(\sum_{s\in [10)}\sum_{j\in \mathcal{J}_s} M_j  \Big)+\sum_{j\in \mathcal{J}_{\ge 10}} M_j,\] where  
\[ \mathcal{J}_{s} = \{j \in [J) \,:\,  a(j+1) =  2^{s}a(j) \},\]
\[\mathcal{J}_{\ge  10} = \{j\in[J)\,:\, a(j) 2^{10} \le  a(j+1) \}.\]
Let $\F\in \mathfrak{F}$. By the triangle inequality, 
\[|\Lambda_1(\sum_{j\in [J)} M_j)(\F)| \leq  \Big(\sum_{s\in [10)} |\Lambda_1(\sum_{j\in \mathcal{J}_s} M_j )(\F)|\Big) + |\Lambda_1(\sum_{j\in \mathcal{J}_{\ge 10}} M_j)(\F)|\]

Proposition \ref{prop:1-on-nonW-gap}  gives
\[\|(M_j)_{j\in \mathcal{J}_{\ge  10}}\|_{\rm M(1)}\le C_{ \ref{prop:1-on-nonW-gap},n }.\]

To estimate $(M_{j})_{j\in \mathcal{J}_s}$,   we will apply Proposition \ref{prop:1-on-W-cor}, and finally sum in    $s\in [10)$.
If  $j\in \mathcal{J}_s$, then 
\[\supp(\widehat{\phi_{0}} - \widehat{(\phi_1)_{(2^s)}}) \subset  [-1,-2^{-11}]\cup [2^{-11},1] \subset \textup{Ann}_1(1, 2^{18}). \]
Proposition \ref{prop:1-on-W-cor} now gives   (with $\psi$ in Proposition \ref{prop:1-on-W-cor} being $\phi_0 - (\phi_1)_{(2^s)}$)
\[\|(M_j)_{j\in \mathcal{J}_s}\|_{\rm M(1)} \le C_{\ref{prop:1-on-W-cor},n}\] for each $s\in [10)$.
% \pd{Thus, with $\Psi = \psi\otimes \psi$, 
% \[\supp(\widehat{\Psi})\subset \textup{Ann}_1(0,2^{20})^2\]
% and for all $u\in \R^2$,   
% \[|\Psi(u)|\leq \langle u_0\rangle^{2M} \langle u_1 \rangle^{2M} \leq \langle u_0+u_1\rangle^{M}\langle u_0-u_1\rangle^{M}\]
% Moreover,
% \[\Psi(u_0,u_1) = \psi(u_0)\psi(u_1) = \Psi(u_1,u_0)\]
% and  the  positivity   \eqref{p-1-on-W_pos} holds by \eqref{form_pos_psipsi}.}
% % in Lemma \ref{lem:form_pos}. 
\end{proof}

\section{Other lemmas not needed - will eventually be deleted}
\begin{lemma}[Derivative of sqrt] 
\label{derivative of sqrt}
Let $\xi\in \R$ and $\phi$ a Schwartz function. Assume there is $\lambda>0$ such that  
\begin{equation}
    \label{dersqrt_asmpt}
    |(\phi^2)^{(l)}(\xi)|\le \lambda^{l}
\end{equation}
for all $0\le l \le m$. 
Assume also that there exists a constant $0< c_0\le 1$ such that  
\[|\phi(\xi)|\ge c_0\]
Then    for each integer  $m\ge 0$, 
\[| \phi^{(m)}(\xi)|\le C_{\ref{derivative of sqrt},m,c_0}\lambda^{m} \]
with  $C_{\ref{derivative of sqrt},m,c_0} = \max(1,2^{2^{m+1}-2} c_0^{-2m+1})$. 
\end{lemma}
\begin{proof}
If $m=0$, we have $|\phi(\xi)|\le 1$ by  the assumption \eqref{dersqrt_asmpt} with $l=0$. So from now on we assume $m\ge 1$. 

Let $1\le a\le m$. 
We compute using the Leibniz rule
\[  (\phi^2)^{(a)}(\xi) = \sum_{l=0}^a {\binom{a}{l}} \phi^{(l)}(\xi) \phi^{(a-l)}(\xi)\]
       \[= 2 \phi(\xi) \phi^{(a)}(\xi) + \sum_{l=1}^{a-1} {\binom{a}{l}} \phi^{(l)}(\xi) \phi^{(a-l)}(\xi)\]
Therefore, using that    $(2|\phi(\xi)|)^{-1}\le c_0^{-1}$ and \eqref{dersqrt_asmpt},
          \begin{equation}
          \label{phia_bd}
           |\phi^{(a)}(\xi)| \le  c_0^{-1}\Big(\lambda^a +  \sum_{l=1}^{a-1} {\binom{a}{l}} 
         |\phi^{(l)}(\xi)| |\phi^{(a-l)}(\xi)|\Big)
        \end{equation}
        
% We claim that for each $1\le a \le a_0$,
% \[|\phi^{(a)}(\xi)| \le  2^{2^{a+1}-2}B^{2a-1} \lambda^{a} \]
% This bound is maximized for  $a=a_0$,  in which case it is bounded by $ 2^{2^{a_0+1}}c_0^{-2a_0+1}\lambda^{a_0}$,  as desired. 

  
Let $B=c_0^{-1}$, so $B\ge 1$. 
If $m=1$,  we use \eqref{phia_bd} with $a=1$, 
giving 
  \[|\phi'(\xi)| \le B \le 2^{2^{2}-2}B\lambda =4B\lambda.\]   
If    $m\ge 2$ we will inductively show that if 
\[
    |\phi^{(l)}(\xi)| \le 2^{2^{l+1}-2}B^{2l-1}\lambda^l
\]
     holds for all $1\le l\le a-1$ where $2\le a\le m$,  that it also holds  for $l=a$. The claim then follows by setting $a=m$.  
 
For $l=a$
we estimate using \eqref{phia_bd} 
   \[|\phi^{(a)}(\xi)|  \le B\Big( \lambda^a +  \sum_{l=1}^{a-1} {\binom{a}{l}} 
         2^{2^{l+1}-2}B^{2l-1}\lambda^l 2^{2^{a-l+1}-2}B^{2(a-l)-1} \lambda^{a-l}
           \Big) \]
           \[=B\lambda^a\Big( 1 +  B^{2(a-1)} \sum_{l=1}^{a-1} {\binom{a}{l}} 
         2^{2^{l+1} + 2^{a-l+1}-4}
           \Big) \]
           The function $l\mapsto 2^{l+1} + 2^{a-l+1}-4$ achieves its maximum for $l=a-1$ on $[1,a-1]$, at which point it equals $2^{a}$. Using $\sum_{l=1}^{a-1}{a \choose l} = 2^a-2$
        \[\le  B\lambda^a(1+  B^{2a-2} (2^{a}-2)2^{2^{a}})  = \lambda^a(B+ B^{2a-1}(2^a-2)2^{2^a})\]
        Since $B\ge 1$, this is 
        \[\le \lambda^a B^{2a-1} (1+2^{a+2^a}-2^{2^a+1})\le\lambda^a  B^{2a-1}2^{a+2^{a}}\le \lambda^a B^{2a-1}2^{2^{a+1}-2},\]
        as desired.
    In the last inequality we used   for $a\ge 2$, 
        \[a+2^{a} \le 2^{a+1}-2.\]
\end{proof}

\begin{theorem}
\label{induct positive terms theorem-old}
For every $k\in\N$ with $1\le k\le n$, $\InductPositiveTerms{k,C_{\ref{induct positive terms theorem-old}}n^2}$ holds,
where $C_{\ref{induct positive terms theorem-old}}=2^{20} C_{\ref{induct positive terms imply increase data}}$.
\end{theorem}

\begin{proof}
This is proved by (reverse) induction on $k$ with the base case being $k=n$.
We claim that
$\InductPositiveTerms{k,C_k}$ holds for all $1\le k\le n$ where
\[ C_k = \left\{\begin{array}{ll}
2+n, & \text{if}\;k=n,\\
2+(n-1)2^{10}C_{\ref{induct positive terms imply increase data}} C_n, & \text{if}\;k=n-1,\\
2+k 2^{10} (C_{\ref{induct positive terms imply increase data}}C_{k+1})^{\frac12}, & \text{if}\;1\le k\le n-2.
\end{array}
\right. \]
The claim holds for $k=n$ by Proposition \ref{vanishing diagonal implies induct positive terms} with $k=n$ and Proposition \ref{vanishing kernel integral}.
Next assume $1\le k\le n-1$ and assume that $\InductPositiveTerms{k+1,C_{k+1}}$ holds.
Then $\InductPositiveTerms{k,C_k}$
follows by applying Proposition \ref{induct positive terms imply increase data}, Proposition \ref{increase data implies diagonal band}, Proposition \ref{diagonal band implies vanishing diagonal} and finally Proposition \ref{vanishing diagonal implies induct positive terms}.
Finally, $C_k\le 2^{20} C_{\ref{induct positive terms imply increase data}} n^2$ follows by induction using the definition of $C_k$.
\end{proof}

\fi
\end{document}